%% main.tex — Tsinghua PL seminar, 90 minutes, 16:9.
%%
%% "Normal Forms and Complete Relations for Circuits in Agda"
%%
%% Build: see the Makefile (WSL): `make` runs Agda's LaTeX backend on
%% sections/*.lagda.tex (from the repository root, so that the library
%% resolves) and then lualatex.  Frames whose title carries the green
%% badge are the POPL submission (tag popl-revision1); the orange badge
%% marks developments on other branches since then.
\documentclass[aspectratio=169,9pt]{beamer}
\usetheme[progressbar=frametitle,numbering=fraction,sectionpage=progressbar]{metropolis}

\usepackage{fontspec}
\usepackage{unicode-math}
\setsansfont{Fira Sans}[BoldFont={Fira Sans SemiBold},
  ItalicFont={Fira Sans Italic},BoldItalicFont={Fira Sans SemiBold Italic}]
\setmathfont{Fira Math}
\setmathfont{STIX Two Math}[range={\wr,\rtimes,\ltimes,\ast,\circledast,\lBrack,\rBrack,\top,\bot,"22EF,"2A01,"2295}]
\newfontfamily\agdafont{JuliaMono}[Scale=MatchLowercase]
\AtBeginDocument{\providecommand{\llbracket}{\lBrack}\providecommand{\rrbracket}{\rBrack}}
\newfontfamily\agdafontfallback{DejaVu Sans}

\usepackage{booktabs}
\usepackage{tabularx}
\usepackage{multirow}
\usepackage{xcolor}
\usepackage{tikz}
\usetikzlibrary{arrows.meta,positioning,calc,fit,backgrounds,shapes,
  shapes.geometric,shapes.misc,decorations.pathreplacing,quantikz2,
  matrix,shadows}
\usepackage{tikz-cd}
\usepackage{pgfplots}
\pgfplotsset{compat=1.18}
\usepackage[most]{tcolorbox}

%% Agda's own style file (the copy Agda 2.8 generates); code blocks in
%% sections/*.lagda.tex are typeset by `agda --latex'.
\usepackage{agda}
\renewcommand{\AgdaFontStyle}[1]{{\agdafont #1}}
\renewcommand{\AgdaKeywordFontStyle}[1]{{\agdafont #1}}
\renewcommand{\AgdaStringFontStyle}[1]{{\agdafont #1}}
\renewcommand{\AgdaCommentFontStyle}[1]{{\agdafont\itshape #1}}
\renewcommand{\AgdaBoundFontStyle}[1]{{\agdafont #1}}
\AtBeginDocument{\renewcommand{\checkmark}{{\agdafontfallback\symbol{"2713}}}}
\renewcommand{\AgdaCodeStyle}{\footnotesize}
\setlength{\mathindent}{0pt}

%% TikZiT figures from the paper (generated by Cir2Tikz) -------------
\pgfdeclarelayer{edgelayer}
\pgfdeclarelayer{nodelayer}
\pgfsetlayers{background,edgelayer,nodelayer,main}
\pgfkeys{/tikz/tikzit fill/.initial=0}
\pgfkeys{/tikz/tikzit draw/.initial=0}
\pgfkeys{/tikz/tikzit shape/.initial=0}
\pgfkeys{/tikz/tikzit category/.initial=0}
\tikzstyle{none}=[inner sep=0mm]
\tikzstyle{tikzfig}=[baseline=-0.25em,scale=0.5]
\newcommand{\tikzfig}[1]{{\tikzstyle{every picture}=[tikzfig]\input{#1.tikz}}}
\definecolor{normGreen}{RGB}{214, 243, 221}
\definecolor{normRed}{RGB}{249, 196, 196}
\definecolor{normYellow}{RGB}{255, 245, 200}
\definecolor{zxgreen}{RGB}{216, 248, 216}
\definecolor{zxred}{RGB}{232, 165, 165}
\definecolor{zxhad}{RGB}{255, 255, 191}
\definecolor{zxpaulibox}{RGB}{138, 159, 198}
\input{circuits.tikzstyles}
\newcommand{\ufig}[2][1]{$\vcenter{\hbox{\scalebox{#1}{\tikzfig{figures/#2}}}}$}
%% Figures from the QPL'26 deck (branch qupit) were drawn for a tikzfig
%% scale of 0.38; \qfig reproduces that, then scales by the factor.
\newcommand{\qfig}[2][0.7]{{\tikzstyle{tikzfig}=[baseline=-0.25em,scale=0.38]%
  $\vcenter{\hbox{\scalebox{#1}{\tikzfig{figures/qupit/#2}}}}$}}

%% Colours -------------------------------------------------------------
\definecolor{paperC}{HTML}{2E7D32}   % in the POPL submission
\definecolor{newC}{HTML}{D35400}     % new since the submission
\definecolor{accent}{HTML}{1565C0}
\definecolor{soft}{HTML}{EEF3FA}
\definecolor{softnew}{HTML}{FDF0E6}
\definecolor{softpaper}{HTML}{EAF5EA}
\definecolor{wirec}{HTML}{37474F}
\setbeamercolor{alerted text}{fg=newC}

%% Parameterised TikZ styles (a `#1' cannot appear inside a frame body).
\tikzset{
  blob/.style={draw=#1,fill=#1!8,ellipse,minimum width=22mm,minimum height=32mm,thick},
  ly/.style={draw=accent,fill=#1,rounded corners=2pt,minimum height=6.2mm,
    font=\scriptsize,align=center,inner sep=2pt},
  w/.style={draw=#1,fill=#1!10,rounded corners=2pt,font=\small,
    minimum width=34mm,minimum height=8mm,align=center},
  tag/.style={draw=#1,fill=#1!12,rounded corners=2pt,font=\scriptsize,
    align=center,inner sep=3pt},
  cm/.style={circle,fill=#1,inner sep=1.6pt},
}
\setbeamercolor{example text}{fg=paperC}

%% Badges in frame titles ---------------------------------------------
\newcommand{\badge}[2]{\tikz[baseline=(b.base)]\node(b)[fill=#1,text=white,
  rounded corners=2pt,inner sep=2.5pt,font=\scriptsize\bfseries]{#2};}
\newcommand{\InPaper}{\hfill\badge{paperC}{in the paper}}
\newcommand{\New}{\hfill\badge{newC}{new since the paper}}
\newcommand{\Mixed}{\hfill\badge{paperC}{paper}\,\badge{newC}{+ new}}
\newcommand{\Branch}[1]{{\scriptsize\color{newC}\texttt{#1}}}

%% Inline identifiers / boxes -----------------------------------------
\newcommand{\aid}[1]{\mbox{{\agdafont\small #1}}}
\newcommand{\aidop}[1]{\ifmmode\mathbin{\aid{#1}}\else\aid{#1}\fi}
\newcommand{\aidrel}[1]{\ifmmode\mathrel{\aid{#1}}\else\aid{#1}\fi}
\newtcolorbox{keybox}[1][]{colback=soft,colframe=accent,boxrule=0.6pt,
  arc=2pt,left=4pt,right=4pt,top=2pt,bottom=2pt,fonttitle=\bfseries\small,#1}
\newtcolorbox{newdevbox}[1][]{colback=softnew,colframe=newC,boxrule=0.6pt,
  arc=2pt,left=4pt,right=4pt,top=2pt,bottom=2pt,fonttitle=\bfseries\small,#1}
\newtcolorbox{paperbox}[1][]{colback=softpaper,colframe=paperC,boxrule=0.6pt,
  arc=2pt,left=4pt,right=4pt,top=2pt,bottom=2pt,fonttitle=\bfseries\small,#1}
\newcommand{\src}[1]{{\par\hfill\scriptsize\color{gray}\texttt{#1}}}

%% Hand-drawn permutation circuits ------------------------------------
%% \permcirc{n}{p1,p2,...}: n wires (wire 0 at the bottom) and a swap of
%% wires p, p+1 at each successive time step; `x' in the list leaves a
%% gap.  Used for the S_3 normal forms and the normalization trace.
\newcommand{\permcirc}[3][0.32]{%
  \begin{tikzpicture}[x=#1cm,y=#1cm,baseline={([yshift=-0.5ex]current bounding box.center)},
      line width=0.7pt,draw=wirec]
    \def\nw{#2}\def\permgap{x}%
    \def\len{0}%
    \foreach \p [count=\i] in {#3} {\xdef\len{\i}}%
    \pgfmathsetmacro{\L}{2*\len+1}%
    \foreach \w in {1,...,\nw} {\draw (0,{(\w-1)*1.6}) -- (\L,{(\w-1)*1.6});}
    \foreach \p [count=\i] in {#3} {%
      \ifx\p\permgap\else
      \pgfmathsetmacro{\xa}{2*\i-1.45}\pgfmathsetmacro{\xb}{2*\i+0.45}%
      \pgfmathsetmacro{\ya}{\p*1.6}\pgfmathsetmacro{\yb}{(\p+1)*1.6}%
      \fill[white] (\xa,\ya-0.05) rectangle (\xb,\yb+0.05);
      \draw (\xa,\ya) -- (\xb,\yb); \draw (\xa,\yb) -- (\xb,\ya);
      \fi}
  \end{tikzpicture}}

%% Section pages: colour the progress bar by provenance ----------------
\newcommand{\PaperPart}{\setbeamercolor{progress bar}{fg=paperC,bg=paperC!20}}
\newcommand{\NewPart}{\setbeamercolor{progress bar}{fg=newC,bg=newC!20}}

\title{Normal Forms and Complete Relations\\for Circuits in Agda}
\subtitle{Presentations, coset tables, and machine-checked completeness}
\author{Xiaoning Bian}
\institute{Tsinghua University \\[2pt] \textcolor{gray}{Tsinghua PL Seminar}}
\date{October 2026}

\begin{document}

\maketitle

%% Opening: plan of the talk, provenance legend, take-home message.

\begin{frame}{Plan for the next 90 minutes}
\begin{columns}[T]
\begin{column}{0.56\textwidth}
\small
\begin{enumerate}
  \setlength{\itemsep}{2pt}
  \item Why completeness proofs need a machine \hfill{\color{gray}\small 7'}
  \item A complete worked example: permutation circuits \hfill{\color{gray}\small 15'}
  \item Normal forms from subgroup chains \hfill{\color{gray}\small 4'}
  \item The library, layer by layer \hfill{\color{gray}\small 18'}
  \item A catalogue of verified presentations \hfill{\color{gray}\small 10'}
  \item {\color{newC}New since the paper} \hfill{\color{gray}\small 31'}
  \begin{itemize}
    \footnotesize
    \item scalars and \alert{minimality} of relation sets
    \item qupit Clifford circuits, all odd primes
    \item \alert{matrix semantics}: $U_n(\mathbb{Z}[\frac12,i])$, 2-qubit Clifford+$T$,
          Real-Clifford+$CH$ over $\mathbb{Z}[\sqrt2]$
    \item CNOT-dihedral (phase polynomials); path sums (polynomial cost)
  \end{itemize}
  \item Related work, lessons, outlook \hfill{\color{gray}\small 5'}
\end{enumerate}
\end{column}
\begin{column}{0.40\textwidth}
\begin{keybox}[title=How to read the slides]
\footnotesize
Every frame title carries its provenance:\\[4pt]
\badge{paperC}{in the paper}\; the POPL'27 submission
\emph{Normal Forms and Complete Relations for Circuits in Agda}
(tag \texttt{popl-revision1}).\\[4pt]
\badge{newC}{new since the paper}\; developed afterwards on other
branches of the same repository (\Branch{qupit},
\Branch{v-office-Cli-CH}, \Branch{popl}, \ldots).\\[4pt]
All Agda on the slides is typeset by Agda's own \LaTeX\ backend from
literate files that are scope-checked: against the library on this
branch, and against stub declarations for snippets quoted from other
branches.
\end{keybox}
\end{column}
\end{columns}
\end{frame}

\begin{frame}{The talk in one slide}
\begin{columns}[c]
\begin{column}{0.52\textwidth}
\begin{itemize}
  \setlength{\itemsep}{5pt}
  \item A \alert{circuit} is a word of gates; a \alert{rule set} is a datatype of axioms;
        the equational theory is the \alert{inductive congruence} they generate.
  \item A \alert{normal form} is a function on words with a \alert{section}
        picking canonical representatives.
  \item \textbf{Completeness = soundness + unique normal form}, proved
        \emph{once}, generically.
  \item Normal forms are built \alert{one wire at a time} by verified
        \alert{coset tables}: one ``digit'' per wire, a mixed-radix numeral.
  \item Everything is \texttt{--safe}, axiom-free, setoid-based Agda on the
        standard library.
\end{itemize}
\end{column}
\begin{column}{0.45\textwidth}
\centering
\begin{tikzpicture}[node distance=7mm and 5mm,
  box/.style={draw=accent,fill=soft,rounded corners=2pt,align=center,
    font=\small,minimum height=8mm,inner sep=3pt},
  arr/.style={-{Stealth[length=2mm]},thick,accent}]
  \node[box] (sound) {soundness\\{\scriptsize $w\approx v \Rightarrow \llbracket w\rrbracket = \llbracket v\rrbracket$}};
  \node[box,right=of sound] (unf) {unique NF\\{\scriptsize $\llbracket \bar u\rrbracket = \llbracket \bar v\rrbracket \Rightarrow u = v$}};
  \node[box,fill=accent!15,below=12mm of $(sound)!0.5!(unf)$] (compl)
     {\textbf{completeness}\\{\scriptsize $\llbracket w\rrbracket = \llbracket v\rrbracket \Rightarrow w\approx v$}};
  \draw[arr] (sound) -- (compl);
  \draw[arr] (unf) -- (compl);
  \node[font=\scriptsize\agdafont,anchor=west,text=accent] at ($(compl.east)+(1mm,0)$) {by-normalization};
  \node[box,below=9mm of compl,fill=white] (tab)
     {verified coset tables\\{\scriptsize finite data + 5 equations per level}};
  \draw[arr] (tab) -- node[right,font=\scriptsize]{build the NF} (compl);
\end{tikzpicture}
\end{column}
\end{columns}
\end{frame}

%% Part 1 — Why completeness proofs need a machine (paper §1).

\PaperPart
\section{Why completeness proofs need a machine}

\begin{frame}{Circuits are programs; rules are their equational theory \InPaper}
\begin{columns}[c]
\begin{column}{0.55\textwidth}
A circuit optimizer, a compiler pass, a diagram normalizer: the
justification is always the same --- \alert{local rewrite rules} on a
syntactic object, preserving a \alert{semantics}.
\vspace{6pt}
\begin{center}
\begin{quantikz}[row sep=0.2cm,column sep=0.18cm,font=\scriptsize]
  & \gate{H} & \gate{H} & \gate{S} & \ctrl{1} & \\
  & & & & \control{} &
\end{quantikz}
$\;\overset{\scriptscriptstyle H^2=\varepsilon}{\approx}\;$
\begin{quantikz}[row sep=0.2cm,column sep=0.18cm,font=\scriptsize]
  & \gate{S} & \ctrl{1} & \\
  & & \control{} &
\end{quantikz}
$\;\overset{\scriptscriptstyle\text{diagonal}}{\approx}\;$
\begin{quantikz}[row sep=0.2cm,column sep=0.18cm,font=\scriptsize]
  & \ctrl{1} & \gate{S} & \\
  & \control{} & &
\end{quantikz}
\end{center}
\vspace{4pt}
Two questions attach to every rule set:
\begin{itemize}
  \item \textbf{Soundness} --- does every rule preserve the semantics?\\
        {\small\color{gray}finitely many small algebraic identities: easy.}
  \item \textbf{Completeness} --- is every true equation derivable?\\
        {\small\color{gray}a different animal altogether.}
\end{itemize}
\end{column}
\begin{column}{0.42\textwidth}
\centering
\begin{tikzpicture}[
  >={Stealth[length=2mm]}]
  \node[blob=accent] (S) at (0,0) {};
  \node[blob=paperC] (M) at (3.3,0) {};
  \node[font=\small\bfseries,accent] at (0,2.0) {syntax};
  \node[font=\small\bfseries,paperC] at (3.3,2.0) {semantics};
  \node[font=\scriptsize,accent] at (0,1.15) {circuits / $\approx$};
  \node[font=\scriptsize,paperC] at (3.3,1.15) {unitaries, \dots};
  \node[circle,fill=accent,inner sep=1.3pt,label=left:{\small $w$}] (w) at (-0.2,0.35) {};
  \node[circle,fill=accent,inner sep=1.3pt,label=left:{\small $v$}] (v) at (-0.2,-0.55) {};
  \node[circle,fill=paperC,inner sep=1.3pt,label=right:{\small $U$}] (u) at (3.2,-0.1) {};
  \draw[->,thick,gray] (w) to[bend left=12] node[above,font=\scriptsize]{$\llbracket\_\rrbracket$} (u);
  \draw[->,thick,gray] (v) to[bend right=12] (u);
  \draw[<->,thick,newC,dashed] (w) to[bend right=40] node[left,font=\scriptsize,align=right]{derivable\\from rules?} (v);
\end{tikzpicture}

\small soundness: $w\approx v\;\Rightarrow\;\llbracket w\rrbracket=\llbracket v\rrbracket$\\
completeness: $\llbracket w\rrbracket=\llbracket v\rrbracket\;\Rightarrow\;w\approx v$
\end{column}
\end{columns}
\end{frame}

\begin{frame}{How completeness is proved on paper \InPaper}
\begin{columns}[c]
\begin{column}{0.47\textwidth}
The standard strategy has three steps:
\begin{enumerate}
  \setlength{\itemsep}{4pt}
  \item define a \alert{normal form} for circuits;
  \item show every circuit \alert{rewrites to normal form} using the rules;
  \item show \alert{distinct normal forms denote distinct} semantic objects.
\end{enumerate}
\vspace{6pt}
Step 2 is where it explodes: for \emph{every} rule and \emph{every} way a
generator can meet a normal form, the combination must reduce again ---
mechanical, unforgiving, enormous.
\end{column}
\begin{column}{0.50\textwidth}
\centering
\begin{tikzpicture}[x=3.3mm,y=3.3mm]
  % a grid of "cases": generators x normal-form shapes
  \foreach \i in {0,...,13} {
    \foreach \j in {0,...,8} {
      \pgfmathsetmacro{\shade}{mod(\i*7+\j*13,11)}
      \ifdim\shade pt<1pt
        \fill[newC!55] (\i,\j) rectangle ++(0.85,0.85);
      \else
        \fill[paperC!35] (\i,\j) rectangle ++(0.85,0.85);
      \fi
    }
  }
  \node[rotate=90,font=\scriptsize,anchor=south] at (-0.3,4.4) {generator $g$ (incl.\ every wire)};
  \node[font=\scriptsize,anchor=north] at (7,-0.2) {shape of the normal form it meets};
  \node[font=\scriptsize,anchor=south,align=center] at (7,9.1)
     {one case per (rule, coset) or (generator, coset)};
  \fill[paperC!35] (15,6) rectangle ++(0.85,0.85); \node[font=\scriptsize,anchor=west] at (16,6.4) {routine};
  \fill[newC!55] (15,4) rectangle ++(0.85,0.85); \node[font=\scriptsize,anchor=west,align=left] at (16,4.4) {a real\\derivation};
\end{tikzpicture}
\end{column}
\end{columns}
\end{frame}

\begin{frame}{The literature wears the scars \InPaper}
\begin{columns}[c]
\begin{column}{0.50\textwidth}
\small
\begin{itemize}
  \setlength{\itemsep}{3pt}
  \item Selinger (2015): $n$-qubit Clifford --- a substantial rewrite
        system and normal-form analysis.
  \item Li, Mosca, Ross, van de Wetering, Zhao (2025): qutrit Clifford, first
        complete fragment in odd prime dimension --- \alert{$>50$ pages} of proofs.
  \item Clément, Heurtel, Mansfield, Perdrix, Valiron (2023, 2024):
        complete theory for full quantum circuits, and its minimal refinement.
  \item Clément (LICS 2026): Real-Clifford+$CH$ --- a 13-page paper with an
        \alert{appendix of 264 pages} of equational derivations.
  \item Rig-categorical completeness (POPL '22, '24, '26); Blake (2026)
        simplifies the rule sets of six near-Clifford fragments.
\end{itemize}
\end{column}
\begin{column}{0.47\textwidth}
\centering
\begin{tikzpicture}
\begin{axis}[xbar, width=5.0cm, height=4.6cm, xmin=0, xmax=300,
  symbolic y coords={Real-Clifford+CH appendix,Real-Clifford+CH paper,Qutrit Clifford proofs},
  ytick=data, y tick label style={font=\scriptsize,align=right,text width=2.1cm},
  x tick label style={font=\scriptsize}, xlabel={\scriptsize pages},
  nodes near coords, nodes near coords style={font=\scriptsize},
  bar width=9pt, enlarge y limits=0.3, axis lines*=left]
\addplot[fill=newC!70,draw=newC] coordinates
  {(264,Real-Clifford+CH appendix) (13,Real-Clifford+CH paper) (50,Qutrit Clifford proofs)};
\end{axis}
\end{tikzpicture}

{\small Long, regular, mechanical: exactly the proofs a proof assistant
should check --- it does not tire of case splits, and where a case is a
computation, its \alert{evaluator} can carry it out.}
\end{column}
\end{columns}
\end{frame}

\begin{frame}{Why this is a programming-languages problem \InPaper}
\begin{columns}[T]
\begin{column}{0.50\textwidth}
\begin{itemize}
  \setlength{\itemsep}{5pt}
  \item Circuit calculi \emph{are} programming languages; completeness of
        their equational theory is a PL question.
  \item It licenses \alert{normal-form decision procedures}, and tells an
        optimizer author that \alert{no valid rewrite is missing}.
  \item It underlies the rig-categorical completeness theorems for
        reversible and quantum languages (Choudhury--Karwowski--Sabry POPL'22,
        Carette--Heunen--Kaarsgaard--Sabry POPL'24, Fang--Heunen--Kaarsgaard
        POPL'26).
  \item The contribution is a \alert{design}: objects that both
        \emph{compute} and \emph{prove}.
\end{itemize}
\end{column}
\begin{column}{0.46\textwidth}
\begin{keybox}[title={Announced up front: setoids, not quotients}]
\small
Plain, setoid-based Agda on the standard library --- not cubical, not HoTT.
A presented monoid is a \emph{setoid on words}, not a quotient type.
\begin{itemize}
  \item the library \alert{computes with syntax}: coset tables and normal
        forms pattern-match on concrete representatives, obligations close
        by evaluation;
  \item central devices are maps \alert{into} the syntax (sections, coset
        representatives, amalgamation tails);
  \item derivations are \alert{data} one can induct on.
\end{itemize}
\end{keybox}
\end{column}
\end{columns}
\end{frame}

\begin{frame}{Contributions of the paper \InPaper}
\small
\begin{columns}[T]
\begin{column}{0.50\textwidth}
\begin{enumerate}
  \setcounter{enumi}{-1}
  \setlength{\itemsep}{4pt}
  \item \textbf{A library} for circuit normal forms and complete relation
        sets: presentations as inductive congruences; normal forms as stdlib
        function bundles; generic \aid{by-normalization}; a verified
        Reidemeister--Schreier engine; solvers by evaluation.
  \item \textbf{Inductively defined circuits}: wire-indexed generators,
        structural rules supplied once for every gate set.
  \item \textbf{A catalogue} of verified presentations: $\mathbb{Z}_n$,
        $S_n$, $\mathbb{Z}_N\wr S_n$, Pauli $(\mathbb{Z}_p{\times}\mathbb{Z}_p)^n$,
        qubit/qutrit Clifford+$T$, $U_3(\mathbb{Z}[\frac12,i])$.
\end{enumerate}
\end{column}
\begin{column}{0.48\textwidth}
\begin{enumerate}
  \setcounter{enumi}{2}
  \setlength{\itemsep}{4pt}
  \item \textbf{Monoid-level} amalgamation and Reidemeister--Schreier
        (after Bian--Selinger).
  \item \textbf{Presentation theorems} for direct, semidirect, $n$-fold and
        amalgamated products, and group extensions.
  \item \textbf{New stdlib theorems}: semidirect products, amalgamated free
        products of monoids, extensions, epimorphisms, $\mathsf{Perm}(n)$ as a group.
\end{enumerate}
\vspace{4pt}
\begin{paperbox}
Every headline theorem is restated with its full type in
\texttt{MainTheorems.agda}; \texttt{--safe}, \texttt{--cubical-compatible},
no postulates, no holes, no classical axioms.
\end{paperbox}
\end{column}
\end{columns}
\end{frame}

\begin{code}[hide]%
\>[0]\AgdaKeyword{module}\AgdaSpace{}%
\AgdaModule{slides.Tsinghua-PL-seminar.sections.permutations}\AgdaSpace{}%
\AgdaKeyword{where}\<%
\\
%
\\[\AgdaEmptyExtraSkip]%
\>[0]\AgdaKeyword{open}\AgdaSpace{}%
\AgdaKeyword{import}\AgdaSpace{}%
\AgdaModule{Data.Nat}\AgdaSpace{}%
\AgdaKeyword{using}\AgdaSpace{}%
\AgdaSymbol{(}\AgdaDatatype{ℕ}\AgdaSpace{}%
\AgdaSymbol{;}\AgdaSpace{}%
\AgdaInductiveConstructor{suc}\AgdaSymbol{)}\<%
\\
\>[0]\AgdaKeyword{open}\AgdaSpace{}%
\AgdaKeyword{import}\AgdaSpace{}%
\AgdaModule{Data.Product}\AgdaSpace{}%
\AgdaKeyword{using}\AgdaSpace{}%
\AgdaSymbol{(}\AgdaOperator{\AgdaFunction{\AgdaUnderscore{}×\AgdaUnderscore{}}}\AgdaSpace{}%
\AgdaSymbol{;}\AgdaSpace{}%
\AgdaOperator{\AgdaInductiveConstructor{\AgdaUnderscore{},\AgdaUnderscore{}}}\AgdaSpace{}%
\AgdaSymbol{;}\AgdaSpace{}%
\AgdaField{proj₁}\AgdaSpace{}%
\AgdaSymbol{;}\AgdaSpace{}%
\AgdaField{proj₂}\AgdaSymbol{)}\<%
\\
\>[0]\AgdaKeyword{open}\AgdaSpace{}%
\AgdaKeyword{import}\AgdaSpace{}%
\AgdaModule{Data.Unit}\AgdaSpace{}%
\AgdaKeyword{using}\AgdaSpace{}%
\AgdaSymbol{(}\AgdaRecord{⊤}\AgdaSymbol{)}\<%
\\
\>[0]\AgdaKeyword{open}\AgdaSpace{}%
\AgdaKeyword{import}\AgdaSpace{}%
\AgdaModule{Level}\AgdaSpace{}%
\AgdaKeyword{using}\AgdaSpace{}%
\AgdaSymbol{(}\AgdaFunction{0ℓ}\AgdaSymbol{)}\<%
\\
\>[0]\AgdaKeyword{open}\AgdaSpace{}%
\AgdaKeyword{import}\AgdaSpace{}%
\AgdaModule{Relation.Binary}\AgdaSpace{}%
\AgdaKeyword{using}\AgdaSpace{}%
\AgdaSymbol{(}\AgdaFunction{Rel}\AgdaSymbol{)}\<%
\\
%
\\[\AgdaEmptyExtraSkip]%
\>[0]\AgdaKeyword{open}\AgdaSpace{}%
\AgdaKeyword{import}\AgdaSpace{}%
\AgdaModule{Notations}\AgdaSpace{}%
\AgdaKeyword{using}\AgdaSpace{}%
\AgdaSymbol{(}\AgdaInductiveConstructor{₁₊}\AgdaSpace{}%
\AgdaSymbol{;}\AgdaSpace{}%
\AgdaInductiveConstructor{₂₊}\AgdaSpace{}%
\AgdaSymbol{;}\AgdaSpace{}%
\AgdaInductiveConstructor{₃₊}\AgdaSymbol{)}\<%
\\
\>[0]\AgdaKeyword{open}\AgdaSpace{}%
\AgdaKeyword{import}\AgdaSpace{}%
\AgdaModule{Word.Base}\AgdaSpace{}%
\AgdaKeyword{using}\AgdaSpace{}%
\AgdaSymbol{(}\AgdaDatatype{Word}\AgdaSpace{}%
\AgdaSymbol{;}\AgdaSpace{}%
\AgdaFunction{WRel}\AgdaSpace{}%
\AgdaSymbol{;}\AgdaSpace{}%
\AgdaOperator{\AgdaInductiveConstructor{[\AgdaUnderscore{}]ʷ}}\AgdaSpace{}%
\AgdaSymbol{;}\AgdaSpace{}%
\AgdaInductiveConstructor{ε}\AgdaSpace{}%
\AgdaSymbol{;}\AgdaSpace{}%
\AgdaOperator{\AgdaInductiveConstructor{\AgdaUnderscore{}•\AgdaUnderscore{}}}\AgdaSpace{}%
\AgdaSymbol{;}\AgdaSpace{}%
\AgdaFunction{wmap}\AgdaSymbol{)}\<%
\\
\>[0]\AgdaKeyword{open}\AgdaSpace{}%
\AgdaKeyword{import}\AgdaSpace{}%
\AgdaModule{Presentation.Definitions}\AgdaSpace{}%
\AgdaKeyword{using}\AgdaSpace{}%
\AgdaSymbol{(}\AgdaOperator{\AgdaRecord{\AgdaUnderscore{}IsPresentationOf\AgdaUnderscore{}}}\AgdaSymbol{)}\<%
\\
\>[0]\AgdaKeyword{open}\AgdaSpace{}%
\AgdaKeyword{import}\AgdaSpace{}%
\AgdaModule{Examples.Groups.Symmetric.Tight.Semantics}\AgdaSpace{}%
\AgdaKeyword{using}\AgdaSpace{}%
\AgdaSymbol{(}\AgdaFunction{Permutation′-group}\AgdaSymbol{)}\<%
\\
\>[0]\AgdaKeyword{import}\AgdaSpace{}%
\AgdaModule{Examples.Groups.Symmetric.Theorems}\AgdaSpace{}%
\AgdaSymbol{as}\AgdaSpace{}%
\AgdaModule{SymThm}\<%
\\
%
\\[\AgdaEmptyExtraSkip]%
\>[0]\AgdaKeyword{postulate}\<%
\\
\>[0][@{}l@{\AgdaIndent{0}}]%
\>[2]\AgdaPostulate{proof-elided}\AgdaSpace{}%
\AgdaSymbol{:}\AgdaSpace{}%
\AgdaSymbol{∀}\AgdaSpace{}%
\AgdaSymbol{\{}\AgdaBound{ℓ}\AgdaSymbol{\}}\AgdaSpace{}%
\AgdaSymbol{\{}\AgdaBound{A}\AgdaSpace{}%
\AgdaSymbol{:}\AgdaSpace{}%
\AgdaPrimitive{Set}\AgdaSpace{}%
\AgdaBound{ℓ}\AgdaSymbol{\}}\AgdaSpace{}%
\AgdaSymbol{→}\AgdaSpace{}%
\AgdaBound{A}\<%
\\
%
\\[\AgdaEmptyExtraSkip]%
\>[0]\AgdaKeyword{private}\AgdaSpace{}%
\AgdaKeyword{variable}\<%
\\
\>[0][@{}l@{\AgdaIndent{0}}]%
\>[2]\AgdaGeneralizable{n}\AgdaSpace{}%
\AgdaSymbol{:}\AgdaSpace{}%
\AgdaDatatype{ℕ}\<%
\end{code}
%% Part 2 — the worked example (paper §2).
\PaperPart
\section{A complete worked example: permutation circuits}

\begin{frame}{The pipeline we are about to run \InPaper}
\centering
\begin{tikzpicture}[node distance=3mm,
  st/.style={draw=paperC,fill=softpaper,rounded corners=3pt,font=\scriptsize,
    minimum height=10mm,text width=14mm,align=center,inner sep=2pt},
  >={Stealth[length=2mm]}]
  \node[st] (g) {gates \&\\generators};
  \node[st,right=of g] (r) {relations\\{\scriptsize 2 axioms}};
  \node[st,right=of r] (c) {cosets \&\\coset table};
  \node[st,right=of c] (t) {tower\\{\scriptsize one digit/wire}};
  \node[st,right=of t] (s) {semantics\\{\scriptsize two of them}};
  \node[st,right=of s] (u) {uniqueness};
  \node[st,fill=paperC!25,right=of u] (k) {completeness\\\& presentation};
  \foreach \a/\b in {g/r,r/c,c/t,t/s,s/u,u/k} \draw[->,thick,paperC] (\a) -- (\b);
\end{tikzpicture}
\vspace{10pt}

\begin{columns}[T]
\begin{column}{0.45\textwidth}
Circuits built from \alert{one} two-wire gate, the swap $\sigma$.
They present the \alert{symmetric groups} $S_n$, at every width $n$ at once.
\end{column}
\begin{column}{0.45\textwidth}
Every later case study is ``the same thing, with richer gates, longer
towers, and amalgamations where the chain branches''.
\end{column}
\end{columns}
\end{frame}

\begin{frame}{Gates, generators, circuits \InPaper}
\begin{columns}[c]
\begin{column}{0.58\textwidth}
A gate set is a family \aid{Gate : ℕ → Set} (gates by arity). The circuit
layer turns it into \alert{wire-indexed generators}:
\begin{code}%
\>[0]\AgdaKeyword{data}\AgdaSpace{}%
\AgdaDatatype{Gate}\AgdaSpace{}%
\AgdaSymbol{:}\AgdaSpace{}%
\AgdaDatatype{ℕ}\AgdaSpace{}%
\AgdaSymbol{→}\AgdaSpace{}%
\AgdaPrimitive{Set}\AgdaSpace{}%
\AgdaKeyword{where}\<%
\\
\>[0][@{}l@{\AgdaIndent{0}}]%
\>[2]\AgdaInductiveConstructor{σ-gate}\AgdaSpace{}%
\AgdaSymbol{:}\AgdaSpace{}%
\AgdaDatatype{Gate}\AgdaSpace{}%
\AgdaNumber{2}\<%
\\
\>[0]\AgdaKeyword{data}\AgdaSpace{}%
\AgdaDatatype{Gen}\AgdaSpace{}%
\AgdaSymbol{:}\AgdaSpace{}%
\AgdaDatatype{ℕ}\AgdaSpace{}%
\AgdaSymbol{→}\AgdaSpace{}%
\AgdaPrimitive{Set}\AgdaSpace{}%
\AgdaKeyword{where}\<%
\\
\>[0][@{}l@{\AgdaIndent{0}}]%
\>[2]\AgdaInductiveConstructor{gate₁}\AgdaSpace{}%
\AgdaSymbol{:}\AgdaSpace{}%
\AgdaDatatype{Gate}\AgdaSpace{}%
\AgdaNumber{1}\AgdaSpace{}%
\AgdaSymbol{→}\AgdaSpace{}%
\AgdaDatatype{Gen}\AgdaSpace{}%
\AgdaSymbol{(}\AgdaInductiveConstructor{₁₊}\AgdaSpace{}%
\AgdaGeneralizable{n}\AgdaSymbol{)}%
\>[31]\AgdaComment{--\ on\ the\ bottom\ wire}\<%
\\
%
\>[2]\AgdaInductiveConstructor{gate₂}\AgdaSpace{}%
\AgdaSymbol{:}\AgdaSpace{}%
\AgdaDatatype{Gate}\AgdaSpace{}%
\AgdaNumber{2}\AgdaSpace{}%
\AgdaSymbol{→}\AgdaSpace{}%
\AgdaDatatype{Gen}\AgdaSpace{}%
\AgdaSymbol{(}\AgdaInductiveConstructor{₂₊}\AgdaSpace{}%
\AgdaGeneralizable{n}\AgdaSymbol{)}%
\>[31]\AgdaComment{--\ on\ the\ bottom\ two\ wires}\<%
\\
%
\>[2]\AgdaOperator{\AgdaInductiveConstructor{\AgdaUnderscore{}↥}}%
\>[7]\AgdaSymbol{:}\AgdaSpace{}%
\AgdaDatatype{Gen}\AgdaSpace{}%
\AgdaGeneralizable{n}\AgdaSpace{}%
\AgdaSymbol{→}\AgdaSpace{}%
\AgdaDatatype{Gen}\AgdaSpace{}%
\AgdaSymbol{(}\AgdaInductiveConstructor{₁₊}\AgdaSpace{}%
\AgdaGeneralizable{n}\AgdaSymbol{)}%
\>[31]\AgdaComment{--\ one\ wire\ higher}\<%
\end{code}
\begin{code}[hide]%
\>[0]\AgdaFunction{Circuit}\AgdaSpace{}%
\AgdaSymbol{:}\AgdaSpace{}%
\AgdaDatatype{ℕ}\AgdaSpace{}%
\AgdaSymbol{→}\AgdaSpace{}%
\AgdaPrimitive{Set}\<%
\\
\>[0]\AgdaFunction{Circuit}\AgdaSpace{}%
\AgdaBound{n}\AgdaSpace{}%
\AgdaSymbol{=}\AgdaSpace{}%
\AgdaDatatype{Word}\AgdaSpace{}%
\AgdaSymbol{(}\AgdaDatatype{Gen}\AgdaSpace{}%
\AgdaBound{n}\AgdaSymbol{)}\<%
\\
%
\\[\AgdaEmptyExtraSkip]%
\>[0]\AgdaFunction{CRel}\AgdaSpace{}%
\AgdaSymbol{:}\AgdaSpace{}%
\AgdaDatatype{ℕ}\AgdaSpace{}%
\AgdaSymbol{→}\AgdaSpace{}%
\AgdaPrimitive{Set₁}\<%
\\
\>[0]\AgdaFunction{CRel}\AgdaSpace{}%
\AgdaBound{n}\AgdaSpace{}%
\AgdaSymbol{=}\AgdaSpace{}%
\AgdaFunction{Rel}\AgdaSpace{}%
\AgdaSymbol{(}\AgdaFunction{Circuit}\AgdaSpace{}%
\AgdaBound{n}\AgdaSymbol{)}\AgdaSpace{}%
\AgdaFunction{0ℓ}\<%
\\
%
\\[\AgdaEmptyExtraSkip]%
\>[0]\AgdaOperator{\AgdaFunction{\AgdaUnderscore{}↑}}\AgdaSpace{}%
\AgdaSymbol{:}\AgdaSpace{}%
\AgdaFunction{Circuit}\AgdaSpace{}%
\AgdaGeneralizable{n}\AgdaSpace{}%
\AgdaSymbol{→}\AgdaSpace{}%
\AgdaFunction{Circuit}\AgdaSpace{}%
\AgdaSymbol{(}\AgdaInductiveConstructor{₁₊}\AgdaSpace{}%
\AgdaGeneralizable{n}\AgdaSymbol{)}\<%
\\
\>[0]\AgdaOperator{\AgdaFunction{\AgdaUnderscore{}↑}}\AgdaSpace{}%
\AgdaSymbol{=}\AgdaSpace{}%
\AgdaFunction{wmap}\AgdaSpace{}%
\AgdaOperator{\AgdaInductiveConstructor{\AgdaUnderscore{}↥}}\<%
\\
%
\\[\AgdaEmptyExtraSkip]%
\>[0]\AgdaKeyword{pattern}\AgdaSpace{}%
\AgdaInductiveConstructor{σ-gen}\AgdaSpace{}%
\AgdaSymbol{=}\AgdaSpace{}%
\AgdaInductiveConstructor{gate₂}\AgdaSpace{}%
\AgdaInductiveConstructor{σ-gate}\<%
\\
%
\\[\AgdaEmptyExtraSkip]%
\>[0]\AgdaFunction{σ}\AgdaSpace{}%
\AgdaSymbol{:}\AgdaSpace{}%
\AgdaFunction{Circuit}\AgdaSpace{}%
\AgdaSymbol{(}\AgdaInductiveConstructor{₂₊}\AgdaSpace{}%
\AgdaGeneralizable{n}\AgdaSymbol{)}\<%
\\
\>[0]\AgdaFunction{σ}\AgdaSpace{}%
\AgdaSymbol{=}\AgdaSpace{}%
\AgdaOperator{\AgdaInductiveConstructor{[}}\AgdaSpace{}%
\AgdaInductiveConstructor{σ-gen}\AgdaSpace{}%
\AgdaOperator{\AgdaInductiveConstructor{]ʷ}}\<%
\\
%
\\[\AgdaEmptyExtraSkip]%
\>[0]\AgdaKeyword{infix}\AgdaSpace{}%
\AgdaNumber{4}\AgdaSpace{}%
\AgdaOperator{\AgdaDatatype{\AgdaUnderscore{}SRel,\AgdaUnderscore{}===\AgdaUnderscore{}}}\AgdaSpace{}%
\AgdaOperator{\AgdaDatatype{\AgdaUnderscore{}VRel,\AgdaUnderscore{}===\AgdaUnderscore{}}}\AgdaSpace{}%
\AgdaOperator{\AgdaPostulate{\AgdaUnderscore{}≈\AgdaUnderscore{}}}\<%
\\
\>[0]\AgdaKeyword{infixr}\AgdaSpace{}%
\AgdaNumber{10}\AgdaSpace{}%
\AgdaOperator{\AgdaInductiveConstructor{σ•\AgdaUnderscore{}}}\<%
\\
%
\\[\AgdaEmptyExtraSkip]%
\>[0]\AgdaKeyword{private}\AgdaSpace{}%
\AgdaKeyword{variable}\<%
\\
\>[0][@{}l@{\AgdaIndent{0}}]%
\>[2]\AgdaGeneralizable{w}\AgdaSpace{}%
\AgdaGeneralizable{v}\AgdaSpace{}%
\AgdaSymbol{:}\AgdaSpace{}%
\AgdaFunction{Circuit}\AgdaSpace{}%
\AgdaGeneralizable{n}\<%
\\
%
\\[\AgdaEmptyExtraSkip]%
\>[0]\AgdaKeyword{postulate}\<%
\\
\>[0][@{}l@{\AgdaIndent{0}}]%
\>[2]\AgdaOperator{\AgdaPostulate{\AgdaUnderscore{}≈\AgdaUnderscore{}}}\AgdaSpace{}%
\AgdaSymbol{:}\AgdaSpace{}%
\AgdaFunction{Circuit}\AgdaSpace{}%
\AgdaGeneralizable{n}\AgdaSpace{}%
\AgdaSymbol{→}\AgdaSpace{}%
\AgdaFunction{Circuit}\AgdaSpace{}%
\AgdaGeneralizable{n}\AgdaSpace{}%
\AgdaSymbol{→}\AgdaSpace{}%
\AgdaPrimitive{Set}\<%
\end{code}
\vspace{2pt}
\small An $n$-wire circuit is a word of generators,
$\aid{Circuit}\;n = \aid{Word}\;(\aid{Gen}\;n)$:\;
$[\,g\,]^{\mathsf{w}}$,\; $\varepsilon$ ($n$ bare wires),\; $c \aidop{•} c'$.
\aid{₂₊ n} is the pattern for $2+n$, so \aid{gate₂ σ-gate} lives at every width $\ge 2$.
\end{column}
\begin{column}{0.38\textwidth}
\centering
\small
\begin{tabular}{@{}r@{\quad}c@{}}
$\sigma$ & \permcirc{3}{0} \\[10pt]
$\sigma\;{\uparrow}$ & \permcirc{3}{1} \\[10pt]
$\sigma \aidop{•} \sigma\;{\uparrow} \aidop{•} \sigma$ & \permcirc{3}{0,1,0}
\end{tabular}

\vspace{8pt}
{\scriptsize wires drawn bottom-up (wire 0 at the bottom),\\
composition left to right; $c\;{\uparrow} = \aid{wmap}\;\aid{\_↥}\;c$}
\end{column}
\end{columns}
\end{frame}

\begin{frame}{Relations: two axioms \InPaper}
\begin{columns}[c]
\begin{column}{0.60\textwidth}
Gate-specific relations are an inductive family indexed by the wire count;
each axiom is stated \alert{once}, at the bottom of the wires, at every width
where it fits:
\begin{code}%
\>[0]\AgdaKeyword{data}\AgdaSpace{}%
\AgdaOperator{\AgdaDatatype{\AgdaUnderscore{}SRel,\AgdaUnderscore{}===\AgdaUnderscore{}}}\AgdaSpace{}%
\AgdaSymbol{:}\AgdaSpace{}%
\AgdaSymbol{(}\AgdaBound{n}\AgdaSpace{}%
\AgdaSymbol{:}\AgdaSpace{}%
\AgdaDatatype{ℕ}\AgdaSymbol{)}\AgdaSpace{}%
\AgdaSymbol{→}\AgdaSpace{}%
\AgdaFunction{WRel}\AgdaSpace{}%
\AgdaSymbol{(}\AgdaDatatype{Gen}\AgdaSpace{}%
\AgdaBound{n}\AgdaSymbol{)}\AgdaSpace{}%
\AgdaKeyword{where}\<%
\\
\>[0][@{}l@{\AgdaIndent{0}}]%
\>[2]\AgdaInductiveConstructor{order}%
\>[14]\AgdaSymbol{:}\AgdaSpace{}%
\AgdaSymbol{(}\AgdaInductiveConstructor{₂₊}\AgdaSpace{}%
\AgdaGeneralizable{n}\AgdaSymbol{)}\AgdaSpace{}%
\AgdaOperator{\AgdaDatatype{SRel,}}\AgdaSpace{}%
\AgdaFunction{σ}\AgdaSpace{}%
\AgdaOperator{\AgdaInductiveConstructor{•}}\AgdaSpace{}%
\AgdaFunction{σ}\AgdaSpace{}%
\AgdaOperator{\AgdaDatatype{===}}\AgdaSpace{}%
\AgdaInductiveConstructor{ε}\<%
\\
%
\>[2]\AgdaInductiveConstructor{yang-baxter}\AgdaSpace{}%
\AgdaSymbol{:}\AgdaSpace{}%
\AgdaSymbol{(}\AgdaInductiveConstructor{₃₊}\AgdaSpace{}%
\AgdaGeneralizable{n}\AgdaSymbol{)}\AgdaSpace{}%
\AgdaOperator{\AgdaDatatype{SRel,}}\AgdaSpace{}%
\AgdaFunction{σ}\AgdaSpace{}%
\AgdaOperator{\AgdaInductiveConstructor{•}}\AgdaSpace{}%
\AgdaFunction{σ}\AgdaSpace{}%
\AgdaOperator{\AgdaFunction{↑}}\AgdaSpace{}%
\AgdaOperator{\AgdaInductiveConstructor{•}}\AgdaSpace{}%
\AgdaFunction{σ}\AgdaSpace{}%
\AgdaOperator{\AgdaDatatype{===}}\AgdaSpace{}%
\AgdaFunction{σ}\AgdaSpace{}%
\AgdaOperator{\AgdaFunction{↑}}\AgdaSpace{}%
\AgdaOperator{\AgdaInductiveConstructor{•}}\AgdaSpace{}%
\AgdaFunction{σ}\AgdaSpace{}%
\AgdaOperator{\AgdaInductiveConstructor{•}}\AgdaSpace{}%
\AgdaFunction{σ}\AgdaSpace{}%
\AgdaOperator{\AgdaFunction{↑}}\<%
\end{code}
\vspace{4pt}
Two thirds of the \alert{Coxeter presentation} of $S_n$:
$\sigma^2=\varepsilon$ and the braid (Yang--Baxter) relation.

\vspace{4pt}
\small The third Coxeter relation --- swaps on \emph{disjoint} wires commute ---
is not about permutations: every circuit theory needs it.
\end{column}
\begin{column}{0.36\textwidth}
\centering
\aid{order}\\[2pt]
\ufig[0.42]{sn-order}\\[12pt]
\aid{yang-baxter}\\[2pt]
\ufig[0.42]{sn-yb}
\end{column}
\end{columns}
\end{frame}

\begin{frame}{Structure for free: \aid{Lift-Relation} \InPaper}
\begin{columns}[c]
\begin{column}{0.60\textwidth}
Circuits are fixed-width slices of a \alert{PROP}. The monoidal structure leaves
two traces at fixed width: equations survive the \alert{shift} (tensoring an
identity wire), and gates on \alert{disjoint wires commute}.
\begin{code}%
\>[0]\AgdaKeyword{data}\AgdaSpace{}%
\AgdaOperator{\AgdaDatatype{\AgdaUnderscore{}VRel,\AgdaUnderscore{}===\AgdaUnderscore{}}}\AgdaSpace{}%
\AgdaSymbol{:}\AgdaSpace{}%
\AgdaSymbol{(}\AgdaBound{n}\AgdaSpace{}%
\AgdaSymbol{:}\AgdaSpace{}%
\AgdaDatatype{ℕ}\AgdaSymbol{)}\AgdaSpace{}%
\AgdaSymbol{→}\AgdaSpace{}%
\AgdaFunction{CRel}\AgdaSpace{}%
\AgdaBound{n}\AgdaSpace{}%
\AgdaKeyword{where}\<%
\\
\>[0][@{}l@{\AgdaIndent{0}}]%
\>[2]\AgdaInductiveConstructor{srel}%
\>[8]\AgdaSymbol{:}\AgdaSpace{}%
\AgdaGeneralizable{n}\AgdaSpace{}%
\AgdaOperator{\AgdaDatatype{SRel,}}\AgdaSpace{}%
\AgdaGeneralizable{w}\AgdaSpace{}%
\AgdaOperator{\AgdaDatatype{===}}\AgdaSpace{}%
\AgdaGeneralizable{v}\AgdaSpace{}%
\AgdaSymbol{→}\AgdaSpace{}%
\AgdaGeneralizable{n}\AgdaSpace{}%
\AgdaOperator{\AgdaDatatype{VRel,}}\AgdaSpace{}%
\AgdaGeneralizable{w}\AgdaSpace{}%
\AgdaOperator{\AgdaDatatype{===}}\AgdaSpace{}%
\AgdaGeneralizable{v}\<%
\\
%
\>[2]\AgdaInductiveConstructor{cong↑}\AgdaSpace{}%
\AgdaSymbol{:}\AgdaSpace{}%
\AgdaGeneralizable{n}\AgdaSpace{}%
\AgdaOperator{\AgdaDatatype{VRel,}}\AgdaSpace{}%
\AgdaGeneralizable{w}\AgdaSpace{}%
\AgdaOperator{\AgdaDatatype{===}}\AgdaSpace{}%
\AgdaGeneralizable{v}\AgdaSpace{}%
\AgdaSymbol{→}\AgdaSpace{}%
\AgdaSymbol{(}\AgdaInductiveConstructor{₁₊}\AgdaSpace{}%
\AgdaGeneralizable{n}\AgdaSymbol{)}\AgdaSpace{}%
\AgdaOperator{\AgdaDatatype{VRel,}}\AgdaSpace{}%
\AgdaGeneralizable{w}\AgdaSpace{}%
\AgdaOperator{\AgdaFunction{↑}}\AgdaSpace{}%
\AgdaOperator{\AgdaDatatype{===}}\AgdaSpace{}%
\AgdaGeneralizable{v}\AgdaSpace{}%
\AgdaOperator{\AgdaFunction{↑}}\<%
\\
%
\>[2]\AgdaInductiveConstructor{comm₁}\AgdaSpace{}%
\AgdaSymbol{:}%
\>[242I]\AgdaSymbol{(}\AgdaBound{h}\AgdaSpace{}%
\AgdaSymbol{:}\AgdaSpace{}%
\AgdaDatatype{Gate}\AgdaSpace{}%
\AgdaNumber{1}\AgdaSymbol{)}\AgdaSpace{}%
\AgdaSymbol{(}\AgdaBound{g}\AgdaSpace{}%
\AgdaSymbol{:}\AgdaSpace{}%
\AgdaDatatype{Gen}\AgdaSpace{}%
\AgdaGeneralizable{n}\AgdaSymbol{)}\AgdaSpace{}%
\AgdaSymbol{→}\<%
\\
\>[.][@{}l@{}]\<[242I]%
\>[10]\AgdaSymbol{(}\AgdaInductiveConstructor{₁₊}\AgdaSpace{}%
\AgdaGeneralizable{n}\AgdaSymbol{)}\AgdaSpace{}%
\AgdaOperator{\AgdaDatatype{VRel,}}\AgdaSpace{}%
\AgdaOperator{\AgdaInductiveConstructor{[}}\AgdaSpace{}%
\AgdaBound{g}\AgdaSpace{}%
\AgdaOperator{\AgdaInductiveConstructor{↥}}\AgdaSpace{}%
\AgdaOperator{\AgdaInductiveConstructor{]ʷ}}\AgdaSpace{}%
\AgdaOperator{\AgdaInductiveConstructor{•}}\AgdaSpace{}%
\AgdaOperator{\AgdaInductiveConstructor{[}}\AgdaSpace{}%
\AgdaInductiveConstructor{gate₁}\AgdaSpace{}%
\AgdaBound{h}\AgdaSpace{}%
\AgdaOperator{\AgdaInductiveConstructor{]ʷ}}\AgdaSpace{}%
\AgdaOperator{\AgdaDatatype{===}}\AgdaSpace{}%
\AgdaOperator{\AgdaInductiveConstructor{[}}\AgdaSpace{}%
\AgdaInductiveConstructor{gate₁}\AgdaSpace{}%
\AgdaBound{h}\AgdaSpace{}%
\AgdaOperator{\AgdaInductiveConstructor{]ʷ}}\AgdaSpace{}%
\AgdaOperator{\AgdaInductiveConstructor{•}}\AgdaSpace{}%
\AgdaOperator{\AgdaInductiveConstructor{[}}\AgdaSpace{}%
\AgdaBound{g}\AgdaSpace{}%
\AgdaOperator{\AgdaInductiveConstructor{↥}}\AgdaSpace{}%
\AgdaOperator{\AgdaInductiveConstructor{]ʷ}}\<%
\\
%
\>[2]\AgdaInductiveConstructor{comm₂}\AgdaSpace{}%
\AgdaSymbol{:}%
\>[273I]\AgdaSymbol{(}\AgdaBound{h}\AgdaSpace{}%
\AgdaSymbol{:}\AgdaSpace{}%
\AgdaDatatype{Gate}\AgdaSpace{}%
\AgdaNumber{2}\AgdaSymbol{)}\AgdaSpace{}%
\AgdaSymbol{(}\AgdaBound{g}\AgdaSpace{}%
\AgdaSymbol{:}\AgdaSpace{}%
\AgdaDatatype{Gen}\AgdaSpace{}%
\AgdaGeneralizable{n}\AgdaSymbol{)}\AgdaSpace{}%
\AgdaSymbol{→}\<%
\\
\>[.][@{}l@{}]\<[273I]%
\>[10]\AgdaSymbol{(}\AgdaInductiveConstructor{₂₊}\AgdaSpace{}%
\AgdaGeneralizable{n}\AgdaSymbol{)}\AgdaSpace{}%
\AgdaOperator{\AgdaDatatype{VRel,}}\AgdaSpace{}%
\AgdaOperator{\AgdaInductiveConstructor{[}}\AgdaSpace{}%
\AgdaBound{g}\AgdaSpace{}%
\AgdaOperator{\AgdaInductiveConstructor{↥}}\AgdaSpace{}%
\AgdaOperator{\AgdaInductiveConstructor{↥}}\AgdaSpace{}%
\AgdaOperator{\AgdaInductiveConstructor{]ʷ}}\AgdaSpace{}%
\AgdaOperator{\AgdaInductiveConstructor{•}}\AgdaSpace{}%
\AgdaOperator{\AgdaInductiveConstructor{[}}\AgdaSpace{}%
\AgdaInductiveConstructor{gate₂}\AgdaSpace{}%
\AgdaBound{h}\AgdaSpace{}%
\AgdaOperator{\AgdaInductiveConstructor{]ʷ}}\AgdaSpace{}%
\AgdaOperator{\AgdaDatatype{===}}\AgdaSpace{}%
\AgdaOperator{\AgdaInductiveConstructor{[}}\AgdaSpace{}%
\AgdaInductiveConstructor{gate₂}\AgdaSpace{}%
\AgdaBound{h}\AgdaSpace{}%
\AgdaOperator{\AgdaInductiveConstructor{]ʷ}}\AgdaSpace{}%
\AgdaOperator{\AgdaInductiveConstructor{•}}\AgdaSpace{}%
\AgdaOperator{\AgdaInductiveConstructor{[}}\AgdaSpace{}%
\AgdaBound{g}\AgdaSpace{}%
\AgdaOperator{\AgdaInductiveConstructor{↥}}\AgdaSpace{}%
\AgdaOperator{\AgdaInductiveConstructor{↥}}\AgdaSpace{}%
\AgdaOperator{\AgdaInductiveConstructor{]ʷ}}\<%
\end{code}
\small A new circuit theory = a \aid{Gate} family + its gate-specific
relations. Structural rules and their lemmas come for free.
\end{column}
\begin{column}{0.37\textwidth}
\centering
\begin{tikzpicture}[x=4.2mm,y=3.7mm,line width=0.6pt,
  g/.style={draw,fill=white,minimum width=4mm,minimum height=4mm,inner sep=0pt,font=\scriptsize}]
  % cong↑
  \node[font=\scriptsize\agdafont,anchor=west] at (-0.3,10.2) {cong↑};
  \foreach \y in {8,9} \draw (0,\y) -- (2.6,\y);
  \node[g,minimum height=9mm] at (1.3,8.5) {$w$};
  \foreach \y in {7} \draw[gray] (0,\y) -- (2.6,\y);
  \node at (3.4,8) {$\approx$};
  \foreach \y in {8,9} \draw (4.2,\y) -- (6.8,\y);
  \node[g,minimum height=9mm] at (5.5,8.5) {$v$};
  \draw[gray] (4.2,7) -- (6.8,7);
  % comm1
  \node[font=\scriptsize\agdafont,anchor=west] at (-0.3,5.6) {comm₁};
  \foreach \y in {2.8,3.8,4.8} \draw (0,\y) -- (2.6,\y);
  \node[g,minimum height=9mm] at (0.8,4.3) {$g$};
  \node[g] at (1.8,2.8) {$h$};
  \node at (3.4,3.8) {$\approx$};
  \foreach \y in {2.8,3.8,4.8} \draw (4.2,\y) -- (6.8,\y);
  \node[g] at (5.0,2.8) {$h$};
  \node[g,minimum height=9mm] at (6.0,4.3) {$g$};
  % comm2
  \node[font=\scriptsize\agdafont,anchor=west] at (-0.3,1.7) {comm₂};
  \foreach \y in {-2,-1,0,1} \draw (0,\y) -- (2.6,\y);
  \node[g,minimum height=9mm] at (0.8,0.5) {$g$};
  \node[g,minimum height=9mm] at (1.8,-1.5) {$h$};
  \node at (3.4,-0.5) {$\approx$};
  \foreach \y in {-2,-1,0,1} \draw (4.2,\y) -- (6.8,\y);
  \node[g,minimum height=9mm] at (5.0,-1.5) {$h$};
  \node[g,minimum height=9mm] at (6.0,0.5) {$g$};
\end{tikzpicture}

{\scriptsize $\approx$ = the monoid congruence generated by
\aid{n VRel,\_===\_}; deciding it is the word problem.}
\end{column}
\end{columns}
\end{frame}

\begin{frame}{Cosets: what normalization will compute \InPaper}
\begin{columns}[c]
\begin{column}{0.52\textwidth}
Chain $S_1 < S_2 < \dots < S_n < S_{n+1}$; $S_n$ sits in $S_{n+1}$ as the
circuits \alert{not touching the bottom wire}: the circuits $c\;{\uparrow}$.

\vspace{4pt}
Two circuits are in the same right coset iff they bring the \alert{same wire
to the bottom}: $n{+}1$ cosets. Representative: the \alert{descending
staircase} carrying that wire down.

\vspace{4pt}
A coset is a \emph{label}, not a set of circuits:
\begin{code}%
\>[0]\AgdaKeyword{data}\AgdaSpace{}%
\AgdaDatatype{C}\AgdaSpace{}%
\AgdaSymbol{:}\AgdaSpace{}%
\AgdaDatatype{ℕ}\AgdaSpace{}%
\AgdaSymbol{→}\AgdaSpace{}%
\AgdaPrimitive{Set}\AgdaSpace{}%
\AgdaKeyword{where}\<%
\\
\>[0][@{}l@{\AgdaIndent{0}}]%
\>[2]\AgdaInductiveConstructor{ε}%
\>[6]\AgdaSymbol{:}\AgdaSpace{}%
\AgdaDatatype{C}\AgdaSpace{}%
\AgdaGeneralizable{n}\<%
\\
%
\>[2]\AgdaOperator{\AgdaInductiveConstructor{σ•\AgdaUnderscore{}}}\AgdaSpace{}%
\AgdaSymbol{:}\AgdaSpace{}%
\AgdaDatatype{C}\AgdaSpace{}%
\AgdaGeneralizable{n}\AgdaSpace{}%
\AgdaSymbol{→}\AgdaSpace{}%
\AgdaDatatype{C}\AgdaSpace{}%
\AgdaSymbol{(}\AgdaInductiveConstructor{₁₊}\AgdaSpace{}%
\AgdaGeneralizable{n}\AgdaSymbol{)}\<%
\\
\>[0]\AgdaOperator{\AgdaFunction{[\AgdaUnderscore{}]ᶜ}}\AgdaSpace{}%
\AgdaSymbol{:}\AgdaSpace{}%
\AgdaDatatype{C}\AgdaSpace{}%
\AgdaGeneralizable{n}\AgdaSpace{}%
\AgdaSymbol{→}\AgdaSpace{}%
\AgdaFunction{Circuit}\AgdaSpace{}%
\AgdaSymbol{(}\AgdaInductiveConstructor{₁₊}\AgdaSpace{}%
\AgdaGeneralizable{n}\AgdaSymbol{)}\<%
\\
\>[0]\AgdaOperator{\AgdaFunction{[}}\AgdaSpace{}%
\AgdaInductiveConstructor{ε}\AgdaSpace{}%
\AgdaOperator{\AgdaFunction{]ᶜ}}%
\>[11]\AgdaSymbol{=}\AgdaSpace{}%
\AgdaInductiveConstructor{ε}\<%
\\
\>[0]\AgdaOperator{\AgdaFunction{[}}\AgdaSpace{}%
\AgdaOperator{\AgdaInductiveConstructor{σ•}}\AgdaSpace{}%
\AgdaBound{c}\AgdaSpace{}%
\AgdaOperator{\AgdaFunction{]ᶜ}}%
\>[11]\AgdaSymbol{=}\AgdaSpace{}%
\AgdaFunction{σ}\AgdaSpace{}%
\AgdaOperator{\AgdaInductiveConstructor{•}}\AgdaSpace{}%
\AgdaSymbol{(}\AgdaOperator{\AgdaFunction{[}}\AgdaSpace{}%
\AgdaBound{c}\AgdaSpace{}%
\AgdaOperator{\AgdaFunction{]ᶜ}}\AgdaSpace{}%
\AgdaOperator{\AgdaFunction{↑}}\AgdaSymbol{)}\<%
\end{code}
\end{column}
\begin{column}{0.45\textwidth}
\centering
\begin{tabular}{@{}l@{\;\;}ccc@{}}
  $\aid{C}\;2$: & \ufig[0.40]{c2-0} & \ufig[0.40]{c2-1} & \ufig[0.40]{c2-2}\\
  & $\varepsilon$ & $\sigma{\bullet}\,\varepsilon$ &
    $\sigma{\bullet}\,\sigma{\bullet}\,\varepsilon$\\[8pt]
  $\aid{C}\;1$: & \ufig[0.40]{c1-0} & \ufig[0.40]{c1-1} & \\
  & $\varepsilon$ & $\sigma{\bullet}\,\varepsilon$ & \\[8pt]
  $\aid{C}\;0$: & \ufig[0.40]{c0-0} & & \\
  & $\varepsilon$ & &
\end{tabular}

\vspace{6pt}
\small $|\aid{C}\;n| = n+1 = [S_{n+1} : S_n]$
\end{column}
\end{columns}
\end{frame}

\begin{frame}{The coset table \aid{ract} \InPaper}
\begin{columns}[c]
\begin{column}{0.60\textwidth}
Push one generator past a staircase: where do we land, and what is left
over on the wires above?
\begin{code}%
\>[0]\AgdaFunction{ract}\AgdaSpace{}%
\AgdaSymbol{:}\AgdaSpace{}%
\AgdaDatatype{C}\AgdaSpace{}%
\AgdaSymbol{(}\AgdaInductiveConstructor{₁₊}\AgdaSpace{}%
\AgdaGeneralizable{n}\AgdaSymbol{)}\AgdaSpace{}%
\AgdaSymbol{→}\AgdaSpace{}%
\AgdaDatatype{Gen}\AgdaSpace{}%
\AgdaSymbol{(}\AgdaInductiveConstructor{₂₊}\AgdaSpace{}%
\AgdaGeneralizable{n}\AgdaSymbol{)}\AgdaSpace{}%
\AgdaSymbol{→}\AgdaSpace{}%
\AgdaFunction{Circuit}\AgdaSpace{}%
\AgdaSymbol{(}\AgdaInductiveConstructor{₁₊}\AgdaSpace{}%
\AgdaGeneralizable{n}\AgdaSymbol{)}\AgdaSpace{}%
\AgdaOperator{\AgdaFunction{×}}\AgdaSpace{}%
\AgdaDatatype{C}\AgdaSpace{}%
\AgdaSymbol{(}\AgdaInductiveConstructor{₁₊}\AgdaSpace{}%
\AgdaGeneralizable{n}\AgdaSymbol{)}\<%
\\
\>[0]\AgdaFunction{ract}\AgdaSpace{}%
\AgdaInductiveConstructor{ε}\AgdaSpace{}%
\AgdaInductiveConstructor{σ-gen}\AgdaSpace{}%
\AgdaSymbol{=}\AgdaSpace{}%
\AgdaInductiveConstructor{ε}\AgdaSpace{}%
\AgdaOperator{\AgdaInductiveConstructor{,}}\AgdaSpace{}%
\AgdaOperator{\AgdaInductiveConstructor{σ•}}\AgdaSpace{}%
\AgdaInductiveConstructor{ε}%
\>[36]\AgdaComment{--\ grows}\<%
\\
\>[0]\AgdaFunction{ract}\AgdaSpace{}%
\AgdaSymbol{(}\AgdaOperator{\AgdaInductiveConstructor{σ•}}\AgdaSpace{}%
\AgdaInductiveConstructor{ε}\AgdaSymbol{)}\AgdaSpace{}%
\AgdaInductiveConstructor{σ-gen}\AgdaSpace{}%
\AgdaSymbol{=}\AgdaSpace{}%
\AgdaInductiveConstructor{ε}\AgdaSpace{}%
\AgdaOperator{\AgdaInductiveConstructor{,}}\AgdaSpace{}%
\AgdaInductiveConstructor{ε}%
\>[36]\AgdaComment{--\ cancels\ (σ²\ =\ ε)}\<%
\\
\>[0]\AgdaFunction{ract}\AgdaSpace{}%
\AgdaSymbol{(}\AgdaOperator{\AgdaInductiveConstructor{σ•}}\AgdaSpace{}%
\AgdaOperator{\AgdaInductiveConstructor{σ•}}\AgdaSpace{}%
\AgdaBound{c}\AgdaSymbol{)}\AgdaSpace{}%
\AgdaInductiveConstructor{σ-gen}\AgdaSpace{}%
\AgdaSymbol{=}\AgdaSpace{}%
\AgdaFunction{σ}\AgdaSpace{}%
\AgdaOperator{\AgdaInductiveConstructor{,}}\AgdaSpace{}%
\AgdaOperator{\AgdaInductiveConstructor{σ•}}\AgdaSpace{}%
\AgdaOperator{\AgdaInductiveConstructor{σ•}}\AgdaSpace{}%
\AgdaBound{c}%
\>[36]\AgdaComment{--\ passes\ (Yang-Baxter)}\<%
\\
\>[0]\AgdaFunction{ract}\AgdaSpace{}%
\AgdaInductiveConstructor{ε}\AgdaSpace{}%
\AgdaSymbol{(}\AgdaBound{g}\AgdaSpace{}%
\AgdaOperator{\AgdaInductiveConstructor{↥}}\AgdaSymbol{)}\AgdaSpace{}%
\AgdaSymbol{=}\AgdaSpace{}%
\AgdaOperator{\AgdaInductiveConstructor{[}}\AgdaSpace{}%
\AgdaBound{g}\AgdaSpace{}%
\AgdaOperator{\AgdaInductiveConstructor{]ʷ}}\AgdaSpace{}%
\AgdaOperator{\AgdaInductiveConstructor{,}}\AgdaSpace{}%
\AgdaInductiveConstructor{ε}%
\>[36]\AgdaComment{--\ commutes\ past}\<%
\\
\>[0]\AgdaFunction{ract}\AgdaSpace{}%
\AgdaSymbol{(}\AgdaOperator{\AgdaInductiveConstructor{σ•}}\AgdaSpace{}%
\AgdaBound{c}\AgdaSymbol{)}\AgdaSpace{}%
\AgdaSymbol{(}\AgdaBound{g}\AgdaSpace{}%
\AgdaOperator{\AgdaInductiveConstructor{↥}}\AgdaSymbol{)}\AgdaSpace{}%
\AgdaSymbol{=}\AgdaSpace{}%
\AgdaField{proj₁}\AgdaSpace{}%
\AgdaSymbol{(}\AgdaFunction{ract}\AgdaSpace{}%
\AgdaBound{c}\AgdaSpace{}%
\AgdaBound{g}\AgdaSymbol{)}\AgdaSpace{}%
\AgdaOperator{\AgdaFunction{↑}}\AgdaSpace{}%
\AgdaOperator{\AgdaInductiveConstructor{,}}\AgdaSpace{}%
\AgdaOperator{\AgdaInductiveConstructor{σ•}}\AgdaSpace{}%
\AgdaSymbol{(}\AgdaField{proj₂}\AgdaSpace{}%
\AgdaSymbol{(}\AgdaFunction{ract}\AgdaSpace{}%
\AgdaBound{c}\AgdaSpace{}%
\AgdaBound{g}\AgdaSymbol{))}\<%
\end{code}
Each clause is certified by the \alert{soundness law}, one short derivation per case:
\begin{code}%
\>[0]\AgdaFunction{ract-sound}\AgdaSpace{}%
\AgdaSymbol{:}%
\>[405I]\AgdaSymbol{∀}\AgdaSpace{}%
\AgdaBound{c}\AgdaSpace{}%
\AgdaBound{b}\AgdaSpace{}%
\AgdaSymbol{→}\AgdaSpace{}%
\AgdaKeyword{let}\AgdaSpace{}%
\AgdaSymbol{(}\AgdaBound{b'}\AgdaSpace{}%
\AgdaOperator{\AgdaInductiveConstructor{,}}\AgdaSpace{}%
\AgdaBound{c'}\AgdaSymbol{)}\AgdaSpace{}%
\AgdaSymbol{=}\AgdaSpace{}%
\AgdaFunction{ract}\AgdaSpace{}%
\AgdaBound{c}\AgdaSpace{}%
\AgdaBound{b}\AgdaSpace{}%
\AgdaKeyword{in}\<%
\\
\>[.][@{}l@{}]\<[405I]%
\>[13]\AgdaOperator{\AgdaFunction{[}}\AgdaSpace{}%
\AgdaBound{c}\AgdaSpace{}%
\AgdaOperator{\AgdaFunction{]ᶜ}}\AgdaSpace{}%
\AgdaOperator{\AgdaInductiveConstructor{•}}\AgdaSpace{}%
\AgdaOperator{\AgdaInductiveConstructor{[}}\AgdaSpace{}%
\AgdaBound{b}\AgdaSpace{}%
\AgdaOperator{\AgdaInductiveConstructor{]ʷ}}\AgdaSpace{}%
\AgdaOperator{\AgdaPostulate{≈}}\AgdaSpace{}%
\AgdaBound{b'}\AgdaSpace{}%
\AgdaOperator{\AgdaFunction{↑}}\AgdaSpace{}%
\AgdaOperator{\AgdaInductiveConstructor{•}}\AgdaSpace{}%
\AgdaOperator{\AgdaFunction{[}}\AgdaSpace{}%
\AgdaBound{c'}\AgdaSpace{}%
\AgdaOperator{\AgdaFunction{]ᶜ}}\<%
\end{code}
\begin{code}[hide]%
\>[0]\AgdaFunction{ract-sound}\AgdaSpace{}%
\AgdaSymbol{=}\AgdaSpace{}%
\AgdaPostulate{proof-elided}\<%
\end{code}
\end{column}
\begin{column}{0.37\textwidth}
\centering
\begin{tikzpicture}
\node (l) {\ufig[0.42]{ract-l}};
\node[right=0mm of l] (eq) {$\approx$};
\node[right=0mm of eq] (r) {\ufig[0.42]{ract-r}};
\end{tikzpicture}

\vspace{6pt}
\begin{keybox}
\small staircase $c$, then generator $b$\\
\quad$\approx$\\
residual $b'$ (one wire up), then staircase $c'$
\end{keybox}
\small The residual is a \emph{circuit}, not a generator: in richer
gate sets it can be long.
\end{column}
\end{columns}
\end{frame}

\begin{frame}{The four cases, as circuit rewrites \InPaper}
\centering
\renewcommand{\arraystretch}{1.6}
\begin{tabular}{@{}r@{\,}c@{\,}l@{\;\;}l@{\qquad}r@{\,}c@{\,}l@{\;\;}l@{}}
  \ufig[0.38]{push4-l} & $\rightarrow$ & \ufig[0.38]{push4-r} &
    \small grows &
  \ufig[0.38]{push3-l} & $\rightarrow$ & \ufig[0.38]{push3-r} &
    \small cancels ($\sigma^2 = \varepsilon$)\\[12pt]
  \ufig[0.38]{push2-l} & $\rightarrow$ & \ufig[0.38]{push2-r} &
    \small passes (Yang--Baxter) &
  \ufig[0.38]{push1-l} & $\rightarrow$ & \ufig[0.38]{push1-r} &
    \small commutes past
\end{tabular}

\vspace{12pt}
\begin{keybox}
\small Each step of normalization is an instance of \aid{order}, of
\aid{yang-baxter}, or of a structural rule (\aid{comm₂}, lifted by
\aid{cong↑}). Nothing else is ever used.
\end{keybox}
\end{frame}

\begin{frame}{Running the table: normalizing $\sigma \aidop{•} \sigma\;{\uparrow} \aidop{•} \sigma$ \InPaper}
\small
The stateful traversal \aid{\_ᵗ} threads the coset left to right and
concatenates the residuals --- one \alert{level} of normalization:
\vspace{6pt}

\centering
\begin{tikzpicture}[>={Stealth[length=2mm]},
  st/.style={draw=accent,fill=soft,rounded corners=2pt,minimum width=18mm,minimum height=13mm,align=center},
  lab/.style={font=\scriptsize,align=center}]
  \node[st] (s0) at (0,0) {\permcirc[0.22]{3}{x}\\[-1pt]{\scriptsize $\varepsilon$}};
  \node[st] (s1) at (3.6,0) {\permcirc[0.22]{3}{0}\\[-1pt]{\scriptsize $\sigma{\bullet}\varepsilon$}};
  \node[st] (s2) at (7.2,0) {\permcirc[0.22]{3}{0,1}\\[-1pt]{\scriptsize $\sigma{\bullet}\sigma{\bullet}\varepsilon$}};
  \node[st] (s3) at (10.8,0) {\permcirc[0.22]{3}{0,1}\\[-1pt]{\scriptsize $\sigma{\bullet}\sigma{\bullet}\varepsilon$}};
  \draw[->,thick] (s0) -- node[above,lab]{read $\sigma$}node[below,lab]{grows\\residual $\varepsilon$} (s1);
  \draw[->,thick] (s1) -- node[above,lab]{read $\sigma\;{\uparrow}$}node[below,lab]{recurses\\residual $\varepsilon$} (s2);
  \draw[->,thick] (s2) -- node[above,lab]{read $\sigma$}node[below,lab]{passes (YB)\\residual $\sigma$} (s3);
  \node[lab,above=1mm of s0] {coset};
\end{tikzpicture}

\vspace{4pt}
\begin{columns}[c]
\begin{column}{0.55\textwidth}
\small
Level 3 wires: top digit $c_2 = \sigma{\bullet}\sigma{\bullet}\varepsilon$, residual $\sigma$ on the upper 2 wires.\\
Level 2 wires: normalize the residual $\sigma$: digit $c_1 = \sigma{\bullet}\varepsilon$.\\[3pt]
Rebuild: $[c_1]^{\mathsf c}\;{\uparrow} \aidop{•} [c_2]^{\mathsf c}
 \;=\; \sigma\;{\uparrow} \aidop{•} \sigma \aidop{•} \sigma\;{\uparrow}$.
\end{column}
\begin{column}{0.40\textwidth}
\centering
\permcirc[0.26]{3}{0,1,0}\; $\approx$\; \permcirc[0.26]{3}{1,0,1}

{\scriptsize the normal form of the Yang--Baxter left side\\
is its right side --- found by computation}
\end{column}
\end{columns}
\end{frame}

\begin{frame}{The coset tower: one digit per wire \InPaper}
\begin{columns}[c]
\begin{column}{0.48\textwidth}
Iterating one level per wire is the \alert{coset tower}:
\begin{code}%
\>[0]\AgdaFunction{NF}\AgdaSpace{}%
\AgdaSymbol{:}\AgdaSpace{}%
\AgdaDatatype{ℕ}\AgdaSpace{}%
\AgdaSymbol{→}\AgdaSpace{}%
\AgdaPrimitive{Set}\<%
\\
\>[0]\AgdaFunction{NF}\AgdaSpace{}%
\AgdaNumber{0}%
\>[11]\AgdaSymbol{=}\AgdaSpace{}%
\AgdaRecord{⊤}\<%
\\
\>[0]\AgdaFunction{NF}\AgdaSpace{}%
\AgdaSymbol{(}\AgdaInductiveConstructor{suc}\AgdaSpace{}%
\AgdaBound{k}\AgdaSymbol{)}\AgdaSpace{}%
\AgdaSymbol{=}\AgdaSpace{}%
\AgdaFunction{NF}\AgdaSpace{}%
\AgdaBound{k}\AgdaSpace{}%
\AgdaOperator{\AgdaFunction{×}}\AgdaSpace{}%
\AgdaDatatype{C}\AgdaSpace{}%
\AgdaBound{k}%
\>[26]\AgdaComment{--\ ⊤\ ×\ C\ 0\ ×\ C\ 1\ ×\ ⋯\ ×\ C\ k}\<%
\end{code}
so $|\aid{NF}\;n| = 1\cdot 2 \cdots n = n!$: the \alert{factorial number
system}.

\vspace{4pt}
\aid{nf : Circuit n → NF n} runs the tables; its \alert{section}
\aid{inv-nf} rebuilds a circuit from the digits --- a ``staircase of
staircases''.
\vspace{4pt}

\centering
\ufig[0.40]{fig3_staircase}\quad{\scriptsize $(\mathsf{tt},c_1,c_2,c_3)\in\aid{NF}\;4$}
\end{column}
\begin{column}{0.49\textwidth}
\centering
{\small all six normal forms on three wires ($S_3$):}\\[6pt]
\renewcommand{\arraystretch}{1.3}
\begin{tabular}{@{}c@{\quad}c@{\quad}c@{}}
 & $c_2=\varepsilon$ & \\
\permcirc[0.24]{3}{x} & \permcirc[0.24]{3}{0} & \permcirc[0.24]{3}{0,1} \\
{\scriptsize $(\varepsilon,\varepsilon)$} & {\scriptsize $(\varepsilon,\sigma\varepsilon)$} & {\scriptsize $(\varepsilon,\sigma\sigma\varepsilon)$}\\[6pt]
\permcirc[0.24]{3}{1} & \permcirc[0.24]{3}{1,0} & \permcirc[0.24]{3}{1,0,1} \\
{\scriptsize $(\sigma\varepsilon,\varepsilon)$} & {\scriptsize $(\sigma\varepsilon,\sigma\varepsilon)$} & {\scriptsize $(\sigma\varepsilon,\sigma\sigma\varepsilon)$}
\end{tabular}

\vspace{6pt}
{\scriptsize digit $c_1\in\aid{C}\;1$ (2 values) selects the row,
$c_2\in\aid{C}\;2$ (3 values) the column;\\
section $(c_1,c_2)\mapsto [c_1]^{\mathsf c}\,{\uparrow}\aidop{•}[c_2]^{\mathsf c}$}
\end{column}
\end{columns}
\end{frame}

\begin{frame}{Two semantics, uniqueness, completeness \InPaper}
\begin{columns}[c]
\begin{column}{0.52\textwidth}
\begin{itemize}
  \setlength{\itemsep}{4pt}
  \item \textbf{loose}: \aid{⟦\_⟧ : Circuit n → (Fin n → Fin n)}, endofunctions
        --- a monoid that is \emph{not} a group;
  \item \textbf{tight}: into \aid{Permutation′-group n}, stdlib permutations
        bundled as a \aid{Group} (our \texttt{ForStdlib}).
\end{itemize}
\vspace{4pt}
\textbf{Soundness}: small case analysis over the axioms.\\
\textbf{Unique NF}: $\llbracket \aid{inv-nf}\,u\rrbracket = \llbracket\aid{inv-nf}\,v\rrbracket \Rightarrow u\equiv v$,
by induction on the tower: the lower digits fix the last point, so the top digit alone
says where it goes.\\
\textbf{Completeness} then follows generically --- even for the \emph{loose}
semantics: circuits equal as mere functions are provably equal.
\end{column}
\begin{column}{0.45\textwidth}
\centering
\begin{tikzpicture}[>={Stealth[length=2mm]},node distance=8mm]
  \node (w) {$w$};
  \node[right=22mm of w] (v) {$v$};
  \node[below=of w] (nw) {$\aid{inv-nf}(\aid{nf}\,w)$};
  \node[below=of v] (nv) {$\aid{inv-nf}(\aid{nf}\,v)$};
  \draw[double,thick,accent] (w) -- node[left,font=\scriptsize]{NF law} (nw);
  \draw[double,thick,accent] (v) -- node[right,font=\scriptsize]{NF law} (nv);
  \draw[double,thick,newC] (nw) -- node[below,font=\scriptsize]{$\equiv$} (nv);
  \draw[->,dashed,gray] (w) -- node[above,font=\scriptsize]{$\llbracket w\rrbracket=\llbracket v\rrbracket$} (v);
  \node[below=2mm of $(nw)!0.5!(nv)$,font=\scriptsize,align=center,yshift=-3mm]
    {soundness: the two normal forms have equal denotations;\\
     uniqueness: so they are \emph{equal}};
\end{tikzpicture}

\vspace{6pt}
{\small $w \approx \aid{inv-nf}(\aid{nf}\,w) = \aid{inv-nf}(\aid{nf}\,v)\approx v$ ---
a derivation assembled from the two normalization runs, i.e.\ from
\aid{order}, \aid{yang-baxter} and the structural rules only.}
\end{column}
\end{columns}
\end{frame}

\begin{frame}{The theorem \InPaper}
Soundness, completeness and surjectivity assemble into the
\alert{presentation record}: a \aid{Grouplike} witness (so the presented
monoid is a group) and a proof that \aid{⟦\_⟧} is a group isomorphism.
\vspace{4pt}
\begin{code}%
\>[0]\AgdaFunction{symmetric-presentation}\AgdaSpace{}%
\AgdaSymbol{:}\AgdaSpace{}%
\AgdaSymbol{∀}\AgdaSpace{}%
\AgdaBound{n}\AgdaSpace{}%
\AgdaSymbol{→}\AgdaSpace{}%
\AgdaSymbol{(}\AgdaBound{n}\AgdaSpace{}%
\AgdaOperator{\AgdaDatatype{VRel,\AgdaUnderscore{}===\AgdaUnderscore{}}}\AgdaSymbol{)}\AgdaSpace{}%
\AgdaOperator{\AgdaRecord{IsPresentationOf}}\AgdaSpace{}%
\AgdaSymbol{(}\AgdaFunction{Permutation′-group}\AgdaSpace{}%
\AgdaBound{n}\AgdaSymbol{)}\<%
\end{code}
\begin{code}[hide]%
\>[0]\AgdaFunction{symmetric-presentation}\AgdaSpace{}%
\AgdaBound{n}\AgdaSpace{}%
\AgdaSymbol{=}\AgdaSpace{}%
\AgdaFunction{SymThm.Tight.presentation}\AgdaSpace{}%
\AgdaBound{n}\<%
\end{code}
\vspace{2pt}
\begin{columns}[c]
\begin{column}{0.55\textwidth}
\small Read: the monoid of swap circuits modulo \{\aid{order},
\aid{yang-baxter}\} + structural rules \emph{is} the group of permutations of
$\mathsf{Fin}\;n$ --- for every $n$.

\vspace{4pt}
Every presentation theorem in the library has this shape.
\end{column}
\begin{column}{0.42\textwidth}
\centering
\begin{tikzpicture}[>={Stealth[length=2mm]}]
  \node (W) at (0,1.6) {\aid{Word (Gen n)}};
  \node (Q) at (0,0) {\aid{Word (Gen n)}$\,/\!\approx$};
  \node (P) at (3.6,1.6) {$\text{Perm}(\mathsf{Fin}\;n)$};
  \draw[->>,thick] (W) -- node[left,font=\scriptsize]{quotient} (Q);
  \draw[->,thick] (W) -- node[above,font=\scriptsize]{$\llbracket\_\rrbracket$} (P);
  \draw[->,thick,paperC] (Q) -- node[below right,font=\scriptsize]{$\cong$ (group iso)} (P);
\end{tikzpicture}

\vspace{4pt}
{\scriptsize \aid{gl}: the quotient is a group; \aid{iso}: the induced map is an isomorphism}
\end{column}
\end{columns}
\end{frame}

%% Part 3 — normal forms from subgroup chains (paper §3).
\PaperPart
\section{Normal forms from subgroup chains}

\begin{frame}{Transversals are digits \InPaper}
\begin{columns}[c]
\begin{column}{0.50\textwidth}
$H \le G$, a \alert{right transversal} picks one $\bar g$ per coset $Hg$:
\[ g = h\cdot\bar g \quad\text{uniquely},\qquad G \cong H\times (H\backslash G)\ \text{(as sets)}. \]
Fix a chain with a transversal $T_k$ for each step:
\[ 1 = G_0 \le G_1 \le \dots \le G_n = G \]
\[ g = t_1\,t_2\cdots t_n,\qquad t_k\in T_k\quad\text{uniquely.} \]
With indices $r_k = [G_k:G_{k-1}]$: $|G| = r_1 r_2\cdots r_n$ --- a
\alert{mixed-radix numeral}.

\vspace{4pt}
\small $S_1<S_2<\dots<S_n$: radices $2,3,\dots,n$ = the factorial number
system. The conceptual core of \alert{Schreier--Sims}; our chains are
chosen \emph{syntactically}.
\end{column}
\begin{column}{0.47\textwidth}
\centering
\begin{tikzpicture}[x=11mm,y=11mm]
  % odometer wheels
  \foreach \r/\d/\lab [count=\i] in {2/1/T_1,3/2/T_2,4/0/T_3,5/3/T_4} {
    \pgfmathsetmacro{\x}{\i*1.25}
    \draw[accent,fill=soft,rounded corners=2pt] (\x-0.5,-0.2) rectangle (\x+0.5,2.6);
    \foreach \k in {1,...,\r} {
      \pgfmathsetmacro{\yy}{2.6-\k*2.8/(\r+1)}
      \pgfmathtruncatemacro{\kk}{\k-1}
      \ifnum\kk=\d
        \node[fill=newC,text=white,font=\scriptsize,inner sep=1.5pt,rounded corners=1pt] at (\x,\yy) {\kk};
      \else
        \node[font=\scriptsize,gray] at (\x,\yy) {\kk};
      \fi
    }
    \node[font=\small] at (\x,-0.55) {$\lab$};
    \node[font=\scriptsize,gray] at (\x,2.85) {$r_\i{=}\r$};
  }
  \node[font=\small,align=center] at (3.1,-1.3) {one digit per level:\; $5!=120$ elements};
\end{tikzpicture}
\end{column}
\end{columns}
\end{frame}

\begin{frame}{Reidemeister--Schreier: certificates, not search \InPaper}
\begin{columns}[c]
\begin{column}{0.50\textwidth}
Monoid presentation $\langle X\mid R\rangle$: the quotient of $X^*$ by the
least congruence $\approx$ containing $R$; the \alert{word problem} decides
$\approx$. Interpretation $X\to G$ extends to $X^*\to G$; \emph{soundness}
descends it to $X^*/{\approx}\to G$; injective = \emph{complete}.

\vspace{6pt}
Given $H\le G=\langle Y\mid S\rangle$, a \alert{coset table} pushes a
generator $y$ past a representative:
\[ \bar t\,y \;=\; h\cdot\overline{\tau(t,y)},\qquad h\in H. \]
That is exactly \aid{ract}, and the equation is \aid{ract-sound}.
\end{column}
\begin{column}{0.47\textwidth}
\centering
\begin{tikzpicture}[node distance=6mm,
  bx/.style={draw,rounded corners=2pt,align=center,font=\small,minimum height=9mm,text width=30mm},
  >={Stealth[length=2mm]}]
  \node[bx,draw=gray,fill=gray!10] (tc) {Todd--Coxeter / by hand\\{\scriptsize search, unverified}};
  \node[bx,draw=accent,fill=soft,below=of tc] (tab) {claimed coset table\\{\scriptsize finite data}};
  \node[bx,draw=paperC,fill=softpaper,below=of tab] (chk) {Agda checks\\{\scriptsize 5 equations}};
  \node[bx,draw=paperC,fill=paperC!20,below=of chk] (nf) {normal form one\\level up, proved once};
  \draw[->,thick] (tc) -- (tab);
  \draw[->,thick] (tab) -- (chk);
  \draw[->,thick,paperC] (chk) -- node[right,font=\scriptsize]{\aid{nf-transp}} (nf);
\end{tikzpicture}

\vspace{4pt}
{\small The certificate is separated from the search that found it
(pioneered in Agda by Bian--Selinger; here a reusable module).}
\end{column}
\end{columns}
\end{frame}

\begin{frame}{For circuits the chain is handed to us \InPaper}
\begin{columns}[c]
\begin{column}{0.48\textwidth}
The shift $c\mapsto c\;{\uparrow}$ embeds
$\mathsf{Cir}\,k \le \mathsf{Cir}\,(k{+}1)$ as the circuits that do not touch
the bottom wire. So:
\begin{itemize}
  \item each \alert{wire} adds one digit, ranging over a transversal:
        staircases for $S_n$; Selinger's staircase representatives for
        Clifford groups; their qutrit analogues;
  \item a normal form for the whole group is a \alert{tower} ---
        \aid{CosetTower};
  \item where the chain \alert{branches} (two subgroups over a common base, as
        for Clifford+$T$), glue by the \alert{amalgamated free product} and its
        alternating normal form.
\end{itemize}
\end{column}
\begin{column}{0.49\textwidth}
\centering
\begin{tikzpicture}[x=6mm,y=6mm,>={Stealth[length=1.8mm]}]
  % nested tower: largest first, so the smaller ones sit on top
  \foreach \k in {4,3,2,1,0} {
    \pgfmathtruncatemacro{\shade}{6+(4-\k)*7}
    \draw[accent,fill=accent!\shade,rounded corners=2pt]
      (-\k*0.5-0.7,0) rectangle (\k*0.5+0.7,\k*1.0+0.8);
    \node[font=\scriptsize] at (0,\k*1.0+0.45) {$\mathsf{Cir}\,\k$};
  }
  \node[font=\scriptsize,align=center] at (0,-0.6) {tower: nested};
  % branching
  \begin{scope}[xshift=44mm,yshift=4mm]
    \node[draw=paperC,fill=softpaper,rounded corners=2pt,font=\scriptsize] (M) at (0,0) {$M$};
    \node[draw=paperC,fill=softpaper,rounded corners=2pt,font=\scriptsize] (A) at (-1.5,2.4) {$\langle T,M\rangle$};
    \node[draw=paperC,fill=softpaper,rounded corners=2pt,font=\scriptsize] (B) at (1.5,2.4) {$\langle H,M\rangle$};
    \node[draw=newC,fill=softnew,rounded corners=2pt,font=\scriptsize] (P) at (0,4.6) {$\langle T,M\rangle *_M \langle H,M\rangle$};
    \draw[->] (M) -- (A); \draw[->] (M) -- (B);
    \draw[->] (A) -- (P); \draw[->] (B) -- (P);
    \node[font=\scriptsize,align=center] at (0,-1.3) {branching: amalgam};
  \end{scope}
\end{tikzpicture}
\end{column}
\end{columns}
\end{frame}

\begin{code}[hide]%
\>[0]\AgdaKeyword{module}\AgdaSpace{}%
\AgdaModule{slides.Tsinghua-PL-seminar.sections.design}\AgdaSpace{}%
\AgdaKeyword{where}\<%
\\
%
\\[\AgdaEmptyExtraSkip]%
\>[0]\AgdaKeyword{open}\AgdaSpace{}%
\AgdaKeyword{import}\AgdaSpace{}%
\AgdaModule{Data.Nat}\AgdaSpace{}%
\AgdaKeyword{using}\AgdaSpace{}%
\AgdaSymbol{(}\AgdaDatatype{ℕ}\AgdaSpace{}%
\AgdaSymbol{;}\AgdaSpace{}%
\AgdaOperator{\AgdaPrimitive{\AgdaUnderscore{}+\AgdaUnderscore{}}}\AgdaSymbol{)}\<%
\\
\>[0]\AgdaKeyword{open}\AgdaSpace{}%
\AgdaKeyword{import}\AgdaSpace{}%
\AgdaModule{Data.Product}\AgdaSpace{}%
\AgdaKeyword{using}\AgdaSpace{}%
\AgdaSymbol{(}\AgdaOperator{\AgdaFunction{\AgdaUnderscore{}×\AgdaUnderscore{}}}\AgdaSpace{}%
\AgdaSymbol{;}\AgdaSpace{}%
\AgdaOperator{\AgdaInductiveConstructor{\AgdaUnderscore{},\AgdaUnderscore{}}}\AgdaSymbol{)}\<%
\\
\>[0]\AgdaKeyword{open}\AgdaSpace{}%
\AgdaKeyword{import}\AgdaSpace{}%
\AgdaModule{Data.Sum}\AgdaSpace{}%
\AgdaKeyword{using}\AgdaSpace{}%
\AgdaSymbol{(}\AgdaOperator{\AgdaDatatype{\AgdaUnderscore{}⊎\AgdaUnderscore{}}}\AgdaSpace{}%
\AgdaSymbol{;}\AgdaSpace{}%
\AgdaInductiveConstructor{inj₁}\AgdaSpace{}%
\AgdaSymbol{;}\AgdaSpace{}%
\AgdaInductiveConstructor{inj₂}\AgdaSymbol{)}\<%
\\
\>[0]\AgdaKeyword{open}\AgdaSpace{}%
\AgdaKeyword{import}\AgdaSpace{}%
\AgdaModule{Data.List}\AgdaSpace{}%
\AgdaKeyword{using}\AgdaSpace{}%
\AgdaSymbol{(}\AgdaDatatype{List}\AgdaSymbol{)}\<%
\\
\>[0]\AgdaKeyword{open}\AgdaSpace{}%
\AgdaKeyword{import}\AgdaSpace{}%
\AgdaModule{Data.Unit}\AgdaSpace{}%
\AgdaKeyword{using}\AgdaSpace{}%
\AgdaSymbol{(}\AgdaRecord{⊤}\AgdaSpace{}%
\AgdaSymbol{;}\AgdaSpace{}%
\AgdaInductiveConstructor{tt}\AgdaSymbol{)}\<%
\\
\>[0]\AgdaKeyword{open}\AgdaSpace{}%
\AgdaKeyword{import}\AgdaSpace{}%
\AgdaModule{Level}\AgdaSpace{}%
\AgdaKeyword{using}\AgdaSpace{}%
\AgdaSymbol{(}\AgdaPostulate{Level}\AgdaSpace{}%
\AgdaSymbol{;}\AgdaSpace{}%
\AgdaOperator{\AgdaPrimitive{\AgdaUnderscore{}⊔\AgdaUnderscore{}}}\AgdaSpace{}%
\AgdaSymbol{;}\AgdaSpace{}%
\AgdaFunction{0ℓ}\AgdaSymbol{)}\<%
\\
\>[0]\AgdaKeyword{open}\AgdaSpace{}%
\AgdaKeyword{import}\AgdaSpace{}%
\AgdaModule{Relation.Binary}\AgdaSpace{}%
\AgdaKeyword{using}\AgdaSpace{}%
\AgdaSymbol{(}\AgdaFunction{Rel}\AgdaSymbol{)}\<%
\\
\>[0]\AgdaKeyword{open}\AgdaSpace{}%
\AgdaKeyword{import}\AgdaSpace{}%
\AgdaModule{Relation.Binary.Definitions}\AgdaSpace{}%
\AgdaKeyword{using}\AgdaSpace{}%
\AgdaSymbol{(}\AgdaFunction{Decidable}\AgdaSymbol{)}\<%
\\
\>[0]\AgdaKeyword{import}\AgdaSpace{}%
\AgdaModule{Relation.Binary.PropositionalEquality}\AgdaSpace{}%
\AgdaSymbol{as}\AgdaSpace{}%
\AgdaModule{Eq}\<%
\\
\>[0]\AgdaKeyword{open}\AgdaSpace{}%
\AgdaKeyword{import}\AgdaSpace{}%
\AgdaModule{Function.Definitions}\AgdaSpace{}%
\AgdaKeyword{using}\AgdaSpace{}%
\AgdaSymbol{(}\AgdaFunction{Congruent}\AgdaSpace{}%
\AgdaSymbol{;}\AgdaSpace{}%
\AgdaFunction{Injective}\AgdaSymbol{)}\<%
\\
\>[0]\AgdaKeyword{open}\AgdaSpace{}%
\AgdaKeyword{import}\AgdaSpace{}%
\AgdaModule{Algebra.Bundles}\AgdaSpace{}%
\AgdaKeyword{using}\AgdaSpace{}%
\AgdaSymbol{(}\AgdaRecord{Group}\AgdaSymbol{)}\<%
\\
\>[0]\AgdaKeyword{open}\AgdaSpace{}%
\AgdaKeyword{import}\AgdaSpace{}%
\AgdaModule{Algebra.Morphism.Structures}\AgdaSpace{}%
\AgdaKeyword{using}\AgdaSpace{}%
\AgdaSymbol{(}\AgdaKeyword{module}\AgdaSpace{}%
\AgdaModule{GroupMorphisms}\AgdaSymbol{)}\<%
\\
%
\\[\AgdaEmptyExtraSkip]%
\>[0]\AgdaKeyword{open}\AgdaSpace{}%
\AgdaKeyword{import}\AgdaSpace{}%
\AgdaModule{Notations}\AgdaSpace{}%
\AgdaKeyword{using}\AgdaSpace{}%
\AgdaSymbol{(}\AgdaInductiveConstructor{₁₊}\AgdaSpace{}%
\AgdaSymbol{;}\AgdaSpace{}%
\AgdaInductiveConstructor{₂₊}\AgdaSymbol{)}\<%
\\
\>[0]\AgdaKeyword{open}\AgdaSpace{}%
\AgdaKeyword{import}\AgdaSpace{}%
\AgdaModule{Word.Base}\AgdaSpace{}%
\AgdaKeyword{using}\AgdaSpace{}%
\AgdaSymbol{(}\AgdaFunction{wconcatmap}\AgdaSymbol{)}\<%
\\
\>[0]\AgdaKeyword{import}\AgdaSpace{}%
\AgdaModule{Word.Base}\AgdaSpace{}%
\AgdaSymbol{as}\AgdaSpace{}%
\AgdaModule{WB}\<%
\\
\>[0]\AgdaKeyword{open}\AgdaSpace{}%
\AgdaKeyword{import}\AgdaSpace{}%
\AgdaModule{Presentation.GroupLike}\AgdaSpace{}%
\AgdaKeyword{using}\AgdaSpace{}%
\AgdaSymbol{(}\AgdaFunction{Grouplike}\AgdaSpace{}%
\AgdaSymbol{;}\AgdaSpace{}%
\AgdaKeyword{module}\AgdaSpace{}%
\AgdaModule{Group-Lemmas}\AgdaSymbol{)}\<%
\\
\>[0]\AgdaKeyword{open}\AgdaSpace{}%
\AgdaKeyword{import}\AgdaSpace{}%
\AgdaModule{Presentation.Construct.Base}\AgdaSpace{}%
\AgdaKeyword{using}\AgdaSpace{}%
\AgdaSymbol{(}\AgdaOperator{\AgdaFunction{[\AgdaUnderscore{}]ₗ}}\AgdaSpace{}%
\AgdaSymbol{;}\AgdaSpace{}%
\AgdaOperator{\AgdaFunction{[\AgdaUnderscore{}]ᵣ}}\AgdaSymbol{)}\<%
\\
%
\\[\AgdaEmptyExtraSkip]%
\>[0]\AgdaKeyword{postulate}\<%
\\
\>[0][@{}l@{\AgdaIndent{0}}]%
\>[2]\AgdaPostulate{proof-elided}\AgdaSpace{}%
\AgdaSymbol{:}\AgdaSpace{}%
\AgdaSymbol{∀}\AgdaSpace{}%
\AgdaSymbol{\{}\AgdaBound{ℓ}\AgdaSymbol{\}}\AgdaSpace{}%
\AgdaSymbol{\{}\AgdaBound{A}\AgdaSpace{}%
\AgdaSymbol{:}\AgdaSpace{}%
\AgdaPrimitive{Set}\AgdaSpace{}%
\AgdaBound{ℓ}\AgdaSymbol{\}}\AgdaSpace{}%
\AgdaSymbol{→}\AgdaSpace{}%
\AgdaBound{A}\<%
\\
%
\>[2]\AgdaPostulate{X}\AgdaSpace{}%
\AgdaPostulate{Y}\AgdaSpace{}%
\AgdaPostulate{A}\AgdaSpace{}%
\AgdaPostulate{B}\AgdaSpace{}%
\AgdaPostulate{C}\AgdaSpace{}%
\AgdaPostulate{D}%
\>[15]\AgdaSymbol{:}\AgdaSpace{}%
\AgdaPrimitive{Set}\<%
\\
%
\>[2]\AgdaPostulate{a}\AgdaSpace{}%
\AgdaPostulate{ℓ}\AgdaSpace{}%
\AgdaPostulate{c}\AgdaSpace{}%
\AgdaPostulate{d}%
\>[15]\AgdaSymbol{:}\AgdaSpace{}%
\AgdaPostulate{Level}\<%
\\
%
\>[2]\AgdaPostulate{n}\AgdaSpace{}%
\AgdaPostulate{k}%
\>[15]\AgdaSymbol{:}\AgdaSpace{}%
\AgdaDatatype{ℕ}\<%
\end{code}
%% Part 4 — the library, layer by layer (paper §4).
\PaperPart
\section{The library, layer by layer}

\begin{frame}{Five design decisions \InPaper}
\begin{columns}[c]
\begin{column}{0.60\textwidth}
\begin{enumerate}
  \setlength{\itemsep}{2pt}
  \item \textbf{Syntax is unquotiented}; everything out of it is a function
        on words --- normal forms, coset tables, sections are structurally
        recursive programs the evaluator runs.
  \item \textbf{Specifications are stdlib bundles}: a normal form \emph{is}
        an \aid{Injection}, \aid{RightInverse} or \aid{Bijection} out of the
        word setoid. No new theory to maintain.
  \item \textbf{Certificates are finite}: a verified coset table is finite
        data plus five equations; obligations are case splits whose cases
        are closed goals.
  \item \textbf{Indices are chosen for the unifier}: wire counts built from
        successors only.
  \item \textbf{Constructions come in pairs}: a syntactic theorem (normal
        forms combine) and a semantic one (the constructed relation set
        presents the constructed group).
\end{enumerate}
\end{column}
\begin{column}{0.37\textwidth}
\centering
\begin{tikzpicture}[x=10mm,y=7.6mm]
  \node[ly=paperC!25,minimum width=50mm] at (2.5,7) {Examples: $S_n$, Pauli, Clifford+$T$, \ldots};
  \node[ly=soft,minimum width=24mm] at (1.2,6) {Constructions\\[-2pt]{\tiny $\times,\rtimes,\ast_M$, ext.}};
  \node[ly=soft,minimum width=24mm] at (3.8,6) {Circuits\\[-2pt]{\tiny \aid{Lift-Relation}}};
  \node[ly=soft,minimum width=50mm] at (2.5,5) {Normalization: NF bundles, R--S, cosets, towers};
  \node[ly=soft,minimum width=50mm] at (2.5,4) {Presentations: $\_{===}\_$ and its congruence $\_{\approx}\_$};
  \node[ly=soft,minimum width=50mm] at (2.5,3) {Words: the free monoid \aid{Word X}};
  \node[ly=gray!10,minimum width=50mm] at (2.5,2) {\texttt{Notations} \quad\textbar\quad \texttt{ForStdlib} (semantics)};
  \node[ly=gray!5,draw=gray,minimum width=50mm] at (2.5,1) {Agda standard library 2.4};
\end{tikzpicture}
\end{column}
\end{columns}
\end{frame}

\begin{frame}{Words: a syntax tree, not a list \InPaper}
\begin{columns}[c]
\begin{column}{0.56\textwidth}
\begin{code}%
\>[0]\AgdaKeyword{data}\AgdaSpace{}%
\AgdaDatatype{Word}\AgdaSpace{}%
\AgdaSymbol{(}\AgdaBound{X}\AgdaSpace{}%
\AgdaSymbol{:}\AgdaSpace{}%
\AgdaPrimitive{Set}\AgdaSymbol{)}\AgdaSpace{}%
\AgdaSymbol{:}\AgdaSpace{}%
\AgdaPrimitive{Set}\AgdaSpace{}%
\AgdaKeyword{where}\<%
\\
\>[0][@{}l@{\AgdaIndent{0}}]%
\>[2]\AgdaOperator{\AgdaInductiveConstructor{[\AgdaUnderscore{}]ʷ}}\AgdaSpace{}%
\AgdaSymbol{:}\AgdaSpace{}%
\AgdaBound{X}\AgdaSpace{}%
\AgdaSymbol{→}\AgdaSpace{}%
\AgdaDatatype{Word}\AgdaSpace{}%
\AgdaBound{X}\<%
\\
%
\>[2]\AgdaInductiveConstructor{ε}%
\>[7]\AgdaSymbol{:}\AgdaSpace{}%
\AgdaDatatype{Word}\AgdaSpace{}%
\AgdaBound{X}\<%
\\
%
\>[2]\AgdaOperator{\AgdaInductiveConstructor{\AgdaUnderscore{}•\AgdaUnderscore{}}}%
\>[7]\AgdaSymbol{:}\AgdaSpace{}%
\AgdaDatatype{Word}\AgdaSpace{}%
\AgdaBound{X}\AgdaSpace{}%
\AgdaSymbol{→}\AgdaSpace{}%
\AgdaDatatype{Word}\AgdaSpace{}%
\AgdaBound{X}\AgdaSpace{}%
\AgdaSymbol{→}\AgdaSpace{}%
\AgdaDatatype{Word}\AgdaSpace{}%
\AgdaBound{X}\<%
\end{code}
\begin{code}[hide]%
\>[0]\AgdaKeyword{infixr}\AgdaSpace{}%
\AgdaNumber{7}\AgdaSpace{}%
\AgdaOperator{\AgdaInductiveConstructor{\AgdaUnderscore{}•\AgdaUnderscore{}}}\<%
\\
\>[0]\AgdaKeyword{infixl}\AgdaSpace{}%
\AgdaNumber{9}\AgdaSpace{}%
\AgdaOperator{\AgdaFunction{\AgdaUnderscore{}ʷ}}\AgdaSpace{}%
\AgdaOperator{\AgdaFunction{\AgdaUnderscore{}ᵗ}}\<%
\\
\>[0]\AgdaKeyword{infix}\AgdaSpace{}%
\AgdaNumber{4}\AgdaSpace{}%
\AgdaOperator{\AgdaDatatype{\AgdaUnderscore{}≈\AgdaUnderscore{}}}\AgdaSpace{}%
\AgdaOperator{\AgdaPostulate{\AgdaUnderscore{}===\AgdaUnderscore{}}}\AgdaSpace{}%
\AgdaOperator{\AgdaPostulate{\AgdaUnderscore{}≈ₙ\AgdaUnderscore{}}}\AgdaSpace{}%
\AgdaOperator{\AgdaPostulate{\AgdaUnderscore{}≈₂\AgdaUnderscore{}}}\AgdaSpace{}%
\AgdaOperator{\AgdaPostulate{\AgdaUnderscore{}\textasciitilde{}\AgdaUnderscore{}}}\AgdaSpace{}%
\AgdaOperator{\AgdaPostulate{\AgdaUnderscore{}===₁\AgdaUnderscore{}}}\AgdaSpace{}%
\AgdaOperator{\AgdaPostulate{\AgdaUnderscore{}===₂\AgdaUnderscore{}}}\<%
\\
%
\\[\AgdaEmptyExtraSkip]%
%
\\[\AgdaEmptyExtraSkip]%
\>[0]\AgdaFunction{WRel}\AgdaSpace{}%
\AgdaSymbol{:}\AgdaSpace{}%
\AgdaPrimitive{Set}\AgdaSpace{}%
\AgdaSymbol{→}\AgdaSpace{}%
\AgdaPrimitive{Set₁}\<%
\\
\>[0]\AgdaFunction{WRel}\AgdaSpace{}%
\AgdaBound{X}\AgdaSpace{}%
\AgdaSymbol{=}\AgdaSpace{}%
\AgdaFunction{Rel}\AgdaSpace{}%
\AgdaSymbol{(}\AgdaDatatype{Word}\AgdaSpace{}%
\AgdaBound{X}\AgdaSymbol{)}\AgdaSpace{}%
\AgdaFunction{0ℓ}\<%
\\
%
\\[\AgdaEmptyExtraSkip]%
\>[0]\AgdaKeyword{postulate}\<%
\\
\>[0][@{}l@{\AgdaIndent{0}}]%
\>[2]\AgdaPostulate{w}\AgdaSpace{}%
\AgdaPostulate{v}\AgdaSpace{}%
\AgdaPostulate{u}\AgdaSpace{}%
\AgdaPostulate{w'}\AgdaSpace{}%
\AgdaPostulate{v'}\AgdaSpace{}%
\AgdaSymbol{:}\AgdaSpace{}%
\AgdaDatatype{Word}\AgdaSpace{}%
\AgdaPostulate{X}\<%
\\
%
\>[2]\AgdaOperator{\AgdaPostulate{\AgdaUnderscore{}===\AgdaUnderscore{}}}%
\>[14]\AgdaSymbol{:}\AgdaSpace{}%
\AgdaFunction{WRel}\AgdaSpace{}%
\AgdaPostulate{X}\<%
\end{code}
A word remembers its bracketing: terms mirror pen-and-paper calculations,
rules apply to any subterm, and associativity is a \alert{congruence
generator}, discharged by a solver.

\vspace{4pt}
Generator substitutions extend to words (``the unique monoid homomorphism''):
\begin{code}%
\>[0]\AgdaOperator{\AgdaFunction{\AgdaUnderscore{}ʷ}}\AgdaSpace{}%
\AgdaSymbol{:}\AgdaSpace{}%
\AgdaSymbol{(}\AgdaPostulate{X}\AgdaSpace{}%
\AgdaSymbol{→}\AgdaSpace{}%
\AgdaDatatype{Word}\AgdaSpace{}%
\AgdaPostulate{Y}\AgdaSymbol{)}\AgdaSpace{}%
\AgdaSymbol{→}\AgdaSpace{}%
\AgdaSymbol{(}\AgdaDatatype{Word}\AgdaSpace{}%
\AgdaPostulate{X}\AgdaSpace{}%
\AgdaSymbol{→}\AgdaSpace{}%
\AgdaDatatype{Word}\AgdaSpace{}%
\AgdaPostulate{Y}\AgdaSymbol{)}\<%
\\
\>[0]\AgdaOperator{\AgdaFunction{\AgdaUnderscore{}ʷ}}\AgdaSpace{}%
\AgdaSymbol{=}\AgdaSpace{}%
\AgdaFunction{wconcatmap}\<%
\end{code}
and a \alert{coset action} extends to words by threading the coset:
\begin{code}%
\>[0]\AgdaOperator{\AgdaFunction{\AgdaUnderscore{}ᵗ}}\AgdaSpace{}%
\AgdaSymbol{:}\AgdaSpace{}%
\AgdaSymbol{(}\AgdaPostulate{C}\AgdaSpace{}%
\AgdaSymbol{→}\AgdaSpace{}%
\AgdaPostulate{Y}\AgdaSpace{}%
\AgdaSymbol{→}\AgdaSpace{}%
\AgdaDatatype{Word}\AgdaSpace{}%
\AgdaPostulate{X}\AgdaSpace{}%
\AgdaOperator{\AgdaFunction{×}}\AgdaSpace{}%
\AgdaPostulate{C}\AgdaSymbol{)}\AgdaSpace{}%
\AgdaSymbol{→}\AgdaSpace{}%
\AgdaSymbol{(}\AgdaPostulate{C}\AgdaSpace{}%
\AgdaSymbol{→}\AgdaSpace{}%
\AgdaDatatype{Word}\AgdaSpace{}%
\AgdaPostulate{Y}\AgdaSpace{}%
\AgdaSymbol{→}\AgdaSpace{}%
\AgdaDatatype{Word}\AgdaSpace{}%
\AgdaPostulate{X}\AgdaSpace{}%
\AgdaOperator{\AgdaFunction{×}}\AgdaSpace{}%
\AgdaPostulate{C}\AgdaSymbol{)}\<%
\end{code}
\begin{code}[hide]%
\>[0]\AgdaOperator{\AgdaFunction{\AgdaUnderscore{}ᵗ}}\AgdaSpace{}%
\AgdaSymbol{=}\AgdaSpace{}%
\AgdaOperator{\AgdaFunction{WB.\AgdaUnderscore{}ᵗ}}\<%
\end{code}
\end{column}
\begin{column}{0.41\textwidth}
\centering
\begin{tikzpicture}[level distance=9mm,
  level 1/.style={sibling distance=22mm},
  level 2/.style={sibling distance=10mm},
  every node/.style={font=\small},
  op/.style={circle,draw=accent,fill=soft,inner sep=1pt,minimum size=5mm},
  lf/.style={rectangle,draw=paperC,fill=softpaper,inner sep=2pt,rounded corners=1pt}]
  \node[op] {$\bullet$}
    child {node[op] {$\bullet$}
      child {node[lf] {$H$}}
      child {node[lf] {$S$}}}
    child {node[op] {$\bullet$}
      child {node[lf] {$H$}}
      child {node[lf,draw=gray,fill=white] {$\varepsilon$}}};
\end{tikzpicture}

{\scriptsize $(H \aidop{•} S) \aidop{•} (H \aidop{•} \varepsilon)$ --- a tree;\\
\aid{by-assoc} flattens both sides to lists and compares by \aid{refl}}

\vspace{8pt}
\begin{tikzpicture}[>={Stealth[length=1.8mm]}]
  \node[draw=accent,fill=soft,rounded corners=2pt,font=\scriptsize] (c) {$c$};
  \node[draw=accent,fill=soft,rounded corners=2pt,font=\scriptsize,right=12mm of c] (c1) {$c_1$};
  \node[draw=accent,fill=soft,rounded corners=2pt,font=\scriptsize,right=12mm of c1] (c2) {$c_2$};
  \draw[->] (c) -- node[above,font=\scriptsize]{$y_1$} node[below,font=\scriptsize]{emit $w_1$} (c1);
  \draw[->] (c1) -- node[above,font=\scriptsize]{$y_2$} node[below,font=\scriptsize]{emit $w_2$} (c2);
\end{tikzpicture}

{\scriptsize $(h\,\aid{ᵗ})\;c\;(y_1\aidop{•}y_2) = (w_1\aidop{•}w_2,\;c_2)$}
\end{column}
\end{columns}
\end{frame}

\begin{frame}{Presented monoids as inductive congruences \InPaper}
\begin{columns}[c]
\begin{column}{0.52\textwidth}
One convention everywhere: \aid{\_===\_} = the raw axioms $\Gamma$;
\aid{\_≈\_} = the monoid congruence they generate:
\begin{code}%
\>[0]\AgdaKeyword{data}\AgdaSpace{}%
\AgdaOperator{\AgdaDatatype{\AgdaUnderscore{}≈\AgdaUnderscore{}}}\AgdaSpace{}%
\AgdaSymbol{:}\AgdaSpace{}%
\AgdaFunction{WRel}\AgdaSpace{}%
\AgdaPostulate{X}\AgdaSpace{}%
\AgdaKeyword{where}\<%
\\
\>[0][@{}l@{\AgdaIndent{0}}]%
\>[2]\AgdaInductiveConstructor{refl}%
\>[8]\AgdaSymbol{:}\AgdaSpace{}%
\AgdaPostulate{w}\AgdaSpace{}%
\AgdaOperator{\AgdaDatatype{≈}}\AgdaSpace{}%
\AgdaPostulate{w}\<%
\\
%
\>[2]\AgdaInductiveConstructor{sym}%
\>[8]\AgdaSymbol{:}\AgdaSpace{}%
\AgdaPostulate{w}\AgdaSpace{}%
\AgdaOperator{\AgdaDatatype{≈}}\AgdaSpace{}%
\AgdaPostulate{v}\AgdaSpace{}%
\AgdaSymbol{→}\AgdaSpace{}%
\AgdaPostulate{v}\AgdaSpace{}%
\AgdaOperator{\AgdaDatatype{≈}}\AgdaSpace{}%
\AgdaPostulate{w}\<%
\\
%
\>[2]\AgdaInductiveConstructor{trans}\AgdaSpace{}%
\AgdaSymbol{:}\AgdaSpace{}%
\AgdaPostulate{w}\AgdaSpace{}%
\AgdaOperator{\AgdaDatatype{≈}}\AgdaSpace{}%
\AgdaPostulate{v}\AgdaSpace{}%
\AgdaSymbol{→}\AgdaSpace{}%
\AgdaPostulate{v}\AgdaSpace{}%
\AgdaOperator{\AgdaDatatype{≈}}\AgdaSpace{}%
\AgdaPostulate{u}\AgdaSpace{}%
\AgdaSymbol{→}\AgdaSpace{}%
\AgdaPostulate{w}\AgdaSpace{}%
\AgdaOperator{\AgdaDatatype{≈}}\AgdaSpace{}%
\AgdaPostulate{u}\<%
\\
%
\>[2]\AgdaInductiveConstructor{cong}%
\>[8]\AgdaSymbol{:}\AgdaSpace{}%
\AgdaPostulate{w}\AgdaSpace{}%
\AgdaOperator{\AgdaDatatype{≈}}\AgdaSpace{}%
\AgdaPostulate{w'}\AgdaSpace{}%
\AgdaSymbol{→}\AgdaSpace{}%
\AgdaPostulate{v}\AgdaSpace{}%
\AgdaOperator{\AgdaDatatype{≈}}\AgdaSpace{}%
\AgdaPostulate{v'}\AgdaSpace{}%
\AgdaSymbol{→}\AgdaSpace{}%
\AgdaPostulate{w}\AgdaSpace{}%
\AgdaOperator{\AgdaInductiveConstructor{•}}\AgdaSpace{}%
\AgdaPostulate{v}\AgdaSpace{}%
\AgdaOperator{\AgdaDatatype{≈}}\AgdaSpace{}%
\AgdaPostulate{w'}\AgdaSpace{}%
\AgdaOperator{\AgdaInductiveConstructor{•}}\AgdaSpace{}%
\AgdaPostulate{v'}\<%
\\
%
\>[2]\AgdaInductiveConstructor{assoc}%
\>[13]\AgdaSymbol{:}\AgdaSpace{}%
\AgdaSymbol{(}\AgdaPostulate{w}\AgdaSpace{}%
\AgdaOperator{\AgdaInductiveConstructor{•}}\AgdaSpace{}%
\AgdaPostulate{v}\AgdaSymbol{)}\AgdaSpace{}%
\AgdaOperator{\AgdaInductiveConstructor{•}}\AgdaSpace{}%
\AgdaPostulate{u}\AgdaSpace{}%
\AgdaOperator{\AgdaDatatype{≈}}\AgdaSpace{}%
\AgdaPostulate{w}\AgdaSpace{}%
\AgdaOperator{\AgdaInductiveConstructor{•}}\AgdaSpace{}%
\AgdaSymbol{(}\AgdaPostulate{v}\AgdaSpace{}%
\AgdaOperator{\AgdaInductiveConstructor{•}}\AgdaSpace{}%
\AgdaPostulate{u}\AgdaSymbol{)}\<%
\\
%
\>[2]\AgdaInductiveConstructor{left-unit}%
\>[13]\AgdaSymbol{:}\AgdaSpace{}%
\AgdaInductiveConstructor{ε}\AgdaSpace{}%
\AgdaOperator{\AgdaInductiveConstructor{•}}\AgdaSpace{}%
\AgdaPostulate{w}\AgdaSpace{}%
\AgdaOperator{\AgdaDatatype{≈}}\AgdaSpace{}%
\AgdaPostulate{w}\<%
\\
%
\>[2]\AgdaInductiveConstructor{right-unit}\AgdaSpace{}%
\AgdaSymbol{:}\AgdaSpace{}%
\AgdaPostulate{w}\AgdaSpace{}%
\AgdaOperator{\AgdaInductiveConstructor{•}}\AgdaSpace{}%
\AgdaInductiveConstructor{ε}\AgdaSpace{}%
\AgdaOperator{\AgdaDatatype{≈}}\AgdaSpace{}%
\AgdaPostulate{w}\<%
\\
%
\>[2]\AgdaInductiveConstructor{axiom}%
\>[13]\AgdaSymbol{:}\AgdaSpace{}%
\AgdaPostulate{w}\AgdaSpace{}%
\AgdaOperator{\AgdaPostulate{===}}\AgdaSpace{}%
\AgdaPostulate{v}\AgdaSpace{}%
\AgdaSymbol{→}\AgdaSpace{}%
\AgdaPostulate{w}\AgdaSpace{}%
\AgdaOperator{\AgdaDatatype{≈}}\AgdaSpace{}%
\AgdaPostulate{v}\<%
\end{code}
\end{column}
\begin{column}{0.45\textwidth}
\begin{itemize}
  \setlength{\itemsep}{4pt}
  \item $\approx$ \emph{is} the word problem of $\langle X\mid\Gamma\rangle$.
  \item A derivation is a \alert{first-class object}: compute with it,
        induct on it, translate it (congruence-lifting lemmas are inductions
        on this type).
  \item Packaged as the \emph{word setoid} and the presented monoid
        \aid{•-ε-monoid}; with a \aid{Grouplike} witness (a left inverse per
        generator) the group \aid{•-ε-group}.
\end{itemize}
\vspace{2pt}
{\small Contrast: mathlib's \aid{PresentedGroup} (quotient), cubical Agda (HIT).}
\end{column}
\end{columns}
\end{frame}

\begin{frame}{What ``$\Gamma$ presents $G$'' means \InPaper}
\begin{code}%
\>[0]\AgdaKeyword{record}\AgdaSpace{}%
\AgdaOperator{\AgdaRecord{\AgdaUnderscore{}IsPresentationOf\AgdaUnderscore{}}}\AgdaSpace{}%
\AgdaSymbol{(}\AgdaOperator{\AgdaBound{\AgdaUnderscore{}===\AgdaUnderscore{}}}\AgdaSpace{}%
\AgdaSymbol{:}\AgdaSpace{}%
\AgdaFunction{WRel}\AgdaSpace{}%
\AgdaPostulate{X}\AgdaSymbol{)}\AgdaSpace{}%
\AgdaSymbol{(}\AgdaBound{G}\AgdaSpace{}%
\AgdaSymbol{:}\AgdaSpace{}%
\AgdaRecord{Group}\AgdaSpace{}%
\AgdaPostulate{a}\AgdaSpace{}%
\AgdaPostulate{ℓ}\AgdaSymbol{)}\AgdaSpace{}%
\AgdaSymbol{:}\AgdaSpace{}%
\AgdaPrimitive{Set}\AgdaSpace{}%
\AgdaSymbol{(}\AgdaPostulate{a}\AgdaSpace{}%
\AgdaOperator{\AgdaPrimitive{⊔}}\AgdaSpace{}%
\AgdaPostulate{ℓ}\AgdaSymbol{)}\AgdaSpace{}%
\AgdaKeyword{where}\<%
\\
\>[0][@{}l@{\AgdaIndent{0}}]%
\>[2]\AgdaKeyword{field}\AgdaSpace{}%
\AgdaField{gl}\AgdaSpace{}%
\AgdaSymbol{:}\AgdaSpace{}%
\AgdaFunction{Grouplike}\AgdaSpace{}%
\AgdaOperator{\AgdaBound{\AgdaUnderscore{}===\AgdaUnderscore{}}}\<%
\\
%
\>[2]\AgdaKeyword{module}\AgdaSpace{}%
\AgdaModule{GL}\AgdaSpace{}%
\AgdaSymbol{=}\AgdaSpace{}%
\AgdaModule{Group-Lemmas}\AgdaSpace{}%
\AgdaOperator{\AgdaBound{\AgdaUnderscore{}===\AgdaUnderscore{}}}\AgdaSpace{}%
\AgdaField{gl}\<%
\\
%
\>[2]\AgdaKeyword{open}\AgdaSpace{}%
\AgdaModule{GroupMorphisms}\AgdaSpace{}%
\AgdaSymbol{(}\AgdaFunction{Group.rawGroup}\AgdaSpace{}%
\AgdaPostulate{GL.•-ε-group}\AgdaSymbol{)}\AgdaSpace{}%
\AgdaSymbol{(}\AgdaFunction{Group.rawGroup}\AgdaSpace{}%
\AgdaBound{G}\AgdaSymbol{)}\<%
\\
%
\>[2]\AgdaKeyword{field}%
\>[302I]\AgdaOperator{\AgdaField{⟦\AgdaUnderscore{}⟧}}\AgdaSpace{}%
\AgdaSymbol{:}\AgdaSpace{}%
\AgdaDatatype{Word}\AgdaSpace{}%
\AgdaPostulate{X}\AgdaSpace{}%
\AgdaSymbol{→}\AgdaSpace{}%
\AgdaField{Group.Carrier}\AgdaSpace{}%
\AgdaBound{G}\<%
\\
\>[.][@{}l@{}]\<[302I]%
\>[8]\AgdaField{iso}\AgdaSpace{}%
\AgdaSymbol{:}\AgdaSpace{}%
\AgdaPostulate{IsGroupIsomorphism}\AgdaSpace{}%
\AgdaOperator{\AgdaField{⟦\AgdaUnderscore{}⟧}}\<%
\end{code}
\begin{code}[hide]%
\>[0]\AgdaKeyword{postulate}\<%
\\
\>[0][@{}l@{\AgdaIndent{0}}]%
\>[2]\AgdaPostulate{|NF|}%
\>[13]\AgdaSymbol{:}\AgdaSpace{}%
\AgdaPrimitive{Set}\<%
\\
%
\>[2]\AgdaPostulate{nf}%
\>[13]\AgdaSymbol{:}\AgdaSpace{}%
\AgdaDatatype{Word}\AgdaSpace{}%
\AgdaPostulate{X}\AgdaSpace{}%
\AgdaSymbol{→}\AgdaSpace{}%
\AgdaPostulate{|NF|}\<%
\\
%
\>[2]\AgdaOperator{\AgdaPostulate{\AgdaUnderscore{}≈ₙ\AgdaUnderscore{}}}%
\>[13]\AgdaSymbol{:}\AgdaSpace{}%
\AgdaPostulate{|NF|}\AgdaSpace{}%
\AgdaSymbol{→}\AgdaSpace{}%
\AgdaPostulate{|NF|}\AgdaSpace{}%
\AgdaSymbol{→}\AgdaSpace{}%
\AgdaPrimitive{Set}\<%
\\
%
\>[2]\AgdaPostulate{NormalForm}\AgdaSpace{}%
\AgdaSymbol{:}\AgdaSpace{}%
\AgdaPrimitive{Set}\<%
\\
%
\>[2]\AgdaPostulate{inv-nf}%
\>[13]\AgdaSymbol{:}\AgdaSpace{}%
\AgdaPostulate{|NF|}\AgdaSpace{}%
\AgdaSymbol{→}\AgdaSpace{}%
\AgdaDatatype{Word}\AgdaSpace{}%
\AgdaPostulate{X}\<%
\\
%
\>[2]\AgdaPostulate{Sem}%
\>[13]\AgdaSymbol{:}\AgdaSpace{}%
\AgdaPrimitive{Set}\<%
\\
%
\>[2]\AgdaOperator{\AgdaPostulate{⟦\AgdaUnderscore{}⟧}}%
\>[13]\AgdaSymbol{:}\AgdaSpace{}%
\AgdaDatatype{Word}\AgdaSpace{}%
\AgdaPostulate{X}\AgdaSpace{}%
\AgdaSymbol{→}\AgdaSpace{}%
\AgdaPostulate{Sem}\<%
\\
%
\>[2]\AgdaOperator{\AgdaPostulate{\AgdaUnderscore{}≈₂\AgdaUnderscore{}}}%
\>[13]\AgdaSymbol{:}\AgdaSpace{}%
\AgdaPostulate{Sem}\AgdaSpace{}%
\AgdaSymbol{→}\AgdaSpace{}%
\AgdaPostulate{Sem}\AgdaSpace{}%
\AgdaSymbol{→}\AgdaSpace{}%
\AgdaPrimitive{Set}\<%
\end{code}
\begin{columns}[c]
\begin{column}{0.54\textwidth}
\small
\begin{itemize}
  \item weaker records: \aid{\_IsSubPresentationOf\_} (monomorphism), monoid versions;
  \item upgrade lemmas: surjective sub-presentation $\Rightarrow$ presentation;
        grouplike monoid presentation $\Rightarrow$ group presentation;
  \item proofs split into soundness, injectivity (by normalization),
        surjectivity; the records assemble the isomorphism.
\end{itemize}
\end{column}
\begin{column}{0.43\textwidth}
\begin{keybox}[title=Monoid vs group: a real distinction]
\small Words have no formal inverses, so \emph{every} presentation is at
bottom a monoid presentation. $\langle\, T \mid T^{0} = \varepsilon\,\rangle$
read as a group presentation is $\mathbb{Z}$; as a \alert{monoid} it is
$(\mathbb{N},+,0)$ --- and \aid{Grouplike} is exactly what fails.
\end{keybox}
\end{column}
\end{columns}
\end{frame}

\begin{frame}{Normal forms are standard-library bundles \InPaper}
\begin{columns}[c]
\begin{column}{0.50\textwidth}
\small From the \alert{word setoid} to a carrier $\mathit{NF}$, on the nose:\\[4pt]
\begin{tabular}{@{\quad}l@{\;\;=\;\;}l@{}}
\aid{NormalFormInjective} & \aid{Injection} \\
\aid{NormalForm}          & \aid{RightInverse} \\
\aid{BijectiveNormalForm} & \aid{Bijection}
\end{tabular}
\vspace{6pt}


\aid{Injection}: $w\approx v \iff \aid{nf}\,w \approx_{\mathit{NF}} \aid{nf}\,v$, and with it
\begin{code}%
\>[0]\AgdaFunction{by-equal-nf}\AgdaSpace{}%
\AgdaSymbol{:}\AgdaSpace{}%
\AgdaPostulate{nf}\AgdaSpace{}%
\AgdaPostulate{w}\AgdaSpace{}%
\AgdaOperator{\AgdaPostulate{≈ₙ}}\AgdaSpace{}%
\AgdaPostulate{nf}\AgdaSpace{}%
\AgdaPostulate{v}\AgdaSpace{}%
\AgdaSymbol{→}\AgdaSpace{}%
\AgdaPostulate{w}\AgdaSpace{}%
\AgdaOperator{\AgdaDatatype{≈}}\AgdaSpace{}%
\AgdaPostulate{v}\<%
\\
\>[0]\AgdaFunction{≈-dec}\AgdaSpace{}%
\AgdaSymbol{:}\AgdaSpace{}%
\AgdaFunction{Decidable}\AgdaSpace{}%
\AgdaOperator{\AgdaPostulate{\AgdaUnderscore{}≈ₙ\AgdaUnderscore{}}}\AgdaSpace{}%
\AgdaSymbol{→}\AgdaSpace{}%
\AgdaFunction{Decidable}\AgdaSpace{}%
\AgdaOperator{\AgdaDatatype{\AgdaUnderscore{}≈\AgdaUnderscore{}}}\<%
\end{code}
\begin{code}[hide]%
\>[0]\AgdaFunction{by-equal-nf}\AgdaSpace{}%
\AgdaSymbol{=}\AgdaSpace{}%
\AgdaPostulate{proof-elided}\<%
\\
\>[0]\AgdaFunction{≈-dec}\AgdaSpace{}%
\AgdaSymbol{=}\AgdaSpace{}%
\AgdaPostulate{proof-elided}\<%
\end{code}
{\small prove word equations by \alert{computing}; a decidable carrier
\alert{decides the word problem}.}
\end{column}
\begin{column}{0.47\textwidth}
\centering
\begin{tikzpicture}[>={Stealth[length=2mm]}]
  \node[draw=accent,fill=soft,ellipse,minimum width=26mm,minimum height=34mm] (W) at (0,0) {};
  \node[draw=paperC,fill=softpaper,ellipse,minimum width=18mm,minimum height=28mm] (N) at (4,0) {};
  \node[font=\scriptsize\bfseries,accent] at (0,1.95) {word setoid};
  \node[font=\scriptsize\bfseries,paperC] at (4,1.95) {$\mathit{NF}$};
  \foreach \y in {0.9,0.35,-0.2} \node[circle,fill=accent,inner sep=1pt] at (-0.4,\y) {};
  \draw[accent!60,rounded corners] (-0.8,-0.45) rectangle (0.1,1.15);
  \node[circle,fill=paperC,inner sep=1.3pt] (u) at (4,0.45) {};
  \draw[->,thick] (0.15,0.45) to[bend left=10] node[above,font=\scriptsize\agdafont]{nf} (u);
  \node[circle,fill=newC,inner sep=1.3pt] (r) at (-0.4,-1.0) {};
  \draw[->,thick,newC] (u) to[bend left=25] node[below right,font=\scriptsize\agdafont]{inv-nf} (r);
  \node[font=\tiny,newC,align=center] at (-0.4,-1.35) {canonical\\representative};
  \node[circle,fill=gray,inner sep=1pt] at (4.3,-1.0) {};
  \node[font=\tiny,gray,align=center] at (4.3,-1.35) {junk allowed};
\end{tikzpicture}

{\small \aid{RightInverse} adds the \alert{section} \aid{inv-nf} with
$\aid{inv-nf}\circ\aid{nf}\approx\mathrm{id}$. Exactness
$\aid{nf}\circ\aid{inv-nf}\approx\mathrm{id}$ is kept separate.}
\end{column}
\end{columns}
\end{frame}

\begin{frame}{Where constructivity bites: the section is data \InPaper}
\begin{columns}[c]
\begin{column}{0.55\textwidth}
\begin{itemize}
  \setlength{\itemsep}{5pt}
  \item Classically, an injective \aid{nf} has a section on its image for
        free: \emph{choose} a preimage.
  \item Constructively, choosing is not free: an \aid{Injection} gives
        \alert{no algorithm} producing a word that realizes a normal form.
  \item That algorithm is exactly the extra field of \aid{RightInverse} ---
        and it is what completeness consumes.
  \item So the two are \alert{different strengths} in Agda, though
        classically equal. Use the weak one for word-problem reasoning, the
        strong one for completeness and presentations.
  \item A fourth, crude witness, \aid{WeakNormalForm}: an injective
        \emph{invariant}, not even congruent.
\end{itemize}
\end{column}
\begin{column}{0.42\textwidth}
\centering
\begin{tikzpicture}[x=12mm,y=10mm]
  \node[w=gray] (weak) at (0,0) {\aid{WeakNormalForm}};
  \node[w=accent] (inj) at (0,1.3) {\aid{NormalFormInjective}};
  \node[w=paperC] (ri) at (0,2.6) {\aid{NormalForm}\\[-2pt]{\scriptsize + section (data)}};
  \node[w=newC] (bij) at (0,3.9) {\aid{BijectiveNormalForm}\\[-2pt]{\scriptsize surjectivity as a property}};
  \draw[-{Stealth},thick] (ri) -- (inj);
  \draw[-{Stealth},thick] (inj) -- (weak);
  \draw[-{Stealth},thick] (bij) -- (ri);
  \node[font=\scriptsize,anchor=west,align=left] at (1.55,1.95) {strictly stronger\\as \emph{data}};
\end{tikzpicture}
\end{column}
\end{columns}
\end{frame}

\begin{frame}{Completeness, once and for all \InPaper}
A \aid{NormalForm} is \alert{unique} for a semantics \aid{⟦\_⟧} into any setoid
when canonical representatives with equal denotations are equal:
\begin{code}%
\>[0]\AgdaKeyword{record}\AgdaSpace{}%
\AgdaRecord{UniqueNormalForm}\AgdaSpace{}%
\AgdaSymbol{(}\AgdaBound{normalForm}\AgdaSpace{}%
\AgdaSymbol{:}\AgdaSpace{}%
\AgdaPostulate{NormalForm}\AgdaSymbol{)}\AgdaSpace{}%
\AgdaSymbol{:}\AgdaSpace{}%
\AgdaPrimitive{Set}\AgdaSpace{}%
\AgdaSymbol{(}\AgdaPostulate{c}\AgdaSpace{}%
\AgdaOperator{\AgdaPrimitive{⊔}}\AgdaSpace{}%
\AgdaPostulate{d}\AgdaSymbol{)}\AgdaSpace{}%
\AgdaKeyword{where}\<%
\\
\>[0][@{}l@{\AgdaIndent{0}}]%
\>[2]\AgdaKeyword{field}\AgdaSpace{}%
\AgdaField{unique}\AgdaSpace{}%
\AgdaSymbol{:}\AgdaSpace{}%
\AgdaSymbol{∀}\AgdaSpace{}%
\AgdaSymbol{\{}\AgdaBound{u}\AgdaSpace{}%
\AgdaBound{v}\AgdaSpace{}%
\AgdaSymbol{:}\AgdaSpace{}%
\AgdaPostulate{|NF|}\AgdaSymbol{\}}\AgdaSpace{}%
\AgdaSymbol{→}\AgdaSpace{}%
\AgdaOperator{\AgdaPostulate{⟦}}\AgdaSpace{}%
\AgdaPostulate{inv-nf}\AgdaSpace{}%
\AgdaBound{u}\AgdaSpace{}%
\AgdaOperator{\AgdaPostulate{⟧}}\AgdaSpace{}%
\AgdaOperator{\AgdaPostulate{≈₂}}\AgdaSpace{}%
\AgdaOperator{\AgdaPostulate{⟦}}\AgdaSpace{}%
\AgdaPostulate{inv-nf}\AgdaSpace{}%
\AgdaBound{v}\AgdaSpace{}%
\AgdaOperator{\AgdaPostulate{⟧}}\AgdaSpace{}%
\AgdaSymbol{→}\AgdaSpace{}%
\AgdaBound{u}\AgdaSpace{}%
\AgdaOperator{\AgdaPostulate{≈ₙ}}\AgdaSpace{}%
\AgdaBound{v}\<%
\\
\>[0]\AgdaFunction{by-normalization}\AgdaSpace{}%
\AgdaSymbol{:}\AgdaSpace{}%
\AgdaSymbol{\{}\AgdaBound{normalForm}\AgdaSpace{}%
\AgdaSymbol{:}\AgdaSpace{}%
\AgdaPostulate{NormalForm}\AgdaSymbol{\}}\AgdaSpace{}%
\AgdaSymbol{→}\AgdaSpace{}%
\AgdaRecord{UniqueNormalForm}\AgdaSpace{}%
\AgdaBound{normalForm}\AgdaSpace{}%
\AgdaSymbol{→}\<%
\\
\>[0][@{}l@{\AgdaIndent{0}}]%
\>[2]\AgdaFunction{Congruent}\AgdaSpace{}%
\AgdaOperator{\AgdaDatatype{\AgdaUnderscore{}≈\AgdaUnderscore{}}}\AgdaSpace{}%
\AgdaOperator{\AgdaPostulate{\AgdaUnderscore{}≈₂\AgdaUnderscore{}}}\AgdaSpace{}%
\AgdaOperator{\AgdaPostulate{⟦\AgdaUnderscore{}⟧}}\AgdaSpace{}%
\AgdaSymbol{→}\AgdaSpace{}%
\AgdaFunction{Injective}\AgdaSpace{}%
\AgdaOperator{\AgdaDatatype{\AgdaUnderscore{}≈\AgdaUnderscore{}}}\AgdaSpace{}%
\AgdaOperator{\AgdaPostulate{\AgdaUnderscore{}≈₂\AgdaUnderscore{}}}\AgdaSpace{}%
\AgdaOperator{\AgdaPostulate{⟦\AgdaUnderscore{}⟧}}\<%
\end{code}
\begin{code}[hide]%
\>[0]\AgdaFunction{by-normalization}\AgdaSpace{}%
\AgdaSymbol{=}\AgdaSpace{}%
\AgdaPostulate{proof-elided}\<%
\end{code}
\begin{columns}[c]
\begin{column}{0.55\textwidth}
\small
\textbf{Soundness + unique NF $\Rightarrow$ completeness.} Five lines of proof;
its value: every case study only ever proves \aid{unique} --- a statement about
\alert{normal forms}, i.e.\ concrete finite data, never about arbitrary words.
\end{column}
\begin{column}{0.42\textwidth}
\small
\begin{itemize}
  \item \aid{by-normalization-wsurj}: surjectivity on canonical forms suffices;
  \item \aid{by-completeness}: the converse, given an exact section --- how
        uniqueness lifts through products.
\end{itemize}
\end{column}
\end{columns}
\end{frame}

\begin{frame}{Inductively defined circuits: indices for the unifier \InPaper}
\begin{columns}[c]
\begin{column}{0.52\textwidth}
\aid{\_+\_} recurses on its \emph{first} argument:
$2 + n \;\leadsto\; \mathsf{suc}\,(\mathsf{suc}\,n)$, while $n + 2$ is stuck.

\vspace{4pt}
So the indices are \alert{left-biased} too:
\begin{itemize}
  \item \aid{gate₂ : Gate 2 → Gen (₂₊ n)}, not \aid{Gen k} with $k\ge 2$;
  \item $k$-fold shift: \aid{Gen n → Gen (k + n)}.
\end{itemize}
Every index equation the coverage checker meets is solved by
\alert{injectivity of \aid{₁₊}}, never by arithmetic. That is why coset tables
are written by direct pattern matching.

\vspace{4pt}
\small \aid{lemma-cong↑} lifts the \emph{entire} congruence one wire up, by
induction over the eight constructors of $\approx$.
\end{column}
\begin{column}{0.45\textwidth}
\centering
\begin{tikzpicture}[
  ok/.style={draw=paperC,fill=softpaper,rounded corners=2pt,font=\small,align=left,text width=50mm,inner sep=4pt},
  bad/.style={draw=newC,fill=softnew,rounded corners=2pt,font=\small,align=left,text width=50mm,inner sep=4pt}]
  \node[ok] (a) {\textbf{left-biased} \hfill\checkmark\\[2pt]
   \aid{Gen (₂₊ n)} $\doteq$ \aid{Gen (₁₊ m)}\\
   $\leadsto\ m := \aid{₁₊}\,n$ \quad by unification};
  \node[bad,below=6mm of a] (b) {\textbf{right-biased} \hfill$\times$\\[2pt]
   \aid{Gen (n + 2)} $\doteq$ \aid{Gen (₁₊ m)}\\
   $\leadsto\ n + 2 \doteq \mathsf{suc}\,m$ \quad stuck:\\ a Diophantine equation};
\end{tikzpicture}
\end{column}
\end{columns}
\end{frame}

\begin{frame}{The Reidemeister--Schreier engine \InPaper}
\begin{code}%
\>[0]\AgdaKeyword{module}\AgdaSpace{}%
\AgdaModule{SingleLevel}\AgdaSpace{}%
\AgdaSymbol{(}\AgdaBound{Γ}\AgdaSpace{}%
\AgdaSymbol{:}\AgdaSpace{}%
\AgdaFunction{WRel}\AgdaSpace{}%
\AgdaPostulate{X}\AgdaSymbol{)}\AgdaSpace{}%
\AgdaSymbol{(}\AgdaBound{Δ}\AgdaSpace{}%
\AgdaSymbol{:}\AgdaSpace{}%
\AgdaFunction{WRel}\AgdaSpace{}%
\AgdaPostulate{Y}\AgdaSymbol{)}\AgdaSpace{}%
\AgdaSymbol{(}\AgdaBound{C}\AgdaSpace{}%
\AgdaSymbol{:}\AgdaSpace{}%
\AgdaPrimitive{Set}\AgdaSymbol{)}\AgdaSpace{}%
\AgdaSymbol{(}\AgdaBound{I}\AgdaSpace{}%
\AgdaSymbol{:}\AgdaSpace{}%
\AgdaBound{C}\AgdaSymbol{)}\<%
\\
\>[0][@{}l@{\AgdaIndent{0}}]%
\>[2]\AgdaSymbol{(}\AgdaBound{f}\AgdaSpace{}%
\AgdaSymbol{:}\AgdaSpace{}%
\AgdaPostulate{X}\AgdaSpace{}%
\AgdaSymbol{→}\AgdaSpace{}%
\AgdaDatatype{Word}\AgdaSpace{}%
\AgdaPostulate{Y}\AgdaSymbol{)}\AgdaSpace{}%
\AgdaSymbol{(}\AgdaBound{h}\AgdaSpace{}%
\AgdaSymbol{:}\AgdaSpace{}%
\AgdaBound{C}\AgdaSpace{}%
\AgdaSymbol{→}\AgdaSpace{}%
\AgdaPostulate{Y}\AgdaSpace{}%
\AgdaSymbol{→}\AgdaSpace{}%
\AgdaDatatype{Word}\AgdaSpace{}%
\AgdaPostulate{X}\AgdaSpace{}%
\AgdaOperator{\AgdaFunction{×}}\AgdaSpace{}%
\AgdaBound{C}\AgdaSymbol{)}\AgdaSpace{}%
\AgdaSymbol{(}\AgdaOperator{\AgdaBound{[\AgdaUnderscore{}]}}\AgdaSpace{}%
\AgdaSymbol{:}\AgdaSpace{}%
\AgdaBound{C}\AgdaSpace{}%
\AgdaSymbol{→}\AgdaSpace{}%
\AgdaDatatype{Word}\AgdaSpace{}%
\AgdaPostulate{Y}\AgdaSymbol{)}\<%
\end{code}
\begin{code}[hide]%
%
\>[2]\AgdaKeyword{where}\<%
\\
%
\\[\AgdaEmptyExtraSkip]%
\>[0]\AgdaKeyword{postulate}\<%
\\
\>[0][@{}l@{\AgdaIndent{0}}]%
\>[2]\AgdaPostulate{I}%
\>[9]\AgdaSymbol{:}\AgdaSpace{}%
\AgdaPostulate{C}\<%
\\
%
\>[2]\AgdaPostulate{f}%
\>[9]\AgdaSymbol{:}\AgdaSpace{}%
\AgdaPostulate{X}\AgdaSpace{}%
\AgdaSymbol{→}\AgdaSpace{}%
\AgdaDatatype{Word}\AgdaSpace{}%
\AgdaPostulate{Y}\<%
\\
%
\>[2]\AgdaPostulate{h}%
\>[9]\AgdaSymbol{:}\AgdaSpace{}%
\AgdaPostulate{C}\AgdaSpace{}%
\AgdaSymbol{→}\AgdaSpace{}%
\AgdaPostulate{Y}\AgdaSpace{}%
\AgdaSymbol{→}\AgdaSpace{}%
\AgdaDatatype{Word}\AgdaSpace{}%
\AgdaPostulate{X}\AgdaSpace{}%
\AgdaOperator{\AgdaFunction{×}}\AgdaSpace{}%
\AgdaPostulate{C}\<%
\\
%
\>[2]\AgdaOperator{\AgdaPostulate{[\AgdaUnderscore{}]}}%
\>[9]\AgdaSymbol{:}\AgdaSpace{}%
\AgdaPostulate{C}\AgdaSpace{}%
\AgdaSymbol{→}\AgdaSpace{}%
\AgdaDatatype{Word}\AgdaSpace{}%
\AgdaPostulate{Y}\<%
\\
%
\>[2]\AgdaOperator{\AgdaPostulate{\AgdaUnderscore{}\textasciitilde{}\AgdaUnderscore{}}}%
\>[9]\AgdaSymbol{:}\AgdaSpace{}%
\AgdaDatatype{Word}\AgdaSpace{}%
\AgdaPostulate{X}\AgdaSpace{}%
\AgdaOperator{\AgdaFunction{×}}\AgdaSpace{}%
\AgdaPostulate{C}\AgdaSpace{}%
\AgdaSymbol{→}\AgdaSpace{}%
\AgdaDatatype{Word}\AgdaSpace{}%
\AgdaPostulate{X}\AgdaSpace{}%
\AgdaOperator{\AgdaFunction{×}}\AgdaSpace{}%
\AgdaPostulate{C}\AgdaSpace{}%
\AgdaSymbol{→}\AgdaSpace{}%
\AgdaPrimitive{Set}\<%
\\
%
\>[2]\AgdaOperator{\AgdaPostulate{\AgdaUnderscore{}===₁\AgdaUnderscore{}}}\AgdaSpace{}%
\AgdaSymbol{:}\AgdaSpace{}%
\AgdaFunction{WRel}\AgdaSpace{}%
\AgdaPostulate{X}\<%
\\
%
\>[2]\AgdaOperator{\AgdaPostulate{\AgdaUnderscore{}===₂\AgdaUnderscore{}}}\AgdaSpace{}%
\AgdaSymbol{:}\AgdaSpace{}%
\AgdaFunction{WRel}\AgdaSpace{}%
\AgdaPostulate{Y}\<%
\end{code}
\vspace{-2pt}
{\small $\Gamma$ presents a subgroup $H$ of the group $G$ presented by $\Delta$;
$C$ = cosets, $I$ = identity coset, $f$ = embedding of generators,
$[\_]$ = Schreier section, $h$ = the coset table.}
\begin{code}%
\>[0]\AgdaKeyword{module}\AgdaSpace{}%
\AgdaModule{Transfer}\<%
\\
\>[0][@{}l@{\AgdaIndent{0}}]%
\>[2]\AgdaSymbol{(}\AgdaBound{h=⁻¹f-gen}\AgdaSpace{}%
\AgdaSymbol{:}\AgdaSpace{}%
\AgdaSymbol{∀}\AgdaSpace{}%
\AgdaSymbol{(}\AgdaBound{x}\AgdaSpace{}%
\AgdaSymbol{:}\AgdaSpace{}%
\AgdaPostulate{X}\AgdaSymbol{)}\AgdaSpace{}%
\AgdaSymbol{→}\AgdaSpace{}%
\AgdaSymbol{(}\AgdaOperator{\AgdaInductiveConstructor{[}}\AgdaSpace{}%
\AgdaBound{x}\AgdaSpace{}%
\AgdaOperator{\AgdaInductiveConstructor{]ʷ}}\AgdaSpace{}%
\AgdaOperator{\AgdaInductiveConstructor{,}}\AgdaSpace{}%
\AgdaPostulate{I}\AgdaSymbol{)}\AgdaSpace{}%
\AgdaOperator{\AgdaPostulate{\textasciitilde{}}}\AgdaSpace{}%
\AgdaSymbol{((}\AgdaPostulate{h}\AgdaSpace{}%
\AgdaOperator{\AgdaFunction{ᵗ}}\AgdaSymbol{)}\AgdaSpace{}%
\AgdaPostulate{I}\AgdaSpace{}%
\AgdaSymbol{(}\AgdaPostulate{f}\AgdaSpace{}%
\AgdaBound{x}\AgdaSymbol{)))}%
\>[83]\AgdaComment{--\ 1}\<%
\\
%
\>[2]\AgdaSymbol{(}\AgdaBound{h-wd-ax}\AgdaSpace{}%
\AgdaSymbol{:}\AgdaSpace{}%
\AgdaSymbol{∀}\AgdaSpace{}%
\AgdaSymbol{(}\AgdaBound{c}\AgdaSpace{}%
\AgdaSymbol{:}\AgdaSpace{}%
\AgdaPostulate{C}\AgdaSymbol{)\{}\AgdaBound{u}\AgdaSpace{}%
\AgdaBound{t}\AgdaSpace{}%
\AgdaSymbol{:}\AgdaSpace{}%
\AgdaDatatype{Word}\AgdaSpace{}%
\AgdaPostulate{Y}\AgdaSymbol{\}}\AgdaSpace{}%
\AgdaSymbol{→}\AgdaSpace{}%
\AgdaBound{u}\AgdaSpace{}%
\AgdaOperator{\AgdaPostulate{===₂}}\AgdaSpace{}%
\AgdaBound{t}\AgdaSpace{}%
\AgdaSymbol{→}\AgdaSpace{}%
\AgdaSymbol{((}\AgdaPostulate{h}\AgdaSpace{}%
\AgdaOperator{\AgdaFunction{ᵗ}}\AgdaSymbol{)}\AgdaSpace{}%
\AgdaBound{c}\AgdaSpace{}%
\AgdaBound{u}\AgdaSymbol{)}\AgdaSpace{}%
\AgdaOperator{\AgdaPostulate{\textasciitilde{}}}\AgdaSpace{}%
\AgdaSymbol{((}\AgdaPostulate{h}\AgdaSpace{}%
\AgdaOperator{\AgdaFunction{ᵗ}}\AgdaSymbol{)}\AgdaSpace{}%
\AgdaBound{c}\AgdaSpace{}%
\AgdaBound{t}\AgdaSymbol{))}%
\>[83]\AgdaComment{--\ 2}\<%
\\
%
\>[2]\AgdaSymbol{(}\AgdaBound{f-wd-ax}\AgdaSpace{}%
\AgdaSymbol{:}\AgdaSpace{}%
\AgdaSymbol{∀}\AgdaSpace{}%
\AgdaSymbol{\{}\AgdaBound{w}\AgdaSpace{}%
\AgdaBound{v}\AgdaSymbol{\}}\AgdaSpace{}%
\AgdaSymbol{→}\AgdaSpace{}%
\AgdaBound{w}\AgdaSpace{}%
\AgdaOperator{\AgdaPostulate{===₁}}\AgdaSpace{}%
\AgdaBound{v}\AgdaSpace{}%
\AgdaSymbol{→}\AgdaSpace{}%
\AgdaSymbol{(}\AgdaPostulate{f}\AgdaSpace{}%
\AgdaOperator{\AgdaFunction{ʷ}}\AgdaSymbol{)}\AgdaSpace{}%
\AgdaBound{w}\AgdaSpace{}%
\AgdaOperator{\AgdaPostulate{≈₂}}\AgdaSpace{}%
\AgdaSymbol{(}\AgdaPostulate{f}\AgdaSpace{}%
\AgdaOperator{\AgdaFunction{ʷ}}\AgdaSymbol{)}\AgdaSpace{}%
\AgdaBound{v}\AgdaSymbol{)}%
\>[83]\AgdaComment{--\ 3}\<%
\\
%
\>[2]\AgdaSymbol{(}\AgdaBound{[I]≈ε}\AgdaSpace{}%
\AgdaSymbol{:}\AgdaSpace{}%
\AgdaOperator{\AgdaPostulate{[}}\AgdaSpace{}%
\AgdaPostulate{I}\AgdaSpace{}%
\AgdaOperator{\AgdaPostulate{]}}\AgdaSpace{}%
\AgdaOperator{\AgdaPostulate{≈₂}}\AgdaSpace{}%
\AgdaInductiveConstructor{ε}\AgdaSymbol{)}%
\>[83]\AgdaComment{--\ 4}\<%
\\
%
\>[2]\AgdaSymbol{(}\AgdaBound{h=ract}\AgdaSpace{}%
\AgdaSymbol{:}\AgdaSpace{}%
\AgdaSymbol{∀}\AgdaSpace{}%
\AgdaBound{c}\AgdaSpace{}%
\AgdaBound{b}\AgdaSpace{}%
\AgdaSymbol{→}\AgdaSpace{}%
\AgdaKeyword{let}\AgdaSpace{}%
\AgdaSymbol{(}\AgdaBound{b'}\AgdaSpace{}%
\AgdaOperator{\AgdaInductiveConstructor{,}}\AgdaSpace{}%
\AgdaBound{c'}\AgdaSymbol{)}\AgdaSpace{}%
\AgdaSymbol{=}\AgdaSpace{}%
\AgdaPostulate{h}\AgdaSpace{}%
\AgdaBound{c}\AgdaSpace{}%
\AgdaBound{b}\AgdaSpace{}%
\AgdaKeyword{in}\AgdaSpace{}%
\AgdaOperator{\AgdaPostulate{[}}\AgdaSpace{}%
\AgdaBound{c}\AgdaSpace{}%
\AgdaOperator{\AgdaPostulate{]}}\AgdaSpace{}%
\AgdaOperator{\AgdaInductiveConstructor{•}}\AgdaSpace{}%
\AgdaOperator{\AgdaInductiveConstructor{[}}\AgdaSpace{}%
\AgdaBound{b}\AgdaSpace{}%
\AgdaOperator{\AgdaInductiveConstructor{]ʷ}}\AgdaSpace{}%
\AgdaOperator{\AgdaPostulate{≈₂}}\AgdaSpace{}%
\AgdaSymbol{(}\AgdaPostulate{f}\AgdaSpace{}%
\AgdaOperator{\AgdaFunction{ʷ}}\AgdaSymbol{)}\AgdaSpace{}%
\AgdaBound{b'}\AgdaSpace{}%
\AgdaOperator{\AgdaInductiveConstructor{•}}\AgdaSpace{}%
\AgdaOperator{\AgdaPostulate{[}}\AgdaSpace{}%
\AgdaBound{c'}\AgdaSpace{}%
\AgdaOperator{\AgdaPostulate{]}}\AgdaSymbol{)}%
\>[83]\AgdaComment{--\ 5}\<%
\end{code}
\begin{code}[hide]%
%
\>[2]\AgdaKeyword{where}\<%
\end{code}
\vspace{-4pt}
{\small (1) the table inverts the embedding at $I$;\; (2) it respects $\Delta$'s
axioms, coset by coset;\; (3) $f$ respects $\Gamma$'s axioms;\; (4) $[I]\approx\varepsilon$;\;
(5) the soundness law, as \aid{ract-sound}.}
\end{frame}

\begin{frame}{What the engine gives back \InPaper}
\begin{columns}[c]
\begin{column}{0.52\textwidth}
From the five hypotheses, \aid{Transfer} derives well-definedness of
$\aid{nf} = (h\,\aid{ᵗ})\,I$, the factorization of whole words, injectivity,
and the payoff --- the \alert{normal-form transport}:
\[
  \aid{nf-transp} : \aid{NormalForm}\;\Gamma\;\mathit{NF}
    \to \aid{NormalForm}\;\Delta\;(\mathit{NF}\times C)
\]
a normal form for the subgroup yields one for the group, \alert{with one
more coset digit}.

\vspace{4pt}
\small \aid{PackedCosetTable} bundles data + hypotheses into one value;
\aid{CosetTower} iterates up an $\mathbb{N}$-indexed family:\\[2pt]
\quad\aid{tower-carrier NF₀ (k+1) = tower-carrier NF₀ k × Cₖ}
\end{column}
\begin{column}{0.45\textwidth}
\centering
\begin{tikzpicture}[x=9mm,y=9mm,>={Stealth[length=1.8mm]},
  lv/.style={draw=accent,fill=soft,rounded corners=2pt,font=\scriptsize,minimum width=26mm,minimum height=6mm},
  dg/.style={draw=newC,fill=softnew,rounded corners=2pt,font=\scriptsize,minimum width=8mm,minimum height=6mm}]
  \node[lv] (l0) at (0,0) {$\mathit{NF}_0$ (base)};
  \node[lv] (l1) at (0,1.2) {level 1};
  \node[lv] (l2) at (0,2.4) {level 2};
  \node[lv] (l3) at (0,3.6) {level $k$};
  \node[dg] (d1) at (2.6,1.2) {$C_1$};
  \node[dg] (d2) at (2.6,2.4) {$C_2$};
  \node[dg] (d3) at (2.6,3.6) {$C_k$};
  \draw[->,thick,paperC] (l0) -- node[left,font=\tiny]{table 1} (l1);
  \draw[->,thick,paperC] (l1) -- node[left,font=\tiny]{table 2} (l2);
  \draw[->,thick,paperC,dashed] (l2) -- (l3);
  \foreach \i in {1,2,3} \draw[-,gray] (l\i) -- (d\i);
  \node[font=\scriptsize,align=center] at (1.3,-1.0) {carrier $\mathit{NF}_0\times C_1\times C_2\times\dots\times C_k$};
\end{tikzpicture}
\end{column}
\end{columns}
\end{frame}


\begin{code}[hide]%
\>[0]\AgdaKeyword{module}\AgdaSpace{}%
\AgdaModule{slides.Tsinghua-PL-seminar.sections.design2}\AgdaSpace{}%
\AgdaKeyword{where}\<%
\\
%
\\[\AgdaEmptyExtraSkip]%
\>[0]\AgdaKeyword{open}\AgdaSpace{}%
\AgdaKeyword{import}\AgdaSpace{}%
\AgdaModule{Data.Nat}\AgdaSpace{}%
\AgdaKeyword{using}\AgdaSpace{}%
\AgdaSymbol{(}\AgdaDatatype{ℕ}\AgdaSymbol{)}\<%
\\
\>[0]\AgdaKeyword{open}\AgdaSpace{}%
\AgdaKeyword{import}\AgdaSpace{}%
\AgdaModule{Data.Product}\AgdaSpace{}%
\AgdaKeyword{using}\AgdaSpace{}%
\AgdaSymbol{(}\AgdaOperator{\AgdaFunction{\AgdaUnderscore{}×\AgdaUnderscore{}}}\AgdaSpace{}%
\AgdaSymbol{;}\AgdaSpace{}%
\AgdaOperator{\AgdaInductiveConstructor{\AgdaUnderscore{},\AgdaUnderscore{}}}\AgdaSymbol{)}\<%
\\
\>[0]\AgdaKeyword{open}\AgdaSpace{}%
\AgdaKeyword{import}\AgdaSpace{}%
\AgdaModule{Data.Sum}\AgdaSpace{}%
\AgdaKeyword{using}\AgdaSpace{}%
\AgdaSymbol{(}\AgdaOperator{\AgdaDatatype{\AgdaUnderscore{}⊎\AgdaUnderscore{}}}\AgdaSymbol{)}\<%
\\
\>[0]\AgdaKeyword{open}\AgdaSpace{}%
\AgdaKeyword{import}\AgdaSpace{}%
\AgdaModule{Data.List}\AgdaSpace{}%
\AgdaKeyword{using}\AgdaSpace{}%
\AgdaSymbol{(}\AgdaDatatype{List}\AgdaSymbol{)}\<%
\\
\>[0]\AgdaKeyword{open}\AgdaSpace{}%
\AgdaKeyword{import}\AgdaSpace{}%
\AgdaModule{Data.Unit}\AgdaSpace{}%
\AgdaKeyword{using}\AgdaSpace{}%
\AgdaSymbol{(}\AgdaRecord{⊤}\AgdaSymbol{)}\<%
\\
\>[0]\AgdaKeyword{open}\AgdaSpace{}%
\AgdaKeyword{import}\AgdaSpace{}%
\AgdaModule{Level}\AgdaSpace{}%
\AgdaKeyword{using}\AgdaSpace{}%
\AgdaSymbol{(}\AgdaFunction{0ℓ}\AgdaSymbol{)}\<%
\\
\>[0]\AgdaKeyword{import}\AgdaSpace{}%
\AgdaModule{Relation.Binary.PropositionalEquality}\AgdaSpace{}%
\AgdaSymbol{as}\AgdaSpace{}%
\AgdaModule{Eq}\<%
\\
%
\\[\AgdaEmptyExtraSkip]%
\>[0]\AgdaKeyword{open}\AgdaSpace{}%
\AgdaKeyword{import}\AgdaSpace{}%
\AgdaModule{Word.Base}\AgdaSpace{}%
\AgdaKeyword{using}\AgdaSpace{}%
\AgdaSymbol{(}\AgdaDatatype{Word}\AgdaSpace{}%
\AgdaSymbol{;}\AgdaSpace{}%
\AgdaFunction{WRel}\AgdaSpace{}%
\AgdaSymbol{;}\AgdaSpace{}%
\AgdaOperator{\AgdaInductiveConstructor{[\AgdaUnderscore{}]ʷ}}\AgdaSpace{}%
\AgdaSymbol{;}\AgdaSpace{}%
\AgdaInductiveConstructor{ε}\AgdaSpace{}%
\AgdaSymbol{;}\AgdaSpace{}%
\AgdaOperator{\AgdaInductiveConstructor{\AgdaUnderscore{}•\AgdaUnderscore{}}}\AgdaSymbol{)}\<%
\\
\>[0]\AgdaKeyword{open}\AgdaSpace{}%
\AgdaKeyword{import}\AgdaSpace{}%
\AgdaModule{Presentation.Construct.Base}\AgdaSpace{}%
\AgdaKeyword{using}\AgdaSpace{}%
\AgdaSymbol{(}\AgdaOperator{\AgdaFunction{[\AgdaUnderscore{}]ₗ}}\AgdaSpace{}%
\AgdaSymbol{;}\AgdaSpace{}%
\AgdaOperator{\AgdaFunction{[\AgdaUnderscore{}]ᵣ}}\AgdaSymbol{)}\<%
\\
%
\\[\AgdaEmptyExtraSkip]%
\>[0]\AgdaKeyword{infix}\AgdaSpace{}%
\AgdaNumber{5}\AgdaSpace{}%
\AgdaOperator{\AgdaDatatype{\AgdaUnderscore{}⋄\AgdaUnderscore{}⋄\AgdaUnderscore{}}}\<%
\\
%
\\[\AgdaEmptyExtraSkip]%
\>[0]\AgdaKeyword{postulate}\<%
\\
\>[0][@{}l@{\AgdaIndent{0}}]%
\>[2]\AgdaPostulate{A}\AgdaSpace{}%
\AgdaPostulate{B}\AgdaSpace{}%
\AgdaSymbol{:}\AgdaSpace{}%
\AgdaPrimitive{Set}\<%
\\
%
\>[2]\AgdaPostulate{u}\AgdaSpace{}%
\AgdaPostulate{v}\AgdaSpace{}%
\AgdaSymbol{:}\AgdaSpace{}%
\AgdaDatatype{Word}\AgdaSpace{}%
\AgdaPostulate{A}\<%
\end{code}
%% Part 4, continued — a real coset level, products, stdlib (paper §4).
\begin{frame}{A real coset level, in full: adjoining $X$ to $\langle\omega\rangle\times\langle S\rangle$ \InPaper}
\begin{code}[hide]%
\>[0]\AgdaKeyword{module}\AgdaSpace{}%
\AgdaModule{Sω}\AgdaSpace{}%
\AgdaKeyword{where}\<%
\\
\>[0][@{}l@{\AgdaIndent{0}}]%
\>[2]\AgdaKeyword{postulate}\AgdaSpace{}%
\AgdaPostulate{ω}\AgdaSpace{}%
\AgdaPostulate{S}\AgdaSpace{}%
\AgdaSymbol{:}\AgdaSpace{}%
\AgdaDatatype{Word}\AgdaSpace{}%
\AgdaRecord{⊤}\<%
\\
\>[0]\AgdaKeyword{module}\AgdaSpace{}%
\AgdaModule{CT}\AgdaSpace{}%
\AgdaKeyword{where}\<%
\\
\>[0][@{}l@{\AgdaIndent{0}}]%
\>[2]\AgdaFunction{Sω}\AgdaSpace{}%
\AgdaSymbol{=}\AgdaSpace{}%
\AgdaRecord{⊤}\<%
\\
%
\>[2]\AgdaKeyword{data}\AgdaSpace{}%
\AgdaDatatype{XSω}\AgdaSpace{}%
\AgdaSymbol{:}\AgdaSpace{}%
\AgdaPrimitive{Set}\AgdaSpace{}%
\AgdaKeyword{where}\<%
\\
\>[2][@{}l@{\AgdaIndent{0}}]%
\>[4]\AgdaInductiveConstructor{X-gen}\AgdaSpace{}%
\AgdaInductiveConstructor{S-gen}\AgdaSpace{}%
\AgdaInductiveConstructor{ω-gen}\AgdaSpace{}%
\AgdaSymbol{:}\AgdaSpace{}%
\AgdaDatatype{XSω}\<%
\\
%
\>[2]\AgdaKeyword{data}\AgdaSpace{}%
\AgdaDatatype{Cr}\AgdaSpace{}%
\AgdaSymbol{:}\AgdaSpace{}%
\AgdaPrimitive{Set}\AgdaSpace{}%
\AgdaKeyword{where}\<%
\\
\>[2][@{}l@{\AgdaIndent{0}}]%
\>[4]\AgdaInductiveConstructor{X-cr}\AgdaSpace{}%
\AgdaInductiveConstructor{ε-cr}\AgdaSpace{}%
\AgdaSymbol{:}\AgdaSpace{}%
\AgdaDatatype{Cr}\<%
\\
%
\>[2]\AgdaKeyword{infixr}\AgdaSpace{}%
\AgdaNumber{8}\AgdaSpace{}%
\AgdaOperator{\AgdaPostulate{\AgdaUnderscore{}\textasciicircum{}\AgdaUnderscore{}}}\<%
\\
%
\>[2]\AgdaKeyword{postulate}\<%
\\
\>[2][@{}l@{\AgdaIndent{0}}]%
\>[4]\AgdaOperator{\AgdaPostulate{\AgdaUnderscore{}\textasciicircum{}\AgdaUnderscore{}}}\AgdaSpace{}%
\AgdaSymbol{:}\AgdaSpace{}%
\AgdaSymbol{∀}\AgdaSpace{}%
\AgdaSymbol{\{}\AgdaBound{A}\AgdaSpace{}%
\AgdaSymbol{:}\AgdaSpace{}%
\AgdaPrimitive{Set}\AgdaSymbol{\}}\AgdaSpace{}%
\AgdaSymbol{→}\AgdaSpace{}%
\AgdaDatatype{Word}\AgdaSpace{}%
\AgdaBound{A}\AgdaSpace{}%
\AgdaSymbol{→}\AgdaSpace{}%
\AgdaDatatype{ℕ}\AgdaSpace{}%
\AgdaSymbol{→}\AgdaSpace{}%
\AgdaDatatype{Word}\AgdaSpace{}%
\AgdaBound{A}\<%
\\
%
\>[4]\AgdaOperator{\AgdaPostulate{\AgdaUnderscore{}===′\AgdaUnderscore{}}}\AgdaSpace{}%
\AgdaSymbol{:}\AgdaSpace{}%
\AgdaFunction{WRel}\AgdaSpace{}%
\AgdaDatatype{XSω}\<%
\\
%
\>[4]\AgdaOperator{\AgdaPostulate{\AgdaUnderscore{}\textasciitilde{}′\AgdaUnderscore{}}}\AgdaSpace{}%
\AgdaSymbol{:}\AgdaSpace{}%
\AgdaDatatype{Word}\AgdaSpace{}%
\AgdaFunction{Sω}\AgdaSpace{}%
\AgdaOperator{\AgdaFunction{×}}\AgdaSpace{}%
\AgdaDatatype{Cr}\AgdaSpace{}%
\AgdaSymbol{→}\AgdaSpace{}%
\AgdaDatatype{Word}\AgdaSpace{}%
\AgdaFunction{Sω}\AgdaSpace{}%
\AgdaOperator{\AgdaFunction{×}}\AgdaSpace{}%
\AgdaDatatype{Cr}\AgdaSpace{}%
\AgdaSymbol{→}\AgdaSpace{}%
\AgdaPrimitive{Set}\<%
\\
%
\>[4]\AgdaPostulate{by-equal-nf′}\AgdaSpace{}%
\AgdaSymbol{:}\AgdaSpace{}%
\AgdaSymbol{∀}\AgdaSpace{}%
\AgdaSymbol{\{}\AgdaBound{A}\AgdaSpace{}%
\AgdaSymbol{:}\AgdaSpace{}%
\AgdaPrimitive{Set}\AgdaSymbol{\}}\AgdaSpace{}%
\AgdaSymbol{→}\AgdaSpace{}%
\AgdaBound{A}\<%
\\
%
\>[2]\AgdaKeyword{open}\AgdaSpace{}%
\AgdaModule{Eq}\AgdaSpace{}%
\AgdaKeyword{using}\AgdaSpace{}%
\AgdaSymbol{(}\AgdaInductiveConstructor{refl}\AgdaSymbol{)}\<%
\\
%
\>[2]\AgdaKeyword{data}\AgdaSpace{}%
\AgdaOperator{\AgdaDatatype{\AgdaUnderscore{}==\AgdaUnderscore{}}}\AgdaSpace{}%
\AgdaSymbol{:}\AgdaSpace{}%
\AgdaFunction{WRel}\AgdaSpace{}%
\AgdaDatatype{XSω}\AgdaSpace{}%
\AgdaKeyword{where}\<%
\\
\>[2][@{}l@{\AgdaIndent{0}}]%
\>[4]\AgdaInductiveConstructor{order-ω}\AgdaSpace{}%
\AgdaInductiveConstructor{order-S}\AgdaSpace{}%
\AgdaInductiveConstructor{order-X}\AgdaSpace{}%
\AgdaInductiveConstructor{order-SX}\AgdaSpace{}%
\AgdaSymbol{:}\AgdaSpace{}%
\AgdaInductiveConstructor{ε}\AgdaSpace{}%
\AgdaOperator{\AgdaDatatype{==}}\AgdaSpace{}%
\AgdaInductiveConstructor{ε}\<%
\\
%
\>[4]\AgdaInductiveConstructor{comm}\AgdaSpace{}%
\AgdaSymbol{:}\AgdaSpace{}%
\AgdaSymbol{∀}\AgdaSpace{}%
\AgdaSymbol{\{}\AgdaBound{gen}\AgdaSpace{}%
\AgdaSymbol{:}\AgdaSpace{}%
\AgdaDatatype{XSω}\AgdaSymbol{\}}\AgdaSpace{}%
\AgdaSymbol{→}\AgdaSpace{}%
\AgdaInductiveConstructor{ε}\AgdaSpace{}%
\AgdaOperator{\AgdaDatatype{==}}\AgdaSpace{}%
\AgdaInductiveConstructor{ε}\<%
\end{code}
\begin{columns}[T]
\begin{column}{0.47\textwidth}
\small Base $\langle\omega\rangle\times\langle S\rangle$ (radices $8\cdot 4$);
adjoin $X$ with $X^2=\varepsilon$, $(SX)^2=\omega^2$, $\omega$ central.
Two cosets, $\varepsilon$ and $X$. The whole table:
\begin{code}%
%
\>[2]\AgdaFunction{h}\AgdaSpace{}%
\AgdaSymbol{:}\AgdaSpace{}%
\AgdaDatatype{Cr}\AgdaSpace{}%
\AgdaSymbol{→}\AgdaSpace{}%
\AgdaDatatype{XSω}\AgdaSpace{}%
\AgdaSymbol{→}\AgdaSpace{}%
\AgdaDatatype{Word}\AgdaSpace{}%
\AgdaFunction{Sω}\AgdaSpace{}%
\AgdaOperator{\AgdaFunction{×}}\AgdaSpace{}%
\AgdaDatatype{Cr}\<%
\\
%
\>[2]\AgdaFunction{h}\AgdaSpace{}%
\AgdaInductiveConstructor{X-cr}\AgdaSpace{}%
\AgdaInductiveConstructor{X-gen}\AgdaSpace{}%
\AgdaSymbol{=}\AgdaSpace{}%
\AgdaInductiveConstructor{ε}\AgdaSpace{}%
\AgdaOperator{\AgdaInductiveConstructor{,}}\AgdaSpace{}%
\AgdaInductiveConstructor{ε-cr}\<%
\\
%
\>[2]\AgdaFunction{h}\AgdaSpace{}%
\AgdaInductiveConstructor{X-cr}\AgdaSpace{}%
\AgdaInductiveConstructor{S-gen}\AgdaSpace{}%
\AgdaSymbol{=}\AgdaSpace{}%
\AgdaSymbol{(}\AgdaPostulate{Sω.ω}\AgdaSpace{}%
\AgdaOperator{\AgdaPostulate{\textasciicircum{}}}\AgdaSpace{}%
\AgdaNumber{2}\AgdaSpace{}%
\AgdaOperator{\AgdaInductiveConstructor{•}}\AgdaSpace{}%
\AgdaPostulate{Sω.S}\AgdaSpace{}%
\AgdaOperator{\AgdaPostulate{\textasciicircum{}}}\AgdaSpace{}%
\AgdaNumber{3}\AgdaSymbol{)}\AgdaSpace{}%
\AgdaOperator{\AgdaInductiveConstructor{,}}\AgdaSpace{}%
\AgdaInductiveConstructor{X-cr}\<%
\\
%
\>[2]\AgdaFunction{h}\AgdaSpace{}%
\AgdaInductiveConstructor{X-cr}\AgdaSpace{}%
\AgdaInductiveConstructor{ω-gen}\AgdaSpace{}%
\AgdaSymbol{=}\AgdaSpace{}%
\AgdaPostulate{Sω.ω}\AgdaSpace{}%
\AgdaOperator{\AgdaInductiveConstructor{,}}\AgdaSpace{}%
\AgdaInductiveConstructor{X-cr}\<%
\\
%
\>[2]\AgdaFunction{h}\AgdaSpace{}%
\AgdaInductiveConstructor{ε-cr}\AgdaSpace{}%
\AgdaInductiveConstructor{X-gen}\AgdaSpace{}%
\AgdaSymbol{=}\AgdaSpace{}%
\AgdaInductiveConstructor{ε}\AgdaSpace{}%
\AgdaOperator{\AgdaInductiveConstructor{,}}\AgdaSpace{}%
\AgdaInductiveConstructor{X-cr}\<%
\\
%
\>[2]\AgdaFunction{h}\AgdaSpace{}%
\AgdaInductiveConstructor{ε-cr}\AgdaSpace{}%
\AgdaInductiveConstructor{S-gen}\AgdaSpace{}%
\AgdaSymbol{=}\AgdaSpace{}%
\AgdaPostulate{Sω.S}\AgdaSpace{}%
\AgdaOperator{\AgdaInductiveConstructor{,}}\AgdaSpace{}%
\AgdaInductiveConstructor{ε-cr}\<%
\\
%
\>[2]\AgdaFunction{h}\AgdaSpace{}%
\AgdaInductiveConstructor{ε-cr}\AgdaSpace{}%
\AgdaInductiveConstructor{ω-gen}\AgdaSpace{}%
\AgdaSymbol{=}\AgdaSpace{}%
\AgdaPostulate{Sω.ω}\AgdaSpace{}%
\AgdaOperator{\AgdaInductiveConstructor{,}}\AgdaSpace{}%
\AgdaInductiveConstructor{ε-cr}\<%
\end{code}
{\scriptsize from \texttt{Examples/Amalgamations/CliffordT1.agda} (coset type renamed \aid{Cr})}
\end{column}
\begin{column}{0.50\textwidth}
\small Hypothesis (2), \emph{every} case closed by evaluation:
\begin{code}%
%
\>[2]\AgdaFunction{h-wd-ax}\AgdaSpace{}%
\AgdaSymbol{:}\AgdaSpace{}%
\AgdaSymbol{∀}\AgdaSpace{}%
\AgdaBound{c}\AgdaSpace{}%
\AgdaSymbol{\{}\AgdaBound{u}\AgdaSpace{}%
\AgdaBound{t}\AgdaSymbol{\}}\AgdaSpace{}%
\AgdaSymbol{→}\AgdaSpace{}%
\AgdaBound{u}\AgdaSpace{}%
\AgdaOperator{\AgdaDatatype{==}}\AgdaSpace{}%
\AgdaBound{t}\AgdaSpace{}%
\AgdaSymbol{→}\AgdaSpace{}%
\AgdaPrimitive{Set}\<%
\\
%
\>[2]\AgdaFunction{h-wd-ax}\AgdaSpace{}%
\AgdaInductiveConstructor{X-cr}\AgdaSpace{}%
\AgdaInductiveConstructor{order-ω}\AgdaSpace{}%
\AgdaSymbol{=}\AgdaSpace{}%
\AgdaPostulate{by-equal-nf′}\AgdaSpace{}%
\AgdaInductiveConstructor{refl}\AgdaSpace{}%
\AgdaOperator{\AgdaInductiveConstructor{,}}\AgdaSpace{}%
\AgdaInductiveConstructor{refl}\<%
\\
%
\>[2]\AgdaFunction{h-wd-ax}\AgdaSpace{}%
\AgdaInductiveConstructor{X-cr}\AgdaSpace{}%
\AgdaInductiveConstructor{order-S}\AgdaSpace{}%
\AgdaSymbol{=}\AgdaSpace{}%
\AgdaPostulate{by-equal-nf′}\AgdaSpace{}%
\AgdaInductiveConstructor{refl}\AgdaSpace{}%
\AgdaOperator{\AgdaInductiveConstructor{,}}\AgdaSpace{}%
\AgdaInductiveConstructor{refl}\<%
\\
%
\>[2]\AgdaFunction{h-wd-ax}\AgdaSpace{}%
\AgdaInductiveConstructor{X-cr}\AgdaSpace{}%
\AgdaInductiveConstructor{order-X}\AgdaSpace{}%
\AgdaSymbol{=}\AgdaSpace{}%
\AgdaPostulate{by-equal-nf′}\AgdaSpace{}%
\AgdaInductiveConstructor{refl}\AgdaSpace{}%
\AgdaOperator{\AgdaInductiveConstructor{,}}\AgdaSpace{}%
\AgdaInductiveConstructor{refl}\<%
\\
%
\>[2]\AgdaFunction{h-wd-ax}\AgdaSpace{}%
\AgdaInductiveConstructor{X-cr}\AgdaSpace{}%
\AgdaInductiveConstructor{order-SX}\AgdaSpace{}%
\AgdaSymbol{=}\AgdaSpace{}%
\AgdaPostulate{by-equal-nf′}\AgdaSpace{}%
\AgdaInductiveConstructor{refl}\AgdaSpace{}%
\AgdaOperator{\AgdaInductiveConstructor{,}}\AgdaSpace{}%
\AgdaInductiveConstructor{refl}\<%
\\
%
\>[2]\AgdaFunction{h-wd-ax}\AgdaSpace{}%
\AgdaInductiveConstructor{X-cr}\AgdaSpace{}%
\AgdaSymbol{(}\AgdaInductiveConstructor{comm}\AgdaSpace{}%
\AgdaSymbol{\{}\AgdaInductiveConstructor{X-gen}\AgdaSymbol{\})}\AgdaSpace{}%
\AgdaSymbol{=}\AgdaSpace{}%
\AgdaPostulate{by-equal-nf′}\AgdaSpace{}%
\AgdaInductiveConstructor{refl}\AgdaSpace{}%
\AgdaOperator{\AgdaInductiveConstructor{,}}\AgdaSpace{}%
\AgdaInductiveConstructor{refl}\<%
\\
%
\>[2]\AgdaFunction{h-wd-ax}\AgdaSpace{}%
\AgdaInductiveConstructor{X-cr}\AgdaSpace{}%
\AgdaSymbol{(}\AgdaInductiveConstructor{comm}\AgdaSpace{}%
\AgdaSymbol{\{}\AgdaInductiveConstructor{S-gen}\AgdaSymbol{\})}\AgdaSpace{}%
\AgdaSymbol{=}\AgdaSpace{}%
\AgdaPostulate{by-equal-nf′}\AgdaSpace{}%
\AgdaInductiveConstructor{refl}\AgdaSpace{}%
\AgdaOperator{\AgdaInductiveConstructor{,}}\AgdaSpace{}%
\AgdaInductiveConstructor{refl}\<%
\\
%
\>[2]\AgdaFunction{h-wd-ax}\AgdaSpace{}%
\AgdaInductiveConstructor{X-cr}\AgdaSpace{}%
\AgdaSymbol{(}\AgdaInductiveConstructor{comm}\AgdaSpace{}%
\AgdaSymbol{\{}\AgdaInductiveConstructor{ω-gen}\AgdaSymbol{\})}\AgdaSpace{}%
\AgdaSymbol{=}\AgdaSpace{}%
\AgdaPostulate{by-equal-nf′}\AgdaSpace{}%
\AgdaInductiveConstructor{refl}\AgdaSpace{}%
\AgdaOperator{\AgdaInductiveConstructor{,}}\AgdaSpace{}%
\AgdaInductiveConstructor{refl}\<%
\\
%
\>[2]\AgdaComment{--\ …\ and\ the\ same\ seven\ for\ ε-cr}\<%
\end{code}
{\scriptsize abridged: \aid{(by-equal-nf Eq.refl) , Eq.refl} in the source}
\end{column}
\end{columns}
\begin{keybox}
\small The typechecker runs the table on both sides of the axiom: the cosets
coincide and the residues have the same normal form one level down.
Qubit Clifford+$T$: \textbf{59} such obligations; qutrit's nine-coset level
alone: \textbf{72} --- every one \aid{refl}.
\end{keybox}
\end{frame}

\begin{frame}{The one case that is a real derivation \InPaper}
\begin{columns}[c]
\begin{column}{0.62\textwidth}
\small Hypothesis (5) at coset $X$, letter $S$: the table says
$X\cdot S = (\omega^2 S^3)\cdot X$. The proof, verbatim:
\vspace{2pt}

{\agdafont\scriptsize
\begin{tabular}{@{}l@{}}
\textcolor{AgdaFunction}{lemma-XS} : \textcolor{AgdaFunction}{X} \textcolor{AgdaInductiveConstructor}{•} \textcolor{AgdaFunction}{S} \textcolor{AgdaDatatype}{≈} (\textcolor{AgdaFunction}{ω} \textcolor{AgdaFunction}{\textasciicircum} 2 \textcolor{AgdaInductiveConstructor}{•} \textcolor{AgdaFunction}{S} \textcolor{AgdaFunction}{\textasciicircum} 3) \textcolor{AgdaInductiveConstructor}{•} \textcolor{AgdaFunction}{X}\\
\textcolor{AgdaFunction}{lemma-XS} = \textcolor{AgdaFunction}{begin}\\
\quad \textcolor{AgdaFunction}{X} \textcolor{AgdaInductiveConstructor}{•} \textcolor{AgdaFunction}{S} \textcolor{AgdaFunction}{≈⟨} \textcolor{AgdaInductiveConstructor}{trans} (\textcolor{AgdaInductiveConstructor}{sym} \textcolor{AgdaInductiveConstructor}{right-unit}) (\textcolor{AgdaInductiveConstructor}{sym} (\textcolor{AgdaInductiveConstructor}{cong} \textcolor{AgdaInductiveConstructor}{refl} (\textcolor{AgdaInductiveConstructor}{axiom} \textcolor{AgdaInductiveConstructor}{order-X}))) \textcolor{AgdaFunction}{⟩}\\
\quad (\textcolor{AgdaFunction}{X} \textcolor{AgdaInductiveConstructor}{•} \textcolor{AgdaFunction}{S}) \textcolor{AgdaInductiveConstructor}{•} (\textcolor{AgdaFunction}{X} \textcolor{AgdaInductiveConstructor}{•} \textcolor{AgdaFunction}{X}) \textcolor{AgdaFunction}{≈⟨} \textcolor{AgdaInductiveConstructor}{trans} (\textcolor{AgdaInductiveConstructor}{sym} \textcolor{AgdaInductiveConstructor}{left-unit}) (\textcolor{AgdaInductiveConstructor}{sym} (\textcolor{AgdaInductiveConstructor}{cong} (\textcolor{AgdaInductiveConstructor}{axiom} \textcolor{AgdaInductiveConstructor}{order-S}) \textcolor{AgdaInductiveConstructor}{refl})) \textcolor{AgdaFunction}{⟩}\\
\quad (\textcolor{AgdaFunction}{S} \textcolor{AgdaFunction}{\textasciicircum} 4) \textcolor{AgdaInductiveConstructor}{•} (\textcolor{AgdaFunction}{X} \textcolor{AgdaInductiveConstructor}{•} \textcolor{AgdaFunction}{S}) \textcolor{AgdaInductiveConstructor}{•} (\textcolor{AgdaFunction}{X} \textcolor{AgdaInductiveConstructor}{•} \textcolor{AgdaFunction}{X}) \textcolor{AgdaFunction}{≈⟨} \textcolor{AgdaFunction}{by-assoc} \textcolor{AgdaInductiveConstructor}{Eq.refl} \textcolor{AgdaFunction}{⟩}\\
\quad (\textcolor{AgdaFunction}{S} \textcolor{AgdaFunction}{\textasciicircum} 3) \textcolor{AgdaInductiveConstructor}{•} (\textcolor{AgdaFunction}{S} \textcolor{AgdaInductiveConstructor}{•} \textcolor{AgdaFunction}{X}) \textcolor{AgdaFunction}{\textasciicircum} 2 \textcolor{AgdaInductiveConstructor}{•} \textcolor{AgdaFunction}{X} \textcolor{AgdaFunction}{≈⟨} \textcolor{AgdaInductiveConstructor}{cong} \textcolor{AgdaInductiveConstructor}{refl} (\textcolor{AgdaInductiveConstructor}{cong} (\textcolor{AgdaInductiveConstructor}{axiom} \textcolor{AgdaInductiveConstructor}{order-SX}) \textcolor{AgdaInductiveConstructor}{refl}) \textcolor{AgdaFunction}{⟩}\\
\quad (\textcolor{AgdaFunction}{S} \textcolor{AgdaFunction}{\textasciicircum} 3) \textcolor{AgdaInductiveConstructor}{•} \textcolor{AgdaFunction}{ω} \textcolor{AgdaFunction}{\textasciicircum} 2 \textcolor{AgdaInductiveConstructor}{•} \textcolor{AgdaFunction}{X} \textcolor{AgdaFunction}{≈⟨} \textcolor{AgdaInductiveConstructor}{sym} \textcolor{AgdaInductiveConstructor}{assoc} \textcolor{AgdaFunction}{⟩}\\
\quad (\textcolor{AgdaFunction}{S} \textcolor{AgdaFunction}{\textasciicircum} 3 \textcolor{AgdaInductiveConstructor}{•} \textcolor{AgdaFunction}{ω} \textcolor{AgdaFunction}{\textasciicircum} 2) \textcolor{AgdaInductiveConstructor}{•} \textcolor{AgdaFunction}{X} \textcolor{AgdaFunction}{≈⟨} \textcolor{AgdaInductiveConstructor}{cong} (\textcolor{AgdaFunction}{lemma-ω\textasciicircum{}n} 2 (\textcolor{AgdaFunction}{S} \textcolor{AgdaFunction}{\textasciicircum} 3)) \textcolor{AgdaInductiveConstructor}{refl} \textcolor{AgdaFunction}{⟩}\\
\quad (\textcolor{AgdaFunction}{ω} \textcolor{AgdaFunction}{\textasciicircum} 2 \textcolor{AgdaInductiveConstructor}{•} \textcolor{AgdaFunction}{S} \textcolor{AgdaFunction}{\textasciicircum} 3) \textcolor{AgdaInductiveConstructor}{•} \textcolor{AgdaFunction}{X} \textcolor{AgdaFunction}{∎}
\end{tabular}}
\end{column}
\begin{column}{0.35\textwidth}
\begin{keybox}[title=Division of labour]
\small
\begin{itemize}
  \item \textbf{coverage checker}: no case is missed;
  \item \textbf{evaluator}: closes the cases that are computation;
  \item \textbf{human}: writes the few real derivations, which read
        line by line like the pen-and-paper ones.
\end{itemize}
\end{keybox}
\vspace{4pt}
{\scriptsize (typeset by hand in Agda's colours: this block needs the full
module context of \texttt{CliffordT1.agda})}
\end{column}
\end{columns}
\end{frame}

\begin{frame}{Products of presentations \InPaper}
\begin{columns}[c]
\begin{column}{0.56\textwidth}
Relation sets over $A \uplus B$ from one primitive: a three-way join whose
\alert{mixed} component is the design parameter.
\begin{code}%
\>[0]\AgdaKeyword{data}\AgdaSpace{}%
\AgdaOperator{\AgdaDatatype{\AgdaUnderscore{}⋄\AgdaUnderscore{}⋄\AgdaUnderscore{}}}%
\>[260I]\AgdaSymbol{(}\AgdaBound{Γ₁}\AgdaSpace{}%
\AgdaSymbol{:}\AgdaSpace{}%
\AgdaFunction{WRel}\AgdaSpace{}%
\AgdaPostulate{A}\AgdaSymbol{)}\AgdaSpace{}%
\AgdaSymbol{(}\AgdaBound{Γ₂}\AgdaSpace{}%
\AgdaSymbol{:}\AgdaSpace{}%
\AgdaFunction{WRel}\AgdaSpace{}%
\AgdaPostulate{B}\AgdaSymbol{)}\AgdaSpace{}%
\AgdaSymbol{(}\AgdaBound{Γ₃}\AgdaSpace{}%
\AgdaSymbol{:}\AgdaSpace{}%
\AgdaFunction{WRel}\AgdaSpace{}%
\AgdaSymbol{(}\AgdaPostulate{A}\AgdaSpace{}%
\AgdaOperator{\AgdaDatatype{⊎}}\AgdaSpace{}%
\AgdaPostulate{B}\AgdaSymbol{))}\<%
\\
\>[.][@{}l@{}]\<[260I]%
\>[11]\AgdaSymbol{:}\AgdaSpace{}%
\AgdaFunction{WRel}\AgdaSpace{}%
\AgdaSymbol{(}\AgdaPostulate{A}\AgdaSpace{}%
\AgdaOperator{\AgdaDatatype{⊎}}\AgdaSpace{}%
\AgdaPostulate{B}\AgdaSymbol{)}\AgdaSpace{}%
\AgdaKeyword{where}\<%
\\
\>[0][@{}l@{\AgdaIndent{0}}]%
\>[2]\AgdaInductiveConstructor{left}%
\>[8]\AgdaSymbol{:}\AgdaSpace{}%
\AgdaBound{Γ₁}\AgdaSpace{}%
\AgdaPostulate{u}\AgdaSpace{}%
\AgdaPostulate{v}\AgdaSpace{}%
\AgdaSymbol{→}\AgdaSpace{}%
\AgdaSymbol{(}\AgdaBound{Γ₁}\AgdaSpace{}%
\AgdaOperator{\AgdaDatatype{⋄}}\AgdaSpace{}%
\AgdaBound{Γ₂}\AgdaSpace{}%
\AgdaOperator{\AgdaDatatype{⋄}}\AgdaSpace{}%
\AgdaBound{Γ₃}\AgdaSymbol{)}\AgdaSpace{}%
\AgdaOperator{\AgdaFunction{[}}\AgdaSpace{}%
\AgdaPostulate{u}\AgdaSpace{}%
\AgdaOperator{\AgdaFunction{]ₗ}}\AgdaSpace{}%
\AgdaOperator{\AgdaFunction{[}}\AgdaSpace{}%
\AgdaPostulate{v}\AgdaSpace{}%
\AgdaOperator{\AgdaFunction{]ₗ}}\<%
\\
%
\>[2]\AgdaInductiveConstructor{right}\AgdaSpace{}%
\AgdaSymbol{:}\AgdaSpace{}%
\AgdaBound{Γ₂}\AgdaSpace{}%
\AgdaPostulate{u}\AgdaSpace{}%
\AgdaPostulate{v}\AgdaSpace{}%
\AgdaSymbol{→}\AgdaSpace{}%
\AgdaSymbol{(}\AgdaBound{Γ₁}\AgdaSpace{}%
\AgdaOperator{\AgdaDatatype{⋄}}\AgdaSpace{}%
\AgdaBound{Γ₂}\AgdaSpace{}%
\AgdaOperator{\AgdaDatatype{⋄}}\AgdaSpace{}%
\AgdaBound{Γ₃}\AgdaSymbol{)}\AgdaSpace{}%
\AgdaOperator{\AgdaFunction{[}}\AgdaSpace{}%
\AgdaPostulate{u}\AgdaSpace{}%
\AgdaOperator{\AgdaFunction{]ᵣ}}\AgdaSpace{}%
\AgdaOperator{\AgdaFunction{[}}\AgdaSpace{}%
\AgdaPostulate{v}\AgdaSpace{}%
\AgdaOperator{\AgdaFunction{]ᵣ}}\<%
\\
%
\>[2]\AgdaInductiveConstructor{mid}%
\>[8]\AgdaSymbol{:}\AgdaSpace{}%
\AgdaBound{Γ₃}\AgdaSpace{}%
\AgdaPostulate{u}\AgdaSpace{}%
\AgdaPostulate{v}\AgdaSpace{}%
\AgdaSymbol{→}\AgdaSpace{}%
\AgdaSymbol{(}\AgdaBound{Γ₁}\AgdaSpace{}%
\AgdaOperator{\AgdaDatatype{⋄}}\AgdaSpace{}%
\AgdaBound{Γ₂}\AgdaSpace{}%
\AgdaOperator{\AgdaDatatype{⋄}}\AgdaSpace{}%
\AgdaBound{Γ₃}\AgdaSymbol{)}\AgdaSpace{}%
\AgdaPostulate{u}\AgdaSpace{}%
\AgdaPostulate{v}\<%
\end{code}
\vspace{2pt}
{\small Congruence-lifting lemmas (\aid{lefts}, \aid{rights}) embed the factors;
normal forms transport across.}
\end{column}
\begin{column}{0.41\textwidth}
\small
\renewcommand{\arraystretch}{1.25}
\begin{tabular}{@{}ll@{}}
\toprule
mixed relation $\Gamma_3$ & construction\\
\midrule
\aid{EmptyRel} & free product\\
\aid{CommRel} & direct product $\times$\\
\aid{ConjRel}, \aid{ConjRelʷ} & semidirect $\rtimes$\\
\aid{AmalgRel} $f_1\,f_2$ & amalgamated $\ast_M$\\
\aid{SugarRel} & definitional extension\\
$\Gamma$ \aid{⊕\textasciicircum} $n$ & $n$-fold power\\
\bottomrule
\end{tabular}

\vspace{6pt}
{\scriptsize \aid{ConjRelʷ}: moving $h$ past $n$ conjugates $n$ by a
\emph{word} --- \aid{conj : H → N → Word N}.}
\end{column}
\end{columns}
\end{frame}

\begin{frame}{Constructions come in pairs \InPaper}
\begin{columns}[c]
\begin{column}{0.50\textwidth}
\textbf{Syntactic}: normal-form lifting.
\begin{itemize}
  \item direct / semidirect: carrier $\mathit{NF}_1\times\mathit{NF}_2$;
  \item amalgam: a normal form for the \emph{common subgroup} plus one verified
        coset table per factor:
\end{itemize}
\begin{code}[hide]%
\>[0]\AgdaFunction{amalNFC}\AgdaSpace{}%
\AgdaSymbol{:}\AgdaSpace{}%
\AgdaPrimitive{Set}\AgdaSpace{}%
\AgdaSymbol{→}\AgdaSpace{}%
\AgdaPrimitive{Set}\AgdaSpace{}%
\AgdaSymbol{→}\AgdaSpace{}%
\AgdaPrimitive{Set}\<%
\end{code}
\begin{code}%
\>[0]\AgdaFunction{amalNFC}\AgdaSpace{}%
\AgdaBound{C}\AgdaSpace{}%
\AgdaBound{D}\AgdaSpace{}%
\AgdaSymbol{=}\AgdaSpace{}%
\AgdaSymbol{(}\AgdaBound{D}\AgdaSpace{}%
\AgdaOperator{\AgdaDatatype{⊎}}\AgdaSpace{}%
\AgdaRecord{⊤}\AgdaSymbol{)}\AgdaSpace{}%
\AgdaOperator{\AgdaFunction{×}}\AgdaSpace{}%
\AgdaDatatype{List}\AgdaSpace{}%
\AgdaSymbol{(}\AgdaBound{C}\AgdaSpace{}%
\AgdaOperator{\AgdaFunction{×}}\AgdaSpace{}%
\AgdaBound{D}\AgdaSymbol{)}\AgdaSpace{}%
\AgdaOperator{\AgdaFunction{×}}\AgdaSpace{}%
\AgdaSymbol{(}\AgdaBound{C}\AgdaSpace{}%
\AgdaOperator{\AgdaDatatype{⊎}}\AgdaSpace{}%
\AgdaRecord{⊤}\AgdaSymbol{)}\<%
\end{code}
\textbf{Semantic}: presentation theorems against concrete groups
\[
\begin{array}{l}
(\Gamma \aidop{⋄} \Delta \aidop{⋄} \aid{CommRel})\ \aid{IsPresentationOf}\ (G_1 \times G_2) \\[2pt]
(\Gamma \aidop{⋄} \Delta \aidop{⋄} \aid{ConjRelʷ}\;\mathit{conj})\ \aid{IsPresentationOf}\ (G_1 \rtimes G_2) \\[2pt]
(\Gamma \mathrel{\aid{⊕\textasciicircum}} n)\ \aid{IsPresentationOf}\ (\otimes\aid{-group}\;n) \\[2pt]
(P_1 \mathrel{\aid{*}} P_2 \mathrel{\aid{⋆}} f_1 \mathrel{\aid{⋆}} f_2)\ \aid{IsPresentationOf}\ (G_1 \ast_{G_0} G_2)
\end{array}
\]
\end{column}
\begin{column}{0.47\textwidth}
\centering
{\small alternating normal form of an amalgam:}\\[4pt]
\begin{tikzpicture}[x=8.5mm,
  b0/.style={draw=gray,fill=gray!15,minimum height=7mm,minimum width=8mm,font=\scriptsize},
  bd/.style={draw=paperC,fill=softpaper,minimum height=7mm,minimum width=8mm,font=\scriptsize},
  bc/.style={draw=accent,fill=soft,minimum height=7mm,minimum width=8mm,font=\scriptsize}]
  \node[b0] at (0,0) {$\mathcal{M}$};
  \node[bd] at (1,0) {$d_0$};
  \node[bc] at (2,0) {$c_1$};
  \node[bd] at (3,0) {$d_1$};
  \node at (4,0) {$\cdots$};
  \node[bc] at (5,0) {$c_k$};
  \node[bd] at (6,0) {$d_k$};
  \node[bc] at (7,0) {$c$};
  \draw[decorate,decoration={brace,amplitude=3pt,mirror}] (0.55,-0.5) -- (1.45,-0.5) node[midway,below=2pt,font=\tiny]{$D\uplus\top$};
  \draw[decorate,decoration={brace,amplitude=3pt,mirror}] (1.55,-0.5) -- (6.45,-0.5) node[midway,below=2pt,font=\tiny]{\aid{List (C × D)}};
  \draw[decorate,decoration={brace,amplitude=3pt,mirror}] (6.55,-0.5) -- (7.45,-0.5) node[midway,below=2pt,font=\tiny]{$C\uplus\top$};
  \draw[decorate,decoration={brace,amplitude=3pt}] (-0.45,0.5) -- (0.45,0.5) node[midway,above=2pt,font=\tiny]{$\mathit{NF}_0$};
\end{tikzpicture}

\vspace{8pt}
{\small Group \alert{extensions}: given $1\to N\to G\to Q\to 1$ with the
conjugation action and the corrections by which lifted relators of $Q$ fail,
$N$'s relations + conjugation rules + corrected lifts present $G$.}

\vspace{4pt}
\begin{tikzpicture}[>={Stealth[length=1.6mm]},x=11mm]
  \foreach \x/\l in {0/1,1/N,2/G,3/Q,4/1} \node (n\x) at (\x,0) {$\l$};
  \foreach \a/\b in {0/1,1/2,2/3,3/4} \draw[->] (n\a) -- (n\b);
  \draw[->,dashed,newC] (n3) to[bend right=40] node[above,font=\scriptsize]{lift} (n2);
\end{tikzpicture}
\end{column}
\end{columns}
\end{frame}

\begin{frame}{The semantic side: contributions to the standard library \InPaper}
\begin{columns}[c]
\begin{column}{0.55\textwidth}
A presentation theorem cannot be stated without the group it presents --- and
stdlib had the structures but not the \alert{constructions}.
\begin{itemize}
  \setlength{\itemsep}{2pt}
  \item \texttt{Algebra.Construct.SemiDirectProduct}: $N\rtimes H$ via a
        bundled \aid{Action} (six laws);
  \item \texttt{Algebra.Construct.Amalgamation}: $M\ast_C N$ of
        \emph{monoids}, as a setoid --- the pushout;
  \item \texttt{Algebra.Construct.Extension}, \texttt{.SplitExtension}:
        short exact sequences;
  \item \texttt{Algebra.Morphism.Consequences}: monoid homomorphisms between
        groups are group homomorphisms;
  \item \aid{IsMonoidEpimorphism};\; \aid{Permutation′ n} as a \aid{Group}.
\end{itemize}
\small Seven modules, 865 lines --- being prepared for upstreaming.
\end{column}
\begin{column}{0.42\textwidth}
\begin{keybox}[title={Solvers by evaluation, not reflection}]
\small
\begin{itemize}
  \item \aid{by-assoc}: same letter sequence $\Rightarrow$ $w\approx v$;
  \item \aid{by-passoc}: re-bracket toward a pattern with holes;
  \item \aid{by-equal-nf}: same computed normal form.
\end{itemize}
Each maps both sides into a canonical data type by a function proved
correct once, then compares by \aid{refl}: the evaluator does the work of
a reflective tactic.
\end{keybox}
\end{column}
\end{columns}
\end{frame}

\begin{code}[hide]%
\>[0]\AgdaKeyword{module}\AgdaSpace{}%
\AgdaModule{slides.Tsinghua-PL-seminar.sections.catalogue}\AgdaSpace{}%
\AgdaKeyword{where}\<%
\\
%
\\[\AgdaEmptyExtraSkip]%
\>[0]\AgdaKeyword{open}\AgdaSpace{}%
\AgdaKeyword{import}\AgdaSpace{}%
\AgdaModule{Algebra.Bundles}\AgdaSpace{}%
\AgdaKeyword{using}\AgdaSpace{}%
\AgdaSymbol{(}\AgdaRecord{Group}\AgdaSpace{}%
\AgdaSymbol{;}\AgdaSpace{}%
\AgdaRecord{Monoid}\AgdaSymbol{)}\<%
\\
\>[0]\AgdaKeyword{open}\AgdaSpace{}%
\AgdaKeyword{import}\AgdaSpace{}%
\AgdaModule{Algebra.Morphism.Structures}\AgdaSpace{}%
\AgdaKeyword{using}\AgdaSpace{}%
\AgdaSymbol{(}\AgdaKeyword{module}\AgdaSpace{}%
\AgdaModule{MonoidMorphisms}\AgdaSymbol{)}\<%
\\
\>[0]\AgdaKeyword{open}\AgdaSpace{}%
\AgdaKeyword{import}\AgdaSpace{}%
\AgdaModule{Data.Nat}\AgdaSpace{}%
\AgdaKeyword{using}\AgdaSpace{}%
\AgdaSymbol{(}\AgdaDatatype{ℕ}\AgdaSpace{}%
\AgdaSymbol{;}\AgdaSpace{}%
\AgdaInductiveConstructor{suc}\AgdaSpace{}%
\AgdaSymbol{;}\AgdaSpace{}%
\AgdaInductiveConstructor{2+}\AgdaSymbol{)}\<%
\\
\>[0]\AgdaKeyword{open}\AgdaSpace{}%
\AgdaKeyword{import}\AgdaSpace{}%
\AgdaModule{Data.Nat.Primality}\AgdaSpace{}%
\AgdaKeyword{using}\AgdaSpace{}%
\AgdaSymbol{(}\AgdaRecord{Prime}\AgdaSymbol{)}\<%
\\
\>[0]\AgdaKeyword{open}\AgdaSpace{}%
\AgdaKeyword{import}\AgdaSpace{}%
\AgdaModule{Level}\AgdaSpace{}%
\AgdaKeyword{using}\AgdaSpace{}%
\AgdaSymbol{(}\AgdaFunction{0ℓ}\AgdaSymbol{)}\<%
\\
%
\\[\AgdaEmptyExtraSkip]%
\>[0]\AgdaKeyword{open}\AgdaSpace{}%
\AgdaKeyword{import}\AgdaSpace{}%
\AgdaModule{Notations}\AgdaSpace{}%
\AgdaKeyword{using}\AgdaSpace{}%
\AgdaSymbol{(}\AgdaInductiveConstructor{₁₊}\AgdaSymbol{)}\<%
\\
\>[0]\AgdaKeyword{open}\AgdaSpace{}%
\AgdaKeyword{import}\AgdaSpace{}%
\AgdaModule{Word.Base}\AgdaSpace{}%
\AgdaKeyword{using}\AgdaSpace{}%
\AgdaSymbol{(}\AgdaOperator{\AgdaFunction{\AgdaUnderscore{}ʷ}}\AgdaSymbol{)}\<%
\\
\>[0]\AgdaKeyword{open}\AgdaSpace{}%
\AgdaKeyword{import}\AgdaSpace{}%
\AgdaModule{Presentation.Definitions}\AgdaSpace{}%
\AgdaKeyword{using}\AgdaSpace{}%
\AgdaSymbol{(}\AgdaOperator{\AgdaRecord{\AgdaUnderscore{}IsPresentationOf\AgdaUnderscore{}}}\AgdaSymbol{)}\<%
\\
\>[0]\AgdaKeyword{open}\AgdaSpace{}%
\AgdaKeyword{import}\AgdaSpace{}%
\AgdaModule{Presentation.Construct.Base}\AgdaSpace{}%
\AgdaKeyword{using}\AgdaSpace{}%
\AgdaSymbol{(}\AgdaOperator{\AgdaDatatype{\AgdaUnderscore{}⋄\AgdaUnderscore{}⋄\AgdaUnderscore{}}}\AgdaSpace{}%
\AgdaSymbol{;}\AgdaSpace{}%
\AgdaDatatype{CommRel}\AgdaSpace{}%
\AgdaSymbol{;}\AgdaSpace{}%
\AgdaDatatype{ConjRelʷ}\AgdaSpace{}%
\AgdaSymbol{;}\AgdaSpace{}%
\AgdaOperator{\AgdaFunction{\AgdaUnderscore{}⊕\textasciicircum{}\AgdaUnderscore{}}}\AgdaSymbol{)}\<%
\\
\>[0]\AgdaKeyword{import}\AgdaSpace{}%
\AgdaModule{Presentation.Properties}\AgdaSpace{}%
\AgdaSymbol{as}\AgdaSpace{}%
\AgdaModule{PP}\<%
\\
\>[0]\AgdaKeyword{open}\AgdaSpace{}%
\AgdaKeyword{import}\AgdaSpace{}%
\AgdaModule{Examples.Groups.Cyclic.Normalization}\AgdaSpace{}%
\AgdaKeyword{using}\AgdaSpace{}%
\AgdaSymbol{(}\AgdaOperator{\AgdaDatatype{\AgdaUnderscore{}Cn,\AgdaUnderscore{}===\AgdaUnderscore{}}}\AgdaSymbol{)}\<%
\\
\>[0]\AgdaKeyword{open}\AgdaSpace{}%
\AgdaKeyword{import}\AgdaSpace{}%
\AgdaModule{Examples.Groups.Symmetric.Syntactics}\AgdaSpace{}%
\AgdaKeyword{using}\AgdaSpace{}%
\AgdaSymbol{(}\AgdaOperator{\AgdaDatatype{\AgdaUnderscore{}VRel,\AgdaUnderscore{}===\AgdaUnderscore{}}}\AgdaSymbol{)}\<%
\\
\>[0]\AgdaKeyword{import}\AgdaSpace{}%
\AgdaModule{Examples.Groups.Pauli.Presentation}\AgdaSpace{}%
\AgdaSymbol{as}\AgdaSpace{}%
\AgdaModule{Pauli}\<%
\\
\>[0]\AgdaKeyword{import}\AgdaSpace{}%
\AgdaModule{Examples.Construct.SemiDirectProduct.SnD}\AgdaSpace{}%
\AgdaSymbol{as}\AgdaSpace{}%
\AgdaModule{SnD}\<%
\\
\>[0]\AgdaKeyword{import}\AgdaSpace{}%
\AgdaModule{Examples.Amalgamations.CliffordT1}\AgdaSpace{}%
\AgdaSymbol{as}\AgdaSpace{}%
\AgdaModule{CliffordT1}\<%
\\
%
\\[\AgdaEmptyExtraSkip]%
\>[0]\AgdaKeyword{open}\AgdaSpace{}%
\AgdaModule{SnD}\AgdaSpace{}%
\AgdaKeyword{using}\AgdaSpace{}%
\AgdaSymbol{(}\AgdaFunction{conj}\AgdaSpace{}%
\AgdaSymbol{;}\AgdaSpace{}%
\AgdaKeyword{module}\AgdaSpace{}%
\AgdaModule{Wreath}\AgdaSymbol{)}\<%
\\
\>[0]\AgdaKeyword{open}\AgdaSpace{}%
\AgdaModule{Pauli}\AgdaSpace{}%
\AgdaKeyword{using}\AgdaSpace{}%
\AgdaSymbol{(}\AgdaFunction{Pauli-group}\AgdaSymbol{)}\AgdaSpace{}%
\AgdaKeyword{renaming}\AgdaSpace{}%
\AgdaSymbol{(}\AgdaFunction{Γ-H}\AgdaSpace{}%
\AgdaSymbol{to}\AgdaSpace{}%
\AgdaFunction{Γ-Pauli₁}\AgdaSymbol{)}\<%
\\
\>[0]\AgdaKeyword{open}\AgdaSpace{}%
\AgdaModule{Monoid}\AgdaSpace{}%
\AgdaKeyword{using}\AgdaSpace{}%
\AgdaSymbol{(}\AgdaFunction{rawMonoid}\AgdaSymbol{)}\<%
\\
\>[0]\AgdaKeyword{open}\AgdaSpace{}%
\AgdaModule{MonoidMorphisms}\AgdaSpace{}%
\AgdaKeyword{using}\AgdaSpace{}%
\AgdaSymbol{(}\AgdaRecord{IsMonoidIsomorphism}\AgdaSymbol{)}\<%
\\
\>[0]\AgdaKeyword{open}\AgdaSpace{}%
\AgdaModule{PP}\AgdaSpace{}%
\AgdaKeyword{using}\AgdaSpace{}%
\AgdaSymbol{(}\AgdaFunction{•-ε-monoid}\AgdaSymbol{)}\<%
\\
\>[0]\AgdaKeyword{open}\AgdaSpace{}%
\AgdaModule{CliffordT1.CliffordT1}\AgdaSpace{}%
\AgdaKeyword{using}\AgdaSpace{}%
\AgdaSymbol{(}\AgdaOperator{\AgdaDatatype{\AgdaUnderscore{}===\AgdaUnderscore{}}}\AgdaSpace{}%
\AgdaSymbol{;}\AgdaSpace{}%
\AgdaFunction{amalPres}\AgdaSpace{}%
\AgdaSymbol{;}\AgdaSpace{}%
\AgdaFunction{f}\AgdaSymbol{)}\<%
\end{code}
%% Part 5 — the catalogue (paper §5).
\PaperPart
\section{A catalogue of verified presentations}

\begin{frame}{Each family was chosen to exercise something \InPaper}
\centering
\small
\renewcommand{\arraystretch}{1.15}
\begin{tabular}{@{}l*{7}{c}@{}}
\toprule
 & circuit & coset & tower & $\times$, $\rtimes$ & $n$-fold & amalgam & \aid{refl}-\\
 & layer & tables & & & power & & obligations\\
\midrule
cyclic $\mathbb{Z}_n$, trivial & & \checkmark & base & & & & \\
symmetric $S_n$ & \checkmark & \checkmark & \checkmark & & & & \checkmark\\
wreath $\mathbb{Z}_N\wr S_n$ & \checkmark & & & \checkmark & \checkmark & & \\
Pauli $(\mathbb{Z}_p\times\mathbb{Z}_p)^n$ & & & & \checkmark & \checkmark & & \\
qubit Clifford+$T$ & & \checkmark & \checkmark & \checkmark & & \checkmark & \checkmark\\
qutrit Clifford+$T$ & & \checkmark & \checkmark & \checkmark & & \checkmark & \checkmark\\
$U_3(\mathbb{Z}[\frac12,i])$ & & \checkmark & \checkmark & \checkmark & & \checkmark & \checkmark\\
\bottomrule
\end{tabular}

\vspace{10pt}
\begin{columns}[T]
\begin{column}{0.30\textwidth}\centering\small
\textbf{base cases}\\ where monoid and group presentations come apart
\end{column}
\begin{column}{0.30\textwidth}\centering\small
\textbf{pure composition}\\ the Pauli group: no coset enumeration at all
\end{column}
\begin{column}{0.30\textwidth}\centering\small
\textbf{stress tests}\\ long towers, twisted relations, branching chains
\end{column}
\end{columns}
\end{frame}

\begin{frame}{Concrete presentation theorems \InPaper}
\centering
\small
\begin{tabular}{@{}lll@{}}
\toprule
Relation set & Presents & Index name in \texttt{MainTheorems} \\
\midrule
\aid{EmptyRel} over $\bot$ & trivial group & \aid{trivial-presentation} \\
\aid{TrivialRel} over any $A$ & trivial group & \aid{trivial-presentation′} \\
$T^{\,n+1} = \varepsilon$ & $\mathbb{Z}/(n{+}1)\mathbb{Z}$ & \aid{cyclic-presentation} \\
$T^{\,0} = \varepsilon$ & the \emph{monoid} $(\mathbb{N}, +, 0)$ & \aid{cyclic-monoid-presentation} \\
\aid{order}, \aid{yang-baxter} + structural & permutations of $\mathsf{Fin}\;n$ & \aid{symmetric-presentation} \\
$C_p \aidop{⋄} C_p \aidop{⋄} \aid{CommRel}$ & $\mathbb{Z}_p \times \mathbb{Z}_p$ & (\aid{H-pres}) \\
$(C_p \aidop{⋄} C_p \aidop{⋄} \aid{CommRel}) \aidop{⊕\textasciicircum} n$ & $(\mathbb{Z}_p{\times}\mathbb{Z}_p)^n$ & \aid{pauli-presentation} \\
$(C_N \aidop{⊕\textasciicircum} n) \aidop{⋄} S_n\text{-circuits} \aidop{⋄} \aid{ConjRelʷ}$ & $\mathbb{Z}_N \wr S_n$ & \aid{wreath-presentation} \\
\midrule
qubit Clifford+$T$ gate relations & $\cong$ amalgam $\langle T,M\rangle\ast_M\langle H,M\rangle$ & \aid{clifford+T-qubit-isomorphism}\\
qutrit Clifford+$T$ gate relations & $\cong$ amalgam over $M=\langle S,X,\zeta,H^2\rangle$ & \aid{clifford+T-qutrit-isomorphism}\\
$U_3(\mathbb{Z}[\frac12,i])$, S1--S14 & $\cong (\mathbb{Z}_4\wr S_3)\ast_M\langle M,K_{01}\rangle$ & \aid{U₃-presentation-isomorphism}\\
\bottomrule
\end{tabular}

\vspace{6pt}
{\small Normal-form carriers: $\mathsf{Fin}(n{+}1)$; $\mathbb{N}$; $\top\times C1\times\dots\times C(n{-}1)$;
$\mathit{NF}_1\times\mathit{NF}_2$; $\mathit{NF}_0\times\aid{amalNFC}\,C\,D$; $\mathsf{Fin}\,8\times\mathsf{Fin}\,4$ (Clifford+$T$ base).}
\end{frame}

\begin{frame}{Wreath products: signed permutations and beyond \InPaper}
\begin{columns}[c]
\begin{column}{0.57\textwidth}
$(\mathbb{Z}_N)^n$ by the $n$-fold power; $S_n$ acts by permuting coordinates,
given \alert{syntactically} as a word-valued conjugation; two finite
compatibility lemmas feed the semidirect-product theorem:
\begin{code}%
\>[0]\AgdaFunction{wreath-presentation}\AgdaSpace{}%
\AgdaSymbol{:}\AgdaSpace{}%
\AgdaSymbol{∀}\AgdaSpace{}%
\AgdaBound{n}\AgdaSpace{}%
\AgdaBound{m}\AgdaSpace{}%
\AgdaSymbol{→}\<%
\\
\>[0][@{}l@{\AgdaIndent{0}}]%
\>[2]\AgdaSymbol{(((}\AgdaInductiveConstructor{suc}\AgdaSpace{}%
\AgdaBound{m}\AgdaSpace{}%
\AgdaOperator{\AgdaDatatype{Cn,\AgdaUnderscore{}===\AgdaUnderscore{}}}\AgdaSymbol{)}\AgdaSpace{}%
\AgdaOperator{\AgdaFunction{⊕\textasciicircum{}}}\AgdaSpace{}%
\AgdaBound{n}\AgdaSymbol{)}\AgdaSpace{}%
\AgdaOperator{\AgdaDatatype{⋄}}\AgdaSpace{}%
\AgdaSymbol{(}\AgdaBound{n}\AgdaSpace{}%
\AgdaOperator{\AgdaDatatype{VRel,\AgdaUnderscore{}===\AgdaUnderscore{}}}\AgdaSymbol{)}\AgdaSpace{}%
\AgdaOperator{\AgdaDatatype{⋄}}\AgdaSpace{}%
\AgdaDatatype{ConjRelʷ}\AgdaSpace{}%
\AgdaFunction{conj}\AgdaSymbol{)}\<%
\\
\>[2][@{}l@{\AgdaIndent{0}}]%
\>[4]\AgdaOperator{\AgdaRecord{IsPresentationOf}}\AgdaSpace{}%
\AgdaSymbol{(}\AgdaFunction{Wreath.wreath-group}\AgdaSpace{}%
\AgdaBound{n}\AgdaSpace{}%
\AgdaBound{m}\AgdaSymbol{)}\<%
\end{code}
\begin{code}[hide]%
\>[0]\AgdaFunction{wreath-presentation}\AgdaSpace{}%
\AgdaBound{n}\AgdaSpace{}%
\AgdaBound{m}\AgdaSpace{}%
\AgdaSymbol{=}\AgdaSpace{}%
\AgdaFunction{Wreath.presentation}\AgdaSpace{}%
\AgdaBound{n}\AgdaSpace{}%
\AgdaBound{m}\<%
\end{code}
\begin{itemize}
  \small
  \item $N=2$: the hyperoctahedral group $B_n$ of \alert{signed permutations};
  \item general $N$: $G(N,1,n)$, monomial matrices with $N$-th roots of unity;
  \item a family of completeness results over \alert{two naturals}, all at once.
\end{itemize}
\end{column}
\begin{column}{0.40\textwidth}
\centering
{\small normal form: residues, then staircases}\\[4pt]
\begin{tikzpicture}[x=5mm,y=6mm,line width=0.7pt]
  \foreach \y in {0,1,2,3} \draw (0,\y) -- (11,\y);
  \foreach \y/\e in {0/2,1/0,2/3,3/1} {
    \node[draw,fill=softnew,font=\scriptsize,minimum width=6mm,minimum height=4mm,inner sep=1pt] at (1.6,\y) {$T^{\e}$};
  }
  % staircases (swap crossings)
  \foreach \x/\p in {4/2,5.5/1,7/2,8.5/0,10/1} {
    \fill[white] (\x-0.55,\p-0.05) rectangle (\x+0.55,\p+1.05);
    \draw (\x-0.55,\p) -- (\x+0.55,\p+1);
    \draw (\x-0.55,\p+1) -- (\x+0.55,\p);
  }
  \draw[decorate,decoration={brace,amplitude=3pt,mirror}] (0.6,-0.6) -- (2.6,-0.6) node[midway,below=2pt,font=\tiny]{$(\mathbb{Z}_N)^n$};
  \draw[decorate,decoration={brace,amplitude=3pt,mirror}] (3.3,-0.6) -- (10.7,-0.6) node[midway,below=2pt,font=\tiny]{$S_n$ (factorial digits)};
\end{tikzpicture}

\vspace{8pt}
$\begin{pmatrix} 0 & 0 & -1 & 0\\ 1 & 0 & 0 & 0\\ 0 & 0 & 0 & -1\\ 0 & -1 & 0 & 0\end{pmatrix}\in B_4$

{\scriptsize a signed permutation ($N=2$): one $\pm1$ per row and column}
\end{column}
\end{columns}
\end{frame}

\begin{frame}{Pauli groups in all prime dimensions: pure composition \InPaper}
\begin{columns}[c]
\begin{column}{0.56\textwidth}
\begin{code}%
\>[0]\AgdaFunction{pauli-presentation}\AgdaSpace{}%
\AgdaSymbol{:}\AgdaSpace{}%
\AgdaSymbol{∀}\AgdaSpace{}%
\AgdaBound{m}\AgdaSpace{}%
\AgdaSymbol{(}\AgdaBound{prime}\AgdaSpace{}%
\AgdaSymbol{:}\AgdaSpace{}%
\AgdaRecord{Prime}\AgdaSpace{}%
\AgdaSymbol{(}\AgdaInductiveConstructor{2+}\AgdaSpace{}%
\AgdaBound{m}\AgdaSymbol{))}\AgdaSpace{}%
\AgdaBound{n}\AgdaSpace{}%
\AgdaSymbol{→}\<%
\\
\>[0][@{}l@{\AgdaIndent{0}}]%
\>[2]\AgdaSymbol{(}\AgdaFunction{Γ-Pauli₁}\AgdaSpace{}%
\AgdaBound{m}\AgdaSpace{}%
\AgdaBound{prime}\AgdaSpace{}%
\AgdaOperator{\AgdaFunction{⊕\textasciicircum{}}}\AgdaSpace{}%
\AgdaBound{n}\AgdaSymbol{)}\AgdaSpace{}%
\AgdaOperator{\AgdaRecord{IsPresentationOf}}\AgdaSpace{}%
\AgdaSymbol{(}\AgdaFunction{Pauli-group}\AgdaSpace{}%
\AgdaBound{m}\AgdaSpace{}%
\AgdaBound{prime}\AgdaSpace{}%
\AgdaBound{n}\AgdaSymbol{)}\<%
\end{code}
\begin{code}[hide]%
\>[0]\AgdaFunction{pauli-presentation}\AgdaSpace{}%
\AgdaBound{m}\AgdaSpace{}%
\AgdaBound{prime}\AgdaSpace{}%
\AgdaSymbol{=}\AgdaSpace{}%
\AgdaFunction{Pauli.Pauli-presentation}\AgdaSpace{}%
\AgdaBound{m}\AgdaSpace{}%
\AgdaBound{prime}\<%
\end{code}
\begin{itemize}
  \small
  \item $p = 2+m$ prime: a \alert{qupit};
  \item \aid{Γ-Pauli₁}: direct-product join of two order-$p$ cyclic
        presentations ($X$, $Z$; they commute modulo phases);
  \item $n$ qupits: the $n$-fold power;
  \item the proof is \alert{three construction theorems composed} --- three
        lines of Agda, no coset enumeration;
  \item normal form: $X^{a_1}Z^{b_1}\cdots X^{a_n}Z^{b_n}$,
        $(a_i,b_i)\in\mathsf{Fin}\,p\times\mathsf{Fin}\,p$.
\end{itemize}
\end{column}
\begin{column}{0.41\textwidth}
\centering
\begin{tikzpicture}[node distance=5mm,>={Stealth[length=1.8mm]},
  bx/.style={draw=paperC,fill=softpaper,rounded corners=2pt,font=\scriptsize,align=center,minimum height=8mm,text width=26mm}]
  \node[bx] (cx) {$\langle X\mid X^p\rangle$\\cyclic};
  \node[bx,right=4mm of cx,text width=18mm] (cz) {$\langle Z\mid Z^p\rangle$\\cyclic};
  \node[bx,below=8mm of $(cx)!0.5!(cz)$] (p1) {$\mathbb{Z}_p\times\mathbb{Z}_p$\\binary product};
  \node[bx,below=8mm of p1] (pn) {$(\mathbb{Z}_p\times\mathbb{Z}_p)^n$\\$n$-fold product};
  \draw[->,thick] (cx) -- (p1); \draw[->,thick] (cz) -- (p1); \draw[->,thick] (p1) -- (pn);
\end{tikzpicture}

\vspace{6pt}
\begin{quantikz}[row sep=0.18cm,column sep=0.25cm,font=\scriptsize]
 & \gate{X^{a_3}} & \gate{Z^{b_3}} & \\
 & \gate{X^{a_2}} & \gate{Z^{b_2}} & \\
 & \gate{X^{a_1}} & \gate{Z^{b_1}} &
\end{quantikz}
\end{column}
\end{columns}
\end{frame}

\begin{frame}{Single-qubit Clifford+$T$ as an amalgamated free product \InPaper}
\begin{columns}[c]
\begin{column}{0.50\textwidth}
Generators $T, H, S, \omega$; $X := HSSH$:
\[
  \omega^8 = \varepsilon,\;
  T^2 = S,\;
  S^4 = \varepsilon,\;
  H^2 = \varepsilon,\]
\[
  (SH)^3 = \omega,\;
  (TX)^2 = \omega,\;
  \omega\ \text{central}.
\]
Verified structure:
\[
  \mathrm{Clifford}{+}T \;\cong\;
  \langle T, M \rangle \,\ast_{M}\, \langle H, M \rangle,\quad
  M = \langle X, S, \omega \rangle .
\]
\begin{code}%
\>[0]\AgdaFunction{clifford+T-qubit-isomorphism}\AgdaSpace{}%
\AgdaSymbol{:}\AgdaSpace{}%
\AgdaRecord{IsMonoidIsomorphism}\<%
\\
\>[0][@{}l@{\AgdaIndent{0}}]%
\>[2]\AgdaSymbol{(}\AgdaFunction{rawMonoid}\AgdaSpace{}%
\AgdaSymbol{(}\AgdaFunction{•-ε-monoid}\AgdaSpace{}%
\AgdaOperator{\AgdaDatatype{\AgdaUnderscore{}===\AgdaUnderscore{}}}\AgdaSymbol{))}\AgdaSpace{}%
\AgdaSymbol{(}\AgdaFunction{rawMonoid}\AgdaSpace{}%
\AgdaSymbol{(}\AgdaFunction{•-ε-monoid}\AgdaSpace{}%
\AgdaFunction{amalPres}\AgdaSymbol{))}\AgdaSpace{}%
\AgdaSymbol{(}\AgdaFunction{f}\AgdaSpace{}%
\AgdaOperator{\AgdaFunction{ʷ}}\AgdaSymbol{)}\<%
\end{code}
\begin{code}[hide]%
\>[0]\AgdaFunction{clifford+T-qubit-isomorphism}\AgdaSpace{}%
\AgdaSymbol{=}\AgdaSpace{}%
\AgdaFunction{CliffordT1.CliffordT1.CliffordT1-isomorphism}\<%
\end{code}
\end{column}
\begin{column}{0.47\textwidth}
\centering
\begin{tikzpicture}[>={Stealth[length=1.8mm]},
  bx/.style={draw=paperC,fill=softpaper,rounded corners=2pt,font=\scriptsize,align=center,minimum height=7mm},
  lab/.style={font=\tiny,align=center,color=newC}]
  \node[bx] (b) at (0,0) {$\langle\omega\rangle\times\langle S\rangle$\\{\tiny radices $8\cdot 4$}};
  \node[bx] (m) at (0,1.6) {$M=\langle X,S,\omega\rangle$};
  \node[bx] (t) at (-1.9,3.3) {$\langle T,M\rangle$};
  \node[bx] (h) at (1.9,3.3) {$\langle H,M\rangle$\\{\tiny Clifford}};
  \node[bx,draw=newC,fill=softnew] (a) at (0,5.0) {$\langle T,M\rangle\ast_M\langle H,M\rangle$};
  \draw[->,thick] (b) -- node[right,lab]{2 cosets: $\varepsilon$, $X$} (m);
  \draw[->,thick] (m) -- node[left,lab]{2 cosets:\\$\varepsilon$, $T$} (t);
  \draw[->,thick] (m) -- node[right,lab]{3 cosets:\\$\varepsilon$, $H$, $HS$} (h);
  \draw[->,thick] (t) -- (a); \draw[->,thick] (h) -- (a);
\end{tikzpicture}

\vspace{2pt}
{\small 59 well-definedness obligations, every one \aid{(by-equal-nf Eq.refl) , Eq.refl}.}
\end{column}
\end{columns}
\end{frame}

\begin{frame}{The amalgam rediscovers Matsumoto--Amano \InPaper}
With $C=\{T\}$ and $D=\{H, HS\}$, the section of the alternating normal form spells, left to right,
\[
  \mathcal{M}\;(\varepsilon \mid H \mid HS)\;(TH \mid THS)^{*}\;(\varepsilon \mid T)
\]
and read right to left:
\[
  (\varepsilon \mid T)\;(HT \mid SHT)^{*}\;(\varepsilon \mid H \mid SH)\;\mathcal{M}
\]
--- the \alert{Matsumoto--Amano normal form}
$(\varepsilon\mid T)(HT\mid SHT)^*\,\mathcal{C}$, with its Clifford factor split
into a coset representative and an element of $M$.
\vspace{6pt}

\centering
\begin{tikzpicture}[x=9mm,
  s/.style={draw,minimum height=6mm,minimum width=8.2mm,font=\scriptsize,inner sep=1pt}]
  \node[s,fill=gray!15] at (0,0) {$\mathcal{M}$};
  \node[s,fill=softpaper] at (1,0) {$HS$};
  \node[s,fill=soft] at (2,0) {$T$}; \node[s,fill=softpaper] at (3,0) {$H$};
  \node[s,fill=soft] at (4,0) {$T$}; \node[s,fill=softpaper] at (5,0) {$HS$};
  \node[s,fill=soft] at (6,0) {$T$}; \node[s,fill=softpaper] at (7,0) {$H$};
  \node[s,fill=soft] at (8,0) {$T$};
  \node[font=\scriptsize,anchor=west] at (8.7,0) {a typical normal form: $T$-count 4};
\end{tikzpicture}

\vspace{4pt}
{\small The syllable structure falls out of the group theory alone (as observed in the
Bian--Selinger line).}
\end{frame}

\begin{frame}{Qutrit Clifford+$T$ and $U_3(\mathbb{Z}[\frac12,i])$ \InPaper}
\begin{columns}[T]
\begin{column}{0.50\textwidth}
\textbf{Qutrit} ($\zeta$ a 9th root; $X,Z$ abbreviations):
\[
  \zeta^9 = \varepsilon,\;
  S^3 = \zeta^6,\;
  H^4 = \varepsilon,\;
  T^3 = Z,\;
  (SH)^3 = \varepsilon,
\]
\[
  (THH)^2 = \varepsilon,\;
  HHSHHS = SHHSHH,\;
  TS = ST,
\]
\[
  TX = \zeta^3 S^2 X\, T,\;\;
  \zeta\ \text{central}.
\]
Tower $\langle\zeta\rangle \to \langle S,\zeta\rangle\ (3) \to \langle S,X,\zeta\rangle\
(\alert{9}) \to M\ (2)$, then two extensions over $M = \langle S,X,\zeta,H^2\rangle$.
Normal form = Glaudell--Ross--Taylor's $T$-optimal form.
The nine-coset level: \alert{72} obligations, all \aid{refl}.
\end{column}
\begin{column}{0.47\textwidth}
\textbf{$U_3(\mathbb{Z}[\frac12,i])$} (Bian--Selinger, $n=3$):
\[
  U_3(\mathbb{Z}[\tfrac12,i]) \cong
  (\mathbb{Z}_4 \wr S_3) \ast_{M} \langle M, K_{01} \rangle,
\]
\[ M = (\mathbb{Z}_4 \wr S_2) \times \mathbb{Z}_4 . \]
The left factor is taken \alert{verbatim} from the wreath development.
\[
  \mathcal{M}\,(\varepsilon \mid K_{01} \mid K_{01} i_0)\,
  \bigl((X_{12} \mid X_{12}X_{01})(K_{01} \mid K_{01} i_0)\bigr)^{*}
\]
\[ (\varepsilon \mid X_{12} \mid X_{12}X_{01}) \]
\begin{newdevbox}
\small A \alert{new canonical form} for $U_3(\mathbb{Z}[\frac12,i])$ with
Matsumoto--Amano-style syllables --- for free from the amalgam.
\end{newdevbox}
\end{column}
\end{columns}
\end{frame}

\begin{frame}{Size and cost \InPaper}
\begin{columns}[c]
\begin{column}{0.48\textwidth}
\centering
\begin{tikzpicture}
\begin{axis}[xbar stacked, width=7.0cm, height=3.6cm, xmin=0, xmax=11500,
  symbolic y coords={examples,core}, ytick=data,
  xlabel={\scriptsize lines of Agda}, x tick label style={font=\scriptsize},
  y tick label style={font=\small}, bar width=16pt, enlarge y limits=0.45,
  legend style={font=\tiny,at={(0.5,-0.42)},anchor=north,legend columns=2,draw=none,fill=none},
  scaled x ticks=false, axis lines*=left]
\addplot[fill=accent!70] coordinates {(1000,core) (0,examples)};
\addplot[fill=accent!40] coordinates {(4300,core) (0,examples)};
\addplot[fill=accent!20] coordinates {(600,core) (0,examples)};
\addplot[fill=gray!30] coordinates {(5100,core) (0,examples)};
\addplot[fill=paperC!60] coordinates {(0,core) (4300,examples)};
\addplot[fill=paperC!25] coordinates {(0,core) (6700,examples)};
\legend{presentation core,product constructions,solvers,rest of core,Clifford+$T_3$ \& $U_3$,other examples}
\end{axis}
\end{tikzpicture}

{\scriptsize approximate split of $\approx$22\,000 lines (74 files, Sept.\ 2026)}
\end{column}
\begin{column}{0.49\textwidth}
\begin{itemize}
  \setlength{\itemsep}{4pt}
  \item about \textbf{22\,000 lines} in 74 files: half reusable core, half examples;
  \item of 7\,000 lines under \texttt{Presentation}: $\approx$1\,000 core proper,
        4\,300 the five product constructions (amalgamation alone 1\,800),
        600 the solvers;
  \item whole development from scratch, Agda 2.8.0: \textbf{1\,min\,42\,s}, peak
        3.5\,GB; a third of it the amalgamated product (31\,s), the three
        amalgamation studies 19\,s, every other module $<4$\,s.
\end{itemize}
\begin{keybox}
\small \textbf{Limitations stated in the paper}: amalgamation studies are
syntax-to-syntax (no matrix semantics); continuous-parameter gate sets and
infinite-index steps are outside the engine.
\end{keybox}
\end{column}
\end{columns}
\end{frame}

\begin{code}[hide]%
\>[0]\AgdaKeyword{module}\AgdaSpace{}%
\AgdaModule{slides.Tsinghua-PL-seminar.sections.newcore}\AgdaSpace{}%
\AgdaKeyword{where}\<%
\\
%
\\[\AgdaEmptyExtraSkip]%
\>[0]\AgdaKeyword{open}\AgdaSpace{}%
\AgdaKeyword{import}\AgdaSpace{}%
\AgdaModule{Data.Nat}\AgdaSpace{}%
\AgdaKeyword{using}\AgdaSpace{}%
\AgdaSymbol{(}\AgdaDatatype{ℕ}\AgdaSymbol{)}\<%
\\
\>[0]\AgdaKeyword{open}\AgdaSpace{}%
\AgdaKeyword{import}\AgdaSpace{}%
\AgdaModule{Data.Product}\AgdaSpace{}%
\AgdaKeyword{using}\AgdaSpace{}%
\AgdaSymbol{(}\AgdaRecord{Σ}\AgdaSpace{}%
\AgdaSymbol{;}\AgdaSpace{}%
\AgdaOperator{\AgdaFunction{\AgdaUnderscore{}×\AgdaUnderscore{}}}\AgdaSpace{}%
\AgdaSymbol{;}\AgdaSpace{}%
\AgdaOperator{\AgdaInductiveConstructor{\AgdaUnderscore{},\AgdaUnderscore{}}}\AgdaSymbol{)}\<%
\\
\>[0]\AgdaKeyword{open}\AgdaSpace{}%
\AgdaKeyword{import}\AgdaSpace{}%
\AgdaModule{Level}\AgdaSpace{}%
\AgdaKeyword{using}\AgdaSpace{}%
\AgdaSymbol{(}\AgdaFunction{0ℓ}\AgdaSymbol{)}\<%
\\
\>[0]\AgdaKeyword{open}\AgdaSpace{}%
\AgdaKeyword{import}\AgdaSpace{}%
\AgdaModule{Relation.Binary}\AgdaSpace{}%
\AgdaKeyword{using}\AgdaSpace{}%
\AgdaSymbol{(}\AgdaFunction{Rel}\AgdaSymbol{)}\<%
\\
\>[0]\AgdaKeyword{open}\AgdaSpace{}%
\AgdaKeyword{import}\AgdaSpace{}%
\AgdaModule{Relation.Nullary}\AgdaSpace{}%
\AgdaKeyword{using}\AgdaSpace{}%
\AgdaSymbol{(}\AgdaOperator{\AgdaFunction{¬\AgdaUnderscore{}}}\AgdaSymbol{)}\<%
\\
\>[0]\AgdaKeyword{open}\AgdaSpace{}%
\AgdaKeyword{import}\AgdaSpace{}%
\AgdaModule{Relation.Binary.PropositionalEquality}\AgdaSpace{}%
\AgdaKeyword{using}\AgdaSpace{}%
\AgdaSymbol{(}\AgdaOperator{\AgdaFunction{\AgdaUnderscore{}≢\AgdaUnderscore{}}}\AgdaSymbol{)}\<%
\\
\>[0]\AgdaKeyword{open}\AgdaSpace{}%
\AgdaKeyword{import}\AgdaSpace{}%
\AgdaModule{Algebra.Bundles}\AgdaSpace{}%
\AgdaKeyword{using}\AgdaSpace{}%
\AgdaSymbol{(}\AgdaRecord{Monoid}\AgdaSymbol{)}\<%
\\
%
\\[\AgdaEmptyExtraSkip]%
\>[0]\AgdaKeyword{open}\AgdaSpace{}%
\AgdaKeyword{import}\AgdaSpace{}%
\AgdaModule{Notations}\AgdaSpace{}%
\AgdaKeyword{using}\AgdaSpace{}%
\AgdaSymbol{(}\AgdaInductiveConstructor{₁₊}\AgdaSpace{}%
\AgdaSymbol{;}\AgdaSpace{}%
\AgdaInductiveConstructor{₂₊}\AgdaSymbol{)}\<%
\\
\>[0]\AgdaKeyword{open}\AgdaSpace{}%
\AgdaKeyword{import}\AgdaSpace{}%
\AgdaModule{Word.Base}\AgdaSpace{}%
\AgdaKeyword{using}\AgdaSpace{}%
\AgdaSymbol{(}\AgdaDatatype{Word}\AgdaSpace{}%
\AgdaSymbol{;}\AgdaSpace{}%
\AgdaFunction{WRel}\AgdaSpace{}%
\AgdaSymbol{;}\AgdaSpace{}%
\AgdaOperator{\AgdaInductiveConstructor{[\AgdaUnderscore{}]ʷ}}\AgdaSpace{}%
\AgdaSymbol{;}\AgdaSpace{}%
\AgdaInductiveConstructor{ε}\AgdaSpace{}%
\AgdaSymbol{;}\AgdaSpace{}%
\AgdaOperator{\AgdaInductiveConstructor{\AgdaUnderscore{}•\AgdaUnderscore{}}}\AgdaSpace{}%
\AgdaSymbol{;}\AgdaSpace{}%
\AgdaFunction{wmap}\AgdaSymbol{)}\<%
\\
%
\\[\AgdaEmptyExtraSkip]%
\>[0]\AgdaKeyword{postulate}\<%
\\
\>[0][@{}l@{\AgdaIndent{0}}]%
\>[2]\AgdaPostulate{proof-elided}\AgdaSpace{}%
\AgdaSymbol{:}\AgdaSpace{}%
\AgdaSymbol{∀}\AgdaSpace{}%
\AgdaSymbol{\{}\AgdaBound{ℓ}\AgdaSymbol{\}}\AgdaSpace{}%
\AgdaSymbol{\{}\AgdaBound{A}\AgdaSpace{}%
\AgdaSymbol{:}\AgdaSpace{}%
\AgdaPrimitive{Set}\AgdaSpace{}%
\AgdaBound{ℓ}\AgdaSymbol{\}}\AgdaSpace{}%
\AgdaSymbol{→}\AgdaSpace{}%
\AgdaBound{A}\<%
\\
%
\>[2]\AgdaPostulate{Gate}\AgdaSpace{}%
\AgdaSymbol{:}\AgdaSpace{}%
\AgdaDatatype{ℕ}\AgdaSpace{}%
\AgdaSymbol{→}\AgdaSpace{}%
\AgdaPrimitive{Set}\<%
\\
%
\>[2]\AgdaPostulate{X}\AgdaSpace{}%
\AgdaSymbol{:}\AgdaSpace{}%
\AgdaPrimitive{Set}\<%
\\
%
\\[\AgdaEmptyExtraSkip]%
\>[0]\AgdaKeyword{private}\AgdaSpace{}%
\AgdaKeyword{variable}\<%
\\
\>[0][@{}l@{\AgdaIndent{0}}]%
\>[2]\AgdaGeneralizable{n}\AgdaSpace{}%
\AgdaSymbol{:}\AgdaSpace{}%
\AgdaDatatype{ℕ}\<%
\end{code}
%% Part 6a — what happened after the submission; the core library.
\NewPart
\section{New since the paper}

\begin{frame}{What happened after the submission \New}
\begin{columns}[c]
\begin{column}{0.58\textwidth}
\centering
\begin{tikzpicture}[x=5.3mm,y=6.6mm,>={Stealth[length=1.6mm]},
  lb/.style={font=\scriptsize,anchor=west,align=left}]
  % trunk
  \draw[thick,paperC] (0,0) -- (3,0);
  \node[cm=paperC,label={[font=\scriptsize,align=center]below:tag\\[-2pt]\texttt{popl-revision1}}] at (0,0) {};
  \node[font=\tiny,gray] at (0,0.45) {Jul};
  \draw[thick,accent] (3,0) -- (6,0);
  \node[cm=accent,label={[font=\scriptsize]below:\texttt{main}}] at (6,0) {};
  \node[font=\tiny,gray] at (6,0.45) {Sep 2};
  % branches
  \draw[thick,newC] (6,0) .. controls (7,0) and (7,2.4) .. (9.5,2.4) node[cm=newC] {} node[lb,right=2pt]{\texttt{popl}\\[-1pt]path sums};
  \draw[thick,newC] (6,0) .. controls (7,0) and (7,0.8) .. (9.5,0.8) node[cm=newC] {} node[lb,right=2pt]{\texttt{v-office}\\[-1pt]unitary groups};
  \draw[thick,newC] (6,0) .. controls (7,0) and (7,-0.8) .. (9.5,-0.8) node[cm=newC] {} node[lb,right=2pt]{\texttt{qupit}\\[-1pt]minimality, CNOT-dihedral};
  \draw[thick,newC] (6,0) .. controls (7,0) and (7,-2.4) .. (9.5,-2.4) node[cm=newC] {} node[lb,right=2pt]{\texttt{v-office-Cli-CH}\\[-1pt]Real-Clifford+$CH$};
  \node[font=\scriptsize,align=center,accent] at (4.5,-1.6) {qupit Clifford,\\$Sp(2n,\mathbb{Z}_p)$,\\exact Clifford};
  \draw[->,gray] (4.5,-1.0) -- (4.5,-0.15);
  \node[font=\tiny,gray] at (9.5,3.1) {late Sep -- Oct 2026};
\end{tikzpicture}
\end{column}
\begin{column}{0.40\textwidth}
\centering
\begin{tikzpicture}
\begin{axis}[xbar, width=5.4cm, height=5.2cm, xmin=0, xmax=230,
  symbolic y coords={v-office-Cli-CH,qupit,v-office,popl,main,tag},
  ytick=data, y tick label style={font=\scriptsize\ttfamily},
  x tick label style={font=\scriptsize}, xlabel={\scriptsize thousand lines of Agda},
  nodes near coords, nodes near coords style={font=\tiny},
  bar width=7pt, enlarge y limits=0.12, axis lines*=left]
\addplot[fill=paperC!70,draw=paperC] coordinates {(22.2,tag)};
\addplot[fill=newC!60,draw=newC] coordinates
  {(112.8,main) (193.9,popl) (155.2,v-office) (139.9,qupit) (187.1,v-office-Cli-CH)};
\end{axis}
\end{tikzpicture}

{\scriptsize \texttt{*.agda} outside \texttt{paper/}; 74 files at the tag,
634 on \texttt{v-office-Cli-CH}}
\end{column}
\end{columns}
\end{frame}

\begin{frame}{Research papers formalized since the submission \New}
\centering\scriptsize
\renewcommand{\arraystretch}{1.18}
\begin{tabular}{@{}p{4.0cm}p{2.7cm}p{2.6cm}p{4.0cm}@{}}
\toprule
paper formalized & gate set / group & semantics & status\\
\midrule
Bian, Li, Ross, van de Wetering, Zhao (QPL'26): qupit Clifford, all odd primes
  & $H,S,CZ$; $\text{Pauli}\rtimes Sp(2n,\mathbb{Z}_p)$ & symplectic maps, Pauli action & \textbf{complete}, every $n$, every odd prime $p$\\
Selinger (LMCS 2015): $n$-qubit Clifford & $\omega,H,S,CZ$ & $Sp(2n,2)$ + phases & \textbf{complete} (modulo scalars; scalars as extension)\\
Bian--Selinger (QPL'21): $U_n(\mathbb{Z}[\frac12,i])$ & two-level $X,K,i$ & \alert{matrices} over $\mathbb{Z}[\frac12,i]$ & \textbf{complete}, every $n$\\
Greylyn (2014); Bian--Selinger (QPL'22): 2-qubit Clifford+$T$ & $\omega,H_i,S_i,T_i,CZ$ & \alert{matrices} over $\mathbb{Z}[\frac1{\sqrt2},i]$ & \textbf{complete} ($n\le4$; faithful into $U_4$)\\
Bian--Selinger (QPL'23): 3-qubit Clifford+$CS$ & $S_i, CS, K_i, i$ & matrices over $\mathbb{Z}[\frac12,i]$ & in progress (steps 1 and 4 of the R--S chain)\\
Clément (2026): Real-Clifford+$CH$ & $H,Z,CZ,CH$ & \alert{matrices} over $\mathbb{Z}[\sqrt2]$ & sound; complete on $\le1$ qubit; $\ge2$ modulo Thm 4.4 (imported by the paper too) + Lemma 8.8 at 3--4 qubits\\
Amy--Chen--Ross (QPL'17): CNOT-dihedral & $\omega,X,T,CNOT$ & phase polynomials & sound; complete on 1 qubit; $n$ modulo imported lemmas\\
Amy (QPL'18): path sums & $H,S,CZ$ / $H,CNOT,R_k$ & exact sums over $\mathbb{Z}[\zeta]$ & Clifford equivalence in \alert{polynomial cost}\\
\bottomrule
\end{tabular}

\vspace{4pt}
{\scriptsize plus, in the core: 0-ary gates (scalars), independence of axioms, central and
factor-set extensions, generic circuit coset towers}
\end{frame}

\begin{frame}{Scalars: 0-ary gates, and a rule that became a theorem \New}
\begin{columns}[c]
\begin{column}{0.55\textwidth}
\small Global phases ($\omega$, $\zeta$, \ldots) live on no wire, at every width:
\begin{code}%
\>[0]\AgdaKeyword{data}\AgdaSpace{}%
\AgdaDatatype{Gen}\AgdaSpace{}%
\AgdaSymbol{:}\AgdaSpace{}%
\AgdaDatatype{ℕ}\AgdaSpace{}%
\AgdaSymbol{→}\AgdaSpace{}%
\AgdaPrimitive{Set}\AgdaSpace{}%
\AgdaKeyword{where}\<%
\\
\>[0][@{}l@{\AgdaIndent{0}}]%
\>[2]\AgdaInductiveConstructor{gate₀}\AgdaSpace{}%
\AgdaSymbol{:}\AgdaSpace{}%
\AgdaPostulate{Gate}\AgdaSpace{}%
\AgdaNumber{0}\AgdaSpace{}%
\AgdaSymbol{→}\AgdaSpace{}%
\AgdaDatatype{Gen}\AgdaSpace{}%
\AgdaGeneralizable{n}\<%
\\
%
\>[2]\AgdaInductiveConstructor{gate₁}\AgdaSpace{}%
\AgdaSymbol{:}\AgdaSpace{}%
\AgdaPostulate{Gate}\AgdaSpace{}%
\AgdaNumber{1}\AgdaSpace{}%
\AgdaSymbol{→}\AgdaSpace{}%
\AgdaDatatype{Gen}\AgdaSpace{}%
\AgdaSymbol{(}\AgdaInductiveConstructor{₁₊}\AgdaSpace{}%
\AgdaGeneralizable{n}\AgdaSymbol{)}\<%
\\
%
\>[2]\AgdaInductiveConstructor{gate₂}\AgdaSpace{}%
\AgdaSymbol{:}\AgdaSpace{}%
\AgdaPostulate{Gate}\AgdaSpace{}%
\AgdaNumber{2}\AgdaSpace{}%
\AgdaSymbol{→}\AgdaSpace{}%
\AgdaDatatype{Gen}\AgdaSpace{}%
\AgdaSymbol{(}\AgdaInductiveConstructor{₂₊}\AgdaSpace{}%
\AgdaGeneralizable{n}\AgdaSymbol{)}\<%
\\
%
\>[2]\AgdaOperator{\AgdaInductiveConstructor{\AgdaUnderscore{}↥}}%
\>[7]\AgdaSymbol{:}\AgdaSpace{}%
\AgdaDatatype{Gen}\AgdaSpace{}%
\AgdaGeneralizable{n}\AgdaSpace{}%
\AgdaSymbol{→}\AgdaSpace{}%
\AgdaDatatype{Gen}\AgdaSpace{}%
\AgdaSymbol{(}\AgdaInductiveConstructor{₁₊}\AgdaSpace{}%
\AgdaGeneralizable{n}\AgdaSymbol{)}\<%
\end{code}
\begin{code}[hide]%
\>[0]\AgdaFunction{Circuit}\AgdaSpace{}%
\AgdaSymbol{:}\AgdaSpace{}%
\AgdaDatatype{ℕ}\AgdaSpace{}%
\AgdaSymbol{→}\AgdaSpace{}%
\AgdaPrimitive{Set}\<%
\\
\>[0]\AgdaFunction{Circuit}\AgdaSpace{}%
\AgdaBound{n}\AgdaSpace{}%
\AgdaSymbol{=}\AgdaSpace{}%
\AgdaDatatype{Word}\AgdaSpace{}%
\AgdaSymbol{(}\AgdaDatatype{Gen}\AgdaSpace{}%
\AgdaBound{n}\AgdaSymbol{)}\<%
\\
\>[0]\AgdaFunction{CRel}\AgdaSpace{}%
\AgdaSymbol{:}\AgdaSpace{}%
\AgdaDatatype{ℕ}\AgdaSpace{}%
\AgdaSymbol{→}\AgdaSpace{}%
\AgdaPrimitive{Set₁}\<%
\\
\>[0]\AgdaFunction{CRel}\AgdaSpace{}%
\AgdaBound{n}\AgdaSpace{}%
\AgdaSymbol{=}\AgdaSpace{}%
\AgdaFunction{Rel}\AgdaSpace{}%
\AgdaSymbol{(}\AgdaFunction{Circuit}\AgdaSpace{}%
\AgdaBound{n}\AgdaSymbol{)}\AgdaSpace{}%
\AgdaFunction{0ℓ}\<%
\\
\>[0]\AgdaOperator{\AgdaFunction{\AgdaUnderscore{}↑}}\AgdaSpace{}%
\AgdaSymbol{:}\AgdaSpace{}%
\AgdaFunction{Circuit}\AgdaSpace{}%
\AgdaGeneralizable{n}\AgdaSpace{}%
\AgdaSymbol{→}\AgdaSpace{}%
\AgdaFunction{Circuit}\AgdaSpace{}%
\AgdaSymbol{(}\AgdaInductiveConstructor{₁₊}\AgdaSpace{}%
\AgdaGeneralizable{n}\AgdaSymbol{)}\<%
\\
\>[0]\AgdaOperator{\AgdaFunction{\AgdaUnderscore{}↑}}\AgdaSpace{}%
\AgdaSymbol{=}\AgdaSpace{}%
\AgdaFunction{wmap}\AgdaSpace{}%
\AgdaOperator{\AgdaInductiveConstructor{\AgdaUnderscore{}↥}}\<%
\\
\>[0]\AgdaKeyword{infix}\AgdaSpace{}%
\AgdaNumber{4}\AgdaSpace{}%
\AgdaOperator{\AgdaPostulate{\AgdaUnderscore{}SRel,\AgdaUnderscore{}===\AgdaUnderscore{}}}\AgdaSpace{}%
\AgdaOperator{\AgdaDatatype{\AgdaUnderscore{}VRel,\AgdaUnderscore{}===\AgdaUnderscore{}}}\AgdaSpace{}%
\AgdaOperator{\AgdaPostulate{\AgdaUnderscore{}≈\AgdaUnderscore{}}}\<%
\\
\>[0]\AgdaKeyword{postulate}\<%
\\
\>[0][@{}l@{\AgdaIndent{0}}]%
\>[2]\AgdaOperator{\AgdaPostulate{\AgdaUnderscore{}SRel,\AgdaUnderscore{}===\AgdaUnderscore{}}}\AgdaSpace{}%
\AgdaSymbol{:}\AgdaSpace{}%
\AgdaSymbol{(}\AgdaBound{n}\AgdaSpace{}%
\AgdaSymbol{:}\AgdaSpace{}%
\AgdaDatatype{ℕ}\AgdaSymbol{)}\AgdaSpace{}%
\AgdaSymbol{→}\AgdaSpace{}%
\AgdaFunction{WRel}\AgdaSpace{}%
\AgdaSymbol{(}\AgdaDatatype{Gen}\AgdaSpace{}%
\AgdaBound{n}\AgdaSymbol{)}\<%
\\
%
\>[2]\AgdaOperator{\AgdaPostulate{\AgdaUnderscore{}≈\AgdaUnderscore{}}}\AgdaSpace{}%
\AgdaSymbol{:}\AgdaSpace{}%
\AgdaFunction{Circuit}\AgdaSpace{}%
\AgdaGeneralizable{n}\AgdaSpace{}%
\AgdaSymbol{→}\AgdaSpace{}%
\AgdaFunction{Circuit}\AgdaSpace{}%
\AgdaGeneralizable{n}\AgdaSpace{}%
\AgdaSymbol{→}\AgdaSpace{}%
\AgdaPrimitive{Set}\<%
\\
\>[0]\AgdaKeyword{private}\AgdaSpace{}%
\AgdaKeyword{variable}\<%
\\
\>[0][@{}l@{\AgdaIndent{0}}]%
\>[2]\AgdaGeneralizable{w}\AgdaSpace{}%
\AgdaGeneralizable{v}\AgdaSpace{}%
\AgdaSymbol{:}\AgdaSpace{}%
\AgdaFunction{Circuit}\AgdaSpace{}%
\AgdaGeneralizable{n}\<%
\end{code}
\vspace{-6pt}
\small Four structural rules now (\Branch{qupit}, \texttt{Circuit/Base.agda}):
\begin{code}%
\>[0]\AgdaKeyword{data}\AgdaSpace{}%
\AgdaOperator{\AgdaDatatype{\AgdaUnderscore{}VRel,\AgdaUnderscore{}===\AgdaUnderscore{}}}\AgdaSpace{}%
\AgdaSymbol{:}\AgdaSpace{}%
\AgdaSymbol{(}\AgdaBound{n}\AgdaSpace{}%
\AgdaSymbol{:}\AgdaSpace{}%
\AgdaDatatype{ℕ}\AgdaSymbol{)}\AgdaSpace{}%
\AgdaSymbol{→}\AgdaSpace{}%
\AgdaFunction{CRel}\AgdaSpace{}%
\AgdaBound{n}\AgdaSpace{}%
\AgdaKeyword{where}\<%
\\
\>[0][@{}l@{\AgdaIndent{0}}]%
\>[2]\AgdaInductiveConstructor{srel}%
\>[8]\AgdaSymbol{:}\AgdaSpace{}%
\AgdaGeneralizable{n}\AgdaSpace{}%
\AgdaOperator{\AgdaPostulate{SRel,}}\AgdaSpace{}%
\AgdaGeneralizable{w}\AgdaSpace{}%
\AgdaOperator{\AgdaPostulate{===}}\AgdaSpace{}%
\AgdaGeneralizable{v}\AgdaSpace{}%
\AgdaSymbol{→}\AgdaSpace{}%
\AgdaGeneralizable{n}\AgdaSpace{}%
\AgdaOperator{\AgdaDatatype{VRel,}}\AgdaSpace{}%
\AgdaGeneralizable{w}\AgdaSpace{}%
\AgdaOperator{\AgdaDatatype{===}}\AgdaSpace{}%
\AgdaGeneralizable{v}\<%
\\
%
\>[2]\AgdaInductiveConstructor{cong↑}\AgdaSpace{}%
\AgdaSymbol{:}\AgdaSpace{}%
\AgdaGeneralizable{n}\AgdaSpace{}%
\AgdaOperator{\AgdaDatatype{VRel,}}\AgdaSpace{}%
\AgdaGeneralizable{w}\AgdaSpace{}%
\AgdaOperator{\AgdaDatatype{===}}\AgdaSpace{}%
\AgdaGeneralizable{v}\AgdaSpace{}%
\AgdaSymbol{→}\AgdaSpace{}%
\AgdaSymbol{(}\AgdaInductiveConstructor{₁₊}\AgdaSpace{}%
\AgdaGeneralizable{n}\AgdaSymbol{)}\AgdaSpace{}%
\AgdaOperator{\AgdaDatatype{VRel,}}\AgdaSpace{}%
\AgdaGeneralizable{w}\AgdaSpace{}%
\AgdaOperator{\AgdaFunction{↑}}\AgdaSpace{}%
\AgdaOperator{\AgdaDatatype{===}}\AgdaSpace{}%
\AgdaGeneralizable{v}\AgdaSpace{}%
\AgdaOperator{\AgdaFunction{↑}}\<%
\\
%
\>[2]\AgdaInductiveConstructor{comm₁}\AgdaSpace{}%
\AgdaSymbol{:}%
\>[194I]\AgdaSymbol{(}\AgdaBound{h}\AgdaSpace{}%
\AgdaSymbol{:}\AgdaSpace{}%
\AgdaPostulate{Gate}\AgdaSpace{}%
\AgdaNumber{1}\AgdaSymbol{)}\AgdaSpace{}%
\AgdaSymbol{(}\AgdaBound{g}\AgdaSpace{}%
\AgdaSymbol{:}\AgdaSpace{}%
\AgdaDatatype{Gen}\AgdaSpace{}%
\AgdaGeneralizable{n}\AgdaSymbol{)}\AgdaSpace{}%
\AgdaSymbol{→}\AgdaSpace{}%
\AgdaSymbol{(}\AgdaInductiveConstructor{₁₊}\AgdaSpace{}%
\AgdaGeneralizable{n}\AgdaSymbol{)}\AgdaSpace{}%
\AgdaOperator{\AgdaDatatype{VRel,}}\<%
\\
\>[.][@{}l@{}]\<[194I]%
\>[10]\AgdaOperator{\AgdaInductiveConstructor{[}}\AgdaSpace{}%
\AgdaBound{g}\AgdaSpace{}%
\AgdaOperator{\AgdaInductiveConstructor{↥}}\AgdaSpace{}%
\AgdaOperator{\AgdaInductiveConstructor{]ʷ}}\AgdaSpace{}%
\AgdaOperator{\AgdaInductiveConstructor{•}}\AgdaSpace{}%
\AgdaOperator{\AgdaInductiveConstructor{[}}\AgdaSpace{}%
\AgdaInductiveConstructor{gate₁}\AgdaSpace{}%
\AgdaBound{h}\AgdaSpace{}%
\AgdaOperator{\AgdaInductiveConstructor{]ʷ}}\AgdaSpace{}%
\AgdaOperator{\AgdaDatatype{===}}\AgdaSpace{}%
\AgdaOperator{\AgdaInductiveConstructor{[}}\AgdaSpace{}%
\AgdaInductiveConstructor{gate₁}\AgdaSpace{}%
\AgdaBound{h}\AgdaSpace{}%
\AgdaOperator{\AgdaInductiveConstructor{]ʷ}}\AgdaSpace{}%
\AgdaOperator{\AgdaInductiveConstructor{•}}\AgdaSpace{}%
\AgdaOperator{\AgdaInductiveConstructor{[}}\AgdaSpace{}%
\AgdaBound{g}\AgdaSpace{}%
\AgdaOperator{\AgdaInductiveConstructor{↥}}\AgdaSpace{}%
\AgdaOperator{\AgdaInductiveConstructor{]ʷ}}\<%
\\
%
\>[2]\AgdaInductiveConstructor{comm₂}\AgdaSpace{}%
\AgdaSymbol{:}%
\>[225I]\AgdaSymbol{(}\AgdaBound{h}\AgdaSpace{}%
\AgdaSymbol{:}\AgdaSpace{}%
\AgdaPostulate{Gate}\AgdaSpace{}%
\AgdaNumber{2}\AgdaSymbol{)}\AgdaSpace{}%
\AgdaSymbol{(}\AgdaBound{g}\AgdaSpace{}%
\AgdaSymbol{:}\AgdaSpace{}%
\AgdaDatatype{Gen}\AgdaSpace{}%
\AgdaGeneralizable{n}\AgdaSymbol{)}\AgdaSpace{}%
\AgdaSymbol{→}\AgdaSpace{}%
\AgdaSymbol{(}\AgdaInductiveConstructor{₂₊}\AgdaSpace{}%
\AgdaGeneralizable{n}\AgdaSymbol{)}\AgdaSpace{}%
\AgdaOperator{\AgdaDatatype{VRel,}}\<%
\\
\>[.][@{}l@{}]\<[225I]%
\>[10]\AgdaOperator{\AgdaInductiveConstructor{[}}\AgdaSpace{}%
\AgdaBound{g}\AgdaSpace{}%
\AgdaOperator{\AgdaInductiveConstructor{↥}}\AgdaSpace{}%
\AgdaOperator{\AgdaInductiveConstructor{↥}}\AgdaSpace{}%
\AgdaOperator{\AgdaInductiveConstructor{]ʷ}}\AgdaSpace{}%
\AgdaOperator{\AgdaInductiveConstructor{•}}\AgdaSpace{}%
\AgdaOperator{\AgdaInductiveConstructor{[}}\AgdaSpace{}%
\AgdaInductiveConstructor{gate₂}\AgdaSpace{}%
\AgdaBound{h}\AgdaSpace{}%
\AgdaOperator{\AgdaInductiveConstructor{]ʷ}}\AgdaSpace{}%
\AgdaOperator{\AgdaDatatype{===}}\AgdaSpace{}%
\AgdaOperator{\AgdaInductiveConstructor{[}}\AgdaSpace{}%
\AgdaInductiveConstructor{gate₂}\AgdaSpace{}%
\AgdaBound{h}\AgdaSpace{}%
\AgdaOperator{\AgdaInductiveConstructor{]ʷ}}\AgdaSpace{}%
\AgdaOperator{\AgdaInductiveConstructor{•}}\AgdaSpace{}%
\AgdaOperator{\AgdaInductiveConstructor{[}}\AgdaSpace{}%
\AgdaBound{g}\AgdaSpace{}%
\AgdaOperator{\AgdaInductiveConstructor{↥}}\AgdaSpace{}%
\AgdaOperator{\AgdaInductiveConstructor{↥}}\AgdaSpace{}%
\AgdaOperator{\AgdaInductiveConstructor{]ʷ}}\<%
\\
%
\>[2]\AgdaInductiveConstructor{ω↑=ω}\AgdaSpace{}%
\AgdaSymbol{:}\AgdaSpace{}%
\AgdaSymbol{(}\AgdaBound{ω}\AgdaSpace{}%
\AgdaSymbol{:}\AgdaSpace{}%
\AgdaPostulate{Gate}\AgdaSpace{}%
\AgdaNumber{0}\AgdaSymbol{)}\AgdaSpace{}%
\AgdaSymbol{→}\AgdaSpace{}%
\AgdaSymbol{(}\AgdaInductiveConstructor{₁₊}\AgdaSpace{}%
\AgdaGeneralizable{n}\AgdaSymbol{)}\AgdaSpace{}%
\AgdaOperator{\AgdaDatatype{VRel,}}\AgdaSpace{}%
\AgdaOperator{\AgdaInductiveConstructor{[}}\AgdaSpace{}%
\AgdaInductiveConstructor{gate₀}\AgdaSpace{}%
\AgdaBound{ω}\AgdaSpace{}%
\AgdaOperator{\AgdaInductiveConstructor{]ʷ}}\AgdaSpace{}%
\AgdaOperator{\AgdaFunction{↑}}\AgdaSpace{}%
\AgdaOperator{\AgdaDatatype{===}}\AgdaSpace{}%
\AgdaOperator{\AgdaInductiveConstructor{[}}\AgdaSpace{}%
\AgdaInductiveConstructor{gate₀}\AgdaSpace{}%
\AgdaBound{ω}\AgdaSpace{}%
\AgdaOperator{\AgdaInductiveConstructor{]ʷ}}\<%
\end{code}
\end{column}
\begin{column}{0.42\textwidth}
\renewcommand{\AgdaCodeStyle}{\scriptsize}
\small There is \alert{no} \aid{comm₀} rule: centrality of a scalar is a
\emph{theorem}, from the other rules plus one hypothesis --- that scalars
commute with each other:
\begin{code}%
\>[0]\AgdaKeyword{module}\AgdaSpace{}%
\AgdaModule{Central-Scalars}\<%
\\
\>[0][@{}l@{\AgdaIndent{0}}]%
\>[2]\AgdaSymbol{(}\AgdaBound{scalars-comm}\AgdaSpace{}%
\AgdaSymbol{:}\AgdaSpace{}%
\AgdaSymbol{∀}\AgdaSpace{}%
\AgdaSymbol{\{}\AgdaBound{n}\AgdaSymbol{\}}\AgdaSpace{}%
\AgdaSymbol{(}\AgdaBound{h}\AgdaSpace{}%
\AgdaBound{h'}\AgdaSpace{}%
\AgdaSymbol{:}\AgdaSpace{}%
\AgdaPostulate{Gate}\AgdaSpace{}%
\AgdaNumber{0}\AgdaSymbol{)}\AgdaSpace{}%
\AgdaSymbol{→}\<%
\\
\>[2][@{}l@{\AgdaIndent{0}}]%
\>[5]\AgdaOperator{\AgdaInductiveConstructor{[}}%
\>[286I]\AgdaInductiveConstructor{gate₀}\AgdaSpace{}%
\AgdaSymbol{\{}\AgdaBound{n}\AgdaSymbol{\}}\AgdaSpace{}%
\AgdaBound{h'}\AgdaSpace{}%
\AgdaOperator{\AgdaInductiveConstructor{]ʷ}}\AgdaSpace{}%
\AgdaOperator{\AgdaInductiveConstructor{•}}\AgdaSpace{}%
\AgdaOperator{\AgdaInductiveConstructor{[}}\AgdaSpace{}%
\AgdaInductiveConstructor{gate₀}\AgdaSpace{}%
\AgdaBound{h}\AgdaSpace{}%
\AgdaOperator{\AgdaInductiveConstructor{]ʷ}}\<%
\\
\>[.][@{}l@{}]\<[286I]%
\>[7]\AgdaOperator{\AgdaPostulate{≈}}\AgdaSpace{}%
\AgdaOperator{\AgdaInductiveConstructor{[}}\AgdaSpace{}%
\AgdaInductiveConstructor{gate₀}\AgdaSpace{}%
\AgdaBound{h}\AgdaSpace{}%
\AgdaOperator{\AgdaInductiveConstructor{]ʷ}}\AgdaSpace{}%
\AgdaOperator{\AgdaInductiveConstructor{•}}\AgdaSpace{}%
\AgdaOperator{\AgdaInductiveConstructor{[}}\AgdaSpace{}%
\AgdaInductiveConstructor{gate₀}\AgdaSpace{}%
\AgdaBound{h'}\AgdaSpace{}%
\AgdaOperator{\AgdaInductiveConstructor{]ʷ}}\AgdaSymbol{)}\<%
\\
%
\>[2]\AgdaKeyword{where}\<%
\\
%
\>[2]\AgdaFunction{comm₀}\AgdaSpace{}%
\AgdaSymbol{:}%
\>[305I]\AgdaSymbol{∀}\AgdaSpace{}%
\AgdaSymbol{(}\AgdaBound{h}\AgdaSpace{}%
\AgdaSymbol{:}\AgdaSpace{}%
\AgdaPostulate{Gate}\AgdaSpace{}%
\AgdaNumber{0}\AgdaSymbol{)}\AgdaSpace{}%
\AgdaSymbol{(}\AgdaBound{g}\AgdaSpace{}%
\AgdaSymbol{:}\AgdaSpace{}%
\AgdaDatatype{Gen}\AgdaSpace{}%
\AgdaGeneralizable{n}\AgdaSymbol{)}\AgdaSpace{}%
\AgdaSymbol{→}\<%
\\
\>[.][@{}l@{}]\<[305I]%
\>[10]\AgdaOperator{\AgdaInductiveConstructor{[}}\AgdaSpace{}%
\AgdaBound{g}\AgdaSpace{}%
\AgdaOperator{\AgdaInductiveConstructor{]ʷ}}\AgdaSpace{}%
\AgdaOperator{\AgdaInductiveConstructor{•}}\AgdaSpace{}%
\AgdaOperator{\AgdaInductiveConstructor{[}}\AgdaSpace{}%
\AgdaInductiveConstructor{gate₀}\AgdaSpace{}%
\AgdaBound{h}\AgdaSpace{}%
\AgdaOperator{\AgdaInductiveConstructor{]ʷ}}\AgdaSpace{}%
\AgdaOperator{\AgdaPostulate{≈}}\AgdaSpace{}%
\AgdaOperator{\AgdaInductiveConstructor{[}}\AgdaSpace{}%
\AgdaInductiveConstructor{gate₀}\AgdaSpace{}%
\AgdaBound{h}\AgdaSpace{}%
\AgdaOperator{\AgdaInductiveConstructor{]ʷ}}\AgdaSpace{}%
\AgdaOperator{\AgdaInductiveConstructor{•}}\AgdaSpace{}%
\AgdaOperator{\AgdaInductiveConstructor{[}}\AgdaSpace{}%
\AgdaBound{g}\AgdaSpace{}%
\AgdaOperator{\AgdaInductiveConstructor{]ʷ}}\<%
\end{code}
\begin{code}[hide]%
%
\>[2]\AgdaFunction{comm₀}\AgdaSpace{}%
\AgdaSymbol{=}\AgdaSpace{}%
\AgdaPostulate{proof-elided}\<%
\end{code}
{\scriptsize (the \aid{let open PB …} of the source elided)}

\vspace{4pt}
\begin{tikzpicture}[x=5mm,y=5mm,line width=0.7pt]
  \foreach \y in {0,1,2} \draw (0,\y) -- (4,\y);
  \node[draw,fill=softnew,circle,inner sep=1pt,font=\scriptsize] at (1,2.9) {$\omega$};
  \node[draw,fill=white,font=\scriptsize,minimum size=4mm,inner sep=0pt] at (2.6,0) {$h$};
  \node at (5,1) {$\approx$};
  \foreach \y in {0,1,2} \draw (6,\y) -- (10,\y);
  \node[draw,fill=softnew,circle,inner sep=1pt,font=\scriptsize] at (8.6,2.9) {$\omega$};
  \node[draw,fill=white,font=\scriptsize,minimum size=4mm,inner sep=0pt] at (7,0) {$h$};
\end{tikzpicture}
\end{column}
\end{columns}
\end{frame}

\begin{frame}{Minimality: when is a rule \emph{not} derivable from the others? \New}
\begin{code}[hide]%
\>[0]\AgdaKeyword{module}\AgdaSpace{}%
\AgdaModule{Indep}\AgdaSpace{}%
\AgdaKeyword{where}\<%
\\
\>[0][@{}l@{\AgdaIndent{0}}]%
\>[2]\AgdaKeyword{infixl}\AgdaSpace{}%
\AgdaNumber{8}\AgdaSpace{}%
\AgdaOperator{\AgdaFunction{\AgdaUnderscore{}∣\AgdaUnderscore{}}}\<%
\\
%
\>[2]\AgdaKeyword{postulate}\<%
\\
\>[2][@{}l@{\AgdaIndent{0}}]%
\>[4]\AgdaPostulate{Derivable}\AgdaSpace{}%
\AgdaSymbol{:}\AgdaSpace{}%
\AgdaSymbol{∀}\AgdaSpace{}%
\AgdaSymbol{\{}\AgdaBound{Γ}\AgdaSpace{}%
\AgdaSymbol{:}\AgdaSpace{}%
\AgdaFunction{WRel}\AgdaSpace{}%
\AgdaPostulate{X}\AgdaSymbol{\}}\AgdaSpace{}%
\AgdaSymbol{\{}\AgdaBound{P}\AgdaSpace{}%
\AgdaSymbol{:}\AgdaSpace{}%
\AgdaPrimitive{Set}\AgdaSymbol{\}}\AgdaSpace{}%
\AgdaSymbol{→}\AgdaSpace{}%
\AgdaPrimitive{Set}\<%
\\
%
\>[4]\AgdaOperator{\AgdaPostulate{⟦\AgdaUnderscore{}⟧}}\AgdaSpace{}%
\AgdaSymbol{:}\AgdaSpace{}%
\AgdaDatatype{Word}\AgdaSpace{}%
\AgdaPostulate{X}\AgdaSpace{}%
\AgdaSymbol{→}\AgdaSpace{}%
\AgdaPostulate{X}\<%
\\
%
\>[4]\AgdaOperator{\AgdaPostulate{\AgdaUnderscore{}≈ᴹ\AgdaUnderscore{}}}\AgdaSpace{}%
\AgdaSymbol{:}\AgdaSpace{}%
\AgdaPostulate{X}\AgdaSpace{}%
\AgdaSymbol{→}\AgdaSpace{}%
\AgdaPostulate{X}\AgdaSpace{}%
\AgdaSymbol{→}\AgdaSpace{}%
\AgdaPrimitive{Set}\<%
\end{code}
\begin{columns}[c]
\begin{column}{0.56\textwidth}
\small \texttt{Presentation/Independence.agda} (\Branch{qupit}): select some
axiom \emph{instances}; is another one derivable from them?
\begin{code}%
%
\>[2]\AgdaFunction{Select}\AgdaSpace{}%
\AgdaSymbol{:}\AgdaSpace{}%
\AgdaFunction{WRel}\AgdaSpace{}%
\AgdaPostulate{X}\AgdaSpace{}%
\AgdaSymbol{→}\AgdaSpace{}%
\AgdaPrimitive{Set₁}\<%
\\
%
\>[2]\AgdaFunction{Select}\AgdaSpace{}%
\AgdaBound{Γ}\AgdaSpace{}%
\AgdaSymbol{=}\AgdaSpace{}%
\AgdaSymbol{∀}\AgdaSpace{}%
\AgdaSymbol{\{}\AgdaBound{w}\AgdaSpace{}%
\AgdaBound{v}\AgdaSymbol{\}}\AgdaSpace{}%
\AgdaSymbol{→}\AgdaSpace{}%
\AgdaBound{Γ}\AgdaSpace{}%
\AgdaBound{w}\AgdaSpace{}%
\AgdaBound{v}\AgdaSpace{}%
\AgdaSymbol{→}\AgdaSpace{}%
\AgdaPrimitive{Set}\<%
\\
%
\\[\AgdaEmptyExtraSkip]%
%
\>[2]\AgdaOperator{\AgdaFunction{\AgdaUnderscore{}∣\AgdaUnderscore{}}}\AgdaSpace{}%
\AgdaSymbol{:}\AgdaSpace{}%
\AgdaSymbol{(}\AgdaBound{Γ}\AgdaSpace{}%
\AgdaSymbol{:}\AgdaSpace{}%
\AgdaFunction{WRel}\AgdaSpace{}%
\AgdaPostulate{X}\AgdaSymbol{)}\AgdaSpace{}%
\AgdaSymbol{→}\AgdaSpace{}%
\AgdaFunction{Select}\AgdaSpace{}%
\AgdaBound{Γ}\AgdaSpace{}%
\AgdaSymbol{→}\AgdaSpace{}%
\AgdaFunction{WRel}\AgdaSpace{}%
\AgdaPostulate{X}\<%
\\
%
\>[2]\AgdaSymbol{(}\AgdaBound{Γ}\AgdaSpace{}%
\AgdaOperator{\AgdaFunction{∣}}\AgdaSpace{}%
\AgdaBound{P}\AgdaSymbol{)}\AgdaSpace{}%
\AgdaBound{w}\AgdaSpace{}%
\AgdaBound{v}\AgdaSpace{}%
\AgdaSymbol{=}\AgdaSpace{}%
\AgdaRecord{Σ}\AgdaSpace{}%
\AgdaSymbol{(}\AgdaBound{Γ}\AgdaSpace{}%
\AgdaBound{w}\AgdaSpace{}%
\AgdaBound{v}\AgdaSymbol{)}\AgdaSpace{}%
\AgdaBound{P}\<%
\end{code}
\begin{code}[hide]%
%
\>[2]\AgdaKeyword{postulate}\<%
\\
\>[2][@{}l@{\AgdaIndent{0}}]%
\>[4]\AgdaPostulate{Independent}\AgdaSpace{}%
\AgdaSymbol{:}\AgdaSpace{}%
\AgdaSymbol{∀}\AgdaSpace{}%
\AgdaSymbol{(}\AgdaBound{Γ}\AgdaSpace{}%
\AgdaSymbol{:}\AgdaSpace{}%
\AgdaFunction{WRel}\AgdaSpace{}%
\AgdaPostulate{X}\AgdaSymbol{)}\AgdaSpace{}%
\AgdaSymbol{(}\AgdaBound{P}\AgdaSpace{}%
\AgdaSymbol{:}\AgdaSpace{}%
\AgdaFunction{Select}\AgdaSpace{}%
\AgdaBound{Γ}\AgdaSymbol{)}\AgdaSpace{}%
\AgdaSymbol{\{}\AgdaBound{w}\AgdaSpace{}%
\AgdaBound{v}\AgdaSymbol{\}}\AgdaSpace{}%
\AgdaSymbol{→}\AgdaSpace{}%
\AgdaBound{Γ}\AgdaSpace{}%
\AgdaBound{w}\AgdaSpace{}%
\AgdaBound{v}\AgdaSpace{}%
\AgdaSymbol{→}\AgdaSpace{}%
\AgdaPrimitive{Set}\<%
\\
%
\>[2]\AgdaKeyword{module}\AgdaSpace{}%
\AgdaModule{Model}\AgdaSpace{}%
\AgdaSymbol{(}\AgdaBound{Γ}\AgdaSpace{}%
\AgdaSymbol{:}\AgdaSpace{}%
\AgdaFunction{WRel}\AgdaSpace{}%
\AgdaPostulate{X}\AgdaSymbol{)}\AgdaSpace{}%
\AgdaSymbol{(}\AgdaBound{P}\AgdaSpace{}%
\AgdaSymbol{:}\AgdaSpace{}%
\AgdaFunction{Select}\AgdaSpace{}%
\AgdaBound{Γ}\AgdaSymbol{)}\AgdaSpace{}%
\AgdaSymbol{(}\AgdaBound{mon}\AgdaSpace{}%
\AgdaSymbol{:}\AgdaSpace{}%
\AgdaRecord{Monoid}\AgdaSpace{}%
\AgdaFunction{0ℓ}\AgdaSpace{}%
\AgdaFunction{0ℓ}\AgdaSymbol{)}\AgdaSpace{}%
\AgdaSymbol{(}\AgdaOperator{\AgdaBound{⟦\AgdaUnderscore{}⟧₀}}\AgdaSpace{}%
\AgdaSymbol{:}\AgdaSpace{}%
\AgdaPostulate{X}\AgdaSpace{}%
\AgdaSymbol{→}\AgdaSpace{}%
\AgdaField{Monoid.Carrier}\AgdaSpace{}%
\AgdaBound{mon}\AgdaSymbol{)}\AgdaSpace{}%
\AgdaKeyword{where}\<%
\\
\>[2][@{}l@{\AgdaIndent{0}}]%
\>[3]\AgdaKeyword{module}\AgdaSpace{}%
\AgdaModule{Sound}\AgdaSpace{}%
\AgdaKeyword{where}\<%
\end{code}
\par\medskip
\small A \alert{model}: a monoid where the selected axioms hold. It
separates the instance's two sides $\Rightarrow$ the instance is independent:
\begin{code}%
\>[3][@{}l@{\AgdaIndent{1}}]%
\>[4]\AgdaFunction{independent}\AgdaSpace{}%
\AgdaSymbol{:}\AgdaSpace{}%
\AgdaSymbol{(}\AgdaBound{a}\AgdaSpace{}%
\AgdaSymbol{:}\AgdaSpace{}%
\AgdaBound{Γ}\AgdaSpace{}%
\AgdaGeneralizable{w}\AgdaSpace{}%
\AgdaGeneralizable{v}\AgdaSymbol{)}\AgdaSpace{}%
\AgdaSymbol{→}\AgdaSpace{}%
\AgdaOperator{\AgdaFunction{¬}}\AgdaSpace{}%
\AgdaSymbol{(}\AgdaOperator{\AgdaPostulate{⟦}}\AgdaSpace{}%
\AgdaGeneralizable{w}\AgdaSpace{}%
\AgdaOperator{\AgdaPostulate{⟧}}\AgdaSpace{}%
\AgdaOperator{\AgdaPostulate{≈ᴹ}}\AgdaSpace{}%
\AgdaOperator{\AgdaPostulate{⟦}}\AgdaSpace{}%
\AgdaGeneralizable{v}\AgdaSpace{}%
\AgdaOperator{\AgdaPostulate{⟧}}\AgdaSymbol{)}\AgdaSpace{}%
\AgdaSymbol{→}\AgdaSpace{}%
\AgdaPostulate{Independent}\AgdaSpace{}%
\AgdaBound{Γ}\AgdaSpace{}%
\AgdaBound{P}\AgdaSpace{}%
\AgdaBound{a}\<%
\end{code}
\begin{code}[hide]%
%
\>[4]\AgdaFunction{independent}\AgdaSpace{}%
\AgdaSymbol{=}\AgdaSpace{}%
\AgdaPostulate{proof-elided}\<%
\end{code}
{\scriptsize (\aid{\_≈ᴹ\_} is the monoid's \aid{\_≈\_} in the source)}
\end{column}
\begin{column}{0.41\textwidth}
\centering
\begin{tikzpicture}[>={Stealth[length=1.8mm]}]
  \node[draw=accent,fill=soft,rounded corners=3pt,minimum width=40mm,minimum height=26mm] (R) at (0,0) {};
  \node[font=\scriptsize\bfseries,accent,anchor=north west] at (R.north west) {axiom instances of $\Gamma$};
  \foreach \x/\y in {-1.4/0.3,-0.8/-0.4,-0.2/0.2,0.4/-0.5,1.0/0.1,-1.1/-0.9,0.1/-1.0} \node[circle,fill=paperC,inner sep=1.4pt] at (\x,\y) {};
  \node[circle,fill=newC,inner sep=2pt,label={[font=\scriptsize,newC]right:$a$}] (a) at (1.25,-0.75) {};
  \draw[dashed,paperC,rounded corners] (-1.75,-1.2) rectangle (1.3,0.55);
  \node[font=\tiny,paperC] at (-0.3,0.7) {selected $P$};
  \node[draw=gray,fill=white,rounded corners=2pt,font=\scriptsize,align=center] (M) at (0,-2.4) {model $\mathcal{M}$: all of $P$ hold,\\ $\llbracket w\rrbracket\neq\llbracket v\rrbracket$ for $a : w = v$};
  \draw[->,thick,gray] (M) -- (0,-1.32);
\end{tikzpicture}

\vspace{6pt}
{\small soundness extends from axioms to the whole congruence of $\Gamma\mid P$
by \aid{StarInterp}: a separated instance is not derivable.}
\end{column}
\end{columns}
\end{frame}

\begin{frame}{The structural rules are mutually independent \New}
\begin{code}[hide]%
\>[0]\AgdaKeyword{data}\AgdaSpace{}%
\AgdaDatatype{Rule}\AgdaSpace{}%
\AgdaSymbol{:}\AgdaSpace{}%
\AgdaPrimitive{Set}\AgdaSpace{}%
\AgdaKeyword{where}\<%
\\
\>[0][@{}l@{\AgdaIndent{0}}]%
\>[2]\AgdaInductiveConstructor{cong↑}\AgdaSpace{}%
\AgdaInductiveConstructor{comm₁}\AgdaSpace{}%
\AgdaInductiveConstructor{comm₂}\AgdaSpace{}%
\AgdaInductiveConstructor{ω↑=ω}\AgdaSpace{}%
\AgdaSymbol{:}\AgdaSpace{}%
\AgdaDatatype{Rule}\<%
\\
\>[0]\AgdaKeyword{infix}\AgdaSpace{}%
\AgdaNumber{4}\AgdaSpace{}%
\AgdaOperator{\AgdaPostulate{\AgdaUnderscore{}IndependentOf\AgdaUnderscore{}}}\<%
\\
\>[0]\AgdaKeyword{postulate}\<%
\\
\>[0][@{}l@{\AgdaIndent{0}}]%
\>[2]\AgdaOperator{\AgdaPostulate{\AgdaUnderscore{}IndependentOf\AgdaUnderscore{}}}\AgdaSpace{}%
\AgdaSymbol{:}\AgdaSpace{}%
\AgdaSymbol{∀}\AgdaSpace{}%
\AgdaSymbol{\{}\AgdaBound{w}\AgdaSpace{}%
\AgdaBound{v}\AgdaSpace{}%
\AgdaSymbol{:}\AgdaSpace{}%
\AgdaFunction{Circuit}\AgdaSpace{}%
\AgdaGeneralizable{n}\AgdaSymbol{\}}\AgdaSpace{}%
\AgdaSymbol{→}\AgdaSpace{}%
\AgdaGeneralizable{n}\AgdaSpace{}%
\AgdaOperator{\AgdaDatatype{VRel,}}\AgdaSpace{}%
\AgdaBound{w}\AgdaSpace{}%
\AgdaOperator{\AgdaDatatype{===}}\AgdaSpace{}%
\AgdaBound{v}\AgdaSpace{}%
\AgdaSymbol{→}\AgdaSpace{}%
\AgdaDatatype{Rule}\AgdaSpace{}%
\AgdaSymbol{→}\AgdaSpace{}%
\AgdaPrimitive{Set}\<%
\\
\>[0]\AgdaKeyword{module}\AgdaSpace{}%
\AgdaModule{Pure}\AgdaSpace{}%
\AgdaKeyword{where}\<%
\end{code}
\begin{columns}[c]
\begin{column}{0.56\textwidth}
\renewcommand{\AgdaCodeStyle}{\scriptsize}
\small Over \emph{any} gate family, for the pure structural theory
(\texttt{Circuit/Independence.agda}), each rule fails in a model of the other three:
\begin{code}%
\>[0][@{}l@{\AgdaIndent{1}}]%
\>[2]\AgdaFunction{comm₁-independent}\AgdaSpace{}%
\AgdaSymbol{:}%
\>[490I]\AgdaSymbol{(}\AgdaBound{h}\AgdaSpace{}%
\AgdaSymbol{:}\AgdaSpace{}%
\AgdaPostulate{Gate}\AgdaSpace{}%
\AgdaNumber{1}\AgdaSymbol{)}\AgdaSpace{}%
\AgdaSymbol{→}\<%
\\
\>[.][@{}l@{}]\<[490I]%
\>[22]\AgdaInductiveConstructor{comm₁}\AgdaSpace{}%
\AgdaBound{h}\AgdaSpace{}%
\AgdaSymbol{(}\AgdaInductiveConstructor{gate₁}\AgdaSpace{}%
\AgdaSymbol{\{}\AgdaGeneralizable{n}\AgdaSymbol{\}}\AgdaSpace{}%
\AgdaBound{h}\AgdaSymbol{)}\AgdaSpace{}%
\AgdaOperator{\AgdaPostulate{IndependentOf}}\AgdaSpace{}%
\AgdaInductiveConstructor{comm₁}\<%
\\
%
\>[2]\AgdaFunction{ω↑=ω-independent}\AgdaSpace{}%
\AgdaSymbol{:}\AgdaSpace{}%
\AgdaSymbol{(}\AgdaBound{ω}\AgdaSpace{}%
\AgdaSymbol{:}\AgdaSpace{}%
\AgdaPostulate{Gate}\AgdaSpace{}%
\AgdaNumber{0}\AgdaSymbol{)}\AgdaSpace{}%
\AgdaSymbol{→}\AgdaSpace{}%
\AgdaInductiveConstructor{ω↑=ω}\AgdaSpace{}%
\AgdaSymbol{\{}\AgdaGeneralizable{n}\AgdaSymbol{\}}\AgdaSpace{}%
\AgdaBound{ω}\AgdaSpace{}%
\AgdaOperator{\AgdaPostulate{IndependentOf}}\AgdaSpace{}%
\AgdaInductiveConstructor{ω↑=ω}\<%
\\
%
\>[2]\AgdaFunction{cong↑-independent₁}\AgdaSpace{}%
\AgdaSymbol{:}%
\>[513I]\AgdaSymbol{(}\AgdaBound{h}\AgdaSpace{}%
\AgdaSymbol{:}\AgdaSpace{}%
\AgdaPostulate{Gate}\AgdaSpace{}%
\AgdaNumber{1}\AgdaSymbol{)}\AgdaSpace{}%
\AgdaSymbol{→}\<%
\\
\>[.][@{}l@{}]\<[513I]%
\>[23]\AgdaInductiveConstructor{cong↑}\AgdaSpace{}%
\AgdaSymbol{(}\AgdaInductiveConstructor{comm₁}\AgdaSpace{}%
\AgdaBound{h}\AgdaSpace{}%
\AgdaSymbol{(}\AgdaInductiveConstructor{gate₁}\AgdaSpace{}%
\AgdaSymbol{\{}\AgdaGeneralizable{n}\AgdaSymbol{\}}\AgdaSpace{}%
\AgdaBound{h}\AgdaSymbol{))}\AgdaSpace{}%
\AgdaOperator{\AgdaPostulate{IndependentOf}}\AgdaSpace{}%
\AgdaInductiveConstructor{cong↑}\<%
\end{code}
\begin{code}[hide]%
%
\>[2]\AgdaFunction{comm₁-independent}\AgdaSpace{}%
\AgdaSymbol{=}\AgdaSpace{}%
\AgdaPostulate{proof-elided}\<%
\\
%
\>[2]\AgdaFunction{ω↑=ω-independent}\AgdaSpace{}%
\AgdaSymbol{=}\AgdaSpace{}%
\AgdaPostulate{proof-elided}\<%
\\
%
\>[2]\AgdaFunction{cong↑-independent₁}\AgdaSpace{}%
\AgdaSymbol{=}\AgdaSpace{}%
\AgdaPostulate{proof-elided}\<%
\end{code}
\small and the hypothesis of \aid{Central-Scalars} is \alert{needed}:
\begin{code}%
%
\>[2]\AgdaFunction{scalars-comm-needed}\AgdaSpace{}%
\AgdaSymbol{:}\AgdaSpace{}%
\AgdaSymbol{(}\AgdaBound{ω}\AgdaSpace{}%
\AgdaBound{ω'}\AgdaSpace{}%
\AgdaSymbol{:}\AgdaSpace{}%
\AgdaPostulate{Gate}\AgdaSpace{}%
\AgdaNumber{0}\AgdaSymbol{)}\AgdaSpace{}%
\AgdaSymbol{→}\AgdaSpace{}%
\AgdaBound{ω}\AgdaSpace{}%
\AgdaOperator{\AgdaFunction{≢}}\AgdaSpace{}%
\AgdaBound{ω'}\AgdaSpace{}%
\AgdaSymbol{→}\<%
\\
\>[2][@{}l@{\AgdaIndent{0}}]%
\>[4]\AgdaOperator{\AgdaFunction{¬}}\AgdaSpace{}%
\AgdaSymbol{(}\AgdaOperator{\AgdaInductiveConstructor{[}}%
\>[543I]\AgdaInductiveConstructor{gate₀}\AgdaSpace{}%
\AgdaSymbol{\{}\AgdaGeneralizable{n}\AgdaSymbol{\}}\AgdaSpace{}%
\AgdaBound{ω'}\AgdaSpace{}%
\AgdaOperator{\AgdaInductiveConstructor{]ʷ}}\AgdaSpace{}%
\AgdaOperator{\AgdaInductiveConstructor{•}}\AgdaSpace{}%
\AgdaOperator{\AgdaInductiveConstructor{[}}\AgdaSpace{}%
\AgdaInductiveConstructor{gate₀}\AgdaSpace{}%
\AgdaBound{ω}\AgdaSpace{}%
\AgdaOperator{\AgdaInductiveConstructor{]ʷ}}\<%
\\
\>[.][@{}l@{}]\<[543I]%
\>[9]\AgdaOperator{\AgdaPostulate{≈}}\AgdaSpace{}%
\AgdaOperator{\AgdaInductiveConstructor{[}}\AgdaSpace{}%
\AgdaInductiveConstructor{gate₀}\AgdaSpace{}%
\AgdaBound{ω}\AgdaSpace{}%
\AgdaOperator{\AgdaInductiveConstructor{]ʷ}}\AgdaSpace{}%
\AgdaOperator{\AgdaInductiveConstructor{•}}\AgdaSpace{}%
\AgdaOperator{\AgdaInductiveConstructor{[}}\AgdaSpace{}%
\AgdaInductiveConstructor{gate₀}\AgdaSpace{}%
\AgdaBound{ω'}\AgdaSpace{}%
\AgdaOperator{\AgdaInductiveConstructor{]ʷ}}\AgdaSymbol{)}\<%
\end{code}
\begin{code}[hide]%
%
\>[2]\AgdaFunction{scalars-comm-needed}\AgdaSpace{}%
\AgdaSymbol{=}\AgdaSpace{}%
\AgdaPostulate{proof-elided}\<%
\end{code}
\end{column}
\begin{column}{0.41\textwidth}
\scriptsize
\renewcommand{\arraystretch}{1.15}
\begin{tabular}{@{}p{1.25cm}p{3.9cm}@{}}
\toprule
rule & model separating it\\
\midrule
\aid{ω↑=ω} & letter $\mapsto 2^{\mathit{depth}}$ in $(\mathbb{N},+)$; shift doubles\\
\aid{comm₁} & depths of 1-wire gates in \aid{List ℕ}; shift $=$ \aid{map suc}\\
\aid{comm₂} & the same, for 2-wire gates\\
\aid{cong↑} & read gates only at depth $\ge1$: no shift homomorphism\\
\aid{comm₀} & scalars in order, in \aid{List (Gate 0)}\\
\bottomrule
\end{tabular}

\vspace{6pt}
\begin{tikzpicture}[x=5mm,y=4.4mm,line width=0.6pt]
  \foreach \y in {0,1,2} \draw (0,\y) -- (6,\y);
  \node[draw,fill=white,font=\tiny,minimum size=3mm,inner sep=0pt] at (1.5,0) {$1$};
  \node[draw,fill=white,font=\tiny,minimum size=3mm,inner sep=0pt] at (3,1) {$2$};
  \node[draw,fill=white,font=\tiny,minimum size=3mm,inner sep=0pt] at (4.5,2) {$4$};
  \node[font=\scriptsize,anchor=west] at (6.3,1) {weights $2^{\mathit{depth}}$};
\end{tikzpicture}

\vspace{2pt}
{\scriptsize Scope today: the \emph{structural} rules. No gate-specific rule set is
shown minimal yet --- but several were \emph{reduced} by derivations
(next slides).}
\end{column}
\end{columns}
\end{frame}

\begin{code}[hide]%
\>[0]\AgdaKeyword{module}\AgdaSpace{}%
\AgdaModule{slides.Tsinghua-PL-seminar.sections.newqupit}\AgdaSpace{}%
\AgdaKeyword{where}\<%
\\
%
\\[\AgdaEmptyExtraSkip]%
\>[0]\AgdaKeyword{open}\AgdaSpace{}%
\AgdaKeyword{import}\AgdaSpace{}%
\AgdaModule{Data.Nat}\AgdaSpace{}%
\AgdaKeyword{using}\AgdaSpace{}%
\AgdaSymbol{(}\AgdaDatatype{ℕ}\AgdaSymbol{)}\<%
\\
\>[0]\AgdaKeyword{open}\AgdaSpace{}%
\AgdaKeyword{import}\AgdaSpace{}%
\AgdaModule{Level}\AgdaSpace{}%
\AgdaKeyword{using}\AgdaSpace{}%
\AgdaSymbol{(}\AgdaFunction{0ℓ}\AgdaSymbol{)}\<%
\\
\>[0]\AgdaKeyword{open}\AgdaSpace{}%
\AgdaKeyword{import}\AgdaSpace{}%
\AgdaModule{Algebra.Bundles}\AgdaSpace{}%
\AgdaKeyword{using}\AgdaSpace{}%
\AgdaSymbol{(}\AgdaRecord{Group}\AgdaSymbol{)}\<%
\\
\>[0]\AgdaKeyword{open}\AgdaSpace{}%
\AgdaKeyword{import}\AgdaSpace{}%
\AgdaModule{Notations}\AgdaSpace{}%
\AgdaKeyword{using}\AgdaSpace{}%
\AgdaSymbol{(}\AgdaInductiveConstructor{₁₊}\AgdaSpace{}%
\AgdaSymbol{;}\AgdaSpace{}%
\AgdaInductiveConstructor{₂₊}\AgdaSpace{}%
\AgdaSymbol{;}\AgdaSpace{}%
\AgdaInductiveConstructor{₃₊}\AgdaSymbol{)}\<%
\\
\>[0]\AgdaKeyword{open}\AgdaSpace{}%
\AgdaKeyword{import}\AgdaSpace{}%
\AgdaModule{Word.Base}\AgdaSpace{}%
\AgdaKeyword{using}\AgdaSpace{}%
\AgdaSymbol{(}\AgdaDatatype{Word}\AgdaSpace{}%
\AgdaSymbol{;}\AgdaSpace{}%
\AgdaFunction{WRel}\AgdaSpace{}%
\AgdaSymbol{;}\AgdaSpace{}%
\AgdaInductiveConstructor{ε}\AgdaSpace{}%
\AgdaSymbol{;}\AgdaSpace{}%
\AgdaOperator{\AgdaInductiveConstructor{\AgdaUnderscore{}•\AgdaUnderscore{}}}\AgdaSymbol{)}\<%
\\
\>[0]\AgdaKeyword{open}\AgdaSpace{}%
\AgdaKeyword{import}\AgdaSpace{}%
\AgdaModule{Presentation.Definitions}\AgdaSpace{}%
\AgdaKeyword{using}\AgdaSpace{}%
\AgdaSymbol{(}\AgdaOperator{\AgdaRecord{\AgdaUnderscore{}IsPresentationOf\AgdaUnderscore{}}}\AgdaSymbol{)}\<%
\\
%
\\[\AgdaEmptyExtraSkip]%
\>[0]\AgdaKeyword{infix}\AgdaSpace{}%
\AgdaNumber{4}\AgdaSpace{}%
\AgdaOperator{\AgdaDatatype{\AgdaUnderscore{}SRel,\AgdaUnderscore{}===\AgdaUnderscore{}}}\AgdaSpace{}%
\AgdaOperator{\AgdaPostulate{\AgdaUnderscore{}QRel,\AgdaUnderscore{}===\AgdaUnderscore{}}}\<%
\\
\>[0]\AgdaKeyword{infixr}\AgdaSpace{}%
\AgdaNumber{8}\AgdaSpace{}%
\AgdaOperator{\AgdaPostulate{\AgdaUnderscore{}\textasciicircum{}\AgdaUnderscore{}}}\<%
\\
\>[0]\AgdaKeyword{infix}\AgdaSpace{}%
\AgdaNumber{9}\AgdaSpace{}%
\AgdaOperator{\AgdaPostulate{S\textasciicircum{}\AgdaUnderscore{}}}\AgdaSpace{}%
\AgdaOperator{\AgdaPostulate{Z\textasciicircum{}\AgdaUnderscore{}}}\AgdaSpace{}%
\AgdaOperator{\AgdaPostulate{CZ\textasciicircum{}\AgdaUnderscore{}}}\AgdaSpace{}%
\AgdaOperator{\AgdaPostulate{XMg\textasciicircum{}\AgdaUnderscore{}}}\AgdaSpace{}%
\AgdaOperator{\AgdaPostulate{g\textasciicircum{}\AgdaUnderscore{}}}\<%
\\
\>[0]\AgdaKeyword{infixl}\AgdaSpace{}%
\AgdaNumber{10}\AgdaSpace{}%
\AgdaOperator{\AgdaPostulate{\AgdaUnderscore{}↑}}\AgdaSpace{}%
\AgdaOperator{\AgdaPostulate{\AgdaUnderscore{}↓}}\<%
\\
\>[0]\AgdaKeyword{infixl}\AgdaSpace{}%
\AgdaNumber{7}\AgdaSpace{}%
\AgdaOperator{\AgdaPostulate{\AgdaUnderscore{}*\AgdaUnderscore{}}}\<%
\\
\>[0]\AgdaKeyword{infixl}\AgdaSpace{}%
\AgdaNumber{6}\AgdaSpace{}%
\AgdaOperator{\AgdaPostulate{\AgdaUnderscore{}+\AgdaUnderscore{}}}\<%
\\
\>[0]\AgdaKeyword{infix}\AgdaSpace{}%
\AgdaNumber{8}\AgdaSpace{}%
\AgdaOperator{\AgdaPostulate{-\AgdaUnderscore{}}}\<%
\\
%
\\[\AgdaEmptyExtraSkip]%
\>[0]\AgdaKeyword{postulate}\<%
\\
\>[0][@{}l@{\AgdaIndent{0}}]%
\>[2]\AgdaPostulate{Gen}\AgdaSpace{}%
\AgdaSymbol{:}\AgdaSpace{}%
\AgdaDatatype{ℕ}\AgdaSpace{}%
\AgdaSymbol{→}\AgdaSpace{}%
\AgdaPrimitive{Set}\<%
\\
%
\>[2]\AgdaPostulate{p}\AgdaSpace{}%
\AgdaPostulate{p-1}\AgdaSpace{}%
\AgdaPostulate{g}\AgdaSpace{}%
\AgdaPostulate{k}\AgdaSpace{}%
\AgdaPostulate{1/2}\AgdaSpace{}%
\AgdaSymbol{:}\AgdaSpace{}%
\AgdaDatatype{ℕ}\<%
\\
%
\>[2]\AgdaPostulate{S}\AgdaSpace{}%
\AgdaPostulate{H}\AgdaSpace{}%
\AgdaPostulate{Z}\AgdaSpace{}%
\AgdaPostulate{CZ}\AgdaSpace{}%
\AgdaPostulate{Ex}\AgdaSpace{}%
\AgdaPostulate{CX}\AgdaSpace{}%
\AgdaPostulate{CZ02}\AgdaSpace{}%
\AgdaPostulate{XMg}\AgdaSpace{}%
\AgdaPostulate{XM₋₁}\AgdaSpace{}%
\AgdaSymbol{:}\AgdaSpace{}%
\AgdaSymbol{∀}\AgdaSpace{}%
\AgdaSymbol{\{}\AgdaBound{n}\AgdaSymbol{\}}\AgdaSpace{}%
\AgdaSymbol{→}\AgdaSpace{}%
\AgdaDatatype{Word}\AgdaSpace{}%
\AgdaSymbol{(}\AgdaPostulate{Gen}\AgdaSpace{}%
\AgdaBound{n}\AgdaSymbol{)}\<%
\\
%
\>[2]\AgdaPostulate{XM}\AgdaSpace{}%
\AgdaSymbol{:}\AgdaSpace{}%
\AgdaDatatype{ℕ}\AgdaSpace{}%
\AgdaSymbol{→}\AgdaSpace{}%
\AgdaSymbol{∀}\AgdaSpace{}%
\AgdaSymbol{\{}\AgdaBound{n}\AgdaSymbol{\}}\AgdaSpace{}%
\AgdaSymbol{→}\AgdaSpace{}%
\AgdaDatatype{Word}\AgdaSpace{}%
\AgdaSymbol{(}\AgdaPostulate{Gen}\AgdaSpace{}%
\AgdaBound{n}\AgdaSymbol{)}\<%
\\
%
\>[2]\AgdaOperator{\AgdaPostulate{\AgdaUnderscore{}\textasciicircum{}\AgdaUnderscore{}}}\AgdaSpace{}%
\AgdaSymbol{:}\AgdaSpace{}%
\AgdaSymbol{∀}\AgdaSpace{}%
\AgdaSymbol{\{}\AgdaBound{n}\AgdaSymbol{\}}\AgdaSpace{}%
\AgdaSymbol{→}\AgdaSpace{}%
\AgdaDatatype{Word}\AgdaSpace{}%
\AgdaSymbol{(}\AgdaPostulate{Gen}\AgdaSpace{}%
\AgdaBound{n}\AgdaSymbol{)}\AgdaSpace{}%
\AgdaSymbol{→}\AgdaSpace{}%
\AgdaDatatype{ℕ}\AgdaSpace{}%
\AgdaSymbol{→}\AgdaSpace{}%
\AgdaDatatype{Word}\AgdaSpace{}%
\AgdaSymbol{(}\AgdaPostulate{Gen}\AgdaSpace{}%
\AgdaBound{n}\AgdaSymbol{)}\<%
\\
%
\>[2]\AgdaOperator{\AgdaPostulate{S\textasciicircum{}\AgdaUnderscore{}}}\AgdaSpace{}%
\AgdaOperator{\AgdaPostulate{Z\textasciicircum{}\AgdaUnderscore{}}}\AgdaSpace{}%
\AgdaOperator{\AgdaPostulate{CZ\textasciicircum{}\AgdaUnderscore{}}}\AgdaSpace{}%
\AgdaOperator{\AgdaPostulate{XMg\textasciicircum{}\AgdaUnderscore{}}}\AgdaSpace{}%
\AgdaSymbol{:}\AgdaSpace{}%
\AgdaSymbol{∀}\AgdaSpace{}%
\AgdaSymbol{\{}\AgdaBound{n}\AgdaSymbol{\}}\AgdaSpace{}%
\AgdaSymbol{→}\AgdaSpace{}%
\AgdaDatatype{ℕ}\AgdaSpace{}%
\AgdaSymbol{→}\AgdaSpace{}%
\AgdaDatatype{Word}\AgdaSpace{}%
\AgdaSymbol{(}\AgdaPostulate{Gen}\AgdaSpace{}%
\AgdaBound{n}\AgdaSymbol{)}\<%
\\
%
\>[2]\AgdaOperator{\AgdaPostulate{g\textasciicircum{}\AgdaUnderscore{}}}\AgdaSpace{}%
\AgdaOperator{\AgdaPostulate{-\AgdaUnderscore{}}}\AgdaSpace{}%
\AgdaSymbol{:}\AgdaSpace{}%
\AgdaDatatype{ℕ}\AgdaSpace{}%
\AgdaSymbol{→}\AgdaSpace{}%
\AgdaDatatype{ℕ}\<%
\\
%
\>[2]\AgdaOperator{\AgdaPostulate{\AgdaUnderscore{}*\AgdaUnderscore{}}}\AgdaSpace{}%
\AgdaOperator{\AgdaPostulate{\AgdaUnderscore{}+\AgdaUnderscore{}}}\AgdaSpace{}%
\AgdaSymbol{:}\AgdaSpace{}%
\AgdaDatatype{ℕ}\AgdaSpace{}%
\AgdaSymbol{→}\AgdaSpace{}%
\AgdaDatatype{ℕ}\AgdaSpace{}%
\AgdaSymbol{→}\AgdaSpace{}%
\AgdaDatatype{ℕ}\<%
\\
%
\>[2]\AgdaOperator{\AgdaPostulate{\AgdaUnderscore{}↑}}\AgdaSpace{}%
\AgdaOperator{\AgdaPostulate{\AgdaUnderscore{}↓}}\AgdaSpace{}%
\AgdaSymbol{:}\AgdaSpace{}%
\AgdaSymbol{∀}\AgdaSpace{}%
\AgdaSymbol{\{}\AgdaBound{n}\AgdaSymbol{\}}\AgdaSpace{}%
\AgdaSymbol{→}\AgdaSpace{}%
\AgdaDatatype{Word}\AgdaSpace{}%
\AgdaSymbol{(}\AgdaPostulate{Gen}\AgdaSpace{}%
\AgdaBound{n}\AgdaSymbol{)}\AgdaSpace{}%
\AgdaSymbol{→}\AgdaSpace{}%
\AgdaDatatype{Word}\AgdaSpace{}%
\AgdaSymbol{(}\AgdaPostulate{Gen}\AgdaSpace{}%
\AgdaBound{n}\AgdaSymbol{)}\<%
\\
%
\>[2]\AgdaOperator{\AgdaPostulate{\AgdaUnderscore{}QRel,\AgdaUnderscore{}===\AgdaUnderscore{}}}\AgdaSpace{}%
\AgdaSymbol{:}\AgdaSpace{}%
\AgdaSymbol{(}\AgdaBound{n}\AgdaSpace{}%
\AgdaSymbol{:}\AgdaSpace{}%
\AgdaDatatype{ℕ}\AgdaSymbol{)}\AgdaSpace{}%
\AgdaSymbol{→}\AgdaSpace{}%
\AgdaFunction{WRel}\AgdaSpace{}%
\AgdaSymbol{(}\AgdaPostulate{Gen}\AgdaSpace{}%
\AgdaBound{n}\AgdaSymbol{)}\<%
\\
%
\>[2]\AgdaPostulate{Sp-group}\AgdaSpace{}%
\AgdaPostulate{Pauli⋊Sp}\AgdaSpace{}%
\AgdaSymbol{:}\AgdaSpace{}%
\AgdaDatatype{ℕ}\AgdaSpace{}%
\AgdaSymbol{→}\AgdaSpace{}%
\AgdaRecord{Group}\AgdaSpace{}%
\AgdaFunction{0ℓ}\AgdaSpace{}%
\AgdaFunction{0ℓ}\<%
\\
%
\\[\AgdaEmptyExtraSkip]%
\>[0]\AgdaKeyword{private}\AgdaSpace{}%
\AgdaKeyword{variable}\<%
\\
\>[0][@{}l@{\AgdaIndent{0}}]%
\>[2]\AgdaGeneralizable{n}\AgdaSpace{}%
\AgdaSymbol{:}\AgdaSpace{}%
\AgdaDatatype{ℕ}\<%
\end{code}
%% Part 6b — qupit Clifford circuits (branches main, qupit, ...).

\begin{frame}{Qupit Clifford circuits in all odd prime dimensions \New}
\begin{columns}[c]
\begin{column}{0.40\textwidth}
\small
\textbf{Bian, Li, Ross, van de Wetering, Zhao}, \emph{A Complete and Natural
Rule Set for Multi-Qudit Clifford Circuits in All Odd Prime Dimensions}
(QPL'26). Its verification is a separate paper in preparation.
\begin{itemize}
  \item gates $H$, $S$ (one qupit), $CZ$ (two), any width $n$, any odd prime $p$;
  \item read \alert{modulo scalars}, the circuits present
        \[\text{Pauli}(n)\rtimes Sp(2n,\mathbb{Z}_p);\]
  \item rules over $p$, a \emph{primitive root} $g$, and $\frac12\in\mathbb{Z}_p$;
  \item 16 rules in the paper (Figure 1); 15 here: $(SH)^3=\varepsilon$
        becomes a theorem.
\end{itemize}
\end{column}
\begin{column}{0.58\textwidth}
\renewcommand{\AgdaCodeStyle}{\scriptsize}
\begin{code}%
\>[0]\AgdaKeyword{data}\AgdaSpace{}%
\AgdaOperator{\AgdaDatatype{\AgdaUnderscore{}SRel,\AgdaUnderscore{}===\AgdaUnderscore{}}}\AgdaSpace{}%
\AgdaSymbol{:}\AgdaSpace{}%
\AgdaSymbol{(}\AgdaBound{n}\AgdaSpace{}%
\AgdaSymbol{:}\AgdaSpace{}%
\AgdaDatatype{ℕ}\AgdaSymbol{)}\AgdaSpace{}%
\AgdaSymbol{→}\AgdaSpace{}%
\AgdaFunction{WRel}\AgdaSpace{}%
\AgdaSymbol{(}\AgdaPostulate{Gen}\AgdaSpace{}%
\AgdaBound{n}\AgdaSymbol{)}\AgdaSpace{}%
\AgdaKeyword{where}\<%
\\
\>[0][@{}l@{\AgdaIndent{0}}]%
\>[2]\AgdaInductiveConstructor{order-S}\AgdaSpace{}%
\AgdaSymbol{:}%
\>[18]\AgdaSymbol{(}\AgdaInductiveConstructor{₁₊}\AgdaSpace{}%
\AgdaGeneralizable{n}\AgdaSymbol{)}\AgdaSpace{}%
\AgdaOperator{\AgdaDatatype{SRel,}}%
\>[32]\AgdaPostulate{S}\AgdaSpace{}%
\AgdaOperator{\AgdaPostulate{\textasciicircum{}}}\AgdaSpace{}%
\AgdaPostulate{p}\AgdaSpace{}%
\AgdaOperator{\AgdaDatatype{===}}\AgdaSpace{}%
\AgdaInductiveConstructor{ε}\<%
\\
%
\>[2]\AgdaInductiveConstructor{order-H}\AgdaSpace{}%
\AgdaSymbol{:}%
\>[18]\AgdaSymbol{(}\AgdaInductiveConstructor{₁₊}\AgdaSpace{}%
\AgdaGeneralizable{n}\AgdaSymbol{)}\AgdaSpace{}%
\AgdaOperator{\AgdaDatatype{SRel,}}%
\>[32]\AgdaPostulate{H}\AgdaSpace{}%
\AgdaOperator{\AgdaPostulate{\textasciicircum{}}}\AgdaSpace{}%
\AgdaNumber{2}\AgdaSpace{}%
\AgdaOperator{\AgdaDatatype{===}}\AgdaSpace{}%
\AgdaPostulate{XM₋₁}\<%
\\
%
\>[2]\AgdaInductiveConstructor{M-power}\AgdaSpace{}%
\AgdaSymbol{:}\AgdaSpace{}%
\AgdaSymbol{∀}\AgdaSpace{}%
\AgdaBound{k}\AgdaSpace{}%
\AgdaSymbol{→}\AgdaSpace{}%
\AgdaSymbol{(}\AgdaInductiveConstructor{₁₊}\AgdaSpace{}%
\AgdaGeneralizable{n}\AgdaSymbol{)}\AgdaSpace{}%
\AgdaOperator{\AgdaDatatype{SRel,}}%
\>[32]\AgdaOperator{\AgdaPostulate{XMg\textasciicircum{}}}\AgdaSpace{}%
\AgdaBound{k}\AgdaSpace{}%
\AgdaOperator{\AgdaDatatype{===}}\AgdaSpace{}%
\AgdaPostulate{XM}\AgdaSpace{}%
\AgdaSymbol{(}\AgdaOperator{\AgdaPostulate{g\textasciicircum{}}}\AgdaSpace{}%
\AgdaBound{k}\AgdaSymbol{)}\<%
\\
%
\>[2]\AgdaInductiveConstructor{semi-MR}\AgdaSpace{}%
\AgdaSymbol{:}%
\>[18]\AgdaSymbol{(}\AgdaInductiveConstructor{₁₊}\AgdaSpace{}%
\AgdaGeneralizable{n}\AgdaSymbol{)}\AgdaSpace{}%
\AgdaOperator{\AgdaDatatype{SRel,}}%
\>[32]\AgdaPostulate{XMg}\AgdaSpace{}%
\AgdaOperator{\AgdaInductiveConstructor{•}}\AgdaSpace{}%
\AgdaOperator{\AgdaPostulate{S\textasciicircum{}}}\AgdaSpace{}%
\AgdaSymbol{(}\AgdaPostulate{g}\AgdaSpace{}%
\AgdaOperator{\AgdaPostulate{*}}\AgdaSpace{}%
\AgdaPostulate{g}\AgdaSymbol{)}\AgdaSpace{}%
\AgdaOperator{\AgdaInductiveConstructor{•}}\AgdaSpace{}%
\AgdaOperator{\AgdaPostulate{Z\textasciicircum{}}}\AgdaSpace{}%
\AgdaSymbol{((}\AgdaPostulate{g}\AgdaSpace{}%
\AgdaOperator{\AgdaPostulate{*}}\AgdaSpace{}%
\AgdaPostulate{g}\AgdaSpace{}%
\AgdaOperator{\AgdaPostulate{+}}\AgdaSpace{}%
\AgdaOperator{\AgdaPostulate{-}}\AgdaSpace{}%
\AgdaPostulate{g}\AgdaSymbol{)}\AgdaSpace{}%
\AgdaOperator{\AgdaPostulate{*}}\AgdaSpace{}%
\AgdaPostulate{1/2}\AgdaSymbol{)}\<%
\\
%
\>[32]\AgdaOperator{\AgdaDatatype{===}}\AgdaSpace{}%
\AgdaPostulate{S}\AgdaSpace{}%
\AgdaOperator{\AgdaInductiveConstructor{•}}\AgdaSpace{}%
\AgdaPostulate{XMg}\<%
\\
%
\>[2]\AgdaInductiveConstructor{comm-HHSHHS}\AgdaSpace{}%
\AgdaSymbol{:}%
\>[18]\AgdaSymbol{(}\AgdaInductiveConstructor{₁₊}\AgdaSpace{}%
\AgdaGeneralizable{n}\AgdaSymbol{)}\AgdaSpace{}%
\AgdaOperator{\AgdaDatatype{SRel,}}%
\>[32]\AgdaPostulate{H}\AgdaSpace{}%
\AgdaOperator{\AgdaInductiveConstructor{•}}\AgdaSpace{}%
\AgdaPostulate{H}\AgdaSpace{}%
\AgdaOperator{\AgdaInductiveConstructor{•}}\AgdaSpace{}%
\AgdaPostulate{S}\AgdaSpace{}%
\AgdaOperator{\AgdaInductiveConstructor{•}}\AgdaSpace{}%
\AgdaPostulate{H}\AgdaSpace{}%
\AgdaOperator{\AgdaInductiveConstructor{•}}\AgdaSpace{}%
\AgdaPostulate{H}\AgdaSpace{}%
\AgdaOperator{\AgdaInductiveConstructor{•}}\AgdaSpace{}%
\AgdaPostulate{S}\AgdaSpace{}%
\AgdaOperator{\AgdaDatatype{===}}\AgdaSpace{}%
\AgdaPostulate{S}\AgdaSpace{}%
\AgdaOperator{\AgdaInductiveConstructor{•}}\AgdaSpace{}%
\AgdaPostulate{H}\AgdaSpace{}%
\AgdaOperator{\AgdaInductiveConstructor{•}}\AgdaSpace{}%
\AgdaPostulate{H}\AgdaSpace{}%
\AgdaOperator{\AgdaInductiveConstructor{•}}\AgdaSpace{}%
\AgdaPostulate{S}\AgdaSpace{}%
\AgdaOperator{\AgdaInductiveConstructor{•}}\AgdaSpace{}%
\AgdaPostulate{H}\AgdaSpace{}%
\AgdaOperator{\AgdaInductiveConstructor{•}}\AgdaSpace{}%
\AgdaPostulate{H}\<%
\\
%
\>[2]\AgdaInductiveConstructor{order-CZ}\AgdaSpace{}%
\AgdaSymbol{:}%
\>[18]\AgdaSymbol{(}\AgdaInductiveConstructor{₂₊}\AgdaSpace{}%
\AgdaGeneralizable{n}\AgdaSymbol{)}\AgdaSpace{}%
\AgdaOperator{\AgdaDatatype{SRel,}}%
\>[32]\AgdaPostulate{CZ}\AgdaSpace{}%
\AgdaOperator{\AgdaPostulate{\textasciicircum{}}}\AgdaSpace{}%
\AgdaPostulate{p}\AgdaSpace{}%
\AgdaOperator{\AgdaDatatype{===}}\AgdaSpace{}%
\AgdaInductiveConstructor{ε}\<%
\\
%
\>[2]\AgdaInductiveConstructor{order-Ex}\AgdaSpace{}%
\AgdaSymbol{:}%
\>[18]\AgdaSymbol{(}\AgdaInductiveConstructor{₂₊}\AgdaSpace{}%
\AgdaGeneralizable{n}\AgdaSymbol{)}\AgdaSpace{}%
\AgdaOperator{\AgdaDatatype{SRel,}}%
\>[32]\AgdaPostulate{Ex}\AgdaSpace{}%
\AgdaOperator{\AgdaPostulate{\textasciicircum{}}}\AgdaSpace{}%
\AgdaNumber{2}\AgdaSpace{}%
\AgdaOperator{\AgdaDatatype{===}}\AgdaSpace{}%
\AgdaInductiveConstructor{ε}\<%
\\
%
\>[2]\AgdaInductiveConstructor{comm-CZ-S↑}\AgdaSpace{}%
\AgdaSymbol{:}%
\>[18]\AgdaSymbol{(}\AgdaInductiveConstructor{₂₊}\AgdaSpace{}%
\AgdaGeneralizable{n}\AgdaSymbol{)}\AgdaSpace{}%
\AgdaOperator{\AgdaDatatype{SRel,}}%
\>[32]\AgdaPostulate{CZ}\AgdaSpace{}%
\AgdaOperator{\AgdaInductiveConstructor{•}}\AgdaSpace{}%
\AgdaPostulate{S}\AgdaSpace{}%
\AgdaOperator{\AgdaPostulate{↑}}\AgdaSpace{}%
\AgdaOperator{\AgdaDatatype{===}}\AgdaSpace{}%
\AgdaPostulate{S}\AgdaSpace{}%
\AgdaOperator{\AgdaPostulate{↑}}\AgdaSpace{}%
\AgdaOperator{\AgdaInductiveConstructor{•}}\AgdaSpace{}%
\AgdaPostulate{CZ}\<%
\\
%
\>[2]\AgdaInductiveConstructor{semi-M↑CZ}\AgdaSpace{}%
\AgdaSymbol{:}%
\>[18]\AgdaSymbol{(}\AgdaInductiveConstructor{₂₊}\AgdaSpace{}%
\AgdaGeneralizable{n}\AgdaSymbol{)}\AgdaSpace{}%
\AgdaOperator{\AgdaDatatype{SRel,}}%
\>[32]\AgdaPostulate{XMg}\AgdaSpace{}%
\AgdaOperator{\AgdaPostulate{↑}}\AgdaSpace{}%
\AgdaOperator{\AgdaInductiveConstructor{•}}\AgdaSpace{}%
\AgdaOperator{\AgdaPostulate{CZ\textasciicircum{}}}\AgdaSpace{}%
\AgdaPostulate{g}\AgdaSpace{}%
\AgdaOperator{\AgdaDatatype{===}}\AgdaSpace{}%
\AgdaPostulate{CZ}\AgdaSpace{}%
\AgdaOperator{\AgdaInductiveConstructor{•}}\AgdaSpace{}%
\AgdaPostulate{XMg}\AgdaSpace{}%
\AgdaOperator{\AgdaPostulate{↑}}\<%
\\
%
\>[2]\AgdaInductiveConstructor{semi-Ex-S↑}\AgdaSpace{}%
\AgdaSymbol{:}%
\>[18]\AgdaSymbol{(}\AgdaInductiveConstructor{₂₊}\AgdaSpace{}%
\AgdaGeneralizable{n}\AgdaSymbol{)}\AgdaSpace{}%
\AgdaOperator{\AgdaDatatype{SRel,}}%
\>[32]\AgdaPostulate{Ex}\AgdaSpace{}%
\AgdaOperator{\AgdaInductiveConstructor{•}}\AgdaSpace{}%
\AgdaPostulate{S}\AgdaSpace{}%
\AgdaOperator{\AgdaPostulate{↑}}\AgdaSpace{}%
\AgdaOperator{\AgdaDatatype{===}}\AgdaSpace{}%
\AgdaPostulate{S}\AgdaSpace{}%
\AgdaOperator{\AgdaPostulate{↓}}\AgdaSpace{}%
\AgdaOperator{\AgdaInductiveConstructor{•}}\AgdaSpace{}%
\AgdaPostulate{Ex}\<%
\\
%
\>[2]\AgdaInductiveConstructor{semi-Ex-H↑}\AgdaSpace{}%
\AgdaSymbol{:}%
\>[18]\AgdaSymbol{(}\AgdaInductiveConstructor{₂₊}\AgdaSpace{}%
\AgdaGeneralizable{n}\AgdaSymbol{)}\AgdaSpace{}%
\AgdaOperator{\AgdaDatatype{SRel,}}%
\>[32]\AgdaPostulate{Ex}\AgdaSpace{}%
\AgdaOperator{\AgdaInductiveConstructor{•}}\AgdaSpace{}%
\AgdaPostulate{H}\AgdaSpace{}%
\AgdaOperator{\AgdaPostulate{↑}}\AgdaSpace{}%
\AgdaOperator{\AgdaDatatype{===}}\AgdaSpace{}%
\AgdaPostulate{H}\AgdaSpace{}%
\AgdaOperator{\AgdaPostulate{↓}}\AgdaSpace{}%
\AgdaOperator{\AgdaInductiveConstructor{•}}\AgdaSpace{}%
\AgdaPostulate{Ex}\<%
\\
%
\>[2]\AgdaInductiveConstructor{blake-c12}\AgdaSpace{}%
\AgdaSymbol{:}%
\>[18]\AgdaSymbol{(}\AgdaInductiveConstructor{₂₊}\AgdaSpace{}%
\AgdaGeneralizable{n}\AgdaSymbol{)}\AgdaSpace{}%
\AgdaOperator{\AgdaDatatype{SRel,}}%
\>[32]\AgdaSymbol{(}\AgdaPostulate{S}\AgdaSpace{}%
\AgdaOperator{\AgdaPostulate{\textasciicircum{}}}\AgdaSpace{}%
\AgdaPostulate{p-1}\AgdaSymbol{)}\AgdaSpace{}%
\AgdaOperator{\AgdaPostulate{↑}}\AgdaSpace{}%
\AgdaOperator{\AgdaInductiveConstructor{•}}\AgdaSpace{}%
\AgdaSymbol{(}\AgdaPostulate{S}\AgdaSpace{}%
\AgdaOperator{\AgdaPostulate{\textasciicircum{}}}\AgdaSpace{}%
\AgdaPostulate{p-1}\AgdaSymbol{)}\AgdaSpace{}%
\AgdaOperator{\AgdaPostulate{↓}}\AgdaSpace{}%
\AgdaOperator{\AgdaInductiveConstructor{•}}\AgdaSpace{}%
\AgdaPostulate{CX}\AgdaSpace{}%
\AgdaOperator{\AgdaPostulate{\textasciicircum{}}}\AgdaSpace{}%
\AgdaPostulate{p-1}\AgdaSpace{}%
\AgdaOperator{\AgdaInductiveConstructor{•}}\AgdaSpace{}%
\AgdaPostulate{S}\AgdaSpace{}%
\AgdaOperator{\AgdaPostulate{↓}}\AgdaSpace{}%
\AgdaOperator{\AgdaInductiveConstructor{•}}\AgdaSpace{}%
\AgdaPostulate{CX}\AgdaSpace{}%
\AgdaOperator{\AgdaDatatype{===}}\AgdaSpace{}%
\AgdaPostulate{CZ}\<%
\\
%
\>[2]\AgdaInductiveConstructor{yang-baxter}\AgdaSpace{}%
\AgdaSymbol{:}%
\>[18]\AgdaSymbol{(}\AgdaInductiveConstructor{₃₊}\AgdaSpace{}%
\AgdaGeneralizable{n}\AgdaSymbol{)}\AgdaSpace{}%
\AgdaOperator{\AgdaDatatype{SRel,}}%
\>[32]\AgdaPostulate{Ex}\AgdaSpace{}%
\AgdaOperator{\AgdaPostulate{↑}}\AgdaSpace{}%
\AgdaOperator{\AgdaInductiveConstructor{•}}\AgdaSpace{}%
\AgdaPostulate{Ex}\AgdaSpace{}%
\AgdaOperator{\AgdaPostulate{↓}}\AgdaSpace{}%
\AgdaOperator{\AgdaInductiveConstructor{•}}\AgdaSpace{}%
\AgdaPostulate{Ex}\AgdaSpace{}%
\AgdaOperator{\AgdaPostulate{↑}}\AgdaSpace{}%
\AgdaOperator{\AgdaDatatype{===}}\AgdaSpace{}%
\AgdaPostulate{Ex}\AgdaSpace{}%
\AgdaOperator{\AgdaPostulate{↓}}\AgdaSpace{}%
\AgdaOperator{\AgdaInductiveConstructor{•}}\AgdaSpace{}%
\AgdaPostulate{Ex}\AgdaSpace{}%
\AgdaOperator{\AgdaPostulate{↑}}\AgdaSpace{}%
\AgdaOperator{\AgdaInductiveConstructor{•}}\AgdaSpace{}%
\AgdaPostulate{Ex}\AgdaSpace{}%
\AgdaOperator{\AgdaPostulate{↓}}\<%
\\
%
\>[2]\AgdaInductiveConstructor{cz-slide}\AgdaSpace{}%
\AgdaSymbol{:}%
\>[18]\AgdaSymbol{(}\AgdaInductiveConstructor{₃₊}\AgdaSpace{}%
\AgdaGeneralizable{n}\AgdaSymbol{)}\AgdaSpace{}%
\AgdaOperator{\AgdaDatatype{SRel,}}%
\>[32]\AgdaPostulate{Ex}\AgdaSpace{}%
\AgdaOperator{\AgdaPostulate{↓}}\AgdaSpace{}%
\AgdaOperator{\AgdaInductiveConstructor{•}}\AgdaSpace{}%
\AgdaPostulate{Ex}\AgdaSpace{}%
\AgdaOperator{\AgdaPostulate{↑}}\AgdaSpace{}%
\AgdaOperator{\AgdaInductiveConstructor{•}}\AgdaSpace{}%
\AgdaPostulate{CZ}\AgdaSpace{}%
\AgdaOperator{\AgdaDatatype{===}}\AgdaSpace{}%
\AgdaPostulate{CZ}\AgdaSpace{}%
\AgdaOperator{\AgdaPostulate{↑}}\AgdaSpace{}%
\AgdaOperator{\AgdaInductiveConstructor{•}}\AgdaSpace{}%
\AgdaPostulate{Ex}\AgdaSpace{}%
\AgdaOperator{\AgdaPostulate{↓}}\AgdaSpace{}%
\AgdaOperator{\AgdaInductiveConstructor{•}}\AgdaSpace{}%
\AgdaPostulate{Ex}\AgdaSpace{}%
\AgdaOperator{\AgdaPostulate{↑}}\<%
\\
%
\>[2]\AgdaInductiveConstructor{semi-CX↑-CZ↓}\AgdaSpace{}%
\AgdaSymbol{:}%
\>[18]\AgdaSymbol{(}\AgdaInductiveConstructor{₃₊}\AgdaSpace{}%
\AgdaGeneralizable{n}\AgdaSymbol{)}\AgdaSpace{}%
\AgdaOperator{\AgdaDatatype{SRel,}}%
\>[32]\AgdaPostulate{CZ}\AgdaSpace{}%
\AgdaOperator{\AgdaPostulate{↓}}\AgdaSpace{}%
\AgdaOperator{\AgdaInductiveConstructor{•}}\AgdaSpace{}%
\AgdaPostulate{CX}\AgdaSpace{}%
\AgdaOperator{\AgdaPostulate{↑}}\AgdaSpace{}%
\AgdaOperator{\AgdaDatatype{===}}\AgdaSpace{}%
\AgdaPostulate{CZ02}\AgdaSpace{}%
\AgdaOperator{\AgdaInductiveConstructor{•}}\AgdaSpace{}%
\AgdaPostulate{CX}\AgdaSpace{}%
\AgdaOperator{\AgdaPostulate{↑}}\AgdaSpace{}%
\AgdaOperator{\AgdaInductiveConstructor{•}}\AgdaSpace{}%
\AgdaPostulate{CZ}\AgdaSpace{}%
\AgdaOperator{\AgdaPostulate{↓}}\<%
\end{code}
{\scriptsize \texttt{ProjectiveClifford/Qupit/Paper-V1/Syntactics.agda} (\Branch{qupit}), one comment elided}
\end{column}
\end{columns}
\end{frame}

\begin{frame}{The rules as circuits (one and two qupits) \New}
\begin{columns}[T]
\begin{column}{0.46\textwidth}\tiny
\begin{tabular}{@{}r@{\;\,}r@{\,}c@{\,}l@{}}
order-S & \qfig{pv0-order-S-l} & $\equiv$ & \qfig{pv0-order-S-r}\\[2pt]
order-H & \qfig{pv0-order-H-l} & $\equiv$ & \qfig{pv0-order-H-r}\\[2pt]
M-power & \qfig{pv0-M-power-l} & $\equiv$ & \qfig{pv0-M-power-r}\\[2pt]
semi-MR & \qfig{pv0-semi-MR-l} & $\equiv$ & \qfig{pv0-semi-MR-r}\\[2pt]
order-CZ & \qfig{pv0-order-CZ-l} & $\equiv$ & \qfig{pv0-order-CZ-r}\\[2pt]
order-Ex & \qfig{pv0-order-Ex-l} & $\equiv$ & \qfig{pv0-order-Ex-r}\\[2pt]
order-SH & \qfig{pv0-order-SH-l} & $\equiv$ & \qfig{pv0-order-SH-r}
\end{tabular}
\end{column}
\begin{column}{0.52\textwidth}\tiny
\begin{tabular}{@{}r@{\;\,}l@{}}
comm-CZ-S↑ & \qfig{pv0-comm-CZ-Sup}\\[2pt]
semi-M↑CZ & \qfig{pv0-semi-Mup-CZ}\\[2pt]
semi-Ex-S↑ & \qfig{pv0-semi-Ex-Sup}\\[2pt]
semi-Ex-H↑ & \qfig{pv0-semi-Ex-Hup}\\[2pt]
comm-HHSHHS & \qfig[0.5]{pv0-comm-HHSHHS}
\end{tabular}
\end{column}
\end{columns}
\vspace{2pt}
{\scriptsize Figure 1 of the QPL'26 paper (the multiplier spelled over $R=S\,Z^{1/2}$);
figures generated by Cir2Tikz from the Agda terms.}
\end{frame}

\begin{frame}{\ldots\ three qupits, and the symplectic group underneath \New}
\begin{columns}[c]
\begin{column}{0.50\textwidth}\tiny
\begin{tabular}{@{}r@{\;\,}l@{}}
blake-c12 & \qfig{pv0-blake-c12}\\[3pt]
yang-baxter & \qfig{pv0-yang-baxter}\\[3pt]
cz-slide & \qfig{pv0-cz-slide}\\[3pt]
semi-CX↑-CZ↓ & \qfig{pv0-semi-CXup-CZdn}
\end{tabular}
\end{column}
\begin{column}{0.48\textwidth}
\small $Sp(2n,\mathbb{Z}_p)$ \textbf{by a coset tower} (\texttt{Examples/Groups/Symplectic}, 35k lines):
\begin{itemize}
  \item one level per wire, Selinger-style staircases as digits;
  \item one wire: $\mathbb{Z}_p\times\mathbb{Z}_p^{*}\times(\varepsilon\mid HS^a)$, i.e.
        $p\,(p{-}1)(p{+}1) = |Sp(2,p)|$ normal forms;
  \item \alert{67 $\to$ 18}: normalization uses 67 relations; a small set of
        natural ones derives them all.
\end{itemize}
\begin{code}%
\>[0]\AgdaComment{--\ MainTheorems,\ Symplectic-Simplified-Theorems}\<%
\\
\>[0]\AgdaFunction{simplified-presentation}\AgdaSpace{}%
\AgdaSymbol{:}\<%
\\
\>[0][@{}l@{\AgdaIndent{0}}]%
\>[2]\AgdaSymbol{∀}\AgdaSpace{}%
\AgdaBound{n}\AgdaSpace{}%
\AgdaSymbol{→}\AgdaSpace{}%
\AgdaSymbol{(}\AgdaBound{n}\AgdaSpace{}%
\AgdaOperator{\AgdaPostulate{QRel,\AgdaUnderscore{}===\AgdaUnderscore{}}}\AgdaSymbol{)}\AgdaSpace{}%
\AgdaOperator{\AgdaRecord{IsPresentationOf}}\AgdaSpace{}%
\AgdaSymbol{(}\AgdaPostulate{Sp-group}\AgdaSpace{}%
\AgdaBound{n}\AgdaSymbol{)}\<%
\end{code}
\begin{code}[hide]%
\>[0]\AgdaFunction{simplified-presentation}\AgdaSpace{}%
\AgdaSymbol{=}\AgdaSpace{}%
\AgdaSymbol{λ}\AgdaSpace{}%
\AgdaBound{\AgdaUnderscore{}}\AgdaSpace{}%
\AgdaSymbol{→}\AgdaSpace{}%
\AgdaSymbol{\AgdaUnderscore{}}\<%
\end{code}
\end{column}
\end{columns}
\end{frame}

\begin{frame}{Assembled, not re-proved \New}
\begin{columns}[c]
\begin{column}{0.52\textwidth}
\centering
\begin{tikzpicture}[node distance=5mm and 6mm,>={Stealth[length=1.8mm]},
  bx/.style={draw=accent,fill=soft,rounded corners=2pt,font=\scriptsize,align=center,minimum height=8mm,text width=26mm},
  nb/.style={draw=newC,fill=softnew,rounded corners=2pt,font=\scriptsize,align=center,minimum height=8mm,text width=26mm}]
  \node[bx] (pl) {ProjectivePauli\\{\tiny $(\mathbb{Z}_p\times\mathbb{Z}_p)^n$, the paper's Pauli}};
  \node[bx,right=of pl] (sp) {Symplectic.Simplified\\{\tiny coset tower, 18 rules}};
  \node[nb,below=9mm of $(pl)!0.5!(sp)$] (sd) {SemiDirect\\{\tiny $\text{Pauli}\rtimes Sp$, generic $\rtimes$ theorem}};
  \node[nb,below=of sd] (v1) {Simplified-V1\\{\tiny the Pauli calculus lives here}};
  \node[nb,below=of v1] (p0) {Paper-V0 (16 rules)\\{\tiny Figure 1}};
  \node[nb,below=of p0] (p1) {Paper-V1 (15 rules)};
  \draw[->,thick] (pl) -- (sd); \draw[->,thick] (sp) -- (sd);
  \draw[<->,thick,newC] (sd) -- node[right,font=\tiny]{iso} (v1);
  \draw[<->,thick,newC] (v1) -- node[right,font=\tiny]{iso, exports \aid{v1⇒pap}} (p0);
  \draw[<->,thick,newC] (p0) -- node[right,font=\tiny]{identity on words} (p1);
  \node[font=\tiny,gray,align=left,anchor=west] at ($(sd.east)+(1mm,0)$) {+ two well-definedness\\proofs of the action};
\end{tikzpicture}
\end{column}
\begin{column}{0.46\textwidth}
\small Every layer above the base is \alert{an isomorphism plus a composition};
no rule set re-proves what another has.
\begin{code}%
\>[0]\AgdaComment{--\ MainTheorems,\ Qupit-Clifford-Theorems}\<%
\\
\>[0]\AgdaFunction{clifford-presentation}\AgdaSpace{}%
\AgdaSymbol{:}\<%
\\
\>[0][@{}l@{\AgdaIndent{0}}]%
\>[2]\AgdaSymbol{∀}\AgdaSpace{}%
\AgdaBound{n}\AgdaSpace{}%
\AgdaSymbol{→}\AgdaSpace{}%
\AgdaSymbol{(}\AgdaBound{n}\AgdaSpace{}%
\AgdaOperator{\AgdaPostulate{QRel,\AgdaUnderscore{}===\AgdaUnderscore{}}}\AgdaSymbol{)}\AgdaSpace{}%
\AgdaOperator{\AgdaRecord{IsPresentationOf}}\AgdaSpace{}%
\AgdaSymbol{(}\AgdaPostulate{Pauli⋊Sp}\AgdaSpace{}%
\AgdaBound{n}\AgdaSymbol{)}\<%
\end{code}
\begin{code}[hide]%
\>[0]\AgdaFunction{clifford-presentation}\AgdaSpace{}%
\AgdaSymbol{=}\AgdaSpace{}%
\AgdaSymbol{λ}\AgdaSpace{}%
\AgdaBound{\AgdaUnderscore{}}\AgdaSpace{}%
\AgdaSymbol{→}\AgdaSpace{}%
\AgdaSymbol{\AgdaUnderscore{}}\<%
\end{code}
\vspace{-2pt}
{\scriptsize (\aid{PapSyn.Clifford-Relations.\_QRel,\_===\_} in the source)}
\begin{itemize}
  \item reuses the paper's machinery verbatim: coset towers, the
        word-valued semidirect product, \aid{StarIsomorphism};
  \item derivability results: two Pauli-vs-$CZ$ axioms dropped (18 $\to$ 16);
        $(SH)^3=\varepsilon$ a theorem in V1;
  \item $\approx$22k lines for the four rule sets.
\end{itemize}
\end{column}
\end{columns}
\end{frame}

\begin{frame}{Getting the scalars back: exact Clifford groups \New}
\begin{columns}[c]
\begin{column}{0.50\textwidth}
\small The Clifford group is a \alert{central extension} of the projective one by
$\langle\omega\rangle$. Add $\omega$, $\omega^p=\varepsilon$, centrality, and
\alert{correct} each relation by the scalar it picks up --- only one needs it:
\vspace{4pt}

\centering\qfig[0.85]{sc-order-SH}

{\scriptsize ($\frac18$ the inverse of 8 in $\mathbb{Z}_p$)}
\vspace{4pt}

\raggedright
\small \texttt{Clifford/Qupit}: \aid{presentation-exact}, against the central
extension $A\times_\gamma H$ with $\gamma$ the Weyl cocycle $\frac12\,\text{sform}$;
new \texttt{ForStdlib}: \aid{CentralExtension}, \aid{FactorSetExtension}
(Rotman 9.8), \aid{CentralProduct}.
\end{column}
\begin{column}{0.47\textwidth}
\small \textbf{Qubits, after Selinger} (LMCS 2015), \texttt{ProjectiveClifford/Qubit},
\texttt{Clifford/Qubit}:
\begin{itemize}
  \item Figure 8 modulo scalars presents $V\!Sp(n)$ (symplectic map + $\mathbb{Z}_4$
        phase function); $\cong$ the Clifford action on signed Paulis;
  \item an \alert{extension} presentation $1\to\mathbb{Z}_8\to\text{Exact}(n)\to
        \text{CMS}(n)\to1$, isomorphic to Figure 8;
  \item that $\omega$ has order exactly 8 needs a model: $2^n\times2^n$
        matrices over $\mathbb{Z}/17$ with $i=4$, $\frac1{\sqrt2}=14$,
        $\omega=2$ --- every relation by \aid{refl} on tries.
\end{itemize}
\begin{newdevbox}
\small A first, small \alert{matrix} semantics --- next: real ones.
\end{newdevbox}
\end{column}
\end{columns}
\end{frame}

\begin{code}[hide]%
\>[0]\AgdaKeyword{module}\AgdaSpace{}%
\AgdaModule{slides.Tsinghua-PL-seminar.sections.newmatrix}\AgdaSpace{}%
\AgdaKeyword{where}\<%
\\
%
\\[\AgdaEmptyExtraSkip]%
\>[0]\AgdaKeyword{open}\AgdaSpace{}%
\AgdaKeyword{import}\AgdaSpace{}%
\AgdaModule{Data.Nat}\AgdaSpace{}%
\AgdaKeyword{using}\AgdaSpace{}%
\AgdaSymbol{(}\AgdaDatatype{ℕ}\AgdaSymbol{)}\<%
\\
\>[0]\AgdaKeyword{open}\AgdaSpace{}%
\AgdaKeyword{import}\AgdaSpace{}%
\AgdaModule{Data.Product}\AgdaSpace{}%
\AgdaKeyword{using}\AgdaSpace{}%
\AgdaSymbol{(}\AgdaField{proj₁}\AgdaSymbol{)}\<%
\\
\>[0]\AgdaKeyword{open}\AgdaSpace{}%
\AgdaKeyword{import}\AgdaSpace{}%
\AgdaModule{Level}\AgdaSpace{}%
\AgdaKeyword{using}\AgdaSpace{}%
\AgdaSymbol{(}\AgdaFunction{0ℓ}\AgdaSymbol{)}\<%
\\
\>[0]\AgdaKeyword{open}\AgdaSpace{}%
\AgdaKeyword{import}\AgdaSpace{}%
\AgdaModule{Relation.Binary.PropositionalEquality}\AgdaSpace{}%
\AgdaKeyword{using}\AgdaSpace{}%
\AgdaSymbol{(}\AgdaOperator{\AgdaDatatype{\AgdaUnderscore{}≡\AgdaUnderscore{}}}\AgdaSpace{}%
\AgdaSymbol{;}\AgdaSpace{}%
\AgdaInductiveConstructor{refl}\AgdaSymbol{)}\<%
\\
\>[0]\AgdaKeyword{open}\AgdaSpace{}%
\AgdaKeyword{import}\AgdaSpace{}%
\AgdaModule{Algebra.Bundles}\AgdaSpace{}%
\AgdaKeyword{using}\AgdaSpace{}%
\AgdaSymbol{(}\AgdaRecord{Group}\AgdaSymbol{)}\<%
\\
\>[0]\AgdaKeyword{open}\AgdaSpace{}%
\AgdaKeyword{import}\AgdaSpace{}%
\AgdaModule{Word.Base}\AgdaSpace{}%
\AgdaKeyword{using}\AgdaSpace{}%
\AgdaSymbol{(}\AgdaDatatype{Word}\AgdaSpace{}%
\AgdaSymbol{;}\AgdaSpace{}%
\AgdaFunction{WRel}\AgdaSymbol{)}\<%
\\
\>[0]\AgdaKeyword{open}\AgdaSpace{}%
\AgdaKeyword{import}\AgdaSpace{}%
\AgdaModule{Presentation.Definitions}\AgdaSpace{}%
\AgdaKeyword{using}\AgdaSpace{}%
\AgdaSymbol{(}\AgdaOperator{\AgdaRecord{\AgdaUnderscore{}IsPresentationOf\AgdaUnderscore{}}}\AgdaSpace{}%
\AgdaSymbol{;}\AgdaSpace{}%
\AgdaOperator{\AgdaRecord{\AgdaUnderscore{}IsSubPresentationOf\AgdaUnderscore{}}}\AgdaSymbol{)}\<%
\\
%
\\[\AgdaEmptyExtraSkip]%
\>[0]\AgdaKeyword{infix}\AgdaSpace{}%
\AgdaNumber{4}\AgdaSpace{}%
\AgdaOperator{\AgdaPostulate{\AgdaUnderscore{}≈\AgdaUnderscore{}}}\AgdaSpace{}%
\AgdaOperator{\AgdaPostulate{\AgdaUnderscore{}⊢\AgdaUnderscore{}===\AgdaUnderscore{}}}\<%
\\
%
\\[\AgdaEmptyExtraSkip]%
\>[0]\AgdaKeyword{postulate}\<%
\\
\>[0][@{}l@{\AgdaIndent{0}}]%
\>[2]\AgdaPostulate{Gen}\AgdaSpace{}%
\AgdaSymbol{:}\AgdaSpace{}%
\AgdaDatatype{ℕ}\AgdaSpace{}%
\AgdaSymbol{→}\AgdaSpace{}%
\AgdaPrimitive{Set}\<%
\\
%
\>[2]\AgdaOperator{\AgdaPostulate{\AgdaUnderscore{}===\AgdaUnderscore{}}}\AgdaSpace{}%
\AgdaSymbol{:}\AgdaSpace{}%
\AgdaSymbol{∀}\AgdaSpace{}%
\AgdaSymbol{\{}\AgdaBound{n}\AgdaSymbol{\}}\AgdaSpace{}%
\AgdaSymbol{→}\AgdaSpace{}%
\AgdaFunction{WRel}\AgdaSpace{}%
\AgdaSymbol{(}\AgdaPostulate{Gen}\AgdaSpace{}%
\AgdaBound{n}\AgdaSymbol{)}\<%
\\
%
\>[2]\AgdaPostulate{U}\AgdaSpace{}%
\AgdaSymbol{:}\AgdaSpace{}%
\AgdaDatatype{ℕ}\AgdaSpace{}%
\AgdaSymbol{→}\AgdaSpace{}%
\AgdaRecord{Group}\AgdaSpace{}%
\AgdaFunction{0ℓ}\AgdaSpace{}%
\AgdaFunction{0ℓ}\<%
\\
%
\>[2]\AgdaOperator{\AgdaPostulate{⟦\AgdaUnderscore{}⟧ᵐ}}\AgdaSpace{}%
\AgdaSymbol{:}\AgdaSpace{}%
\AgdaSymbol{∀}\AgdaSpace{}%
\AgdaSymbol{\{}\AgdaBound{n}\AgdaSymbol{\}}\AgdaSpace{}%
\AgdaSymbol{→}\AgdaSpace{}%
\AgdaDatatype{Word}\AgdaSpace{}%
\AgdaSymbol{(}\AgdaPostulate{Gen}\AgdaSpace{}%
\AgdaBound{n}\AgdaSymbol{)}\AgdaSpace{}%
\AgdaSymbol{→}\AgdaSpace{}%
\AgdaDatatype{ℕ}\<%
\\
%
\>[2]\AgdaOperator{\AgdaPostulate{\AgdaUnderscore{}≈\AgdaUnderscore{}}}\AgdaSpace{}%
\AgdaSymbol{:}\AgdaSpace{}%
\AgdaSymbol{∀}\AgdaSpace{}%
\AgdaSymbol{\{}\AgdaBound{n}\AgdaSymbol{\}}\AgdaSpace{}%
\AgdaSymbol{→}\AgdaSpace{}%
\AgdaDatatype{Word}\AgdaSpace{}%
\AgdaSymbol{(}\AgdaPostulate{Gen}\AgdaSpace{}%
\AgdaBound{n}\AgdaSymbol{)}\AgdaSpace{}%
\AgdaSymbol{→}\AgdaSpace{}%
\AgdaDatatype{Word}\AgdaSpace{}%
\AgdaSymbol{(}\AgdaPostulate{Gen}\AgdaSpace{}%
\AgdaBound{n}\AgdaSymbol{)}\AgdaSpace{}%
\AgdaSymbol{→}\AgdaSpace{}%
\AgdaPrimitive{Set}\<%
\\
%
\>[2]\AgdaPostulate{l17}\AgdaSpace{}%
\AgdaPostulate{r17}\AgdaSpace{}%
\AgdaSymbol{:}\AgdaSpace{}%
\AgdaDatatype{Word}\AgdaSpace{}%
\AgdaSymbol{(}\AgdaPostulate{Gen}\AgdaSpace{}%
\AgdaNumber{4}\AgdaSymbol{)}\<%
\\
%
\>[2]\AgdaPostulate{Rel}\AgdaSpace{}%
\AgdaSymbol{:}\AgdaSpace{}%
\AgdaFunction{WRel}\AgdaSpace{}%
\AgdaSymbol{(}\AgdaPostulate{Gen}\AgdaSpace{}%
\AgdaNumber{4}\AgdaSymbol{)}\<%
\\
%
\>[2]\AgdaOperator{\AgdaPostulate{\AgdaUnderscore{}⊢\AgdaUnderscore{}===\AgdaUnderscore{}}}\AgdaSpace{}%
\AgdaSymbol{:}\AgdaSpace{}%
\AgdaFunction{WRel}\AgdaSpace{}%
\AgdaSymbol{(}\AgdaPostulate{Gen}\AgdaSpace{}%
\AgdaNumber{4}\AgdaSymbol{)}\AgdaSpace{}%
\AgdaSymbol{→}\AgdaSpace{}%
\AgdaFunction{WRel}\AgdaSpace{}%
\AgdaSymbol{(}\AgdaPostulate{Gen}\AgdaSpace{}%
\AgdaNumber{4}\AgdaSymbol{)}\<%
\\
%
\\[\AgdaEmptyExtraSkip]%
\>[0]\AgdaKeyword{private}\AgdaSpace{}%
\AgdaKeyword{variable}\<%
\\
\>[0][@{}l@{\AgdaIndent{0}}]%
\>[2]\AgdaGeneralizable{n}\AgdaSpace{}%
\AgdaSymbol{:}\AgdaSpace{}%
\AgdaDatatype{ℕ}\<%
\end{code}
%% Part 6c — matrix semantics for unitary groups (branch v-office).

\begin{frame}{Matrix semantics, for real: $U_n(\mathbb{Z}[\frac12,i])$ \New}
\begin{columns}[c]
\begin{column}{0.47\textwidth}
\small The paper's stated limitation: its $U_3(\mathbb{Z}[\frac12,i])$ and
Clifford+$T$ studies stop at syntax, because the rings and unitary matrices
were missing. \Branch{v-office} supplies them:
\begin{itemize}
  \item rings $\mathbb{Z}[\frac12]$, $\mathbb{Z}[\frac12,i]$, $\mathbb{Z}[\frac1{\sqrt2},i]$ and
        matrices from \textbf{EucDomain} --- an Agda port of Ross--Selinger's
        \texttt{newsynth} (17.6k lines);
  \item generators act as \alert{two-level} matrices:
\end{itemize}
\vspace{2pt}
\centering
\begin{tikzpicture}[x=2.75mm,y=2.75mm,font=\tiny]
  \foreach \M/\x/\lab in {X/0/{X_{[j,k]}},K/9/{K_{[j,k]}},I/18/{i_{[j]}}} {
    \draw[gray] (\x,0) rectangle ++(6,6);
    \node at (\x+3,-0.9) {$\lab$};
  }
  % X: swap rows 1,3 (j=1,k=3)
  \fill[newC!20] (1,4) rectangle (2,5); \fill[newC!20] (3,2) rectangle (4,3);
  \fill[newC!20] (1,2) rectangle (2,3); \fill[newC!20] (3,4) rectangle (4,5);
  \node at (1.5,4.5) {0}; \node at (3.5,4.5) {1}; \node at (1.5,2.5) {1}; \node at (3.5,2.5) {0};
  \foreach \d in {0,2,4,5} \node at (\d+0.5,5.5-\d) {1};
  % K
  \fill[accent!20] (10,4) rectangle (11,5); \fill[accent!20] (12,2) rectangle (13,3);
  \fill[accent!20] (10,2) rectangle (11,3); \fill[accent!20] (12,4) rectangle (13,5);
  \node at (10.5,4.5) {$\gamma$}; \node at (12.5,4.5) {$\gamma$}; \node at (10.5,2.5) {$\gamma$}; \node at (12.5,2.5) {$\!\!-\gamma$};
  \foreach \d in {0,2,4,5} \node at (9+\d+0.5,5.5-\d) {1};
  % i
  \fill[paperC!20] (19,4) rectangle (20,5); \node at (19.5,4.5) {$i$};
  \foreach \d in {0,2,3,4,5} \node at (18+\d+0.5,5.5-\d) {1};
  \node[anchor=west] at (24.3,3) {$\gamma = \frac{1}{1+i}$};
\end{tikzpicture}
\end{column}
\begin{column}{0.51\textwidth}
\renewcommand{\AgdaCodeStyle}{\scriptsize}
\small Bian--Selinger (QPL'21), \textbf{every} $n$, against matrices:
\begin{code}[hide]%
\>[0]\AgdaKeyword{module}\AgdaSpace{}%
\AgdaModule{Clifford+CS-TwoLevel}\AgdaSpace{}%
\AgdaKeyword{where}\<%
\end{code}
\begin{code}%
\>[0][@{}l@{\AgdaIndent{1}}]%
\>[2]\AgdaComment{--\ Presentation.agda:\ the\ relations\ present\ Uₙ(𝔻[i])}\<%
\\
%
\>[2]\AgdaFunction{presentation}\AgdaSpace{}%
\AgdaSymbol{:}\AgdaSpace{}%
\AgdaSymbol{(}\AgdaOperator{\AgdaPostulate{\AgdaUnderscore{}===\AgdaUnderscore{}}}\AgdaSpace{}%
\AgdaSymbol{\{}\AgdaGeneralizable{n}\AgdaSymbol{\})}\AgdaSpace{}%
\AgdaOperator{\AgdaRecord{IsPresentationOf}}\AgdaSpace{}%
\AgdaSymbol{(}\AgdaPostulate{U}\AgdaSpace{}%
\AgdaGeneralizable{n}\AgdaSymbol{)}\<%
\\
%
\>[2]\AgdaComment{--\ MainLemma.agda}\<%
\\
%
\>[2]\AgdaFunction{relations-complete}\AgdaSpace{}%
\AgdaSymbol{:}%
\>[107I]\AgdaSymbol{\{}\AgdaBound{u}\AgdaSpace{}%
\AgdaBound{v}\AgdaSpace{}%
\AgdaSymbol{:}\AgdaSpace{}%
\AgdaDatatype{Word}\AgdaSpace{}%
\AgdaSymbol{(}\AgdaPostulate{Gen}\AgdaSpace{}%
\AgdaGeneralizable{n}\AgdaSymbol{)\}}\AgdaSpace{}%
\AgdaSymbol{→}\<%
\\
\>[.][@{}l@{}]\<[107I]%
\>[23]\AgdaOperator{\AgdaPostulate{⟦}}\AgdaSpace{}%
\AgdaBound{u}\AgdaSpace{}%
\AgdaOperator{\AgdaPostulate{⟧ᵐ}}\AgdaSpace{}%
\AgdaOperator{\AgdaDatatype{≡}}\AgdaSpace{}%
\AgdaOperator{\AgdaPostulate{⟦}}\AgdaSpace{}%
\AgdaBound{v}\AgdaSpace{}%
\AgdaOperator{\AgdaPostulate{⟧ᵐ}}\AgdaSpace{}%
\AgdaSymbol{→}\AgdaSpace{}%
\AgdaBound{u}\AgdaSpace{}%
\AgdaOperator{\AgdaPostulate{≈}}\AgdaSpace{}%
\AgdaBound{v}\<%
\\
%
\>[2]\AgdaComment{--\ Soundness.agda:\ decided\ on\ concrete\ 4\ ×\ 4\ matrices}\<%
\\
%
\>[2]\AgdaFunction{rel-17₄}\AgdaSpace{}%
\AgdaSymbol{:}\AgdaSpace{}%
\AgdaOperator{\AgdaPostulate{⟦}}\AgdaSpace{}%
\AgdaPostulate{l17}\AgdaSpace{}%
\AgdaOperator{\AgdaPostulate{⟧ᵐ}}\AgdaSpace{}%
\AgdaOperator{\AgdaDatatype{≡}}\AgdaSpace{}%
\AgdaOperator{\AgdaPostulate{⟦}}\AgdaSpace{}%
\AgdaPostulate{r17}\AgdaSpace{}%
\AgdaOperator{\AgdaPostulate{⟧ᵐ}}\<%
\\
%
\>[2]\AgdaFunction{rel-17₄}\AgdaSpace{}%
\AgdaSymbol{=}\AgdaSpace{}%
\AgdaInductiveConstructor{refl}\<%
\end{code}
\begin{code}[hide]%
%
\>[2]\AgdaFunction{presentation}\AgdaSpace{}%
\AgdaSymbol{=}\AgdaSpace{}%
\AgdaSymbol{\AgdaUnderscore{}}\<%
\\
%
\>[2]\AgdaFunction{relations-complete}\AgdaSpace{}%
\AgdaSymbol{=}\AgdaSpace{}%
\AgdaSymbol{\AgdaUnderscore{}}\<%
\end{code}
\begin{itemize}
  \small
  \item the normal form is \alert{exact synthesis} (Algorithm 2.10), by
        well-founded recursion on (pivot, lde, \#odd entries);
  \item soundness: commutations at any indices; the rest pulled back to
        dimension $\le4$ and decided by \aid{refl};
  \item Greylyn's $U_n(\mathbb{Z}[\frac1{\sqrt2},i])$ by the same code, generic in
        the ring: complete for $n\le4$.
\end{itemize}
\end{column}
\end{columns}
\end{frame}

\begin{frame}{The authors' own Agda, now on the library's engine \New}
\begin{columns}[c]
\begin{column}{0.53\textwidth}
\renewcommand{\AgdaCodeStyle}{\scriptsize}
\small \textbf{2-qubit Clifford+$T$} (Bian--Selinger, QPL'22), end to end:
\vspace{2pt}

\centering
\begin{tikzpicture}[node distance=3.5mm,>={Stealth[length=1.6mm]},
  bx/.style={draw=accent,fill=soft,rounded corners=2pt,font=\scriptsize,align=center,text width=40mm,minimum height=7mm}]
  \node[bx,draw=newC,fill=softnew] (ct) {Clifford+$T$: $\omega,H_i,S_i,T_i,CZ$\\{\tiny Figure 1, 33 relation constructors}};
  \node[bx,below=of ct] (gs) {Greylyn, simplified\\{\tiny 5 generators, 19 relations}};
  \node[bx,below=of gs] (g) {Greylyn: $U_4(\mathbb{Z}[\frac1{\sqrt2},i])$\\{\tiny 16 generators, 123 relations}};
  \node[bx,below=of g,draw=paperC,fill=softpaper] (m) {two-level matrices, $n\le4$\\{\tiny complete against matrices}};
  \draw[->,thick] (ct) -- node[right,font=\tiny,align=left]{full R--S: 2 cosets $I,\Omega$,\\ $8+2\cdot19=46$ equations} (gs);
  \draw[<->,thick] (gs) -- node[right,font=\tiny]{iso} (g);
  \draw[<->,thick] (g) -- node[right,font=\tiny]{\texttt{Greylyn-TwoLevel}} (m);
\end{tikzpicture}
\begin{code}[hide]%
\>[0]\AgdaKeyword{module}\AgdaSpace{}%
\AgdaModule{Clifford+T-2qubit}\AgdaSpace{}%
\AgdaKeyword{where}\<%
\end{code}
\begin{code}%
\>[0][@{}l@{\AgdaIndent{1}}]%
\>[2]\AgdaComment{--\ Clifford+T-2qubit/Presentation.agda:}\<%
\\
%
\>[2]\AgdaComment{--\ ⟦\AgdaUnderscore{}⟧\ is\ a\ group\ monomorphism\ into\ U₄(ℤ[1/√2,i])}\<%
\\
%
\>[2]\AgdaFunction{presentation}\AgdaSpace{}%
\AgdaSymbol{:}\AgdaSpace{}%
\AgdaPostulate{Rel}\AgdaSpace{}%
\AgdaOperator{\AgdaRecord{IsSubPresentationOf}}\AgdaSpace{}%
\AgdaPostulate{U}\AgdaSpace{}%
\AgdaNumber{4}\<%
\end{code}
\begin{code}[hide]%
%
\>[2]\AgdaFunction{presentation}\AgdaSpace{}%
\AgdaSymbol{=}\AgdaSpace{}%
\AgdaSymbol{\AgdaUnderscore{}}\<%
\end{code}
\end{column}
\begin{column}{0.44\textwidth}
\small \textbf{3-qubit Clifford+$CS$} (QPL'23), in progress: the
two-level $U_8(\mathbb{Z}[\frac12,i])$ side is done; Step 4 of the R--S chain
derives each of \alert{3\,192} relations from \alert{508}:
\vspace{2pt}

\centering
\begin{tikzpicture}[scale=0.95]
  \def\R{1.25}
  % 2486 / 3192 = 0.7788 -> 280.4 deg; 341 -> 38.5 deg; 346 -> 39.0 deg
  \fill[paperC!55] (0,0) -- (90:\R) arc (90:370.4:\R) -- cycle;
  \fill[accent!45] (0,0) -- (370.4:\R) arc (370.4:408.9:\R) -- cycle;
  \fill[newC!55] (0,0) -- (408.9:\R) arc (408.9:447.9:\R) -- cycle;
  \fill[gray!50] (0,0) -- (447.9:\R) arc (447.9:450:\R) -- cycle;
  \node[font=\scriptsize,anchor=west,align=left] at (1.5,0.55) {\textcolor{paperC}{\rule{6pt}{6pt}} 2\,486 by normal-form\\[-2pt]\quad computation};
  \node[font=\scriptsize,anchor=west] at (1.5,-0.15) {\textcolor{accent!70}{\rule{6pt}{6pt}} 341 axioms};
  \node[font=\scriptsize,anchor=west] at (1.5,-0.65) {\textcolor{newC}{\rule{6pt}{6pt}} 346 by rewriting};
  \node[font=\scriptsize,anchor=west] at (1.5,-1.15) {\textcolor{gray}{\rule{6pt}{6pt}} 19 by small lemmas};
\end{tikzpicture}

\vspace{4pt}
\raggedright
\small \texttt{Presentation/Tactics}: a \aid{Γ ⊢ w === v} \emph{view} of our
congruence, and the original paper's R--S theorem \alert{as an instance} of
\texttt{Normalization.Reidemeister-Schreier}. \texttt{Step4/Main}: 10 min, 2.2\,GB.
\end{column}
\end{columns}
\end{frame}

\begin{code}[hide]%
\>[0]\AgdaKeyword{module}\AgdaSpace{}%
\AgdaModule{slides.Tsinghua-PL-seminar.sections.newrcch}\AgdaSpace{}%
\AgdaKeyword{where}\<%
\\
%
\\[\AgdaEmptyExtraSkip]%
\>[0]\AgdaKeyword{open}\AgdaSpace{}%
\AgdaKeyword{import}\AgdaSpace{}%
\AgdaModule{Data.Nat}\AgdaSpace{}%
\AgdaKeyword{using}\AgdaSpace{}%
\AgdaSymbol{(}\AgdaDatatype{ℕ}\AgdaSymbol{)}\<%
\\
\>[0]\AgdaKeyword{open}\AgdaSpace{}%
\AgdaKeyword{import}\AgdaSpace{}%
\AgdaModule{Data.Product}\AgdaSpace{}%
\AgdaKeyword{using}\AgdaSpace{}%
\AgdaSymbol{(}\AgdaOperator{\AgdaFunction{\AgdaUnderscore{}×\AgdaUnderscore{}}}\AgdaSpace{}%
\AgdaSymbol{;}\AgdaSpace{}%
\AgdaField{proj₁}\AgdaSpace{}%
\AgdaSymbol{;}\AgdaSpace{}%
\AgdaField{proj₂}\AgdaSymbol{)}\<%
\\
\>[0]\AgdaKeyword{open}\AgdaSpace{}%
\AgdaKeyword{import}\AgdaSpace{}%
\AgdaModule{Word.Base}\AgdaSpace{}%
\AgdaKeyword{using}\AgdaSpace{}%
\AgdaSymbol{(}\AgdaDatatype{Word}\AgdaSpace{}%
\AgdaSymbol{;}\AgdaSpace{}%
\AgdaFunction{WRel}\AgdaSpace{}%
\AgdaSymbol{;}\AgdaSpace{}%
\AgdaInductiveConstructor{ε}\AgdaSpace{}%
\AgdaSymbol{;}\AgdaSpace{}%
\AgdaOperator{\AgdaInductiveConstructor{\AgdaUnderscore{}•\AgdaUnderscore{}}}\AgdaSymbol{)}\<%
\\
\>[0]\AgdaKeyword{open}\AgdaSpace{}%
\AgdaKeyword{import}\AgdaSpace{}%
\AgdaModule{Notations}\AgdaSpace{}%
\AgdaKeyword{using}\AgdaSpace{}%
\AgdaSymbol{(}\AgdaInductiveConstructor{₁₊}\AgdaSpace{}%
\AgdaSymbol{;}\AgdaSpace{}%
\AgdaInductiveConstructor{₂₊}\AgdaSpace{}%
\AgdaSymbol{;}\AgdaSpace{}%
\AgdaInductiveConstructor{₃₊}\AgdaSpace{}%
\AgdaSymbol{;}\AgdaSpace{}%
\AgdaInductiveConstructor{₄₊}\AgdaSymbol{)}\<%
\\
%
\\[\AgdaEmptyExtraSkip]%
\>[0]\AgdaKeyword{infix}\AgdaSpace{}%
\AgdaNumber{4}\AgdaSpace{}%
\AgdaOperator{\AgdaDatatype{\AgdaUnderscore{}SRel,\AgdaUnderscore{}===\AgdaUnderscore{}}}\AgdaSpace{}%
\AgdaOperator{\AgdaFunction{\AgdaUnderscore{}\textasciitilde{}\AgdaUnderscore{}}}\AgdaSpace{}%
\AgdaOperator{\AgdaPostulate{\AgdaUnderscore{}≐\AgdaUnderscore{}}}\AgdaSpace{}%
\AgdaOperator{\AgdaPostulate{\AgdaUnderscore{}VRel,\AgdaUnderscore{}===\AgdaUnderscore{}}}\<%
\\
\>[0]\AgdaKeyword{infixl}\AgdaSpace{}%
\AgdaNumber{10}\AgdaSpace{}%
\AgdaOperator{\AgdaPostulate{\AgdaUnderscore{}↑}}\AgdaSpace{}%
\AgdaOperator{\AgdaPostulate{\AgdaUnderscore{}↓}}\<%
\\
\>[0]\AgdaKeyword{infixr}\AgdaSpace{}%
\AgdaNumber{8}\AgdaSpace{}%
\AgdaOperator{\AgdaPostulate{\AgdaUnderscore{}\textasciicircum{}\AgdaUnderscore{}}}\<%
\\
\>[0]\AgdaKeyword{infixr}\AgdaSpace{}%
\AgdaNumber{7}\AgdaSpace{}%
\AgdaOperator{\AgdaPostulate{\AgdaUnderscore{}·\AgdaUnderscore{}}}\<%
\\
%
\\[\AgdaEmptyExtraSkip]%
\>[0]\AgdaKeyword{postulate}\<%
\\
\>[0][@{}l@{\AgdaIndent{0}}]%
\>[2]\AgdaPostulate{Gen}\AgdaSpace{}%
\AgdaSymbol{:}\AgdaSpace{}%
\AgdaDatatype{ℕ}\AgdaSpace{}%
\AgdaSymbol{→}\AgdaSpace{}%
\AgdaPrimitive{Set}\<%
\\
%
\>[2]\AgdaPostulate{Bits}\AgdaSpace{}%
\AgdaSymbol{:}\AgdaSpace{}%
\AgdaDatatype{ℕ}\AgdaSpace{}%
\AgdaSymbol{→}\AgdaSpace{}%
\AgdaPrimitive{Set}\<%
\\
%
\>[2]\AgdaPostulate{𝔽}\AgdaSpace{}%
\AgdaSymbol{:}\AgdaSpace{}%
\AgdaPrimitive{Set}\<%
\\
%
\>[2]\AgdaOperator{\AgdaPostulate{\AgdaUnderscore{}\textasciicircum{}\AgdaUnderscore{}}}\AgdaSpace{}%
\AgdaSymbol{:}\AgdaSpace{}%
\AgdaSymbol{∀}\AgdaSpace{}%
\AgdaSymbol{\{}\AgdaBound{n}\AgdaSymbol{\}}\AgdaSpace{}%
\AgdaSymbol{→}\AgdaSpace{}%
\AgdaDatatype{Word}\AgdaSpace{}%
\AgdaSymbol{(}\AgdaPostulate{Gen}\AgdaSpace{}%
\AgdaBound{n}\AgdaSymbol{)}\AgdaSpace{}%
\AgdaSymbol{→}\AgdaSpace{}%
\AgdaDatatype{ℕ}\AgdaSpace{}%
\AgdaSymbol{→}\AgdaSpace{}%
\AgdaDatatype{Word}\AgdaSpace{}%
\AgdaSymbol{(}\AgdaPostulate{Gen}\AgdaSpace{}%
\AgdaBound{n}\AgdaSymbol{)}\<%
\\
%
\>[2]\AgdaOperator{\AgdaPostulate{\AgdaUnderscore{}↑}}\AgdaSpace{}%
\AgdaOperator{\AgdaPostulate{\AgdaUnderscore{}↓}}\AgdaSpace{}%
\AgdaSymbol{:}\AgdaSpace{}%
\AgdaSymbol{∀}\AgdaSpace{}%
\AgdaSymbol{\{}\AgdaBound{n}\AgdaSymbol{\}}\AgdaSpace{}%
\AgdaSymbol{→}\AgdaSpace{}%
\AgdaDatatype{Word}\AgdaSpace{}%
\AgdaSymbol{(}\AgdaPostulate{Gen}\AgdaSpace{}%
\AgdaBound{n}\AgdaSymbol{)}\AgdaSpace{}%
\AgdaSymbol{→}\AgdaSpace{}%
\AgdaDatatype{Word}\AgdaSpace{}%
\AgdaSymbol{(}\AgdaPostulate{Gen}\AgdaSpace{}%
\AgdaBound{n}\AgdaSymbol{)}\<%
\\
%
\>[2]\AgdaPostulate{H}\AgdaSpace{}%
\AgdaPostulate{Z}\AgdaSpace{}%
\AgdaPostulate{CZ}\AgdaSpace{}%
\AgdaPostulate{CH}\AgdaSpace{}%
\AgdaPostulate{°CZ}\AgdaSpace{}%
\AgdaPostulate{°CH}\AgdaSpace{}%
\AgdaPostulate{CZ₂₀}\AgdaSpace{}%
\AgdaPostulate{CH₂₀}\AgdaSpace{}%
\AgdaPostulate{PP}\AgdaSpace{}%
\AgdaSymbol{:}\AgdaSpace{}%
\AgdaSymbol{∀}\AgdaSpace{}%
\AgdaSymbol{\{}\AgdaBound{n}\AgdaSymbol{\}}\AgdaSpace{}%
\AgdaSymbol{→}\AgdaSpace{}%
\AgdaDatatype{Word}\AgdaSpace{}%
\AgdaSymbol{(}\AgdaPostulate{Gen}\AgdaSpace{}%
\AgdaBound{n}\AgdaSymbol{)}\<%
\\
%
\>[2]\AgdaPostulate{Λ□}\AgdaSpace{}%
\AgdaSymbol{:}\AgdaSpace{}%
\AgdaDatatype{ℕ}\AgdaSpace{}%
\AgdaSymbol{→}\AgdaSpace{}%
\AgdaSymbol{∀}\AgdaSpace{}%
\AgdaSymbol{\{}\AgdaBound{n}\AgdaSymbol{\}}\AgdaSpace{}%
\AgdaSymbol{→}\AgdaSpace{}%
\AgdaDatatype{Word}\AgdaSpace{}%
\AgdaSymbol{(}\AgdaPostulate{Gen}\AgdaSpace{}%
\AgdaBound{n}\AgdaSymbol{)}\<%
\\
%
\>[2]\AgdaOperator{\AgdaPostulate{√2\textasciicircum{}\AgdaUnderscore{}}}\AgdaSpace{}%
\AgdaSymbol{:}\AgdaSpace{}%
\AgdaDatatype{ℕ}\AgdaSpace{}%
\AgdaSymbol{→}\AgdaSpace{}%
\AgdaPostulate{𝔽}\<%
\\
%
\>[2]\AgdaOperator{\AgdaPostulate{\AgdaUnderscore{}·\AgdaUnderscore{}}}\AgdaSpace{}%
\AgdaSymbol{:}\AgdaSpace{}%
\AgdaPostulate{𝔽}\AgdaSpace{}%
\AgdaSymbol{→}\AgdaSpace{}%
\AgdaSymbol{∀}\AgdaSpace{}%
\AgdaSymbol{\{}\AgdaBound{n}\AgdaSymbol{\}}\AgdaSpace{}%
\AgdaSymbol{→}\AgdaSpace{}%
\AgdaSymbol{(}\AgdaPostulate{Bits}\AgdaSpace{}%
\AgdaBound{n}\AgdaSpace{}%
\AgdaSymbol{→}\AgdaSpace{}%
\AgdaPostulate{Bits}\AgdaSpace{}%
\AgdaBound{n}\AgdaSpace{}%
\AgdaSymbol{→}\AgdaSpace{}%
\AgdaPostulate{𝔽}\AgdaSymbol{)}\AgdaSpace{}%
\AgdaSymbol{→}\AgdaSpace{}%
\AgdaSymbol{(}\AgdaPostulate{Bits}\AgdaSpace{}%
\AgdaBound{n}\AgdaSpace{}%
\AgdaSymbol{→}\AgdaSpace{}%
\AgdaPostulate{Bits}\AgdaSpace{}%
\AgdaBound{n}\AgdaSpace{}%
\AgdaSymbol{→}\AgdaSpace{}%
\AgdaPostulate{𝔽}\AgdaSymbol{)}\<%
\\
%
\>[2]\AgdaOperator{\AgdaPostulate{\AgdaUnderscore{}≐\AgdaUnderscore{}}}\AgdaSpace{}%
\AgdaSymbol{:}\AgdaSpace{}%
\AgdaSymbol{∀}\AgdaSpace{}%
\AgdaSymbol{\{}\AgdaBound{n}\AgdaSymbol{\}}\AgdaSpace{}%
\AgdaSymbol{→}\AgdaSpace{}%
\AgdaSymbol{(}\AgdaPostulate{Bits}\AgdaSpace{}%
\AgdaBound{n}\AgdaSpace{}%
\AgdaSymbol{→}\AgdaSpace{}%
\AgdaPostulate{Bits}\AgdaSpace{}%
\AgdaBound{n}\AgdaSpace{}%
\AgdaSymbol{→}\AgdaSpace{}%
\AgdaPostulate{𝔽}\AgdaSymbol{)}\AgdaSpace{}%
\AgdaSymbol{→}\AgdaSpace{}%
\AgdaSymbol{(}\AgdaPostulate{Bits}\AgdaSpace{}%
\AgdaBound{n}\AgdaSpace{}%
\AgdaSymbol{→}\AgdaSpace{}%
\AgdaPostulate{Bits}\AgdaSpace{}%
\AgdaBound{n}\AgdaSpace{}%
\AgdaSymbol{→}\AgdaSpace{}%
\AgdaPostulate{𝔽}\AgdaSymbol{)}\AgdaSpace{}%
\AgdaSymbol{→}\AgdaSpace{}%
\AgdaPrimitive{Set}\<%
\\
%
\>[2]\AgdaOperator{\AgdaPostulate{\AgdaUnderscore{}VRel,\AgdaUnderscore{}===\AgdaUnderscore{}}}\AgdaSpace{}%
\AgdaSymbol{:}\AgdaSpace{}%
\AgdaSymbol{(}\AgdaBound{n}\AgdaSpace{}%
\AgdaSymbol{:}\AgdaSpace{}%
\AgdaDatatype{ℕ}\AgdaSymbol{)}\AgdaSpace{}%
\AgdaSymbol{→}\AgdaSpace{}%
\AgdaFunction{WRel}\AgdaSpace{}%
\AgdaSymbol{(}\AgdaPostulate{Gen}\AgdaSpace{}%
\AgdaBound{n}\AgdaSymbol{)}\<%
\\
%
\>[2]\AgdaPostulate{proof-elided}\AgdaSpace{}%
\AgdaSymbol{:}\AgdaSpace{}%
\AgdaSymbol{∀}\AgdaSpace{}%
\AgdaSymbol{\{}\AgdaBound{A}\AgdaSpace{}%
\AgdaSymbol{:}\AgdaSpace{}%
\AgdaPrimitive{Set}\AgdaSymbol{\}}\AgdaSpace{}%
\AgdaSymbol{→}\AgdaSpace{}%
\AgdaBound{A}\<%
\\
%
\\[\AgdaEmptyExtraSkip]%
\>[0]\AgdaFunction{Circuit}\AgdaSpace{}%
\AgdaSymbol{:}\AgdaSpace{}%
\AgdaDatatype{ℕ}\AgdaSpace{}%
\AgdaSymbol{→}\AgdaSpace{}%
\AgdaPrimitive{Set}\<%
\\
\>[0]\AgdaFunction{Circuit}\AgdaSpace{}%
\AgdaBound{n}\AgdaSpace{}%
\AgdaSymbol{=}\AgdaSpace{}%
\AgdaDatatype{Word}\AgdaSpace{}%
\AgdaSymbol{(}\AgdaPostulate{Gen}\AgdaSpace{}%
\AgdaBound{n}\AgdaSymbol{)}\<%
\\
%
\\[\AgdaEmptyExtraSkip]%
\>[0]\AgdaKeyword{private}\AgdaSpace{}%
\AgdaKeyword{variable}\<%
\\
\>[0][@{}l@{\AgdaIndent{0}}]%
\>[2]\AgdaGeneralizable{n}\AgdaSpace{}%
\AgdaSymbol{:}\AgdaSpace{}%
\AgdaDatatype{ℕ}\<%
\end{code}
%% Part 6d — Real-Clifford+CH (branch v-office-Cli-CH).

\begin{frame}{Real-Clifford+$CH$: Clément's complete equational theory \New}
\begin{columns}[c]
\begin{column}{0.53\textwidth}
\renewcommand{\AgdaCodeStyle}{\scriptsize}
\small A.\ Clément, arXiv:2602.06644: a 13-page paper, \alert{264 pages} of appendix.
\begin{code}%
\>[0]\AgdaKeyword{data}\AgdaSpace{}%
\AgdaDatatype{Gate}\AgdaSpace{}%
\AgdaSymbol{:}\AgdaSpace{}%
\AgdaDatatype{ℕ}\AgdaSpace{}%
\AgdaSymbol{→}\AgdaSpace{}%
\AgdaPrimitive{Set}\AgdaSpace{}%
\AgdaKeyword{where}\<%
\\
\>[0][@{}l@{\AgdaIndent{0}}]%
\>[2]\AgdaInductiveConstructor{H-gate}%
\>[10]\AgdaSymbol{:}\AgdaSpace{}%
\AgdaDatatype{Gate}\AgdaSpace{}%
\AgdaNumber{1}\<%
\\
%
\>[2]\AgdaInductiveConstructor{Z-gate}%
\>[10]\AgdaSymbol{:}\AgdaSpace{}%
\AgdaDatatype{Gate}\AgdaSpace{}%
\AgdaNumber{1}\<%
\\
%
\>[2]\AgdaInductiveConstructor{CZ-gate}\AgdaSpace{}%
\AgdaSymbol{:}\AgdaSpace{}%
\AgdaDatatype{Gate}\AgdaSpace{}%
\AgdaNumber{2}\<%
\\
%
\>[2]\AgdaInductiveConstructor{CH-gate}\AgdaSpace{}%
\AgdaSymbol{:}\AgdaSpace{}%
\AgdaDatatype{Gate}\AgdaSpace{}%
\AgdaNumber{2}\<%
\\
%
\\[\AgdaEmptyExtraSkip]%
\>[0]\AgdaFunction{Ex}\AgdaSpace{}%
\AgdaSymbol{:}\AgdaSpace{}%
\AgdaFunction{Circuit}\AgdaSpace{}%
\AgdaSymbol{(}\AgdaInductiveConstructor{₂₊}\AgdaSpace{}%
\AgdaGeneralizable{n}\AgdaSymbol{)}%
\>[22]\AgdaComment{--\ Equation\ (8):\ the\ swap\ is\ a\ definition}\<%
\\
\>[0]\AgdaFunction{Ex}\AgdaSpace{}%
\AgdaSymbol{=}\AgdaSpace{}%
\AgdaPostulate{CZ}\AgdaSpace{}%
\AgdaOperator{\AgdaInductiveConstructor{•}}\AgdaSpace{}%
\AgdaPostulate{H}\AgdaSpace{}%
\AgdaOperator{\AgdaPostulate{↓}}\AgdaSpace{}%
\AgdaOperator{\AgdaInductiveConstructor{•}}\AgdaSpace{}%
\AgdaPostulate{H}\AgdaSpace{}%
\AgdaOperator{\AgdaPostulate{↑}}\AgdaSpace{}%
\AgdaOperator{\AgdaInductiveConstructor{•}}\AgdaSpace{}%
\AgdaPostulate{CZ}\AgdaSpace{}%
\AgdaOperator{\AgdaInductiveConstructor{•}}\AgdaSpace{}%
\AgdaPostulate{H}\AgdaSpace{}%
\AgdaOperator{\AgdaPostulate{↓}}\AgdaSpace{}%
\AgdaOperator{\AgdaInductiveConstructor{•}}\AgdaSpace{}%
\AgdaPostulate{H}\AgdaSpace{}%
\AgdaOperator{\AgdaPostulate{↑}}\AgdaSpace{}%
\AgdaOperator{\AgdaInductiveConstructor{•}}\AgdaSpace{}%
\AgdaPostulate{CZ}\AgdaSpace{}%
\AgdaOperator{\AgdaInductiveConstructor{•}}\AgdaSpace{}%
\AgdaPostulate{H}\AgdaSpace{}%
\AgdaOperator{\AgdaPostulate{↓}}\AgdaSpace{}%
\AgdaOperator{\AgdaInductiveConstructor{•}}\AgdaSpace{}%
\AgdaPostulate{H}\AgdaSpace{}%
\AgdaOperator{\AgdaPostulate{↑}}\<%
\end{code}
\begin{code}[hide]%
\>[0]\<%
\end{code}
\vspace{-6pt}
{\scriptsize 23 rules: (1)--(19) of Figure 4 minus (8), and the 5 swap rules of Remark 1;
a sample:}
\begin{code}%
\>[0]\AgdaKeyword{data}\AgdaSpace{}%
\AgdaOperator{\AgdaDatatype{\AgdaUnderscore{}SRel,\AgdaUnderscore{}===\AgdaUnderscore{}}}\AgdaSpace{}%
\AgdaSymbol{:}\AgdaSpace{}%
\AgdaSymbol{(}\AgdaBound{n}\AgdaSpace{}%
\AgdaSymbol{:}\AgdaSpace{}%
\AgdaDatatype{ℕ}\AgdaSymbol{)}\AgdaSpace{}%
\AgdaSymbol{→}\AgdaSpace{}%
\AgdaFunction{WRel}\AgdaSpace{}%
\AgdaSymbol{(}\AgdaPostulate{Gen}\AgdaSpace{}%
\AgdaBound{n}\AgdaSymbol{)}\AgdaSpace{}%
\AgdaKeyword{where}\<%
\\
\>[0][@{}l@{\AgdaIndent{0}}]%
\>[2]\AgdaInductiveConstructor{order-HZ}\AgdaSpace{}%
\AgdaSymbol{:}\AgdaSpace{}%
\AgdaSymbol{(}\AgdaInductiveConstructor{₁₊}\AgdaSpace{}%
\AgdaGeneralizable{n}\AgdaSymbol{)}\AgdaSpace{}%
\AgdaOperator{\AgdaDatatype{SRel,}}\AgdaSpace{}%
\AgdaSymbol{(}\AgdaPostulate{H}\AgdaSpace{}%
\AgdaOperator{\AgdaInductiveConstructor{•}}\AgdaSpace{}%
\AgdaPostulate{Z}\AgdaSymbol{)}\AgdaSpace{}%
\AgdaOperator{\AgdaPostulate{\textasciicircum{}}}\AgdaSpace{}%
\AgdaNumber{8}\AgdaSpace{}%
\AgdaOperator{\AgdaDatatype{===}}\AgdaSpace{}%
\AgdaInductiveConstructor{ε}\<%
\\
%
\>[2]\AgdaInductiveConstructor{merge-CZ}%
\>[14]\AgdaSymbol{:}\AgdaSpace{}%
\AgdaSymbol{(}\AgdaInductiveConstructor{₂₊}\AgdaSpace{}%
\AgdaGeneralizable{n}\AgdaSymbol{)}\AgdaSpace{}%
\AgdaOperator{\AgdaDatatype{SRel,}}\AgdaSpace{}%
\AgdaPostulate{CZ}\AgdaSpace{}%
\AgdaOperator{\AgdaInductiveConstructor{•}}\AgdaSpace{}%
\AgdaPostulate{°CZ}\AgdaSpace{}%
\AgdaOperator{\AgdaDatatype{===}}\AgdaSpace{}%
\AgdaPostulate{Z}\AgdaSpace{}%
\AgdaOperator{\AgdaPostulate{↓}}\<%
\\
%
\>[2]\AgdaInductiveConstructor{slide-CH}%
\>[14]\AgdaSymbol{:}\AgdaSpace{}%
\AgdaSymbol{(}\AgdaInductiveConstructor{₂₊}\AgdaSpace{}%
\AgdaGeneralizable{n}\AgdaSymbol{)}\AgdaSpace{}%
\AgdaOperator{\AgdaDatatype{SRel,}}%
\>[255I]\AgdaPostulate{CH}\AgdaSpace{}%
\AgdaOperator{\AgdaInductiveConstructor{•}}\AgdaSpace{}%
\AgdaPostulate{H}\AgdaSpace{}%
\AgdaOperator{\AgdaPostulate{↑}}\AgdaSpace{}%
\AgdaOperator{\AgdaInductiveConstructor{•}}\AgdaSpace{}%
\AgdaPostulate{Z}\AgdaSpace{}%
\AgdaOperator{\AgdaPostulate{↑}}\AgdaSpace{}%
\AgdaOperator{\AgdaInductiveConstructor{•}}\AgdaSpace{}%
\AgdaPostulate{CH}\AgdaSpace{}%
\AgdaOperator{\AgdaInductiveConstructor{•}}\AgdaSpace{}%
\AgdaPostulate{H}\AgdaSpace{}%
\AgdaOperator{\AgdaPostulate{↑}}\<%
\\
\>[.][@{}l@{}]\<[255I]%
\>[29]\AgdaOperator{\AgdaDatatype{===}}\AgdaSpace{}%
\AgdaPostulate{H}\AgdaSpace{}%
\AgdaOperator{\AgdaPostulate{↑}}\AgdaSpace{}%
\AgdaOperator{\AgdaInductiveConstructor{•}}\AgdaSpace{}%
\AgdaPostulate{Z}\AgdaSpace{}%
\AgdaOperator{\AgdaPostulate{↑}}\AgdaSpace{}%
\AgdaOperator{\AgdaInductiveConstructor{•}}\AgdaSpace{}%
\AgdaPostulate{CH}\AgdaSpace{}%
\AgdaOperator{\AgdaInductiveConstructor{•}}\AgdaSpace{}%
\AgdaPostulate{H}\AgdaSpace{}%
\AgdaOperator{\AgdaPostulate{↑}}\AgdaSpace{}%
\AgdaOperator{\AgdaInductiveConstructor{•}}\AgdaSpace{}%
\AgdaPostulate{CH}\<%
\\
%
\>[2]\AgdaInductiveConstructor{square-CCZX}%
\>[16]\AgdaSymbol{:}\AgdaSpace{}%
\AgdaSymbol{(}\AgdaInductiveConstructor{₃₊}\AgdaSpace{}%
\AgdaGeneralizable{n}\AgdaSymbol{)}\AgdaSpace{}%
\AgdaOperator{\AgdaDatatype{SRel,}}\AgdaSpace{}%
\AgdaSymbol{(}\AgdaPostulate{CH}\AgdaSpace{}%
\AgdaOperator{\AgdaPostulate{↓}}\AgdaSpace{}%
\AgdaOperator{\AgdaInductiveConstructor{•}}\AgdaSpace{}%
\AgdaPostulate{CZ₂₀}\AgdaSymbol{)}\AgdaSpace{}%
\AgdaOperator{\AgdaPostulate{\textasciicircum{}}}\AgdaSpace{}%
\AgdaNumber{4}\AgdaSpace{}%
\AgdaOperator{\AgdaDatatype{===}}\AgdaSpace{}%
\AgdaPostulate{CZ}\AgdaSpace{}%
\AgdaOperator{\AgdaPostulate{↑}}\<%
\\
%
\>[2]\AgdaInductiveConstructor{conj-PP}%
\>[16]\AgdaSymbol{:}\AgdaSpace{}%
\AgdaSymbol{(}\AgdaInductiveConstructor{₃₊}\AgdaSpace{}%
\AgdaGeneralizable{n}\AgdaSymbol{)}\AgdaSpace{}%
\AgdaOperator{\AgdaDatatype{SRel,}}\AgdaSpace{}%
\AgdaPostulate{CZ}\AgdaSpace{}%
\AgdaOperator{\AgdaPostulate{↑}}\AgdaSpace{}%
\AgdaOperator{\AgdaInductiveConstructor{•}}\AgdaSpace{}%
\AgdaPostulate{PP}\AgdaSpace{}%
\AgdaOperator{\AgdaPostulate{↓}}\AgdaSpace{}%
\AgdaOperator{\AgdaDatatype{===}}\AgdaSpace{}%
\AgdaPostulate{PP}\AgdaSpace{}%
\AgdaOperator{\AgdaPostulate{↓}}\AgdaSpace{}%
\AgdaOperator{\AgdaInductiveConstructor{•}}\AgdaSpace{}%
\AgdaPostulate{CH}\AgdaSpace{}%
\AgdaOperator{\AgdaPostulate{↑}}\<%
\\
%
\>[2]\AgdaInductiveConstructor{box-Z}\AgdaSpace{}%
\AgdaSymbol{:}\AgdaSpace{}%
\AgdaSymbol{∀}\AgdaSpace{}%
\AgdaBound{k}\AgdaSpace{}%
\AgdaSymbol{→}\AgdaSpace{}%
\AgdaSymbol{(}\AgdaInductiveConstructor{₁₊}\AgdaSpace{}%
\AgdaSymbol{(}\AgdaInductiveConstructor{₄₊}\AgdaSpace{}%
\AgdaBound{k}\AgdaSymbol{))}\AgdaSpace{}%
\AgdaOperator{\AgdaDatatype{SRel,}}\AgdaSpace{}%
\AgdaPostulate{Z}\AgdaSpace{}%
\AgdaOperator{\AgdaPostulate{↓}}\AgdaSpace{}%
\AgdaOperator{\AgdaInductiveConstructor{•}}\AgdaSpace{}%
\AgdaPostulate{Λ□}\AgdaSpace{}%
\AgdaSymbol{(}\AgdaInductiveConstructor{₄₊}\AgdaSpace{}%
\AgdaBound{k}\AgdaSymbol{)}\AgdaSpace{}%
\AgdaOperator{\AgdaDatatype{===}}\AgdaSpace{}%
\AgdaPostulate{Λ□}\AgdaSpace{}%
\AgdaSymbol{(}\AgdaInductiveConstructor{₄₊}\AgdaSpace{}%
\AgdaBound{k}\AgdaSymbol{)}\AgdaSpace{}%
\AgdaOperator{\AgdaInductiveConstructor{•}}\AgdaSpace{}%
\AgdaPostulate{Z}\AgdaSpace{}%
\AgdaOperator{\AgdaPostulate{↓}}\<%
\end{code}
\end{column}
\begin{column}{0.44\textwidth}
\centering
{\small the swap, by definition (8):}\\[2pt]
\begin{quantikz}[row sep=0.25cm,column sep=0.14cm,font=\scriptsize]
 & \swap{1} & \\
 & \targX{} &
\end{quantikz}
$=$
\begin{quantikz}[row sep=0.25cm,column sep=0.14cm,font=\scriptsize]
 & \ctrl{1} & \gate{H} & \ctrl{1} & \gate{H} & \ctrl{1} & \gate{H} & \\
 & \control{} & \gate{H} & \control{} & \gate{H} & \control{} & \gate{H} &
\end{quantikz}

\vspace{8pt}
{\small rule (11) \aid{slide-CH}:}\\[2pt]
\begin{quantikz}[row sep=0.25cm,column sep=0.12cm,font=\scriptsize]
 & \ctrl{1} & \gate{H} & \gate{Z} & \ctrl{1} & \gate{H} & \\
 & \gate{H} & & & \gate{H} & &
\end{quantikz}
$=$
\begin{quantikz}[row sep=0.25cm,column sep=0.12cm,font=\scriptsize]
 & \gate{H} & \gate{Z} & \ctrl{1} & \gate{H} & \ctrl{1} & \\
 & & & \gate{H} & & \gate{H} &
\end{quantikz}

\vspace{6pt}
{\scriptsize wire 0 is the bottom wire; $CH$ is controlled by wire 1.\\
(19) is a \alert{schema}: one rule per width $\ge5$, about the multi-controlled
box $\Lambda\square$ of Definition 2.4, defined by recursion.}
\end{column}
\end{columns}
\end{frame}

\begin{frame}{A matrix semantics over $\mathbb{Z}[\sqrt2]$ \New}
\begin{columns}[c]
\begin{column}{0.52\textwidth}
\renewcommand{\AgdaCodeStyle}{\scriptsize}
\small Gate matrices are taken $\times\sqrt2$, so entries lie in
$\mathbb{Z}[\sqrt2]$; the power of $\frac1{\sqrt2}$ is kept apart:
\begin{code}%
\>[0]\AgdaFunction{Op}\AgdaSpace{}%
\AgdaSymbol{:}\AgdaSpace{}%
\AgdaDatatype{ℕ}\AgdaSpace{}%
\AgdaSymbol{→}\AgdaSpace{}%
\AgdaPrimitive{Set}\<%
\\
\>[0]\AgdaFunction{Op}\AgdaSpace{}%
\AgdaBound{n}\AgdaSpace{}%
\AgdaSymbol{=}\AgdaSpace{}%
\AgdaPostulate{Bits}\AgdaSpace{}%
\AgdaBound{n}\AgdaSpace{}%
\AgdaSymbol{→}\AgdaSpace{}%
\AgdaPostulate{Bits}\AgdaSpace{}%
\AgdaBound{n}\AgdaSpace{}%
\AgdaSymbol{→}\AgdaSpace{}%
\AgdaPostulate{𝔽}\<%
\\
%
\\[\AgdaEmptyExtraSkip]%
\>[0]\AgdaFunction{Scaled}\AgdaSpace{}%
\AgdaSymbol{:}\AgdaSpace{}%
\AgdaDatatype{ℕ}\AgdaSpace{}%
\AgdaSymbol{→}\AgdaSpace{}%
\AgdaPrimitive{Set}\<%
\\
\>[0]\AgdaFunction{Scaled}\AgdaSpace{}%
\AgdaBound{n}\AgdaSpace{}%
\AgdaSymbol{=}\AgdaSpace{}%
\AgdaDatatype{ℕ}\AgdaSpace{}%
\AgdaOperator{\AgdaFunction{×}}\AgdaSpace{}%
\AgdaFunction{Op}\AgdaSpace{}%
\AgdaBound{n}%
\>[29]\AgdaComment{--\ (ℓ\ ,\ M)\ \ reads\ \ (1/√2)\textasciicircum{}ℓ\ ·\ M}\<%
\\
%
\\[\AgdaEmptyExtraSkip]%
\>[0]\AgdaOperator{\AgdaFunction{\AgdaUnderscore{}\textasciitilde{}\AgdaUnderscore{}}}\AgdaSpace{}%
\AgdaSymbol{:}\AgdaSpace{}%
\AgdaFunction{Scaled}\AgdaSpace{}%
\AgdaGeneralizable{n}\AgdaSpace{}%
\AgdaSymbol{→}\AgdaSpace{}%
\AgdaFunction{Scaled}\AgdaSpace{}%
\AgdaGeneralizable{n}\AgdaSpace{}%
\AgdaSymbol{→}\AgdaSpace{}%
\AgdaPrimitive{Set}\<%
\\
\>[0]\AgdaBound{s}\AgdaSpace{}%
\AgdaOperator{\AgdaFunction{\textasciitilde{}}}\AgdaSpace{}%
\AgdaBound{t}\AgdaSpace{}%
\AgdaSymbol{=}\AgdaSpace{}%
\AgdaSymbol{((}\AgdaOperator{\AgdaPostulate{√2\textasciicircum{}}}\AgdaSpace{}%
\AgdaField{proj₁}\AgdaSpace{}%
\AgdaBound{t}\AgdaSymbol{)}\AgdaSpace{}%
\AgdaOperator{\AgdaPostulate{·}}\AgdaSpace{}%
\AgdaField{proj₂}\AgdaSpace{}%
\AgdaBound{s}\AgdaSymbol{)}\AgdaSpace{}%
\AgdaOperator{\AgdaPostulate{≐}}\AgdaSpace{}%
\AgdaSymbol{((}\AgdaOperator{\AgdaPostulate{√2\textasciicircum{}}}\AgdaSpace{}%
\AgdaField{proj₁}\AgdaSpace{}%
\AgdaBound{s}\AgdaSymbol{)}\AgdaSpace{}%
\AgdaOperator{\AgdaPostulate{·}}\AgdaSpace{}%
\AgdaField{proj₂}\AgdaSpace{}%
\AgdaBound{t}\AgdaSymbol{)}\<%
\end{code}
\begin{code}[hide]%
\>[0]\AgdaKeyword{postulate}\AgdaSpace{}%
\AgdaOperator{\AgdaPostulate{⟦\AgdaUnderscore{}⟧}}\AgdaSpace{}%
\AgdaSymbol{:}\AgdaSpace{}%
\AgdaSymbol{∀}\AgdaSpace{}%
\AgdaSymbol{\{}\AgdaBound{n}\AgdaSymbol{\}}\AgdaSpace{}%
\AgdaSymbol{→}\AgdaSpace{}%
\AgdaFunction{Circuit}\AgdaSpace{}%
\AgdaBound{n}\AgdaSpace{}%
\AgdaSymbol{→}\AgdaSpace{}%
\AgdaFunction{Scaled}\AgdaSpace{}%
\AgdaBound{n}\<%
\end{code}
\begin{itemize}
  \item matrices are computed on stored \alert{tries} \aid{Mat k};
  \item \textbf{soundness} (Prop.\ 3.1): each relator checked \emph{once}, by
        \aid{refl} on tries at the width it is drawn on; \aid{localise} /
        \aid{by-relator} lift it to \emph{every} width;
  \item the schema (19) at every width by a symbolic lemma: the box is
        $(\frac1{\sqrt2})^\ell$ times a diagonal sign.
\end{itemize}
\end{column}
\begin{column}{0.45\textwidth}
\centering
\begin{tikzpicture}[x=5.5mm,y=5.5mm,font=\scriptsize]
  \node at (-1.6,2) {$\tfrac{1}{\sqrt2}\;\cdot$};
  \draw[thick] (-0.3,-0.2) -- (-0.4,-0.2) -- (-0.4,4.2) -- (-0.3,4.2);
  \draw[thick] (4.3,-0.2) -- (4.4,-0.2) -- (4.4,4.2) -- (4.3,4.2);
  \foreach \r/\row in {0/{1,1,0,0},1/{1,-1,0,0},2/{0,0,1,1},3/{0,0,1,-1}} {
    \foreach \e [count=\c] in \row {
      \node at (\c-0.5,3.5-\r) {$\e$};
    }
  }
  \node[font=\scriptsize,align=center] at (2,-1) {$I\otimes H$ as $(1,M)$, $M$ over $\mathbb{Z}[\sqrt2]$};
\end{tikzpicture}

\vspace{8pt}
\begin{code}%
\>[0]\AgdaFunction{axiom-sound}\AgdaSpace{}%
\AgdaSymbol{:}%
\>[384I]\AgdaSymbol{\{}\AgdaBound{w}\AgdaSpace{}%
\AgdaBound{v}\AgdaSpace{}%
\AgdaSymbol{:}\AgdaSpace{}%
\AgdaFunction{Circuit}\AgdaSpace{}%
\AgdaGeneralizable{n}\AgdaSymbol{\}}\AgdaSpace{}%
\AgdaSymbol{→}\<%
\\
\>[.][@{}l@{}]\<[384I]%
\>[14]\AgdaGeneralizable{n}\AgdaSpace{}%
\AgdaOperator{\AgdaPostulate{VRel,}}\AgdaSpace{}%
\AgdaBound{w}\AgdaSpace{}%
\AgdaOperator{\AgdaPostulate{===}}\AgdaSpace{}%
\AgdaBound{v}\AgdaSpace{}%
\AgdaSymbol{→}\AgdaSpace{}%
\AgdaOperator{\AgdaPostulate{⟦}}\AgdaSpace{}%
\AgdaBound{w}\AgdaSpace{}%
\AgdaOperator{\AgdaPostulate{⟧}}\AgdaSpace{}%
\AgdaOperator{\AgdaFunction{\textasciitilde{}}}\AgdaSpace{}%
\AgdaOperator{\AgdaPostulate{⟦}}\AgdaSpace{}%
\AgdaBound{v}\AgdaSpace{}%
\AgdaOperator{\AgdaPostulate{⟧}}\<%
\end{code}
\begin{code}[hide]%
\>[0]\AgdaFunction{axiom-sound}\AgdaSpace{}%
\AgdaSymbol{=}\AgdaSpace{}%
\AgdaPostulate{proof-elided}\<%
\end{code}
{\scriptsize \texttt{Real-Clifford+CH/Semantics/Scaled.agda}, \texttt{Soundness.agda}}
\end{column}
\end{columns}
\end{frame}

\begin{frame}{Completeness: the paper's back-and-forth argument \New}
\begin{columns}[c]
\begin{column}{0.53\textwidth}
\centering
\begin{tikzpicture}[>={Stealth[length=1.8mm]}]
  \node[draw=newC,fill=softnew,rounded corners=2pt,font=\scriptsize,align=center,minimum height=10mm,text width=24mm] (C) at (0,0) {$n$-qubit circuits\\Figure 4};
  \node[draw=accent,fill=soft,rounded corners=2pt,font=\scriptsize,align=center,minimum height=10mm,text width=24mm] (P) at (4,0) {words over $P$\\Figure 8 (27 rules)};
  \node[draw=paperC,fill=softpaper,rounded corners=2pt,font=\scriptsize,align=center,minimum height=10mm,text width=24mm] (O) at (4,-2.6) {$O_{2^n}(\mathbb{Z}[\frac1{\sqrt2}])$\\Figure 7 (13 rules)};
  \draw[->,thick] (C.10) -- node[above,font=\scriptsize]{encode $e$} (P.170);
  \draw[->,thick] (P.190) -- node[below,font=\scriptsize]{decode $d$} (C.350);
  \draw[->,thick] (P) -- node[right,font=\tiny,align=left]{Thm 4.10\\(R--S, 4 cosets)} (O);
  \node[font=\tiny,align=center,paperC] at (4,-3.6) {Thm 4.4: complete --- imported\\by the paper (Fang--Heunen--Kaarsgaard)};
  \node[font=\tiny,align=left,anchor=north] at (0,-0.9) {E-sem: $e$ preserves semantics\\Lemma 8.7: $d(e\,g)\approx g$\\Lemma 8.8: $d$ respects Figure 8\\Thm 4.10: Figure 8 is complete};
\end{tikzpicture}
\end{column}
\begin{column}{0.44\textwidth}
\small
\renewcommand{\arraystretch}{1.2}
\begin{tabular}{@{}lp{4.3cm}@{}}
\toprule
width & status (\Branch{v-office-Cli-CH}, Sep 29)\\
\midrule
0, 1 & \textbf{complete} outright (1 qubit: 16-state automaton NF)\\
2 & from Theorem 4.4 at $N=4$; Lemmas 7.4, 7.5 proved\\
$\ge3$ & \aid{BackAndForth}: Thm 4.10 \textbf{proved} (from 4.4); $e$'s semantics
         and Lemma 8.7 \textbf{proved} at every width; Lemma 8.8 proved at every
         width $\ge5$; \alert{open at 3, 4 qubits}\\
\bottomrule
\end{tabular}

\vspace{6pt}
{\scriptsize The hypothesis list shrank over one week:
\aid{Theorem-8-9} (4.10, 8.7, 8.8) $\to$ \aid{′} (4.4, 8.7, 8.8) $\to$ \aid{″}
$\to$ \aid{‴} (4.4, and 8.8 at 3 and 4 qubits). No postulates; \texttt{--safe}.}
\end{column}
\end{columns}
\end{frame}

\begin{frame}{Appendix D, the 264 pages: where we are \New}
\begin{columns}[c]
\begin{column}{0.52\textwidth}
\centering
\begin{tikzpicture}
\begin{axis}[xbar stacked, width=7.4cm, height=5.6cm, xmin=0, xmax=105,
  symbolic y coords={D8,D7,D5,D2,D1}, yticklabels={{D.8 (269--307)},{D.7 (249--268)},{D.5 (150--248)},{D.2 (117--147)},{D.1}},
  ytick=data, y tick label style={font=\scriptsize},
  x tick label style={font=\scriptsize}, xlabel={\scriptsize equations},
  bar width=9pt, enlarge y limits=0.15, axis lines*=left,
  legend style={font=\tiny,at={(0.98,0.97)},anchor=north east}]
\addplot[fill=paperC!60,draw=paperC] coordinates
  {(6,D1) (31,D2) (99,D5) (20,D7) (37,D8)};
\addplot[fill=gray!25,draw=gray] coordinates
  {(0,D1) (0,D2) (0,D5) (0,D7) (2,D8)};
\legend{stated and proved in Agda, not yet named}
\end{axis}
\end{tikzpicture}

{\scriptsize D.8 counts named statements, proved inside the width induction from
completeness one width down.}
\end{column}
\begin{column}{0.45\textwidth}
\small
\begin{itemize}
  \setlength{\itemsep}{3pt}
  \item 294 files, \textbf{77\,193 lines}, 48 commits in ten days (Sep 19--29);
  \item an $n$-qubit derivation uses completeness \alert{one width down} as a
        black box --- the paper's ``Lemma 5.1'' is our induction hypothesis;
  \item a 3-wire evaluation \aid{Same u v} becomes a derivation
        (\aid{SemanticSteps}), $\approx$1\,s per gate;
  \item structure instead of pages: two lemmas ``replace some sixty diagram
        steps''; (203) ``replaces the paper's three pages'';
  \item Theorem 4.10 is the library's \alert{generic} R--S engine: Figure 8 over
        Figure 7, 4 cosets, $13\times4$ obligations, uniformly via one embedding.
\end{itemize}
\end{column}
\end{columns}
\end{frame}

\begin{frame}{What the formalization found \New}
\begin{columns}[T]
\begin{column}{0.50\textwidth}
\begin{keybox}[title=In the paper]
\footnotesize
\begin{itemize}
  \setlength{\itemsep}{2pt}
  \item Appendix A.4.1 prints $\mathcal{E}^{1-x}$; a decision procedure
        (\texttt{Checks8}) shows (36)/(37) hold only the other way round, and
        \emph{fail at every instance} as printed.
  \item \texttt{Checks8} found the \alert{missing distinctness side conditions}
        of (30), (34), (42), and the negation parity of $W_1$/$W_2$.
  \item The figure in the proof of (216)--(218) does not parse: as drawn it is
        \alert{not semantically valid}; replaced by a structural argument.
  \item Footnote 13's ``constructed in a similar way'' needed a generic
        construction (six-index words).
\end{itemize}
\end{keybox}
\end{column}
\begin{column}{0.47\textwidth}
\begin{newdevbox}[title=In our own transcription]
\footnotesize
\begin{itemize}
  \setlength{\itemsep}{2pt}
  \item a swap network iterated the wrong way: invisible on 3 qubits, so (36)/(37)
        as transcribed were \alert{unsound from 4 qubits on}; caught only when
        (65) was proved at symbolic width;
  \item (261) misread (a control dot) --- the misread, harder statement was
        proved too; lesson: read the dots at 400\,dpi;
  \item a draft typechecked while silently wrong: \aid{zero} not imported, so
        it became a pattern variable --- caught only by warnings.
\end{itemize}
\end{newdevbox}
\vspace{4pt}
\begin{keybox}
\small Discipline: every claim tested numerically first (a pure-Python
simulator), then decided by evaluation, then proved.
\end{keybox}
\end{column}
\end{columns}
\end{frame}

\begin{frame}{Making evaluation affordable \New}
\begin{columns}[c]
\begin{column}{0.50\textwidth}
\small Operators are functions; Agda decides equality of two of them by
applying and normalizing --- for a product of forty gate matrices, a sum over
every intermediate index, \alert{exponential in the word}. It happens whenever two
readings are not \emph{syntactically} equal, for invisible reasons
(\aid{3} vs \aid{₃₊ 0}).
\begin{itemize}
  \item \aid{⟦\_⟧ₒ} is \alert{\texttt{abstract}}: readings compared by comparing
        words; defining equations exported as lemmas;
  \item $\sqrt2^{\,n}$ evaluated strictly level by level (lazily,
        $\sqrt2^{27}$ alone exhausts memory);
  \item never a lemma taking \aid{Same (… u …)} with a free word: pass a
        record (\aid{Evaluated}) instead;
  \item a 4-wire evaluation is \alert{not an option} (36 gates: no answer in
        10 min) --- use structure.
\end{itemize}
\end{column}
\begin{column}{0.47\textwidth}
\centering
\begin{tikzpicture}
\begin{axis}[xbar, width=6.6cm, height=5.0cm, xmode=log, xmin=0.5, xmax=6000,
  symbolic y coords={80-gate equation,Soundness.Relators,FiveQubit/LemmaD7,FourQubit/LemmaD5},
  ytick=data, y tick label style={font=\scriptsize\ttfamily},
  x tick label style={font=\scriptsize}, xlabel={\scriptsize seconds (log)},
  bar width=5pt, enlarge y limits=0.2, axis lines*=left,
  legend style={font=\tiny,at={(0.98,0.03)},anchor=south east}]
\addplot[fill=gray!40,draw=gray] coordinates
  {(3600,FourQubit/LemmaD5) (300,FiveQubit/LemmaD7) (525.6,Soundness.Relators) (158,80-gate equation)};
\addplot[fill=paperC!60,draw=paperC] coordinates
  {(54,FourQubit/LemmaD5) (26,FiveQubit/LemmaD7) (1.2,Soundness.Relators) (0.9,80-gate equation)};
\legend{before,after}
\end{axis}
\end{tikzpicture}

{\scriptsize typechecking times reported in the commit history (``about an hour'',
``about five minutes'', ``under a second'' rounded)}
\end{column}
\end{columns}
\end{frame}

\begin{code}[hide]%
\>[0]\AgdaKeyword{module}\AgdaSpace{}%
\AgdaModule{slides.Tsinghua-PL-seminar.sections.newothers}\AgdaSpace{}%
\AgdaKeyword{where}\<%
\\
%
\\[\AgdaEmptyExtraSkip]%
\>[0]\AgdaKeyword{open}\AgdaSpace{}%
\AgdaKeyword{import}\AgdaSpace{}%
\AgdaModule{Data.Nat}\AgdaSpace{}%
\AgdaKeyword{using}\AgdaSpace{}%
\AgdaSymbol{(}\AgdaDatatype{ℕ}\AgdaSpace{}%
\AgdaSymbol{;}\AgdaSpace{}%
\AgdaOperator{\AgdaDatatype{\AgdaUnderscore{}≤\AgdaUnderscore{}}}\AgdaSpace{}%
\AgdaSymbol{;}\AgdaSpace{}%
\AgdaOperator{\AgdaPrimitive{\AgdaUnderscore{}*\AgdaUnderscore{}}}\AgdaSpace{}%
\AgdaSymbol{;}\AgdaSpace{}%
\AgdaOperator{\AgdaPrimitive{\AgdaUnderscore{}+\AgdaUnderscore{}}}\AgdaSpace{}%
\AgdaSymbol{;}\AgdaSpace{}%
\AgdaOperator{\AgdaFunction{\AgdaUnderscore{}\textasciicircum{}\AgdaUnderscore{}}}\AgdaSymbol{)}\<%
\\
\>[0]\AgdaKeyword{open}\AgdaSpace{}%
\AgdaKeyword{import}\AgdaSpace{}%
\AgdaModule{Data.Bool}\AgdaSpace{}%
\AgdaKeyword{using}\AgdaSpace{}%
\AgdaSymbol{(}\AgdaInductiveConstructor{true}\AgdaSymbol{)}\<%
\\
\>[0]\AgdaKeyword{open}\AgdaSpace{}%
\AgdaKeyword{import}\AgdaSpace{}%
\AgdaModule{Data.Fin}\AgdaSpace{}%
\AgdaKeyword{using}\AgdaSpace{}%
\AgdaSymbol{(}\AgdaDatatype{Fin}\AgdaSymbol{)}\<%
\\
\>[0]\AgdaKeyword{open}\AgdaSpace{}%
\AgdaKeyword{import}\AgdaSpace{}%
\AgdaModule{Data.Product}\AgdaSpace{}%
\AgdaKeyword{using}\AgdaSpace{}%
\AgdaSymbol{(}\AgdaOperator{\AgdaFunction{\AgdaUnderscore{}×\AgdaUnderscore{}}}\AgdaSymbol{)}\<%
\\
\>[0]\AgdaKeyword{open}\AgdaSpace{}%
\AgdaKeyword{import}\AgdaSpace{}%
\AgdaModule{Relation.Binary.PropositionalEquality}\AgdaSpace{}%
\AgdaKeyword{using}\AgdaSpace{}%
\AgdaSymbol{(}\AgdaOperator{\AgdaDatatype{\AgdaUnderscore{}≡\AgdaUnderscore{}}}\AgdaSymbol{)}\<%
\\
\>[0]\AgdaKeyword{open}\AgdaSpace{}%
\AgdaKeyword{import}\AgdaSpace{}%
\AgdaModule{Function.Bundles}\AgdaSpace{}%
\AgdaKeyword{using}\AgdaSpace{}%
\AgdaSymbol{(}\AgdaOperator{\AgdaFunction{\AgdaUnderscore{}⇔\AgdaUnderscore{}}}\AgdaSymbol{)}\<%
\\
%
\\[\AgdaEmptyExtraSkip]%
\>[0]\AgdaKeyword{infix}\AgdaSpace{}%
\AgdaNumber{4}\AgdaSpace{}%
\AgdaOperator{\AgdaPostulate{\AgdaUnderscore{}≐\AgdaUnderscore{}}}\AgdaSpace{}%
\AgdaOperator{\AgdaPostulate{\AgdaUnderscore{}⊢\AgdaUnderscore{}≈\AgdaUnderscore{}}}\AgdaSpace{}%
\AgdaOperator{\AgdaPostulate{\AgdaUnderscore{}≋\AgdaUnderscore{}}}\<%
\\
%
\\[\AgdaEmptyExtraSkip]%
\>[0]\AgdaKeyword{postulate}\<%
\\
\>[0][@{}l@{\AgdaIndent{0}}]%
\>[2]\AgdaPostulate{Bits}\AgdaSpace{}%
\AgdaSymbol{:}\AgdaSpace{}%
\AgdaDatatype{ℕ}\AgdaSpace{}%
\AgdaSymbol{→}\AgdaSpace{}%
\AgdaPrimitive{Set}\<%
\\
%
\>[2]\AgdaPostulate{ℤ₈}\AgdaSpace{}%
\AgdaSymbol{:}\AgdaSpace{}%
\AgdaPrimitive{Set}\<%
\\
%
\>[2]\AgdaPostulate{Circuit}\AgdaSpace{}%
\AgdaSymbol{:}\AgdaSpace{}%
\AgdaDatatype{ℕ}\AgdaSpace{}%
\AgdaSymbol{→}\AgdaSpace{}%
\AgdaPrimitive{Set}\<%
\\
%
\>[2]\AgdaOperator{\AgdaPostulate{⟦\AgdaUnderscore{}⟧}}\AgdaSpace{}%
\AgdaSymbol{:}\AgdaSpace{}%
\AgdaSymbol{∀}\AgdaSpace{}%
\AgdaSymbol{\{}\AgdaBound{n}\AgdaSymbol{\}}\AgdaSpace{}%
\AgdaSymbol{→}\AgdaSpace{}%
\AgdaPostulate{Circuit}\AgdaSpace{}%
\AgdaBound{n}\AgdaSpace{}%
\AgdaSymbol{→}\AgdaSpace{}%
\AgdaPrimitive{Set}\<%
\\
%
\>[2]\AgdaOperator{\AgdaPostulate{\AgdaUnderscore{}≐\AgdaUnderscore{}}}\AgdaSpace{}%
\AgdaOperator{\AgdaPostulate{\AgdaUnderscore{}≋\AgdaUnderscore{}}}\AgdaSpace{}%
\AgdaSymbol{:}\AgdaSpace{}%
\AgdaPrimitive{Set}\AgdaSpace{}%
\AgdaSymbol{→}\AgdaSpace{}%
\AgdaPrimitive{Set}\AgdaSpace{}%
\AgdaSymbol{→}\AgdaSpace{}%
\AgdaPrimitive{Set}\<%
\\
%
\>[2]\AgdaOperator{\AgdaPostulate{\AgdaUnderscore{}⊢\AgdaUnderscore{}≈\AgdaUnderscore{}}}\AgdaSpace{}%
\AgdaSymbol{:}\AgdaSpace{}%
\AgdaSymbol{(}\AgdaBound{n}\AgdaSpace{}%
\AgdaSymbol{:}\AgdaSpace{}%
\AgdaDatatype{ℕ}\AgdaSymbol{)}\AgdaSpace{}%
\AgdaSymbol{→}\AgdaSpace{}%
\AgdaPostulate{Circuit}\AgdaSpace{}%
\AgdaBound{n}\AgdaSpace{}%
\AgdaSymbol{→}\AgdaSpace{}%
\AgdaPostulate{Circuit}\AgdaSpace{}%
\AgdaBound{n}\AgdaSpace{}%
\AgdaSymbol{→}\AgdaSpace{}%
\AgdaPrimitive{Set}\<%
\\
%
\>[2]\AgdaPostulate{NormalForms}\AgdaSpace{}%
\AgdaSymbol{:}\AgdaSpace{}%
\AgdaDatatype{ℕ}\AgdaSpace{}%
\AgdaSymbol{→}\AgdaSpace{}%
\AgdaPrimitive{Set₁}\<%
\\
%
\>[2]\AgdaPostulate{Poly}\AgdaSpace{}%
\AgdaSymbol{:}\AgdaSpace{}%
\AgdaDatatype{ℕ}\AgdaSpace{}%
\AgdaSymbol{→}\AgdaSpace{}%
\AgdaDatatype{ℕ}\AgdaSpace{}%
\AgdaSymbol{→}\AgdaSpace{}%
\AgdaPrimitive{Set}\<%
\\
%
\>[2]\AgdaPostulate{idPS}\AgdaSpace{}%
\AgdaSymbol{:}\AgdaSpace{}%
\AgdaPrimitive{Set}\<%
\\
%
\>[2]\AgdaPostulate{decideᶜ}\AgdaSpace{}%
\AgdaSymbol{:}\AgdaSpace{}%
\AgdaSymbol{∀}\AgdaSpace{}%
\AgdaSymbol{\{}\AgdaBound{n}\AgdaSymbol{\}}\AgdaSpace{}%
\AgdaSymbol{→}\AgdaSpace{}%
\AgdaPostulate{Circuit}\AgdaSpace{}%
\AgdaBound{n}\AgdaSpace{}%
\AgdaSymbol{→}\AgdaSpace{}%
\AgdaPrimitive{Set}\<%
\\
%
\>[2]\AgdaPostulate{value}\AgdaSpace{}%
\AgdaPostulate{cost}\AgdaSpace{}%
\AgdaSymbol{:}\AgdaSpace{}%
\AgdaPrimitive{Set}\AgdaSpace{}%
\AgdaSymbol{→}\AgdaSpace{}%
\AgdaDatatype{ℕ}\<%
\\
%
\>[2]\AgdaPostulate{decideBound}\AgdaSpace{}%
\AgdaSymbol{:}\AgdaSpace{}%
\AgdaDatatype{ℕ}\AgdaSpace{}%
\AgdaSymbol{→}\AgdaSpace{}%
\AgdaDatatype{ℕ}\AgdaSpace{}%
\AgdaSymbol{→}\AgdaSpace{}%
\AgdaDatatype{ℕ}\<%
\\
%
\>[2]\AgdaPostulate{volumeBound}\AgdaSpace{}%
\AgdaSymbol{:}\AgdaSpace{}%
\AgdaDatatype{ℕ}\AgdaSpace{}%
\AgdaSymbol{→}\AgdaSpace{}%
\AgdaDatatype{ℕ}\<%
\\
%
\>[2]\AgdaPostulate{length}\AgdaSpace{}%
\AgdaSymbol{:}\AgdaSpace{}%
\AgdaSymbol{∀}\AgdaSpace{}%
\AgdaSymbol{\{}\AgdaBound{n}\AgdaSymbol{\}}\AgdaSpace{}%
\AgdaSymbol{→}\AgdaSpace{}%
\AgdaPostulate{Circuit}\AgdaSpace{}%
\AgdaBound{n}\AgdaSpace{}%
\AgdaSymbol{→}\AgdaSpace{}%
\AgdaDatatype{ℕ}\<%
\\
%
\>[2]\AgdaPostulate{proof-elided}\AgdaSpace{}%
\AgdaSymbol{:}\AgdaSpace{}%
\AgdaSymbol{∀}\AgdaSpace{}%
\AgdaSymbol{\{}\AgdaBound{A}\AgdaSpace{}%
\AgdaSymbol{:}\AgdaSpace{}%
\AgdaPrimitive{Set}\AgdaSymbol{\}}\AgdaSpace{}%
\AgdaSymbol{→}\AgdaSpace{}%
\AgdaBound{A}\<%
\\
%
\\[\AgdaEmptyExtraSkip]%
\>[0]\AgdaKeyword{private}\AgdaSpace{}%
\AgdaKeyword{variable}\<%
\\
\>[0][@{}l@{\AgdaIndent{0}}]%
\>[2]\AgdaGeneralizable{n}\AgdaSpace{}%
\AgdaSymbol{:}\AgdaSpace{}%
\AgdaDatatype{ℕ}\<%
\end{code}
%% Part 6e — CNOT-dihedral, path sums, summary.

\begin{frame}{CNOT-dihedral circuits: phase polynomials \New}
\begin{columns}[c]
\begin{column}{0.50\textwidth}
\renewcommand{\AgdaCodeStyle}{\scriptsize}
\small Amy, Chen, Ross (QPL'17): gates $\omega$, $X$, $T$, $CNOT$, SWAP;
relations $R_1$--$R_{13}$ + symmetry laws. An operator is a \alert{phase
polynomial} $|x\rangle\mapsto\omega^{p(x)}|f(x)\rangle$, $f$ affine:
\begin{code}%
\>[0]\AgdaFunction{Op}\AgdaSpace{}%
\AgdaSymbol{:}\AgdaSpace{}%
\AgdaDatatype{ℕ}\AgdaSpace{}%
\AgdaSymbol{→}\AgdaSpace{}%
\AgdaPrimitive{Set}\<%
\\
\>[0]\AgdaFunction{Op}\AgdaSpace{}%
\AgdaBound{n}\AgdaSpace{}%
\AgdaSymbol{=}\AgdaSpace{}%
\AgdaPostulate{Bits}\AgdaSpace{}%
\AgdaBound{n}\AgdaSpace{}%
\AgdaSymbol{→}\AgdaSpace{}%
\AgdaPostulate{Bits}\AgdaSpace{}%
\AgdaBound{n}\AgdaSpace{}%
\AgdaOperator{\AgdaFunction{×}}\AgdaSpace{}%
\AgdaPostulate{ℤ₈}\<%
\end{code}
\begin{itemize}
  \small
  \item composition adds phases: \alert{no sums}; each relation is \aid{refl}
        on a table of $\le16$ entries;
  \item soundness; 24 + 4 derived rules; \alert{complete on one qubit}
        ($\omega^kT^aX^e$, 128 tables);
  \item $n$ qubits: complete given the paper's imported lemmas (Lafont's
        affine normal forms, \ldots) as a record:
\end{itemize}
\begin{code}%
\>[0]\AgdaKeyword{module}\AgdaSpace{}%
\AgdaModule{Complete}\AgdaSpace{}%
\AgdaSymbol{(}\AgdaBound{NFs}\AgdaSpace{}%
\AgdaSymbol{:}\AgdaSpace{}%
\AgdaPostulate{NormalForms}\AgdaSpace{}%
\AgdaGeneralizable{n}\AgdaSymbol{)}\AgdaSpace{}%
\AgdaKeyword{where}\<%
\\
\>[0][@{}l@{\AgdaIndent{0}}]%
\>[2]\AgdaFunction{complete}\AgdaSpace{}%
\AgdaSymbol{:}\AgdaSpace{}%
\AgdaSymbol{∀}\AgdaSpace{}%
\AgdaSymbol{\{}\AgdaBound{w}\AgdaSpace{}%
\AgdaBound{v}\AgdaSpace{}%
\AgdaSymbol{:}\AgdaSpace{}%
\AgdaPostulate{Circuit}\AgdaSpace{}%
\AgdaBound{n}\AgdaSymbol{\}}\AgdaSpace{}%
\AgdaSymbol{→}\AgdaSpace{}%
\AgdaOperator{\AgdaPostulate{⟦}}\AgdaSpace{}%
\AgdaBound{w}\AgdaSpace{}%
\AgdaOperator{\AgdaPostulate{⟧}}\AgdaSpace{}%
\AgdaOperator{\AgdaPostulate{≐}}\AgdaSpace{}%
\AgdaOperator{\AgdaPostulate{⟦}}\AgdaSpace{}%
\AgdaBound{v}\AgdaSpace{}%
\AgdaOperator{\AgdaPostulate{⟧}}\AgdaSpace{}%
\AgdaSymbol{→}\AgdaSpace{}%
\AgdaBound{n}\AgdaSpace{}%
\AgdaOperator{\AgdaPostulate{⊢}}\AgdaSpace{}%
\AgdaBound{w}\AgdaSpace{}%
\AgdaOperator{\AgdaPostulate{≈}}\AgdaSpace{}%
\AgdaBound{v}\<%
\end{code}
\begin{code}[hide]%
%
\>[2]\AgdaFunction{complete}\AgdaSpace{}%
\AgdaSymbol{=}\AgdaSpace{}%
\AgdaPostulate{proof-elided}\<%
\end{code}
\end{column}
\begin{column}{0.47\textwidth}
\centering
\begin{quantikz}[row sep=0.3cm,column sep=0.22cm,font=\scriptsize]
\lstick{$x_1$} & \ctrl{1} & & \ctrl{1} & \rstick[2]{$\omega^{\,x_0\oplus x_1}\,|x_0,x_1\rangle$}\\
\lstick{$x_0$} & \targ{} & \gate{T} & \targ{} &
\end{quantikz}

\vspace{6pt}
{\scriptsize $CNOT\cdot(I\otimes T)\cdot CNOT$: a diagonal, i.e.\ $f=\mathrm{id}$, phase
$p(x)=x_0\oplus x_1$ --- tables, not matrices}

\vspace{10pt}
\begin{keybox}
\small Same circuit layer, same congruence: ω's centrality is the theorem
\aid{Central-Scalars.comm₀}; 2\,829 lines; checks in under a minute.
\end{keybox}
\end{column}
\end{columns}
\end{frame}

\begin{frame}{Path sums, with polynomial bounds \New}
\begin{columns}[c]
\begin{column}{0.52\textwidth}
\renewcommand{\AgdaCodeStyle}{\scriptsize}
\small Amy (QPL'18), \emph{Towards large-scale functional verification of
universal quantum circuits} --- \Branch{popl}, 87k lines, \alert{no assumptions}:
\[
  |x\rangle \;\mapsto\; 2^{-k/2}\!\!\sum_{y\in\mathbb{B}^m}\! e^{2\pi i\,P(x,y)}\,|f(x,y)\rangle
\]
\vspace{-8pt}
\begin{itemize}
  \small
  \item exact amplitudes in $\mathbb{Z}[\zeta]=\mathbb{Z}[X]/(X^H+1)$;
        rules [Elim], [$\omega$], [HH], [Case];
  \item a \alert{cost model}: each algorithm written once in a cost monad;
        value correctness + polynomial cost:
\end{itemize}
\begin{code}%
\>[0]\AgdaKeyword{record}\AgdaSpace{}%
\AgdaRecord{Corollary-4-4}\AgdaSpace{}%
\AgdaSymbol{\{}\AgdaBound{n}\AgdaSpace{}%
\AgdaSymbol{:}\AgdaSpace{}%
\AgdaDatatype{ℕ}\AgdaSymbol{\}}\AgdaSpace{}%
\AgdaSymbol{(}\AgdaBound{C}\AgdaSpace{}%
\AgdaSymbol{:}\AgdaSpace{}%
\AgdaPostulate{Circuit}\AgdaSpace{}%
\AgdaBound{n}\AgdaSymbol{)}\AgdaSpace{}%
\AgdaSymbol{:}\AgdaSpace{}%
\AgdaPrimitive{Set}\AgdaSpace{}%
\AgdaKeyword{where}\<%
\\
\>[0][@{}l@{\AgdaIndent{0}}]%
\>[2]\AgdaKeyword{field}\<%
\\
\>[2][@{}l@{\AgdaIndent{0}}]%
\>[4]\AgdaField{decides}%
\>[15]\AgdaSymbol{:}\AgdaSpace{}%
\AgdaPostulate{value}\AgdaSpace{}%
\AgdaSymbol{(}\AgdaPostulate{decideᶜ}\AgdaSpace{}%
\AgdaBound{C}\AgdaSymbol{)}\AgdaSpace{}%
\AgdaOperator{\AgdaDatatype{≡}}\AgdaSpace{}%
\AgdaInductiveConstructor{true}\AgdaSpace{}%
\AgdaOperator{\AgdaFunction{⇔}}\AgdaSpace{}%
\AgdaOperator{\AgdaPostulate{⟦}}\AgdaSpace{}%
\AgdaBound{C}\AgdaSpace{}%
\AgdaOperator{\AgdaPostulate{⟧}}\AgdaSpace{}%
\AgdaOperator{\AgdaPostulate{≋}}\AgdaSpace{}%
\AgdaPostulate{idPS}\<%
\\
%
\>[4]\AgdaField{polynomial}\AgdaSpace{}%
\AgdaSymbol{:}\AgdaSpace{}%
\AgdaPostulate{cost}\AgdaSpace{}%
\AgdaSymbol{(}\AgdaPostulate{decideᶜ}\AgdaSpace{}%
\AgdaBound{C}\AgdaSymbol{)}\AgdaSpace{}%
\AgdaOperator{\AgdaDatatype{≤}}\AgdaSpace{}%
\AgdaPostulate{decideBound}\AgdaSpace{}%
\AgdaBound{n}\AgdaSpace{}%
\AgdaSymbol{(}\AgdaPostulate{length}\AgdaSpace{}%
\AgdaBound{C}\AgdaSymbol{)}\<%
\\
%
\>[4]\AgdaField{volume}%
\>[15]\AgdaSymbol{:}\AgdaSpace{}%
\AgdaPostulate{cost}\AgdaSpace{}%
\AgdaSymbol{(}\AgdaPostulate{decideᶜ}\AgdaSpace{}%
\AgdaBound{C}\AgdaSymbol{)}\AgdaSpace{}%
\AgdaOperator{\AgdaDatatype{≤}}\AgdaSpace{}%
\AgdaPostulate{volumeBound}\AgdaSpace{}%
\AgdaSymbol{(}\AgdaBound{n}\AgdaSpace{}%
\AgdaOperator{\AgdaPrimitive{*}}\AgdaSpace{}%
\AgdaPostulate{length}\AgdaSpace{}%
\AgdaBound{C}\AgdaSymbol{)}\<%
\\
\>[0]\AgdaFunction{decideBound-def}\AgdaSpace{}%
\AgdaSymbol{:}\AgdaSpace{}%
\AgdaSymbol{∀}\AgdaSpace{}%
\AgdaBound{n}\AgdaSpace{}%
\AgdaBound{ℓ}\AgdaSpace{}%
\AgdaSymbol{→}\AgdaSpace{}%
\AgdaPostulate{decideBound}\AgdaSpace{}%
\AgdaBound{n}\AgdaSpace{}%
\AgdaBound{ℓ}\AgdaSpace{}%
\AgdaOperator{\AgdaDatatype{≡}}\AgdaSpace{}%
\AgdaNumber{313}\AgdaSpace{}%
\AgdaOperator{\AgdaPrimitive{*}}\AgdaSpace{}%
\AgdaSymbol{(}\AgdaNumber{3}\AgdaSpace{}%
\AgdaOperator{\AgdaPrimitive{+}}\AgdaSpace{}%
\AgdaSymbol{(}\AgdaBound{n}\AgdaSpace{}%
\AgdaOperator{\AgdaPrimitive{+}}\AgdaSpace{}%
\AgdaBound{ℓ}\AgdaSymbol{))}\AgdaSpace{}%
\AgdaOperator{\AgdaFunction{\textasciicircum{}}}\AgdaSpace{}%
\AgdaNumber{12}\<%
\end{code}
\begin{code}[hide]%
\>[0]\AgdaFunction{decideBound-def}\AgdaSpace{}%
\AgdaSymbol{=}\AgdaSpace{}%
\AgdaPostulate{proof-elided}\<%
\end{code}
\end{column}
\begin{column}{0.45\textwidth}
\footnotesize
\begin{itemize}
  \setlength{\itemsep}{2pt}
  \item Section 5.2 families for \emph{every} size: QFT, $n$-bit Toffoli, Maslov,
        adder, hidden shift; Table 2 rows reproduced exactly (e.g.\ Toffoli$_{50}$:
        190 paths, $T$-count 665);
  \item footnote 2's reduction from UNSAT, as an equivalence;
  \item beyond Clifford: normal forms \alert{not unique}; a complete but
        exponential decision procedure.
\end{itemize}
\begin{keybox}[title={Errata in the paper, each with a checked counterexample}]
\footnotesize
\begin{itemize}
  \item Lemma 4.2 needs $Q$ odd at some input ($Q=2x_1$ fails);
  \item Prop.\ 2.14's bound is $\max(2,k)$, not $k$;
  \item Prop.\ 2.7: well-formedness is not preserved by composition.
\end{itemize}
\end{keybox}
{\scriptsize independent of the presentation library: a second style of completeness
(decision by rewriting, not normal forms of a group)}
\end{column}
\end{columns}
\end{frame}

\begin{frame}{Summary: in the paper vs.\ new since \Mixed}
\centering\footnotesize
\renewcommand{\arraystretch}{1.15}
\begin{tabular}{@{}p{3.3cm}p{5.0cm}p{5.6cm}@{}}
\toprule
 & \textcolor{paperC}{\textbf{POPL submission} (tag, 22k lines)} & \textcolor{newC}{\textbf{since} (other branches, up to 194k lines each)}\\
\midrule
core & words, inductive congruences, NF bundles, \aid{by-normalization}, R--S engine, towers, $\times$ $\rtimes$ $\ast_M$ ext., circuits (3 structural rules) & 0-ary gates, \aid{Central-Scalars}, 4 structural rules; independence of axioms; central / factor-set extensions; generic circuit towers\\
symmetric \& friends & $S_n$ (2 semantics), $\mathbb{Z}_n$, trivial, $\mathbb{Z}_N\wr S_n$, Pauli & reorganized; projective Pauli\\
Clifford & qubit / qutrit Clifford+$T$ (1 qubit), $U_3(\mathbb{Z}[\frac12,i])$ --- \emph{syntactic} & qupit Clifford (all $n$, odd $p$), $Sp(2n,\mathbb{Z}_p)$, Selinger's qubit Clifford, exact versions with scalars\\
matrices & --- (stated limitation) & $U_n(\mathbb{Z}[\frac12,i])$ all $n$; $U_{\le4}(\mathbb{Z}[\frac1{\sqrt2},i])$; 2-qubit Clifford+$T$; Real-Clifford+$CH$ over $\mathbb{Z}[\sqrt2]$; $\mathbb{Z}/17$ model\\
other semantics & --- & phase polynomials (CNOT-dihedral); path sums with cost bounds\\
in progress & --- & 3-qubit Clifford+$CS$; Real-Clifford+$CH$ Lemma 8.8 at 3--4 qubits; gate-specific minimality\\
\bottomrule
\end{tabular}
\end{frame}

%% Part 7 — related work, lessons, outlook (paper §6–§8, plus new).
\PaperPart
\section{Related work, lessons, outlook}

\begin{frame}{Where this sits \InPaper}
\centering\footnotesize
\setlength{\tabcolsep}{5pt}
\begin{tabular}{@{}llcccc@{}}
\toprule
Development & Prover & Circuits & Pres. & NF & Compl. \\
\midrule
\textbf{This library} & Agda & \checkmark & \checkmark & \checkmark & \checkmark \\
Bian--Selinger developments & Agda & \checkmark & \checkmark & \checkmark & syntactic \\
\textsc{sqir}/\textsc{voqc}, VyZX & Rocq & \checkmark & -- & -- & sound only \\
\textsc{Qwire}, CoqQ, Qbricks, QHLProver, IMD & Rocq/Why3/Isabelle & \checkmark & -- & -- & -- \\
$\Pi$ as free rig groupoid & Agda (HoTT) & \checkmark & categorical & \checkmark & classical \\
$\sqrt{\Pi}$ & Agda (partial) & \checkmark & categorical & -- & 2q Clifford \\
mathlib group theory & Lean & -- & \checkmark & -- & -- \\
MathComp / Odd Order & Rocq & -- & partial & -- & -- \\
Free groups (AFP), IsaFoR completion & Isabelle & -- & \checkmark\,/\,-- & \checkmark & -- \\
\bottomrule
\end{tabular}

\vspace{8pt}
\begin{columns}[T]
\begin{column}{0.47\textwidth}
\small \textbf{vs.\ mathlib}: presentation \emph{theory} (quotients, Coxeter);
no infrastructure to prove a \emph{given} relation set presents a \emph{given}
group --- no coset enumeration, no verified transversals.
\end{column}
\begin{column}{0.47\textwidth}
\small \textbf{vs.\ rewriting}: we never orient or prove confluence; a normal form is
any verified \emph{function} --- sidestepping Squier's obstruction (decidable word
problem $\not\Rightarrow$ finite convergent presentation).
\end{column}
\end{columns}
\end{frame}

\begin{frame}{Lessons for proof engineering \Mixed}
\begin{columns}[T]
\begin{column}{0.48\textwidth}
\begin{paperbox}[title=From the paper]
\small
\begin{enumerate}
  \setlength{\itemsep}{2pt}
  \item \textbf{Setoids over syntax} beat quotients when you need to compute and
        to map \emph{into} the syntax.
  \item \textbf{Certificates, not search}: finite tables + five equations;
        the engine is proved once.
  \item \textbf{Evaluation as a tactic}: map both sides to canonical data, compare by
        \aid{refl} --- no reflection needed.
  \item \textbf{Index engineering}: left-biased successors make coverage a matter
        of unification.
  \item \textbf{Constructions in pairs}: syntax theorem + semantic theorem, against
        stdlib bundles.
\end{enumerate}
\end{paperbox}
\end{column}
\begin{column}{0.48\textwidth}
\begin{newdevbox}[title=Since]
\small
\begin{enumerate}
  \setcounter{enumi}{5}
  \setlength{\itemsep}{2pt}
  \item \textbf{Tame the conversion checker}: \texttt{abstract} semantics,
        strict arithmetic, records instead of unfolded propositions.
  \item \textbf{Simulate, decide, then prove at symbolic width}: a transcription
        bug invisible on three qubits surfaced only in a symbolic-width proof.
  \item \textbf{Formalization finds errata}: in Clément's appendix and in Amy's
        path-sum paper --- each confirmed by a decision procedure or a
        counterexample.
  \item \textbf{Reuse compounds}: Theorem 4.10, the Clifford+$T$ 2-qubit proof and
        the qupit semidirect product all run on the paper's generic engine.
\end{enumerate}
\end{newdevbox}
\end{column}
\end{columns}
\end{frame}

\begin{frame}{Outlook \Mixed}
\begin{columns}[c]
\begin{column}{0.55\textwidth}
\begin{itemize}
  \setlength{\itemsep}{3pt}
  \item \textbf{Upstreaming} the \texttt{ForStdlib} constructions (semidirect,
        amalgamated, central and factor-set extensions) to agda-stdlib.
  \item \textbf{Close the loops}: Lemma 8.8 on 3--4 qubits; Theorem 4.4 itself
        (a presentation of $O_N(\mathbb{Z}[\frac1{\sqrt2}])$, Fang--Heunen--Kaarsgaard);
        the remaining R--S steps of 3-qubit Clifford+$CS$.
  \item \textbf{Connect} the syntactic Clifford+$T$ amalgams of the paper to the
        new matrix semantics.
  \item \textbf{Minimality} of gate-specific rule sets with the independence
        engine --- e.g.\ Blake's simplified near-Clifford presentations.
  \item The ``264-page appendix'' should become \alert{ordinary, reusable,
        machine-checked mathematics}.
\end{itemize}
\end{column}
\begin{column}{0.42\textwidth}
\centering
\begin{tikzpicture}[x=9mm,y=9mm,>={Stealth[length=1.8mm]}]
  \node[tag=paperC,text width=30mm] (a) at (0,3) {presentations as data};
  \node[tag=paperC,text width=30mm] (b) at (0,2) {normal forms as bundles};
  \node[tag=paperC,text width=30mm] (c) at (0,1) {coset tables as certificates};
  \node[tag=newC,text width=30mm] (d) at (0,0) {matrices as tries,\\decided by \aid{refl}};
  \node[tag=accent,text width=30mm,font=\small\bfseries] (e) at (0,-1.4) {complete rule sets,\\machine-checked};
  \draw[->,thick,gray] (d) -- (e);
\end{tikzpicture}
\end{column}
\end{columns}
\end{frame}

\begin{frame}[standout]
Thank you!\\[10pt]
{\normalsize Questions?}\\[16pt]
{\footnotesize Everything stated as a theorem is checked by Agda 2.8 + stdlib 2.4:\\
\texttt{--safe}, no postulates; conditional results take their hypotheses as
module parameters.}
\end{frame}


\end{document}
